% !TeX program = xelatex
% 由扫描件人工逐行转录。每个 \BookLine 对应原书一行，可直接编辑。
\documentclass[10pt]{article}
\usepackage[paperwidth=140mm,paperheight=210mm,margin=16mm]{geometry}
\usepackage{fontspec,xeCJK,graphicx,xcolor,tikz}
\usepackage[absolute,overlay]{textpos}
\usepackage[unicode,pdfpagelayout=TwoPageRight,pdftitle={阿尔刻提斯之歌（扫描部分重排）},pdfauthor={}]{hyperref}
\setmainfont{Liberation Serif}
\IfFileExists{/usr/local/share/fonts/f/FZShuSong_Z01_Regular.ttf}{
  \setCJKmainfont{FZShuSong_Z01_Regular.ttf}[Path=/usr/local/share/fonts/f/]
}{\setCJKmainfont{Noto Serif CJK SC}}
\IfFileExists{/usr/local/share/fonts/f/FZKai_Z03_Regular.ttf}{
  \newCJKfontfamily\kaifont{FZKai_Z03_Regular.ttf}[Path=/usr/local/share/fonts/f/]
}{\newCJKfontfamily\kaifont{FangSong}}
\IfFileExists{/usr/local/share/fonts/f/FZHei_B01_Regular.ttf}{
  \newCJKfontfamily\heifont{FZHei_B01_Regular.ttf}[Path=/usr/local/share/fonts/f/]
}{\newCJKfontfamily\heifont{Noto Sans CJK SC}}
\pagestyle{empty}
\setlength{\parindent}{0pt}
\setlength{\TPHorizModule}{1mm}
\setlength{\TPVertModule}{1mm}
\textblockorigin{0mm}{0mm}
\xeCJKsetup{PunctStyle=plain}
\newsavebox{\rowbox}
% y（毫米）, 缩进（毫米）, 字体, 是否铺满行, 文字。
\newcommand{\BookLine}[5]{%
  \begin{textblock*}{108mm}(16mm,#1mm)
  \fontsize{10.5}{14}\selectfont #3%
  \hspace*{#2mm}%
  \begingroup
  \ifnum#4=1 \renewcommand\CJKglue{\hskip 0pt plus 1fill}\fi
  \sbox{\rowbox}{#5}%
  \ifdim\wd\rowbox>\dimexpr108mm-#2mm\relax
    \resizebox{\dimexpr108mm-#2mm\relax}{!}{\usebox{\rowbox}}%
  \else
    \ifnum#4=1
      \makebox[\dimexpr108mm-#2mm\relax][s]{#5}%
    \else\usebox{\rowbox}\fi
  \fi
  \endgroup
  \end{textblock*}%
}
\newcommand{\Vertical}[5]{%
  \begin{textblock*}{8mm}(#1mm,#2mm)
    #3\fontsize{#4}{#5}\selectfont\centering
}
\newcommand{\EndVertical}{\end{textblock*}}
\newcommand{\BookNumber}[1]{%
  \begin{textblock*}{140mm}(0mm,192mm)
  \centering\fontsize{12}{12}\selectfont\itshape #1
  \end{textblock*}
}
\newcommand{\ActOrnament}[1][一]{%
  \Vertical{66}{28}{}{14}{21}第\\#1\\幕\EndVertical
  \begin{textblock*}{140mm}(0mm,63mm)
    \centering\textcolor{black!65}{\fontsize{10}{10}\selectfont\textbullet}
  \end{textblock*}
  \begin{textblock*}{75mm}(38mm,96mm)
  \begin{tikzpicture}[x=1mm,y=1mm,line width=.35pt,draw=black!65]
    % 后方拱门。
    \draw (16,0) -- (57,0) -- (57,39)
      .. controls (57,48) and (53,55) .. (48,59);
    \draw (16,0) -- (16,5);
    % 前方拱门。
    \draw (0,5) -- (48,5) -- (48,49)
      arc[start angle=0,end angle=180,radius=24] -- cycle;
    % 沿拱门上方逐字排出书名，不依赖扫描图。
    \foreach \letter [count=\i from 0] in {T,h,e,A,l,c,e,s,t,i,a,d}{%
      \pgfmathsetmacro{\a}{153-\i*11.45}
      \node[rotate=\a-90,text=black!80,font=\fontsize{9}{9}\selectfont]
        at ({24+35*cos(\a)},{49+35*sin(\a)}) {\letter};
    }
  \end{tikzpicture}
  \end{textblock*}
}
\newcommand{\WorkOrnament}[1]{%
  \Vertical{58}{26}{}{14}{21}#1\EndVertical
  \begin{textblock*}{140mm}(0mm,63mm)
    \centering\textcolor{black!65}{\fontsize{10}{10}\selectfont\textbullet}
  \end{textblock*}
  \begin{textblock*}{75mm}(38mm,96mm)
  \begin{tikzpicture}[x=1mm,y=1mm,line width=.35pt,draw=black!65]
    \draw (16,0) -- (57,0) -- (57,39) .. controls (57,48) and (53,55) .. (48,59);
    \draw (16,0) -- (16,5);
    \draw (0,5) -- (48,5) -- (48,49) arc[start angle=0,end angle=180,radius=24] -- cycle;
    \foreach \letter [count=\i from 0] in {T,h,e,A,l,c,e,s,t,i,a,d}{%
      \pgfmathsetmacro{\a}{153-\i*11.45}
      \node[rotate=\a-90,text=black!80,font=\fontsize{9}{9}\selectfont]
        at ({24+35*cos(\a)},{49+35*sin(\a)}) {\letter};
    }
  \end{tikzpicture}
  \end{textblock*}
}
\begin{document}

% 原书页码：空白
\null
\newpage
% 原书页码：目录
\null
\Vertical{15}{28}{}{14}{21}目\\录\EndVertical
\BookLine{88.00}{29.52}{}{0}{前言 / 1}
\BookLine{105.00}{29.52}{}{0}{人物 / 20}
\BookLine{122.00}{29.52}{}{0}{场景 / 22}
\BookLine{139.00}{29.52}{}{0}{第一幕 / 23}
\BookLine{156.00}{29.52}{}{0}{第二幕 / 81}
\BookLine{173.00}{29.52}{}{0}{第三幕 / 157}
\BookLine{190.00}{29.52}{}{0}{醉酒的女神 / 227}
\newpage
% 原书页码：空白
\null
\newpage
% 原书页码：1
\null
\BookNumber{1}
\Vertical{15}{28}{}{14}{21}前\\言\EndVertical
\BookLine{105.00}{7.38}{}{0}{人们总是将生命看作一条沿着既定路线奔流的河，而在桑}
\BookLine{113.50}{0.00}{}{0}{顿·怀尔德的小说《第八日》的最后几页，我哥哥让其中一位}
\BookLine{122.00}{0.00}{}{0}{主角反驳了这种看法。相较之下，一个人的人生经历就像一张}
\BookLine{130.50}{0.00}{}{0}{织毯，它的长和宽都不确定，两面的纹路也不尽相同。彩色的}
\BookLine{139.00}{0.00}{}{0}{纱线勾勒出故事的轮廓，设计上也许极不对称，却清晰、真实}
\BookLine{147.50}{0.00}{}{0}{而完整。是记录，也是肖像。}
\BookLine{156.00}{7.38}{}{0}{桑顿的命运织毯起初还很小，七八岁时他生活在威斯康}
\BookLine{164.50}{0.00}{}{0}{星州的麦迪逊市，在阅读布尔凡奇的《寓言时代》时第一次听}
\BookLine{173.00}{0.00}{}{0}{到了阿尔刻提斯这个名字。在前基督教时期的希腊神话和歌}
\newpage
% 原书页码：2
\null
\BookNumber{2}
\Vertical{4}{169}{}{6.2}{8}阿\\尔\\刻\\提\\斯\\之\\歌\EndVertical
\BookLine{30.00}{0.00}{}{0}{谣中，国王佩利阿斯的女儿是公主中的公主。她的故事占据}
\BookLine{38.50}{0.00}{}{0}{了他的想象，也俘获了他的心灵。她的白色纱线嵌进了他的}
\BookLine{47.00}{0.00}{}{0}{命运织物中，阿尔刻提斯成为他潜意识中一份温柔的坚持，}
\BookLine{55.50}{0.00}{}{0}{一个有待发掘的灵感素材。她是一个戏剧雏形，直到多年后}
\BookLine{64.00}{0.00}{}{0}{才得以创作出来。}
\BookLine{72.50}{7.38}{}{0}{当怀尔德一家搬到加利福尼亚州伯克利时，阿尔刻提斯}
\BookLine{81.00}{0.00}{}{0}{之线打了个结。桉树林山坡上建有一座宏伟的希腊剧院，是}
\BookLine{89.50}{0.00}{}{0}{学校和小镇生活中生机勃勃的一部分。古典学系每年都会制作}
\BookLine{98.00}{0.00}{}{0}{几部索福克勒斯、埃斯库罗斯和欧里庇得斯的戏剧作品。我们}
\BookLine{106.50}{0.00}{}{0}{的母亲参与了服装组的志愿工作，用模板印出希腊钥匙的形}
\BookLine{115.00}{0.00}{}{0}{状，或是在漂亮的长袍上印上月桂树叶的图案。她给桑顿做了}
\BookLine{123.50}{0.00}{}{0}{一件有贝壳镶边的蓝色小袍子——给哥哥阿莫斯做了一件绿色}
\BookLine{132.00}{0.00}{}{0}{的——然后让他们去扮演雅典民众。桑顿就这样了解到了“整}
\BookLine{140.50}{0.00}{}{0}{个”剧场和古希腊的黄金时代，在这之前，他的全部剧场经验}
\BookLine{149.00}{0.00}{}{0}{不过是坐在密尔沃基剧院的后排看了一场《皆大欢喜》。}
\BookLine{157.50}{7.38}{}{0}{桑顿十岁了，黑头发，蓝眼睛，充满好奇心，精力旺盛。}
\BookLine{166.00}{0.00}{}{0}{早在十岁之前他就掌握了作家们的工具，驯服了二十六个字母}
\BookLine{174.50}{0.00}{}{0}{为自己服务，将它们组成词汇，在适当的时候用来表达自己。}
\newpage
% 原书页码：3
\null
\BookNumber{3}
\Vertical{129}{11}{}{7}{10}前\\言\EndVertical
\BookLine{30.00}{0.00}{}{0}{他早睡早起写作，随着词汇量的增长，他开始创造对话。他给}
\BookLine{38.50}{0.00}{}{0}{我们和邻居的孩子披上借来的纱，哄着我们一起朗诵他的豪言}
\BookLine{47.00}{0.00}{}{0}{壮语。}
\BookLine{55.50}{7.38}{}{0}{四年之后，他成为中国烟台一所教会学校的寄宿生。在一}
\BookLine{64.00}{0.00}{}{0}{所英国人开办的专为牛津、剑桥入学考试做准备的学校里，他}
\BookLine{72.50}{0.00}{}{0}{开始了对古典文学的学习。}
\BookLine{81.00}{7.38}{}{0}{他大学时代的前两年是在俄亥俄州的欧伯林学院度过的。}
\BookLine{89.50}{0.00}{}{0}{在那里他成了查尔斯·H. A. 瓦格教授的得意门生和朋友，而瓦}
\BookLine{98.00}{0.00}{}{0}{格教授本身也是一名出色的古典文学家。1920 年从耶鲁大学毕}
\BookLine{106.50}{0.00}{}{0}{业后，他在美国设立于罗马的古典学研究院待了八个月，阿尔}
\BookLine{115.00}{0.00}{}{0}{刻提斯之线在这些年里有了更清晰的纹路。尽管不是一名研究}
\BookLine{123.50}{0.00}{}{0}{员，他还是会去听讲座，和考古学家们一起去挖掘现场，仔细}
\BookLine{132.00}{0.00}{}{0}{研读莎草纸上的文字。在意大利的博物馆和美术馆里，他捕捉}
\BookLine{140.50}{0.00}{}{0}{到了阿尔刻提斯的影子。}
\BookLine{149.00}{7.38}{}{0}{在二十世纪二三十年代里，这张命运织毯延展得更大，也}
\BookLine{157.50}{0.00}{}{0}{变得更加密实。他开始从事教学工作——先是在劳伦斯维尔学}
\BookLine{166.00}{0.00}{}{0}{校教授法语，然后在芝加哥大学教授比较文学，其中也会讲到}
\BookLine{174.50}{0.00}{}{0}{古希腊的剧作家们——还有巡回演讲以及在好莱坞进行电影剧}
\newpage
% 原书页码：4
\null
\BookNumber{4}
\Vertical{4}{169}{}{6.2}{8}阿\\尔\\刻\\提\\斯\\之\\歌\EndVertical
\BookLine{30.00}{0.00}{}{0}{本创作。他出版了两部短剧集和四部小说，其中一部《圣路易}
\BookLine{38.50}{0.00}{}{0}{斯雷大桥》为他赢得了普利策奖，在三十岁的这一年使他名扬}
\BookLine{47.00}{0.00}{}{0}{四海。随之而来的是各种特权、荣誉和令人眼花缭乱的机遇，}
\BookLine{55.50}{0.00}{}{0}{同时他也失去了很多隐私，身体、思想和精神都受到了危害。}
\BookLine{64.00}{7.38}{}{0}{这几年里，桑顿的第一部长剧《我们的小镇》在制作演出}
\BookLine{72.50}{0.00}{}{0}{后便取得了成功。第二部剧《扬克斯商人》则以失败告终。安}
\BookLine{81.00}{0.00}{}{0}{德烈·奥比的剧作《卢克丽丝》由凯瑟琳·康奈尔委托给我哥}
\BookLine{89.50}{0.00}{}{0}{哥翻译，这部剧开演后又停演。与此同时，他帮鲁斯·戈登改}
\BookLine{98.00}{0.00}{}{0}{编的易卜生的《玩偶之家》则取得了巨大的成功。}
\BookLine{106.50}{7.38}{}{0}{1939 年夏天，桑顿终于完成了所有的邀约工作，带着完全}
\BookLine{115.00}{0.00}{}{0}{自由的状态和满腹个人创作计划回到家中。他已经写了两部作}
\BookLine{123.50}{0.00}{}{0}{品的开头，两部都是戏剧。一部还没有命名，另一部就是《阿}
\BookLine{132.00}{0.00}{}{0}{尔刻提斯之歌》。九月到了，希特勒的军队横扫了荷兰和比利}
\BookLine{140.50}{0.00}{}{1}{时，桑顿别无选择。《阿尔刻提斯之歌》并不适合一个战乱的世}
\BookLine{149.00}{0.00}{}{0}{界。此时尚未被命名的《九死一生》才更合适。}
\BookLine{157.50}{7.38}{}{0}{珍珠港事件后，桑顿被任命为空军上尉。很巧的是，他}
\BookLine{166.00}{0.00}{}{1}{出国服役前的最后一天，正是《九死一生》在纽约首演的日子。}
\BookLine{174.50}{0.00}{}{0}{我和母亲同他告别的那个早上，他拍了拍自己鼓鼓的行囊说}
\newpage
% 原书页码：5
\null
\BookNumber{5}
\Vertical{129}{11}{}{7}{10}前\\言\EndVertical
\BookLine{30.00}{0.00}{}{0}{道：“我打包了《阿尔刻提斯之歌》的手稿在里面。我一直在琢}
\BookLine{38.50}{0.00}{}{0}{磨把它写成一个不错的独幕剧，很可能只需要第二幕的那一大}
\BookLine{47.00}{0.00}{}{0}{场戏。我想留下一部完成的作品……说不定能派上用场，妈}
\BookLine{55.50}{0.00}{}{0}{妈。《九死一生》能不能成功也说不定呢……”}
\BookLine{64.00}{7.38}{}{0}{两年后，他弄丢了自己的手稿。彼时他已经是一位空军中}
\BookLine{72.50}{0.00}{}{0}{校，正在波士顿附近的军营里等待着退伍。指挥官批准他离开}
\BookLine{81.00}{0.00}{}{0}{三天，桑顿从后湾车站直奔科普利广场的波士顿图书馆，再次}
\BookLine{89.50}{0.00}{}{0}{沉浸到了古希腊的黄金时代之中。几天后，他终于在三年后又}
\BookLine{98.00}{0.00}{}{0}{一次坐在了康涅狄格汉姆登自己的书桌前。}
\BookLine{106.50}{7.38}{}{0}{此刻我面前正放着那时他留下的一张纸，一面用大写字母}
\BookLine{115.00}{0.00}{}{0}{打印着：}
\BookLine{132.00}{14.76}{\kaifont}{0}{《阿尔刻提斯之歌》}
\BookLine{140.50}{14.76}{\kaifont}{0}{一出充满问题的戏剧}
\BookLine{157.50}{7.38}{}{0}{反面是他写的一些小字：}
\BookLine{174.50}{14.76}{\kaifont}{0}{（第一幕戏的草稿，包括有忒瑞西阿斯的场景，都在}
\newpage
% 原书页码：6
\null
\BookNumber{6}
\Vertical{4}{169}{}{6.2}{8}阿\\尔\\刻\\提\\斯\\之\\歌\EndVertical
\BookLine{30.00}{7.38}{\kaifont}{0}{战争前完成，已经遗失）}
\BookLine{38.50}{14.76}{\kaifont}{0}{第一幕初稿，1945 年 5 月和 7 月}
\BookLine{47.00}{14.76}{\kaifont}{0}{第二稿，1945 年 7 月 8 日开始创作}
\BookLine{55.50}{14.76}{\kaifont}{0}{第三稿，12 月开始（在第二幕与第三幕的一半完}
\BookLine{64.00}{7.38}{\kaifont}{0}{成后）}
\BookLine{81.00}{7.38}{}{0}{没写上的日期应该是 1945 年。}
\BookLine{89.50}{7.38}{}{0}{遗失手稿给哥哥带来的影响远不是这几行笔记可以传达}
\BookLine{98.00}{0.00}{}{0}{的。他急切地想要回到和平时期的状态中，这激发了他的热}
\BookLine{106.50}{0.00}{}{0}{情，让他如记录中所写的样子，连续创作了七个月的时间。然}
\BookLine{115.00}{0.00}{}{0}{而，就如同其他上百万的士兵们一样，他回到家时不只是迷茫}
\BookLine{123.50}{0.00}{}{0}{与疲惫，整个人也完全改变了。他不再能够以战争之前那样的}
\BookLine{132.00}{0.00}{}{0}{视角来看待阿尔刻提斯了，他的体力和精力也不足以支撑自己}
\BookLine{140.50}{0.00}{}{0}{保持一个亢奋的状态来完成工作。他发现自己所做的更多是破}
\BookLine{149.00}{0.00}{}{0}{坏而不是创造。好在他意识到了这一点，在造成破坏前就停止}
\BookLine{157.50}{0.00}{}{0}{了自己的工作。}
\BookLine{166.00}{7.38}{}{0}{古代文学故事里满是战场和战士，暴君和英雄。哥哥的思}
\BookLine{174.50}{0.00}{}{0}{绪还围绕着他自己在希特勒和墨索里尼带来的战争中的经历：}
\newpage
% 原书页码：7
\null
\BookNumber{7}
\Vertical{129}{11}{}{7}{10}前\\言\EndVertical
\BookLine{30.00}{0.00}{}{0}{他在北非待了一年，就在那些希腊与罗马帝国留下的废墟与阴}
\BookLine{38.50}{0.00}{}{0}{影里；又在意大利南部待了一年，那里同样充满戏剧性的传}
\BookLine{47.00}{0.00}{}{0}{说，例如恺撒的胜利。桑顿突然有了新作品的灵感，他的小说}
\BookLine{55.50}{0.00}{}{1}{《三月十五》在 1948 年出版。这部小说开启了他回归文明世界}
\BookLine{64.00}{0.00}{}{0}{的道路，他又开始教书，演讲，上剧场，重新思考阿尔刻提斯}
\BookLine{72.50}{0.00}{}{0}{的故事。}
\BookLine{81.00}{7.38}{}{0}{在欧里庇得斯之后，有许多作家想要以自己的方式改}
\BookLine{89.50}{0.00}{}{0}{编阿尔刻提斯的故事。1773 年，克里斯托夫·马丁·维兰德}
\BookLine{98.00}{0.00}{}{0}{以此为题创作了一部歌剧剧本。在十九世纪，奥地利诗人胡}
\BookLine{106.50}{0.00}{}{0}{戈·冯·霍夫曼斯塔尔和英国作家罗伯特·勃朗宁也都选取过}
\BookLine{115.00}{0.00}{}{1}{这个题材。T. S. 艾略特在他成功的剧作《鸡尾酒会》中也有所借}
\BookLine{123.50}{0.00}{}{0}{用。也难怪，这个故事的主旨是夫妻之间的爱，以及妻子愿意}
\BookLine{132.00}{0.00}{}{0}{为了丈夫而死。这对任何想要发掘和赞美人类能动性的作者来}
\BookLine{140.50}{0.00}{}{0}{说，都是充满吸引力且发人深省的挑战。}
\BookLine{149.00}{7.38}{}{0}{1954 年，《媒人》在爱丁堡艺术节上演并获得成功——这}
\BookLine{157.50}{0.00}{}{0}{部有趣的作品由纽约的失败之作《扬克斯商人》而来，后来发}
\BookLine{166.00}{0.00}{}{0}{展成为大热的音乐剧《你好，多莉！》——制作人委托我哥哥}
\BookLine{174.50}{0.00}{}{0}{为下一季创作新作品，这给了他完成《阿尔刻提斯之歌》的动}
\newpage
% 原书页码：8
\null
\BookNumber{8}
\Vertical{4}{169}{}{6.2}{8}阿\\尔\\刻\\提\\斯\\之\\歌\EndVertical
\BookLine{30.00}{0.00}{}{0}{力。他在普罗旺斯的艾克斯闭关创作了很长一段时间。圣诞}
\BookLine{38.50}{0.00}{}{0}{节时，他聚集了包括我在内的一帮朋友，一起在瑞士的维拉}
\BookLine{47.00}{0.00}{}{0}{斯待了两周。在此期间，我依据他的手稿打出了新文本的前}
\BookLine{55.50}{0.00}{}{0}{两幕，他立刻提交给了艺术节组委会。泰隆·格思里接受这}
\BookLine{64.00}{0.00}{}{0}{出戏参与到 1955 年的节日中，春天离开艾克斯之前，他提交}
\BookLine{72.50}{0.00}{}{0}{了第三幕。}
\BookLine{81.00}{7.38}{}{0}{古希腊剧场的演出由内容相关联的悲剧三部曲组成，通常}
\BookLine{89.50}{0.00}{}{0}{还有一出羊人剧\textsuperscript{1}——用前面剧中的人物和情节进行或幽默或尖}
\BookLine{98.00}{0.00}{}{0}{锐的注解。欧里庇得斯是三大对手剧作家中最年轻的一位，他}
\BookLine{106.50}{0.00}{}{0}{有时会打破传统，写一部可以独立上演的剧本，或者该剧可以}
\BookLine{115.00}{0.00}{}{0}{与其他剧目组合成为三部曲。他的《阿尔刻提斯》就是这个类}
\BookLine{123.50}{0.00}{}{0}{型的剧目之一。事实上，由于它的“大团圆结局”——女主角}
\BookLine{132.00}{0.00}{}{1}{被赫拉克勒斯从冥界带回人间，《阿尔刻提斯》偶尔甚至会被用}
\BookLine{140.50}{0.00}{}{0}{作羊人剧演出。}
\BookLine{149.00}{7.38}{}{0}{我哥哥对于他这版改编的第一个想法，是将欧里庇得斯的}
\BookLine{157.50}{0.00}{}{0}{文本进行全新翻译作为中间一幕，然后创作第一幕和第三幕把}
\begin{textblock*}{45mm}(16mm,171mm)\rule{43mm}{.3pt}\end{textblock*}
\begin{textblock*}{108mm}(16mm,177mm)\kaifont\fontsize{7}{10}\selectfont 1　羊人剧，是一种轻松的笑剧，不是喜剧。一般在三部悲剧上演后进行演出，\end{textblock*}
\begin{textblock*}{108mm}(16mm,182mm)\kaifont\fontsize{7}{10}\selectfont 　 作为一种调剂。——译者注\end{textblock*}
\newpage
% 原书页码：9
\null
\BookNumber{9}
\Vertical{129}{11}{}{7}{10}前\\言\EndVertical
\BookLine{30.00}{0.00}{}{0}{第二幕框起来。很快他就放弃了这个计划，因为这不符合他的}
\BookLine{38.50}{0.00}{}{0}{目的。}
\BookLine{47.00}{7.38}{}{0}{在希腊人讲述的传说和几个世纪后布尔凡奇的记录中，阿}
\BookLine{55.50}{0.00}{}{0}{德墨托斯国王知道，如果有人愿意为他而死，他就可以活下}
\BookLine{64.00}{0.00}{}{0}{去。他向他的仆人、昔日的军中同僚和朋友们求救，但无济于}
\BookLine{72.50}{0.00}{}{0}{事。而后，他不遗余力地恳求他的父母，指出既然他们的寿命}
\BookLine{81.00}{0.00}{}{0}{所剩无几，请求他们代死只是小事一桩。但无论是老父亲还是}
\BookLine{89.50}{0.00}{}{0}{老母亲，都不愿意为儿子放弃生命。国王没有请求他的王后做}
\BookLine{98.00}{0.00}{}{0}{出这个牺牲，当阿尔刻提斯带着爱献出她的生命时，他的抗议}
\BookLine{106.50}{0.00}{}{0}{显得不温不火，毫无意义。}
\BookLine{115.00}{7.38}{}{0}{桑顿让阿德墨托斯不曾得知自己可以通过他人的死亡得到}
\BookLine{123.50}{0.00}{}{0}{拯救，在第二幕造就了一个戏剧性和情感的高潮，并深化了人}
\BookLine{132.00}{0.00}{}{0}{物形象和夫妻关系，为这个传说原本简单的故事线赋予了来自}
\BookLine{140.50}{0.00}{}{0}{二十世纪的丰富解读。关于他的这个版本，哥哥说过：}
\BookLine{157.50}{14.76}{\kaifont}{0}{从某个层面上来说，我的剧本描述了一个女人——}
\BookLine{166.00}{7.38}{\kaifont}{0}{许多女人——从一个犹疑的新娘成长为经受着苦考验的}
\BookLine{174.50}{7.38}{\kaifont}{0}{妻子，最后成为不堪重负的老人。在另一个层面上，它}
\newpage
% 原书页码：10
\null
\BookNumber{10}
\Vertical{4}{169}{}{6.2}{8}阿\\尔\\刻\\提\\斯\\之\\歌\EndVertical
\BookLine{30.00}{7.38}{\kaifont}{0}{是一个关于神和人，关于死亡和冥界，关于伟大的爱和}
\BookLine{38.50}{7.38}{\kaifont}{0}{伟大的试炼后的复活，关于篡夺和复仇的疯狂而浪漫的}
\BookLine{47.00}{7.38}{\kaifont}{0}{故事。在又一个层面上，它是一部喜剧……关于天界和}
\BookLine{55.50}{7.38}{\kaifont}{0}{人间沟通上的极端困难，关于人类和神明的不可比性所}
\BookLine{64.00}{7.38}{\kaifont}{0}{造成的误解。\textsuperscript{1}}
\BookLine{81.00}{7.38}{}{0}{在桑顿的其他作品中，1930 年出版的第三部小说《安德罗}
\BookLine{89.50}{0.00}{}{1}{斯的女人》在设定、情感和思想上与《阿尔刻提斯之歌》最为接}
\BookLine{98.00}{0.00}{}{0}{近。故事发生在希腊的一个小岛上，彼时阿尔刻提斯已经成为}
\BookLine{106.50}{0.00}{}{0}{传说。这部作品是一首怀旧的、宁静却也充满激情的散文诗，}
\BookLine{115.00}{0.00}{}{0}{它像一座桥，一端陷在残酷的异教世界的泥沼中，另一端延伸}
\BookLine{123.50}{0.00}{}{0}{向正缓缓到来的解放的黎明。女主角克里西丝漂亮、聪明，受}
\BookLine{132.00}{0.00}{}{0}{过一定的教育，但因为自己交际花的身份被边缘化。然而，她}
\BookLine{140.50}{0.00}{}{0}{是阿尔刻提斯女王精神上的姐妹。}
\BookLine{149.00}{7.38}{}{0}{通过这两个女性角色，桑顿向所有女性表达了敬意，也让}
\BookLine{157.50}{0.00}{}{0}{角色承担起传递自己主旨的重任。克里西丝在她的生和死中都}
\begin{textblock*}{45mm}(16mm,171mm)\rule{43mm}{.3pt}\end{textblock*}
\begin{textblock*}{108mm}(16mm,177mm)\kaifont\fontsize{7}{10}\selectfont 1　引自保罗·M. 库贝塔的《现代戏剧分析》（纽约：霍尔特，莱因哈特和温斯顿公司，\end{textblock*}
\begin{textblock*}{108mm}(16mm,182mm)\kaifont\fontsize{7}{10}\selectfont 　 1955），第 595—596 页。\end{textblock*}
\newpage
% 原书页码：11
\null
\BookNumber{11}
\Vertical{129}{11}{}{7}{10}前\\言\EndVertical
\BookLine{30.00}{0.00}{}{0}{赞颂了全部生命——“光明抑或黑暗”——并教导其他人也要}
\BookLine{38.50}{0.00}{}{0}{这样做。她在不自知的情形下成了一位先知，为人类指明了更}
\BookLine{47.00}{0.00}{}{0}{高的希望所在。}
\BookLine{55.50}{7.38}{}{0}{当然，阿尔刻提斯这个名字，早在桑顿听说她之前，就}
\BookLine{64.00}{0.00}{}{0}{已经承载了许多意义，被月桂花环围绕。在他的剧中，阿尔刻}
\BookLine{72.50}{0.00}{}{0}{提斯从冥界被带回，重新生活在阿波罗的光芒之中，她的结局}
\BookLine{81.00}{0.00}{}{0}{也与人类不同。神把她从死神那里偷了回来，坚持要她永远活}
\BookLine{89.50}{0.00}{}{0}{着。阿波罗将她从阿德墨托斯身边偷走，也夺去了她——以及}
\BookLine{98.00}{0.00}{}{0}{国王为她在皇家陵墓中立碑的机会。阿波罗把她带到自己的长}
\BookLine{106.50}{0.00}{}{0}{青林中，她仍在那里徘徊，成为一个象征，一个神话，一个讲}
\BookLine{115.00}{0.00}{}{0}{给我们所有人的真理：在基督教人文主义精神看来，一个逃离}
\BookLine{123.50}{0.00}{}{0}{了黑暗的人，她自己就是地平线上的一束光。}
\BookLine{132.00}{7.38}{}{0}{一部以古希腊戏剧形式写下的戏剧，理论上应该以传统}
\BookLine{140.50}{0.00}{}{0}{的古典形式在一个大理石建造的露天剧场中演出。桑顿创作}
\BookLine{149.00}{0.00}{}{0}{这版《阿尔刻提斯之歌》，也是为了它能够适应传统的镜框式}
\BookLine{157.50}{0.00}{}{0}{舞台。在爱丁堡，这出戏的演出环境有所不同，是在一个叫}
\BookLine{166.00}{0.00}{}{0}{作苏格兰教会大会堂的巨大谷仓里进行表演的。剧院呈箱形，}
\BookLine{174.50}{0.00}{}{0}{一排排并不舒适的座位构成了一个三面舞台——若是要为}
\newpage
% 原书页码：12
\null
\BookNumber{12}
\Vertical{4}{169}{}{6.2}{8}阿\\尔\\刻\\提\\斯\\之\\歌\EndVertical
\BookLine{30.00}{0.00}{}{0}{教会政治和教会条例而在会议上互相争斗，这倒是个完美的}
\BookLine{38.50}{0.00}{}{0}{环境。}
\BookLine{47.00}{7.38}{}{1}{在这个不大合拍的演出环境中，《阿尔刻提斯之歌》这个简}
\BookLine{55.50}{0.00}{}{0}{单的故事不得不被放大到一个夸张的比例。泰隆·格思里导演}
\BookLine{64.00}{0.00}{}{0}{的处理激动人心却过度戏剧化了，他增加了舞台行动，用群戏}
\BookLine{72.50}{0.00}{}{0}{场面填满舞台和过道。且不说是好是坏，这出戏的本质完全丧}
\BookLine{81.00}{0.00}{}{0}{失了。}
\BookLine{89.50}{7.38}{}{0}{甚至连剧名都改成了《阳光下的一生》，委员会认为这个}
\BookLine{98.00}{0.00}{}{0}{名字“更具票房吸引力”。桑顿从苏格兰给身在麻省剑桥的哥}
\BookLine{106.50}{0.00}{}{0}{哥阿莫斯写的信中提到了这一点：}
\BookLine{123.50}{14.76}{\kaifont}{0}{这就是《阿尔刻提斯之歌》，但老板们不想用这个名}
\BookLine{132.00}{7.38}{\kaifont}{0}{字。对于一出充满了可怖情节的戏来说，《阳光下的一}
\BookLine{140.50}{7.38}{\kaifont}{0}{生》这个名字过于轻松愉快了，但他们说阿尔刻提斯的}
\BookLine{149.00}{7.38}{\kaifont}{0}{生命源于阿波罗，而阿波罗就是太阳，所以就有了这个}
\BookLine{157.50}{7.38}{\kaifont}{0}{名字……}
\BookLine{174.50}{7.38}{}{0}{桑顿还是得到了一些补偿。他不是一个会绝望的人，甚}
\newpage
% 原书页码：13
\null
\BookNumber{13}
\Vertical{129}{11}{}{7}{10}前\\言\EndVertical
\BookLine{30.00}{0.00}{}{0}{至也不会焦虑太久。格思里的才华创造了一些精彩段落，以艾}
\BookLine{38.50}{0.00}{}{0}{琳·沃思为首的是一个强大的演员阵容。尽管批评家们还是老}
\BookLine{47.00}{0.00}{}{0}{样子，先是充满兴趣与理解，随即便是因全知而不屑。而大部}
\BookLine{55.50}{0.00}{}{0}{分的观众尽管坐在很不舒服的位置上，但都给出了积极的反响}
\BookLine{64.00}{0.00}{}{0}{和明确的肯定。}
\BookLine{72.50}{7.38}{}{0}{最重要的是，作者看到自己的作品立体而生动地呈现了}
\BookLine{81.00}{0.00}{}{0}{出来，他知道自己做到了哪些，又有哪些不足。当泰隆·格}
\BookLine{89.50}{0.00}{}{0}{思里非常有雅量地向他吐露“我没有把你的戏演好”时，桑}
\BookLine{98.00}{0.00}{}{0}{顿也得到了安慰，他觉得是被苏格兰教会大会堂这个场地给}
\BookLine{106.50}{0.00}{}{0}{坑了。}
\BookLine{115.00}{7.38}{}{0}{国内外有许多戏剧制作人对这部作品感兴趣，想获得授}
\BookLine{123.50}{0.00}{}{0}{权在纽约制作。哥哥对于他想要在文本上进行的调整也有自}
\BookLine{132.00}{0.00}{}{0}{己的想法。他与一些导演进行了交流，听取了各种建议，开始}
\BookLine{140.50}{0.00}{}{0}{着手修改。这对他来说一直不是件容易的事，他擅长的是活在}
\BookLine{149.00}{0.00}{}{0}{当下，着眼于未来。有太多人围着他给他建议，于是他告诉他}
\BookLine{157.50}{0.00}{}{0}{们：“我需要更长远的考量。”他收回了剧本，拒绝了那些想要}
\BookLine{166.00}{0.00}{}{0}{基于此版本制作的请求。}
\BookLine{174.50}{7.38}{}{0}{当他从这种完全无法适应的压力中解脱出来后，他开始有}
\newpage
% 原书页码：14
\null
\BookNumber{14}
\Vertical{4}{169}{}{6.2}{8}阿\\尔\\刻\\提\\斯\\之\\歌\EndVertical
\BookLine{30.00}{0.00}{}{0}{了新的想法，想要按照希腊戏剧的传统来创作《阿尔刻提斯之}
\BookLine{38.50}{0.00}{}{1}{歌》和羊人剧。很快，他就沉浸在《醉酒的女神》的创作中了。}
\BookLine{47.00}{0.00}{}{0}{这是作家或各种形式的创作者偶尔会经历的状况，他一下就}
\BookLine{55.50}{0.00}{}{0}{“想到”了这部作品，就像一条直线，从头到尾一气呵成。作}
\BookLine{64.00}{0.00}{}{1}{品 1957 年在《大西洋月刊》上发表，和原稿相比几乎没有任何}
\BookLine{72.50}{0.00}{}{0}{改动。}
\BookLine{81.00}{7.38}{}{0}{因为希望《醉酒的女神》能在舞台上呈现，我哥哥签了合}
\BookLine{89.50}{0.00}{}{0}{约，让这两部戏能够在瑞士苏黎世的剧院上演。第二次世界}
\BookLine{98.00}{0.00}{}{1}{大战期间，《我们的小镇》和《九死一生》都是在这个剧院进行}
\BookLine{106.50}{0.00}{}{0}{的欧洲首演。德国和奥地利最知名的导演和演员们都在瑞士避}
\BookLine{115.00}{0.00}{}{0}{难，桑顿的叙事作品和戏剧作品，都为他们所熟知并给予了高}
\BookLine{123.50}{0.00}{}{0}{度评价。}
\BookLine{132.00}{7.38}{}{0}{桑顿与译者赫伯特·E. 赫利奇卡合作，赫利奇卡是维也纳}
\BookLine{140.50}{0.00}{}{0}{人，曾将他早期的许多作品翻译成德文。相较爱丁堡的演出版}
\BookLine{149.00}{0.00}{}{0}{本，桑顿做的修改并不大，他希望看到另一位导演和另一个团}
\BookLine{157.50}{0.00}{}{0}{队来呈现同一个剧本。这个版本的制作在各种意义上都令人激}
\BookLine{166.00}{0.00}{}{0}{动和震撼，观众在剧院里高度专注，在各幕之间进行讨论，演}
\BookLine{174.50}{0.00}{}{0}{出结束时用力鼓掌，离开剧院时活力满满。}
\newpage
% 原书页码：15
\null
\BookNumber{15}
\Vertical{129}{11}{}{7}{10}前\\言\EndVertical
\BookLine{30.00}{7.38}{}{0}{苏黎世首演之后，该剧在全德国大大小小城市的国有剧}
\BookLine{38.50}{0.00}{}{0}{院都进行了完整版的演出，维也纳城堡剧院也上演了该剧。在}
\BookLine{47.00}{0.00}{}{0}{此期间，桑顿做了进一步的修改，为了回应在德国的巨大反}
\BookLine{55.50}{0.00}{}{1}{响，《阿尔刻提斯之歌》在 1960 年由费舍尔出版社出版了德文}
\BookLine{64.00}{0.00}{}{0}{版剧本。}
\BookLine{72.50}{7.38}{}{0}{该剧在报纸和杂志的文学栏目，以及学术期刊和书籍中都}
\BookLine{81.00}{0.00}{}{0}{引起了讨论。凯特·汉伯格在她关于现代文学中的古典人物的}
\BookLine{89.50}{0.00}{}{0}{集子《从索福克勒斯到萨特》\textsuperscript{1}中写道：“怀尔德的作品是现代世}
\BookLine{98.00}{0.00}{}{0}{界文学中对阿尔刻提斯这一题材最重要的表达。”}
\BookLine{106.50}{7.38}{}{0}{尽管桑顿没有继续把《阿尔刻提斯之歌》写成一个英文的}
\BookLine{115.00}{0.00}{}{0}{戏剧剧本，他也没有放任不管。他在为自己非常欣赏的作曲家}
\BookLine{123.50}{0.00}{}{0}{露易丝·塔尔玛寻找歌剧剧本题材时，似乎不可避免地完成了}
\BookLine{132.00}{0.00}{}{0}{这个循环。}
\BookLine{140.50}{7.38}{}{0}{这部歌剧也叫《阿尔刻提斯之歌》，历时六年完成制作，}
\BookLine{149.00}{0.00}{}{1}{1962 年 3 月 2 日在德国法兰克福重建后的极为华丽的歌剧院进}
\BookLine{157.50}{0.00}{}{0}{行了首演。这是一部规模宏大的歌剧，需要动用很大的舞台、}
\begin{textblock*}{45mm}(16mm,171mm)\rule{43mm}{.3pt}\end{textblock*}
\begin{textblock*}{108mm}(16mm,177mm)\kaifont\fontsize{7}{10}\selectfont 1　纽约：昂格尔，1969。编译自 Von Sophokles zu Sartre, Griechische Dramenfiguren\end{textblock*}
\begin{textblock*}{108mm}(16mm,182mm)\kaifont\fontsize{7}{10}\selectfont 　 Antik und Modern, Stuttgart, 1962.\end{textblock*}
\newpage
% 原书页码：16
\null
\BookNumber{16}
\Vertical{4}{169}{}{6.2}{8}阿\\尔\\刻\\提\\斯\\之\\歌\EndVertical
\BookLine{30.00}{0.00}{}{0}{大规模管弦乐队和一大批演员。在补贴的支持下，制作得以满}
\BookLine{38.50}{0.00}{}{0}{足这些要求，英格·博克将阿尔刻提斯演绎得极为出色。大幕}
\BookLine{47.00}{0.00}{}{0}{落下时，剧院里响起雷鸣般的掌声和此起彼伏的叫好声。大家}
\BookLine{55.50}{0.00}{}{0}{进行了三十多次谢幕，塔尔玛小姐、桑顿，每个人都被一次又}
\BookLine{64.00}{0.00}{}{0}{一次地叫上去鞠躬致意。演出结束后的晚宴上，大家都十分雀}
\BookLine{72.50}{0.00}{}{0}{跃。不过，在接下来的几天里和后来的周刊中，评论家们的反}
\BookLine{81.00}{0.00}{}{0}{馈褒贬不一，其中一个隐晦但明确的抱怨是女人不应该写歌}
\BookLine{89.50}{0.00}{}{0}{剧，还有就是音乐过于现代了。}
\BookLine{98.00}{7.38}{}{0}{对桑顿而言，他对配乐完全满意，他认为作品的不足之处}
\BookLine{106.50}{0.00}{}{0}{都源于自己作为词作者的经验不足。这部歌剧后来常以音乐会}
\BookLine{115.00}{0.00}{}{0}{的形式上演，其中的咏叹调也在各种独唱会中被演唱。桑顿从}
\BookLine{123.50}{0.00}{}{0}{未怀疑过，这是一部领先于时代的作品，总有一天会得到充分}
\BookLine{132.00}{0.00}{}{0}{的认可。}
\BookLine{140.50}{7.38}{}{0}{虽然我哥哥很早就开始写作，但他的作品列表并不长：七}
\BookLine{149.00}{0.00}{}{0}{部小说，一些论说性散文和关于自己与他人作品的序言，两部}
\BookLine{157.50}{0.00}{}{1}{短剧集和四部长剧，其中《扬克斯商人》和《媒人》算作一部。}
\BookLine{166.00}{0.00}{}{0}{在作品中，他不怕重复自己，也不怕受到薪火相传的传统大师}
\BookLine{174.50}{0.00}{}{0}{们的影响。}
\newpage
% 原书页码：17
\null
\BookNumber{17}
\Vertical{129}{11}{}{7}{10}前\\言\EndVertical
\BookLine{30.00}{7.38}{}{1}{关于他的戏剧作品，他在集子《戏剧三种》\textsuperscript{1}的序言中写道：}
\BookLine{47.00}{14.76}{\kaifont}{0}{我不是大家所期待的创新剧作家之一。我希望我是。}
\BookLine{55.50}{7.38}{\kaifont}{0}{我希望自己的创作是为他们铺路。我不是一个创新者，而}
\BookLine{64.00}{7.38}{\kaifont}{0}{是重新发现那些被遗忘的事物的人，希望我能清除那些妨}
\BookLine{72.50}{7.38}{\kaifont}{0}{碍着我们的东西。}
\BookLine{89.50}{7.38}{}{0}{但他对形式进行了实验，敢于对现实时间、历史时间和文}
\BookLine{98.00}{0.00}{}{0}{学时间进行创造，从历史中摘取事件，从故事和绘画中摘取人}
\BookLine{106.50}{0.00}{}{0}{物，从神圣传说中摘取角色。他让死亡、天使、异教神灵和动}
\BookLine{115.00}{0.00}{}{0}{物发声。}
\BookLine{123.50}{7.38}{}{0}{在他的三分钟剧本集《搅乱水池的天使及其他作品》\textsuperscript{2}的前}
\BookLine{132.00}{0.00}{}{0}{言中，我哥哥写到一个年轻作家的问题{\kaifont （他了解这些问题，因}}
\BookLine{140.50}{0.00}{}{0}{{\kaifont 为他在十几岁的时候就写了集子中的一些剧本）}。他描述了写}
\BookLine{149.00}{0.00}{}{0}{作时欣喜和沮丧的时刻，以及各种各样的可能性：}
\begin{textblock*}{45mm}(16mm,171mm)\rule{43mm}{.3pt}\end{textblock*}
\begin{textblock*}{108mm}(16mm,177mm)\kaifont\fontsize{7}{10}\selectfont 1　纽约：哈珀兄弟出版公司，1957。\end{textblock*}
\begin{textblock*}{108mm}(16mm,182mm)\kaifont\fontsize{7}{10}\selectfont 2　纽约：科沃德－麦肯出版公司，1928。\end{textblock*}
\newpage
% 原书页码：18
\null
\BookNumber{18}
\Vertical{4}{169}{}{6.2}{8}阿\\尔\\刻\\提\\斯\\之\\歌\EndVertical
\BookLine{30.00}{14.76}{\kaifont}{0}{如果一个人不允许自己自命不凡，那么写作的实践将}
\BookLine{38.50}{7.38}{\kaifont}{0}{会多么不同。有些人别无选择：他们宁愿写一部失败的清}
\BookLine{47.00}{7.38}{\kaifont}{0}{唱剧，也不愿写一首成功的民谣曲。}
\BookLine{64.00}{7.38}{}{0}{桑顿始终坚信，废纸篓是作家最好的朋友。即便如此，现}
\BookLine{72.50}{0.00}{}{0}{存的文档还是证明他写过一些“失败的清唱剧”，包括只有前}
\BookLine{81.00}{0.00}{}{0}{两幕，第三幕从没创作出来的《百货商场》。这部作品以世界}
\BookLine{89.50}{0.00}{}{0}{上最大、最诱人的商店为背景，在那里你可以找到你需要或想}
\BookLine{98.00}{0.00}{}{0}{要的一切。还有一部未完成的史诗剧情片剧本《熔炉》，通过}
\BookLine{106.50}{0.00}{}{0}{几代人的生活反映我们国家的发展。尽管他没有从其他项目中}
\BookLine{115.00}{0.00}{}{0}{抽出时间来修改《阿尔刻提斯之歌》的英文剧本，使其与德文}
\BookLine{123.50}{0.00}{}{0}{版内容保持一致，而且他肯定还会对剧本做出进一步调整，但}
\BookLine{132.00}{0.00}{}{0}{目前的剧本并不算是“失败的清唱剧”。这部作品成功地反映}
\BookLine{140.50}{0.00}{}{0}{了他对人类经验真实性最持久的信念。}
\BookLine{149.00}{7.38}{}{0}{桑顿·怀尔德的织毯已经完成。随着《阿尔刻提斯之歌》}
\BookLine{157.50}{0.00}{}{0}{与《醉酒的女神》的出版和故事的讲述，他的全部设计得以清}
\BookLine{166.00}{0.00}{}{0}{晰展现出来，以自己的文字能力，以自己的天赋，以人生的全}
\BookLine{174.50}{0.00}{}{0}{部岁月，最重要的，以自身的信念和分享的愿望。现在，阿}
\newpage
% 原书页码：19
\null
\BookNumber{19}
\Vertical{129}{11}{}{7}{10}前\\言\EndVertical
\BookLine{30.00}{0.00}{}{0}{尔刻提斯这团巨大的纱线，最后一段也已经解开，最后一针也}
\BookLine{38.50}{0.00}{}{0}{打好了结。在这个故事里，已经再没有谜团需要德尔斐来守}
\BookLine{47.00}{0.00}{}{0}{护了。}
\BookLine{64.00}{66.43}{}{0}{伊莎贝尔·怀尔德\textsuperscript{1}}
\BookLine{72.50}{81.19}{}{0}{玛莎葡萄园}
\BookLine{81.00}{81.19}{}{1}{1977 年 7 月}
\begin{textblock*}{45mm}(16mm,171mm)\rule{43mm}{.3pt}\end{textblock*}
\begin{textblock*}{108mm}(16mm,177mm)\kaifont\fontsize{7}{10}\selectfont 1　作者为桑顿·怀尔德的妹妹，她是一位小说家和传记作家，同时也是桑顿·怀\end{textblock*}
\begin{textblock*}{108mm}(16mm,182mm)\kaifont\fontsize{7}{10}\selectfont 　 尔德的文学经纪人、发言人和传记作者。——译者注\end{textblock*}
\newpage
% 原书页码：20
\null
\Vertical{15}{28}{}{14}{21}人\\物\EndVertical
\BookLine{82.00}{0.00}{}{0}{阿波罗——太阳神}
\BookLine{90.50}{0.00}{}{0}{死神——冥界之神}
\BookLine{99.00}{0.00}{}{0}{第一个看门人——守夜人，中年，阿德墨托斯国王的仆人}
\BookLine{107.50}{0.00}{}{0}{阿尔刻提斯——塞萨利王后}
\BookLine{116.00}{0.00}{}{0}{阿格莱亚——年迈的女仆}
\BookLine{124.50}{0.00}{}{0}{忒瑞西阿斯——盲人，年纪非常大，暴躁易怒，控制欲强，}
\BookLine{133.00}{25.83}{}{0}{马上就要老糊涂的状态}
\BookLine{141.50}{0.00}{}{0}{男孩——忒瑞西阿斯的仆人}
\BookLine{150.00}{0.00}{}{0}{牧人一——牧羊人}
\BookLine{158.50}{0.00}{}{0}{牧人二——牧羊人}
\BookLine{167.00}{0.00}{}{0}{牧人三——牧羊人}
\BookLine{175.50}{0.00}{}{0}{牧人四——牧羊人}
\BookLine{184.00}{0.00}{}{0}{罗德普——王宫中年轻的女仆}
\newpage
% 原书页码：21
\null
\BookNumber{21}
\BookLine{88.00}{0.00}{}{0}{赫拉克勒斯——宙斯之子}
\BookLine{96.50}{0.00}{}{0}{第二个看门人——比第一个看门人年轻，阿吉斯国王的仆人}
\BookLine{105.00}{0.00}{}{0}{伊皮伊尼兹——21 岁，阿尔刻提斯的儿子}
\BookLine{113.50}{0.00}{}{0}{切瑞安德——伊皮伊尼兹的朋友}
\BookLine{122.00}{0.00}{}{0}{阿吉斯——篡位的塞萨利国王}
\BookLine{130.50}{0.00}{}{0}{守卫一——阿吉斯王宫的守卫}
\BookLine{139.00}{0.00}{}{0}{守卫二——阿吉斯王宫的守卫}
\BookLine{147.50}{0.00}{}{0}{守卫三——阿吉斯王宫的守卫}
\BookLine{156.00}{0.00}{}{0}{守卫四——阿吉斯王宫的守卫}
\BookLine{164.50}{0.00}{}{0}{仆人们——阿德墨托斯国王和阿吉斯国王的仆人们}
\BookLine{173.00}{0.00}{}{0}{塞萨利的百姓}
\newpage
% 原书页码：22
\null
\Vertical{15}{28}{}{14}{21}场\\景\EndVertical
\BookLine{82.00}{0.00}{}{0}{塞萨利国王阿德墨托斯王宫的后院，}
\BookLine{90.50}{0.00}{}{0}{希腊黄金时代几个世纪之前。}
\newpage
% 原书页码：23
\null
\ActOrnament[一]
\newpage
% 原书页码：24
\null
\newpage
% 原书页码：25
\null
\BookNumber{25}
\Vertical{129}{11}{}{7}{10}第\\一\\幕\EndVertical
\BookLine{30.00}{7.38}{\kaifont}{0}{除了第三幕结尾和《醉酒的女神》结束后，全剧不设大}
\BookLine{38.50}{0.00}{\kaifont}{0}{幕。《阿尔刻提斯之歌》全剧三幕戏的场景均在塞萨利国王阿}
\BookLine{47.00}{0.00}{\kaifont}{0}{德墨托斯宫殿的后院，时间是希腊黄金时代几个世纪之前。}
\BookLine{55.50}{0.00}{\kaifont}{0}{每一幕戏都开始于一天的黎明，结束于同一日的黄昏。}
\BookLine{64.00}{7.38}{\kaifont}{0}{宫殿是一座由粗粝石块砌成的低矮却宽敞的房屋，有着}
\BookLine{72.50}{0.00}{\kaifont}{0}{平坦的屋顶。巨大的树干仿佛组成了宫殿的门廊。宫殿大门}
\BookLine{81.00}{0.00}{\kaifont}{0}{是木制的，两扇门上各镶有一块镀金的牛骨。门前是平台，}
\BookLine{89.50}{0.00}{\kaifont}{0}{下几级台阶就能走到被松软土壤覆盖的庭院。}
\BookLine{98.00}{7.38}{\kaifont}{0}{宫殿的前区占据了舞台后方左侧四分之三的空间，舞台}
\BookLine{106.50}{0.00}{\kaifont}{0}{其他区域则被泥砖墙围了起来。右边墙上是一扇大木门，直}
\BookLine{115.00}{0.00}{\kaifont}{0}{通外面的大路和费莱城。左边的墙稍低一些，墙上一扇小门}
\BookLine{123.50}{0.00}{\kaifont}{0}{可通往仆人们的住所。}
\BookLine{132.00}{7.38}{\kaifont}{0}{舞台中央前区有一条小路（向右逐渐走低）直接通向剧院}
\BookLine{140.50}{0.00}{\kaifont}{0}{的乐池。这条路的下方是一个“洞穴”——一处有活水流动的}
\BookLine{149.00}{0.00}{\kaifont}{0}{清泉，一扇铜门通往冥界，门被藤条覆盖，却也足够大，能}
\BookLine{157.50}{0.00}{\kaifont}{0}{够让演员从中通过。这儿还有一条（看不到，只存在于想象中}
\BookLine{166.00}{0.00}{\kaifont}{0}{的）蛇——佩索。}
\newpage
% 原书页码：26
\null
\BookNumber{26}
\Vertical{4}{169}{}{6.2}{8}阿\\尔\\刻\\提\\斯\\之\\歌\EndVertical
\BookLine{30.00}{7.38}{\kaifont}{0}{黎明的第一道光线。}
\BookLine{38.50}{7.38}{\kaifont}{0}{一束光渐亮，阿波罗正站在宫殿的房顶上。他穿着金色}
\BookLine{47.00}{0.00}{\kaifont}{0}{的衣服，右肩上披着深蓝色的长斗篷。一束蓝色的光渐渐打}
\BookLine{55.50}{0.00}{\kaifont}{0}{亮泉水下方的冥界入口处，死神——穿着一件黑色的补丁外}
\BookLine{64.00}{0.00}{\kaifont}{0}{套，看上去像一只蝙蝠，或者甲虫——蠕动着沿着小路爬上}
\BookLine{72.50}{0.00}{\kaifont}{0}{来，在左边的门口以及宫殿的门口嗅个不停。}
\BookLine{81.00}{7.38}{\kaifont}{0}{这场戏自始至终，阿波罗一直注视着渐渐升起的太}
\BookLine{89.50}{0.00}{\kaifont}{0}{阳——平静而克制，嘴角带着一丝不易察觉的微笑。}
\BookLine{106.50}{0.00}{\heifont}{0}{阿波罗：}
\BookLine{115.00}{22.14}{}{0}{{\kaifont （像在问早安一样）}死神！}
\BookLine{132.00}{0.00}{\heifont}{0}{死神：}
\BookLine{140.50}{22.14}{}{0}{啊！你来了！阿德墨托斯的宫殿迎来了一位尊}
\BookLine{149.00}{22.14}{}{0}{贵的客人啊！我们这里今天要举行一场婚礼，}
\BookLine{157.50}{22.14}{}{0}{竟有你来做客！你是不是来偷新娘子的啊，大}
\BookLine{166.00}{22.14}{}{0}{名鼎鼎的阿波罗？}
\newpage
% 原书页码：27
\null
\BookNumber{27}
\Vertical{129}{11}{}{7}{10}第\\一\\幕\EndVertical
\BookLine{30.00}{0.00}{\heifont}{0}{阿波罗：}
\BookLine{38.50}{22.14}{}{0}{死神，你活在黑暗之中。}
\BookLine{55.50}{0.00}{\heifont}{0}{死神：}
\BookLine{64.00}{22.14}{}{0}{是啊，是啊。阿波罗大人，你今天来，是不是}
\BookLine{72.50}{22.14}{}{0}{要给我们展现一些伟大的启示，伟大的神迹？}
\BookLine{89.50}{7.38}{\kaifont}{0}{看门人摇着手中的鼓，提着羊皮灯笼绕过宫殿走到舞台}
\BookLine{98.00}{0.00}{\kaifont}{0}{中央。}
\BookLine{115.00}{0.00}{\heifont}{0}{看门人：}
\BookLine{123.50}{22.14}{}{0}{{\kaifont （念词）}黎明时分，照例巡查。热情诚挚的阿德}
\BookLine{132.00}{22.14}{}{0}{墨托斯，塞萨利之王，千金万马的拥有者，宫}
\BookLine{140.50}{22.14}{}{0}{中万事顺遂。{\kaifont （向左侧走下台）}黎明时分，大婚}
\BookLine{149.00}{22.14}{}{0}{之日——将是最伟大的婚礼之一。}
\BookLine{166.00}{7.38}{\kaifont}{0}{他从左侧下场，回到仆人们的住所。}
\newpage
% 原书页码：28
\null
\BookNumber{28}
\Vertical{4}{169}{}{6.2}{8}阿\\尔\\刻\\提\\斯\\之\\歌\EndVertical
\BookLine{30.00}{0.00}{\heifont}{0}{死神：}
\BookLine{38.50}{22.14}{}{0}{我刚问你呢，阿波罗大人，你是不是来给我们}
\BookLine{47.00}{22.14}{}{0}{点奇迹看看的？{\kaifont （停顿）}是？不是？“是。”我希}
\BookLine{55.50}{22.14}{}{0}{望是。当众神靠近人类之时，总是会有人死去。}
\BookLine{64.00}{22.14}{}{0}{我的王国今天是不是会有幸迎来某个尊贵的客}
\BookLine{72.50}{22.14}{}{0}{人？是阿德墨托斯国王还是阿尔刻提斯公主？}
\BookLine{89.50}{0.00}{\heifont}{0}{阿波罗：}
\BookLine{98.00}{22.14}{}{0}{没有。}
\BookLine{115.00}{0.00}{\heifont}{0}{死神：}
\BookLine{123.50}{22.14}{}{0}{让我等着瞧瞧。{\kaifont （他蠕动到场中央）}我能不能}
\BookLine{132.00}{22.14}{}{0}{再问问，你今天是带着哪一种属性和能力出现}
\BookLine{140.50}{22.14}{}{0}{的？治疗者？{\kaifont （停顿）}光明和生命的赐予者？}
\BookLine{149.00}{22.14}{}{0}{{\kaifont （停顿）}歌唱者？}
\BookLine{166.00}{0.00}{\heifont}{0}{阿波罗：}
\BookLine{174.50}{22.14}{}{0}{{\kaifont （仍然注视着太阳，轻松地）}他们是同一个人，}
\newpage
% 原书页码：29
\null
\BookNumber{29}
\Vertical{129}{11}{}{7}{10}第\\一\\幕\EndVertical
\BookLine{30.00}{22.14}{}{0}{是同样的。我今天来，是要开启一曲有情节的}
\BookLine{38.50}{22.14}{}{0}{歌谣——一个故事——}
\BookLine{55.50}{0.00}{\heifont}{0}{死神：}
\BookLine{64.00}{22.14}{}{0}{一个故事！}
\BookLine{81.00}{0.00}{\heifont}{0}{阿波罗：}
\BookLine{89.50}{22.14}{}{0}{一个会被不断传讲的故事。}
\BookLine{106.50}{0.00}{\heifont}{0}{死神：}
\BookLine{115.00}{22.14}{}{0}{啊！那就是一个教导了！其中有给我的教}
\BookLine{123.50}{22.14}{}{0}{导吗？}
\BookLine{140.50}{0.00}{\heifont}{0}{阿波罗：}
\BookLine{149.00}{22.14}{}{0}{是的。}
\BookLine{166.00}{0.00}{\heifont}{0}{死神：}
\BookLine{174.50}{22.14}{}{0}{{\kaifont （用他的手脚拍打地面）}不！}
\newpage
% 原书页码：30
\null
\BookNumber{30}
\Vertical{4}{169}{}{6.2}{8}阿\\尔\\刻\\提\\斯\\之\\歌\EndVertical
\BookLine{30.00}{0.00}{\heifont}{0}{阿波罗：}
\BookLine{38.50}{22.14}{}{0}{是的，你会得到你的教导。}
\BookLine{55.50}{0.00}{\heifont}{0}{死神：}
\BookLine{64.00}{22.14}{}{0}{{\kaifont （勃然大怒）}不！你没资格教导我。我永远都在}
\BookLine{72.50}{22.14}{}{0}{这儿，而且永不改变。你们这些天上的神才需}
\BookLine{81.00}{22.14}{}{0}{要教导。我现在就先给你上一课，{\kaifont （尖声尖气}}
\BookLine{89.50}{22.14}{}{0}{{\kaifont 地）}别管这些人类了，好好待在你的奥林匹斯}
\BookLine{98.00}{22.14}{}{0}{山上，好好过你们的日子。你们的这些蠢念头}
\BookLine{106.50}{22.14}{}{0}{我早就知道了，你们创造了这些生物，然后又}
\BookLine{115.00}{22.14}{}{0}{迷上了他们。结果你们让整个世界乱成了一团，}
\BookLine{123.50}{22.14}{}{0}{一天比一天糟糕。你所做的对他们不过是种折}
\BookLine{132.00}{22.14}{}{0}{磨——谁能比我更了解呢？{\kaifont （他往回蠕动到小}}
\BookLine{140.50}{22.14}{}{0}{{\kaifont 路的最上方，狂暴地晃动着身体）}他们永远不}
\BookLine{149.00}{22.14}{}{0}{会理解你的意思。你越是想要表达，就会把他}
\BookLine{157.50}{22.14}{}{0}{们弄得越疯狂。}
\newpage
% 原书页码：31
\null
\BookNumber{31}
\Vertical{129}{11}{}{7}{10}第\\一\\幕\EndVertical
\BookLine{30.00}{0.00}{\heifont}{0}{阿波罗：}
\BookLine{38.50}{22.14}{}{0}{他们已经开始理解我了。最初他们的确就像野}
\BookLine{47.00}{22.14}{}{0}{兽——野蛮，恐惧。就像是笼中的困兽，他们}
\BookLine{55.50}{22.14}{}{0}{自己就是自己的牢笼。后来有两件事情让他们}
\BookLine{64.00}{22.14}{}{0}{开始有了意识，开始寻索：我父亲的雷电，让}
\BookLine{72.50}{22.14}{}{0}{他们的恐惧变为敬畏；我的阳光，让他们学会}
\BookLine{81.00}{22.14}{}{0}{了感谢。在感谢之中他们发明了语言，我又给}
\BookLine{89.50}{22.14}{}{0}{了他们音乐。这些就是启示，他们理解了。一}
\BookLine{98.00}{22.14}{}{0}{个又一个人，意识到是我在搅动他们的心灵，}
\BookLine{106.50}{22.14}{}{0}{对他们言语。}
\BookLine{123.50}{0.00}{\heifont}{0}{死神：}
\BookLine{132.00}{22.14}{}{0}{是的，他们和以前不一样了：“阿波罗爱塞萨}
\BookLine{140.50}{22.14}{}{0}{利，阿波罗爱费莱。”回你的奥林匹斯山去吧。}
\BookLine{149.00}{22.14}{}{0}{这种爱吧……很难说你们谁更不幸——是你，}
\BookLine{157.50}{22.14}{}{0}{还是这些可怜的造物。你想要介入他们的生活，}
\BookLine{166.00}{22.14}{}{0}{这就像一个巨人走进了小小的房间：你每挪一}
\BookLine{174.50}{22.14}{}{0}{小步，都会弄坏点什么。你今天打算让谁吃这}
\newpage
% 原书页码：32
\null
\BookNumber{32}
\Vertical{4}{169}{}{6.2}{8}阿\\尔\\刻\\提\\斯\\之\\歌\EndVertical
\BookLine{30.00}{22.14}{}{0}{个苦头呢？国王？还是他的新娘？}
\BookLine{47.00}{0.00}{\heifont}{0}{阿波罗：}
\BookLine{55.50}{22.14}{}{0}{你。}
\BookLine{72.50}{0.00}{\heifont}{0}{死神：}
\BookLine{81.00}{22.14}{}{0}{我？我？所以你是打算也爱一爱我喽？不！谢}
\BookLine{89.50}{22.14}{}{0}{谢！{\kaifont （移动起所有手脚，在小路上来回爬着，尖}}
\BookLine{98.00}{22.14}{}{0}{{\kaifont 叫着）}你休想找我麻烦，你也休想教导我。我}
\BookLine{106.50}{22.14}{}{0}{和我的王国会永远存在。你怎么可能给我找麻}
\BookLine{115.00}{22.14}{}{0}{烦呢？}
\BookLine{132.00}{0.00}{\heifont}{0}{阿波罗：}
\BookLine{140.50}{22.14}{}{0}{你生活在黑暗之中，所以看不到万事万物的}
\BookLine{149.00}{22.14}{}{0}{改变。}
\BookLine{166.00}{0.00}{\heifont}{0}{死神：}
\BookLine{174.50}{22.14}{}{0}{{\kaifont （喊叫着）}改变！没什么会变！{\kaifont （恐惧地环视一}}
\newpage
% 原书页码：33
\null
\BookNumber{33}
\Vertical{129}{11}{}{7}{10}第\\一\\幕\EndVertical
\BookLine{30.00}{22.14}{}{0}{{\kaifont 下四周）}天马上要亮了。所以你今天要讲的这}
\BookLine{38.50}{22.14}{}{0}{个故事是关于改变？针对我的改变？}
\BookLine{55.50}{0.00}{\heifont}{0}{阿波罗：}
\BookLine{64.00}{22.14}{}{0}{对你，也对我。}
\BookLine{81.00}{0.00}{\heifont}{0}{死神：}
\BookLine{89.50}{22.14}{}{0}{{\kaifont （消失在自己的洞穴，最后尖叫道）}就对你！}
\BookLine{106.50}{7.38}{\kaifont}{0}{阿波罗身上的光渐渐隐去，他消失在大家的视野中。}
\BookLine{115.00}{7.38}{\kaifont}{0}{看门人回到场上，他摇了摇鼓，把灯吹灭。}
\BookLine{132.00}{0.00}{\heifont}{0}{看门人：}
\BookLine{140.50}{22.14}{}{0}{黎明到来，黎明到来。热情诚挚的阿德墨托}
\BookLine{149.00}{22.14}{}{0}{斯，塞萨利之王，千金万马的拥有者，宫中万}
\BookLine{157.50}{22.14}{}{0}{事顺遂。{\kaifont （走向观众）}大婚之日——将是最伟}
\BookLine{166.00}{22.14}{}{0}{大的婚礼之一，可事情都不顺利。她为什么无}
\BookLine{174.50}{22.14}{}{0}{法入睡——这位公主，这位新娘，我们未来的}
\newpage
% 原书页码：34
\null
\BookNumber{34}
\Vertical{4}{169}{}{6.2}{8}阿\\尔\\刻\\提\\斯\\之\\歌\EndVertical
\BookLine{30.00}{22.14}{}{0}{王后？就在夜里，我曾八次、十次在这里见到}
\BookLine{38.50}{22.14}{}{0}{她——望着天空四处游荡。有时她走到大路上}
\BookLine{47.00}{22.14}{}{0}{去，好像在等着什么信使。她站在这里，举起}
\BookLine{55.50}{22.14}{}{0}{她的双臂低语：“阿波罗！一点启示！一个启}
\BookLine{64.00}{22.14}{}{0}{示！”启示什么？她嫁给阿德墨托斯国王是否}
\BookLine{72.50}{22.14}{}{0}{正确？唉！她想寻求的最清楚明白的启示早已}
\BookLine{81.00}{22.14}{}{0}{写在国王脸上了。哦，我已经有些岁数了，我}
\BookLine{89.50}{22.14}{}{0}{知道每一位新娘在她们的婚礼前夜内心都会充}
\BookLine{98.00}{22.14}{}{0}{满不安，但是何必为我们的阿德墨托斯国王感}
\BookLine{106.50}{22.14}{}{0}{到不安，他赢得她的奇妙故事让整个希腊都为}
\BookLine{115.00}{22.14}{}{0}{之惊讶。哦，朋友们，请记住我这个老看门人}
\BookLine{123.50}{22.14}{}{0}{的建议：不要在凌晨三点的时候思考人生大事。}
\BookLine{132.00}{22.14}{}{0}{在那样的时刻，你的心灵和头脑都没有温度，}
\BookLine{140.50}{22.14}{}{0}{在那样的时刻——嘘——你会看到自己的房子}
\BookLine{149.00}{22.14}{}{0}{着起火，你的孩子挣扎着死在你的脚边。等到}
\BookLine{157.50}{22.14}{}{0}{太阳升起吧，事实都是一样的——人类生活的}
\BookLine{166.00}{22.14}{}{0}{一切事实都是一样的——但阳光赋予了它们意}
\BookLine{174.50}{22.14}{}{0}{义。请记住看门人给你们的建议。现在我要去}
\newpage
% 原书页码：35
\null
\BookNumber{35}
\Vertical{129}{11}{}{7}{10}第\\一\\幕\EndVertical
\BookLine{30.00}{22.14}{}{0}{喝点水。{\kaifont （走到泉水旁边，和蛇佩索打了个招}}
\BookLine{38.50}{22.14}{}{0}{{\kaifont 呼）}日安佩索，我的老朋友。今天你一定会过}
\BookLine{47.00}{22.14}{}{0}{得不错，婚宴少不了你的份儿——一只羊，或}
\BookLine{55.50}{22.14}{}{0}{者是半头牛。过去，好让我进行我的祭礼。{\kaifont （让}}
\BookLine{64.00}{22.14}{}{0}{{\kaifont 手中捧着的水流出，一边念着祷文：“生命之源}}
\BookLine{72.50}{22.14}{}{0}{{\kaifont 啊——大地，空气，火，水……”然后又捧了}}
\BookLine{81.00}{22.14}{}{0}{{\kaifont 一捧水喝了下去。对观众说话）}看，我的朋友}
\BookLine{89.50}{22.14}{}{0}{们。你们看到这些藤条下面的洞穴了吗？这是}
\BookLine{98.00}{22.14}{}{0}{前往冥界的五个入口之一，我们的好佩索正是}
\BookLine{106.50}{22.14}{}{0}{在此守护着入口。没有人进到过那里，也没有}
\BookLine{115.00}{22.14}{}{0}{人从那里出来。就是这样——我们有很多不明}
\BookLine{123.50}{22.14}{}{0}{白的事，这只是万分之一。{\kaifont （又喝了点水，同}}
\BookLine{132.00}{22.14}{}{0}{{\kaifont 时后面传来动物被宰杀时的尖叫）}好啦，大日}
\BookLine{140.50}{22.14}{}{0}{子开始了。他们已经在为宴会宰杀牲畜了。厨}
\BookLine{149.00}{22.14}{}{0}{子们生起了火，草原上满是国王们和首领们的}
\BookLine{157.50}{22.14}{}{0}{帐篷，他们都是来庆祝阿德墨托斯国王和伊俄}
\BookLine{166.00}{22.14}{}{0}{尔科斯国王佩利阿斯之女阿尔刻提斯公主的婚}
\BookLine{174.50}{22.14}{}{0}{礼的。}
\newpage
% 原书页码：36
\null
\BookNumber{36}
\Vertical{4}{169}{}{6.2}{8}阿\\尔\\刻\\提\\斯\\之\\歌\EndVertical
\BookLine{30.00}{7.38}{\kaifont}{0}{他沿着小路往回走着。阿尔刻提斯穿着白色衣服从宫门}
\BookLine{38.50}{0.00}{\kaifont}{0}{缓步走出。可以听到动物们叫唤的声音。}
\BookLine{55.50}{0.00}{\heifont}{0}{看门人：}
\BookLine{64.00}{22.14}{}{0}{{\kaifont （继续说道）}嘘，她又来了！}
\BookLine{81.00}{7.38}{\kaifont}{0}{看门人在路边略低于舞台的位置藏了起来。}
\BookLine{89.50}{7.38}{\kaifont}{0}{阿尔刻提斯走到舞台中央，举起双臂，低语着。}
\BookLine{106.50}{0.00}{\heifont}{0}{阿尔刻提斯：}
\BookLine{115.00}{22.14}{}{0}{阿波罗！一点启示！一个启示！{\kaifont （走出大门，向}}
\BookLine{123.50}{22.14}{}{0}{{\kaifont 右朝大路上走去）}}
\BookLine{140.50}{0.00}{\heifont}{0}{看门人：}
\BookLine{149.00}{22.14}{}{0}{{\kaifont （模仿她，对着观众轻轻地喊着）}阿波罗！一点}
\BookLine{157.50}{22.14}{}{0}{启示！}
\BookLine{174.50}{7.38}{\kaifont}{0}{动物们仍然在叫唤。老奶妈阿格莱亚从宫殿里跑出来。}
\newpage
% 原书页码：37
\null
\BookNumber{37}
\Vertical{129}{11}{}{7}{10}第\\一\\幕\EndVertical
\BookLine{30.00}{0.00}{\kaifont}{0}{她四处张望，看见了看门人。}
\BookLine{47.00}{0.00}{\heifont}{0}{阿格莱亚：}
\BookLine{55.50}{22.14}{}{1}{{\kaifont （小声地）}公主在哪儿？{\kaifont （看向看门人指的方向）}}
\BookLine{64.00}{22.14}{}{0}{一整个晚上——一整晚不眠不休，忧伤满怀！}
\BookLine{72.50}{22.14}{}{0}{“阿波罗！一点启示！”}
\BookLine{89.50}{0.00}{\heifont}{0}{看门人：}
\BookLine{98.00}{22.14}{}{0}{“阿波罗！一点启示！”}
\BookLine{115.00}{7.38}{\kaifont}{0}{阿尔刻提斯从右侧的门走进来，带着一种神经质的果决}
\BookLine{123.50}{0.00}{\kaifont}{0}{情绪。}
\BookLine{140.50}{0.00}{\heifont}{0}{阿尔刻提斯：}
\BookLine{149.00}{22.14}{}{0}{看门人！}
\BookLine{166.00}{0.00}{\heifont}{0}{看门人：}
\BookLine{174.50}{22.14}{}{0}{在，公主？}
\newpage
% 原书页码：38
\null
\BookNumber{38}
\Vertical{4}{169}{}{6.2}{8}阿\\尔\\刻\\提\\斯\\之\\歌\EndVertical
\BookLine{30.00}{0.00}{\heifont}{0}{阿尔刻提斯：}
\BookLine{38.50}{22.14}{}{0}{把我的马夫都找来。吩咐他们把马都备好，我}
\BookLine{47.00}{22.14}{}{0}{要出门。阿格莱亚，把我的女仆都找来，吩咐}
\BookLine{55.50}{22.14}{}{0}{她们把行李都准备好。}
\BookLine{72.50}{0.00}{\heifont}{0}{阿格莱亚：}
\BookLine{81.00}{22.14}{}{0}{您要出门，公主——在您大婚的日子，要}
\BookLine{89.50}{22.14}{}{0}{出门？}
\BookLine{106.50}{0.00}{\heifont}{0}{阿尔刻提斯：}
\BookLine{115.00}{22.14}{}{0}{{\kaifont （迅速走下台阶到她面前）}阿格莱亚，我别无选}
\BookLine{123.50}{22.14}{}{0}{择。我必须离开。原谅我。不，恨我吧，蔑视}
\BookLine{132.00}{22.14}{}{0}{我吧——不过最终，请忘了我。}
\BookLine{149.00}{0.00}{\heifont}{0}{阿格莱亚：}
\BookLine{157.50}{22.14}{}{0}{公主，这是个耻辱，这对阿德墨托斯国王是种}
\BookLine{166.00}{22.14}{}{0}{侮辱。}
\newpage
% 原书页码：39
\null
\BookNumber{39}
\Vertical{129}{11}{}{7}{10}第\\一\\幕\EndVertical
\BookLine{30.00}{0.00}{\heifont}{0}{阿尔刻提斯：}
\BookLine{38.50}{22.14}{}{0}{我都明白，阿格莱亚。阿格莱亚，等我离开之}
\BookLine{47.00}{22.14}{}{0}{后，告诉国王——告诉阿德墨托斯——所有的}
\BookLine{55.50}{22.14}{}{0}{耻辱我独自承担，我不奢求他可以原谅，让他}
\BookLine{64.00}{22.14}{}{0}{蔑视我，然后忘了我吧。}
\BookLine{81.00}{0.00}{\heifont}{0}{阿格莱亚：}
\BookLine{89.50}{22.14}{}{0}{公主，我是个老妇人了，但我不是宫里一个普}
\BookLine{98.00}{22.14}{}{0}{通的仆人。我养育了幼年的阿德墨托斯，还有}
\BookLine{106.50}{22.14}{}{0}{他的父亲。{\kaifont （对看门人）}看门人，请留我们单独}
\BookLine{115.00}{22.14}{}{0}{说话。}
\BookLine{132.00}{7.38}{\kaifont}{0}{看门人下场。}
\BookLine{149.00}{0.00}{\heifont}{0}{阿格莱亚：}
\BookLine{157.50}{22.14}{}{0}{您不了解阿德墨托斯国王。希腊众岛之中，您}
\BookLine{166.00}{22.14}{}{0}{再找不到比他更好的丈夫了。}
\newpage
% 原书页码：40
\null
\BookNumber{40}
\Vertical{4}{169}{}{6.2}{8}阿\\尔\\刻\\提\\斯\\之\\歌\EndVertical
\BookLine{30.00}{0.00}{\heifont}{0}{阿尔刻提斯：}
\BookLine{38.50}{22.14}{}{0}{我知道。}
\BookLine{55.50}{0.00}{\heifont}{0}{阿格莱亚：}
\BookLine{64.00}{22.14}{}{0}{或许有人比他更善战、更进取、更强壮，但不}
\BookLine{72.50}{22.14}{}{0}{会有人比他更——更受人爱戴了。}
\BookLine{89.50}{0.00}{\heifont}{0}{阿尔刻提斯：}
\BookLine{98.00}{22.14}{}{0}{我都知道，阿格莱亚。我，我也爱阿德墨托斯。}
\BookLine{106.50}{22.14}{}{0}{所以我也不知道为什么高兴不起来。可是比起}
\BookLine{115.00}{22.14}{}{0}{他，我更爱另一个。}
\BookLine{132.00}{0.00}{\heifont}{0}{阿格莱亚：}
\BookLine{140.50}{22.14}{}{0}{另一个？另一个男人？比阿德墨托斯更爱？那}
\BookLine{149.00}{22.14}{}{0}{去吧，公主，快去吧。我们错待您了，您不必}
\BookLine{157.50}{22.14}{}{0}{待在这里。既然您有眼无珠，麻木无知，什么}
\BookLine{166.00}{22.14}{}{0}{都看不到。{\kaifont （严厉地）}看门人！看门人！您路上}
\BookLine{174.50}{22.14}{}{0}{所需要的一切都会准备好，公主，快点走吧。}
\newpage
% 原书页码：41
\null
\BookNumber{41}
\Vertical{129}{11}{}{7}{10}第\\一\\幕\EndVertical
\BookLine{30.00}{0.00}{\heifont}{0}{阿尔刻提斯：}
\BookLine{38.50}{22.14}{}{0}{不，阿格莱亚，不是别的男人。比起阿德墨托}
\BookLine{47.00}{22.14}{}{0}{斯，我更爱的那一位是，是一位神祇，是阿}
\BookLine{55.50}{22.14}{}{0}{波罗。}
\BookLine{72.50}{0.00}{\heifont}{0}{阿格莱亚：}
\BookLine{81.00}{22.14}{}{0}{阿波罗？}
\BookLine{98.00}{0.00}{\heifont}{0}{阿尔刻提斯：}
\BookLine{106.50}{22.14}{}{0}{是的。从我还是个小姑娘起，我就只有一个心}
\BookLine{115.00}{22.14}{}{0}{愿——在德尔斐，做他的女祭司。{\kaifont （越说越沮}}
\BookLine{123.50}{22.14}{}{0}{{\kaifont 丧，突然哭了起来，很激动地）}我想过真正的}
\BookLine{132.00}{22.14}{}{0}{生活。只此一生，奉献所有。——不是我们所}
\BookLine{140.50}{22.14}{}{0}{看到的狂热、骄傲、占有的生活，而是真正的}
\BookLine{149.00}{22.14}{}{0}{生活，在德尔斐，在真理所在之处。}
\BookLine{166.00}{0.00}{\heifont}{0}{阿格莱亚：}
\BookLine{174.50}{22.14}{}{0}{可是，神呼召您了吗？{\kaifont （停顿）}神没有差遣过谁}
\newpage
% 原书页码：42
\null
\BookNumber{42}
\Vertical{4}{169}{}{6.2}{8}阿\\尔\\刻\\提\\斯\\之\\歌\EndVertical
\BookLine{30.00}{22.14}{}{0}{来找您吗？}
\BookLine{47.00}{0.00}{\heifont}{0}{阿尔刻提斯：}
\BookLine{55.50}{22.14}{}{0}{{\kaifont （低声，羞愧地）}没有。}
\BookLine{72.50}{0.00}{\heifont}{0}{阿格莱亚：}
\BookLine{81.00}{22.14}{}{0}{所谓真正的——难道成为阿德墨托斯的妻子、}
\BookLine{89.50}{22.14}{}{0}{孩子们的母亲、塞萨利的王后，就不算真正的}
\BookLine{98.00}{22.14}{}{0}{生活吗？}
\BookLine{115.00}{0.00}{\heifont}{0}{阿尔刻提斯：}
\BookLine{123.50}{22.14}{}{0}{每个女人都可以做妻子，做母亲，成百上千人}
\BookLine{132.00}{22.14}{}{0}{做过王后。然后，我的丈夫，我的孩子……生}
\BookLine{140.50}{22.14}{}{0}{活就开始围绕这五六个人打转，我只会在意与}
\BookLine{149.00}{22.14}{}{0}{他们相关的琐事……每一天就这样过去了——}
\BookLine{157.50}{22.14}{}{0}{被塞满——被他们无数的需求塞满，最终带}
\BookLine{166.00}{22.14}{}{0}{着点爱和尊荣躺进坟墓——仍然和出生时一样}
\BookLine{174.50}{22.14}{}{0}{无知——}
\newpage
% 原书页码：43
\null
\BookNumber{43}
\Vertical{129}{11}{}{7}{10}第\\一\\幕\EndVertical
\BookLine{30.00}{0.00}{\heifont}{0}{阿格莱亚：}
\BookLine{38.50}{22.14}{}{0}{无知？}
\BookLine{55.50}{0.00}{\heifont}{0}{阿尔刻提斯：}
\BookLine{64.00}{22.14}{}{0}{对我们为何出生，为何死去——人们为何出生}
\BookLine{72.50}{22.14}{}{0}{死去——一无所知，和我们出生时一样无知。}
\BookLine{89.50}{0.00}{\heifont}{0}{阿格莱亚：}
\BookLine{98.00}{22.14}{}{0}{{\kaifont （不带感情地）}所以，您认为在德尔斐可以知}
\BookLine{106.50}{22.14}{}{0}{道？但神并没有呼召您啊。}
\BookLine{123.50}{0.00}{\heifont}{0}{阿尔刻提斯：}
\BookLine{132.00}{22.14}{}{0}{{\kaifont （羞愧地）}我送去了奉献，带去了信息，献了}
\BookLine{140.50}{22.14}{}{0}{祭……{\kaifont （停顿）}我是父亲最疼爱的女儿，他曾}
\BookLine{149.00}{22.14}{}{0}{希望我终身不嫁，一直陪伴他左右。但追求者}
\BookLine{157.50}{22.14}{}{0}{们从希腊各个地方涌来，于是父亲提出了一项}
\BookLine{166.00}{22.14}{}{0}{不可能完成的任务。他让他们同时拴住一头狮}
\BookLine{174.50}{22.14}{}{0}{子和一只野猪，并驾着它们沿伊俄尔科斯城的}
\newpage
% 原书页码：44
\null
\BookNumber{44}
\Vertical{4}{169}{}{6.2}{8}阿\\尔\\刻\\提\\斯\\之\\歌\EndVertical
\BookLine{30.00}{22.14}{}{0}{城墙绕满三周。人们从希腊各个地方涌来：伊}
\BookLine{38.50}{22.14}{}{0}{阿宋来了，涅斯托耳来了，宙斯之子赫拉克勒}
\BookLine{47.00}{22.14}{}{0}{斯也来了，还有阿特柔斯。可所有人都失败了。}
\BookLine{55.50}{22.14}{}{0}{几个月过去，不断有新的追求者失败甚至丧命}
\BookLine{64.00}{22.14}{}{0}{于此。父亲和我就坐在城门之上，父亲笑着，}
\BookLine{72.50}{22.14}{}{0}{我也笑着。不是因为我渴望留在父亲身边，而}
\BookLine{81.00}{22.14}{}{0}{是因为我所求的那唯一一件事——此生只在德}
\BookLine{89.50}{22.14}{}{0}{尔斐，做阿波罗的女祭司。}
\BookLine{106.50}{0.00}{\heifont}{0}{阿格莱亚：}
\BookLine{115.00}{22.14}{}{0}{然后阿德墨托斯出现了。他驾驭着狮子和野}
\BookLine{123.50}{22.14}{}{0}{猪——就像驾驭着温顺的公牛一样绕城而行，}
\BookLine{132.00}{22.14}{}{0}{并赢得了您，公主。}
\BookLine{149.00}{0.00}{\heifont}{0}{阿尔刻提斯：}
\BookLine{157.50}{22.14}{}{0}{可我更爱阿波罗。}
\newpage
% 原书页码：45
\null
\BookNumber{45}
\Vertical{129}{11}{}{7}{10}第\\一\\幕\EndVertical
\BookLine{30.00}{0.00}{\heifont}{0}{阿格莱亚：}
\BookLine{38.50}{22.14}{}{0}{是，可正是阿波罗安排了这场婚礼。}
\BookLine{55.50}{0.00}{\heifont}{0}{阿尔刻提斯：}
\BookLine{64.00}{22.14}{}{0}{那可说不准。}
\BookLine{81.00}{0.00}{\heifont}{0}{阿格莱亚：}
\BookLine{89.50}{22.14}{}{0}{您要求的启示，公主，就在您的面前——最清}
\BookLine{98.00}{22.14}{}{0}{楚无误的启示。{\kaifont （激烈地）}您从没有被呼召前往}
\BookLine{106.50}{22.14}{}{0}{德尔斐，难道您看不出神最简单的指引吗？}
\BookLine{123.50}{7.38}{\kaifont}{0}{阿尔刻提斯捧起了自己的面庞。}
\BookLine{140.50}{0.00}{\heifont}{0}{阿格莱亚：}
\BookLine{149.00}{22.14}{}{0}{现在让我告诉您：阿德墨托斯能够驾驭狮子与}
\BookLine{157.50}{22.14}{}{0}{野猪，难道您不惊讶吗？阿特柔斯失败了，宙}
\BookLine{166.00}{22.14}{}{0}{斯之子赫拉克勒斯也失败了，他却做到了。我}
\newpage
% 原书页码：46
\null
\BookNumber{46}
\Vertical{4}{169}{}{6.2}{8}阿\\尔\\刻\\提\\斯\\之\\歌\EndVertical
\BookLine{30.00}{22.14}{}{0}{来告诉您他是怎么做到的：阿波罗曾在梦中教}
\BookLine{38.50}{22.14}{}{0}{授他如何同时驾驭一只狮子和一只野猪。}
\BookLine{55.50}{7.38}{\kaifont}{0}{阿尔刻提斯向后退了两步。}
\BookLine{72.50}{0.00}{\heifont}{0}{阿格莱亚：}
\BookLine{81.00}{22.14}{}{0}{起初，他对您一见钟情。他第二次回到伊俄尔}
\BookLine{89.50}{22.14}{}{0}{科斯之前——也就是第一次失败后，他病了。}
\BookLine{98.00}{22.14}{}{0}{爱情与绝望让他濒临死亡，我一直在身旁照顾}
\BookLine{106.50}{22.14}{}{0}{他，有三个晚上他都差点死去。就在第三晚，}
\BookLine{115.00}{22.14}{}{0}{我就坐在他身边——在痛苦与谵妄之中——我}
\BookLine{123.50}{22.14}{}{0}{听到，我看到阿波罗正在梦中教授他，如何同}
\BookLine{132.00}{22.14}{}{0}{时驾驭一只狮子和一只野猪。}
\BookLine{149.00}{7.38}{\kaifont}{0}{阿尔刻提斯盯着她。}
\BookLine{166.00}{0.00}{\heifont}{0}{阿格莱亚：}
\BookLine{174.50}{22.14}{}{0}{这是真的。我发誓这都是真的。}
\newpage
% 原书页码：47
\null
\BookNumber{47}
\Vertical{129}{11}{}{7}{10}第\\一\\幕\EndVertical
\BookLine{30.00}{0.00}{\heifont}{0}{阿尔刻提斯：}
\BookLine{38.50}{22.14}{}{0}{是真的，好的。不过我们听说过许多关于谵妄}
\BookLine{47.00}{22.14}{}{0}{和梦境、发热和幻象的事情。阿格莱亚，我们}
\BookLine{55.50}{22.14}{}{0}{必须百分百确定。神清晰地公然降临——只会}
\BookLine{64.00}{22.14}{}{0}{是在德尔斐。}
\BookLine{81.00}{0.00}{\heifont}{0}{阿格莱亚：}
\BookLine{89.50}{22.14}{}{0}{清晰？公然？即便在德尔斐，女预言家们也都}
\BookLine{98.00}{22.14}{}{0}{处在谵妄之中，她们高声咆哮，魂游象外。谁}
\BookLine{106.50}{22.14}{}{0}{听到过她们清晰明白地讲话？}
\BookLine{123.50}{0.00}{\heifont}{0}{阿尔刻提斯：}
\BookLine{132.00}{22.14}{}{0}{{\kaifont （拖着犹豫不决的步子朝王宫走去，十分沮丧）}}
\BookLine{140.50}{22.14}{}{0}{我孤独一人，只有一人……}
\BookLine{157.50}{0.00}{\heifont}{0}{阿格莱亚：}
\BookLine{166.00}{22.14}{}{0}{{\kaifont （坚定地，动情地）}听我说，公主。回到房间好}
\BookLine{174.50}{22.14}{}{0}{好睡一觉。{\kaifont （朝天看一眼）}还有两个钟头就到正}
\newpage
% 原书页码：48
\null
\BookNumber{48}
\Vertical{4}{169}{}{6.2}{8}阿\\尔\\刻\\提\\斯\\之\\歌\EndVertical
\BookLine{30.00}{22.14}{}{0}{午了。如果休息之后，您的想法还是没有改变，}
\BookLine{38.50}{22.14}{}{0}{您想去哪里就去哪里，没有人会阻拦。{\kaifont （托起阿}}
\BookLine{47.00}{22.14}{}{0}{{\kaifont 尔刻提斯的胳膊，把她领进宫门，像母亲和女}}
\BookLine{55.50}{22.14}{}{0}{{\kaifont 儿闲聊一样）}您希望神明们清晰坦率地对我们}
\BookLine{64.00}{22.14}{}{0}{说话，公主？您把他们当作什么了啊，公主？}
\BookLine{72.50}{22.14}{}{0}{您可不能把他们想象成普通的人啊！}
\BookLine{89.50}{7.38}{\kaifont}{0}{两人都退回宫殿中。舞台上空了一段时间。门外一阵声}
\BookLine{98.00}{0.00}{\kaifont}{0}{响愈加吵闹起来，有人在敲右侧的门。看门人绕过宫殿走到}
\BookLine{106.50}{0.00}{\kaifont}{0}{舞台中央。}
\BookLine{123.50}{0.00}{\heifont}{0}{看门人：}
\BookLine{132.00}{22.14}{}{0}{好吧，又是什么情况啊？吵吵嚷嚷是怎么回}
\BookLine{140.50}{22.14}{}{0}{事？{\kaifont （开门，透过缝隙交谈）}婚礼的客人们请走}
\BookLine{149.00}{22.14}{}{0}{宫殿正门。这是后门。什么？能不能不要同时}
\BookLine{157.50}{22.14}{}{0}{讲话！什么？好吧，好吧，让老人家进来。}
\BookLine{174.50}{7.38}{\kaifont}{0}{忒瑞西阿斯走了进来，他是个瞎子，非常非常老，容易}
\newpage
% 原书页码：49
\null
\BookNumber{49}
\Vertical{129}{11}{}{7}{10}第\\一\\幕\EndVertical
\BookLine{30.00}{0.00}{\kaifont}{0}{动怒，言辞恶毒，盛气凌人，因为高龄，讲起话像是胡言乱}
\BookLine{38.50}{0.00}{\kaifont}{0}{语。他一只手搭在带领着自己的男孩肩上，另一只手则不停}
\BookLine{47.00}{0.00}{\kaifont}{0}{挥舞着手中的手杖。镇上的人们都跟着他走进院子，仆人们}
\BookLine{55.50}{0.00}{\kaifont}{0}{也从宫门和皇宫四周来到院子里。}
\BookLine{72.50}{0.00}{\heifont}{0}{忒瑞西阿斯：}
\BookLine{81.00}{22.14}{}{0}{{\kaifont （声音倒是令人惊讶地响亮）}这是克里特国王米}
\BookLine{89.50}{22.14}{}{0}{诺斯的王宫吗？}
\BookLine{106.50}{7.38}{\kaifont}{0}{大家笑起来，男孩拽了拽他的袖子并对他耳语几句。}
\BookLine{123.50}{0.00}{\heifont}{0}{忒瑞西阿斯：}
\BookLine{132.00}{22.14}{}{0}{我是说，这是忒拜国王俄狄浦斯的王宫吗？{\kaifont （用}}
\BookLine{140.50}{22.14}{}{0}{{\kaifont 手杖敲打男孩）}别推搡我！我知道自己在说什}
\BookLine{149.00}{22.14}{}{0}{么！{\kaifont （躲开男孩的推搡）}蜜蜂！胡蜂！大黄蜂！}
\BookLine{166.00}{0.00}{\heifont}{0}{看门人：}
\BookLine{174.50}{22.14}{}{0}{老人家。这是塞萨利国王，阿德墨托斯的王宫。}
\newpage
% 原书页码：50
\null
\BookNumber{50}
\Vertical{4}{169}{}{6.2}{8}阿\\尔\\刻\\提\\斯\\之\\歌\EndVertical
\BookLine{30.00}{0.00}{\heifont}{0}{忒瑞西阿斯：}
\BookLine{38.50}{22.14}{}{0}{{\kaifont （重复着他的话）}塞萨利国王。对，我是这么说}
\BookLine{47.00}{22.14}{}{0}{的。我就是这个意思。把塞萨利国王阿德墨托}
\BookLine{55.50}{22.14}{}{0}{斯找来，我有个消息带给他。}
\BookLine{72.50}{0.00}{\heifont}{0}{看门人：}
\BookLine{81.00}{22.14}{}{0}{老人家，国王今天要举行婚礼，他此刻正忙着}
\BookLine{89.50}{22.14}{}{0}{招待客人。您先坐在这里晒晒太阳，我们会给}
\BookLine{98.00}{22.14}{}{0}{您洗洗脚。稍后国王会来听您的消息。}
\BookLine{115.00}{0.00}{\heifont}{0}{忒瑞西阿斯：}
\BookLine{123.50}{22.14}{}{0}{{\kaifont （举着手杖威胁道）}结婚……洗脚……结婚和}
\BookLine{132.00}{22.14}{}{0}{洗脚跟我有什么关系？我一分钟都不会多等。}
\BookLine{140.50}{22.14}{}{0}{{\kaifont （跺脚）}把那个什么什么国王给我找来！}
\BookLine{157.50}{7.38}{\kaifont}{0}{阿格莱亚从宫中走出来。}
\newpage
% 原书页码：51
\null
\BookNumber{51}
\Vertical{129}{11}{}{7}{10}第\\一\\幕\EndVertical
\BookLine{30.00}{0.00}{\heifont}{0}{阿格莱亚：}
\BookLine{38.50}{22.14}{}{0}{老人家，请问您是？我去通报国王。}
\BookLine{55.50}{0.00}{\heifont}{0}{忒瑞西阿斯：}
\BookLine{64.00}{22.14}{}{0}{告诉国王，我是德尔斐，忒瑞西阿斯的祭}
\BookLine{72.50}{22.14}{}{0}{司——我是阿波罗，德尔斐的祭司……孩子，}
\BookLine{81.00}{22.14}{}{0}{我在说什么？}
\BookLine{98.00}{7.38}{\kaifont}{0}{男孩耳语几句。}
\BookLine{115.00}{0.00}{\heifont}{0}{忒瑞西阿斯：}
\BookLine{123.50}{22.14}{}{0}{告诉国王我是忒瑞西阿斯，阿波罗的祭司。我}
\BookLine{132.00}{22.14}{}{0}{从德尔斐来，带来了一个消息。我今天还要赶}
\BookLine{140.50}{22.14}{}{0}{回去。}
\BookLine{157.50}{0.00}{\heifont}{0}{阿格莱亚：}
\BookLine{166.00}{22.14}{}{0}{你是忒瑞西阿斯？忒瑞西阿斯！}
\newpage
% 原书页码：52
\null
\BookNumber{52}
\Vertical{4}{169}{}{6.2}{8}阿\\尔\\刻\\提\\斯\\之\\歌\EndVertical
\BookLine{30.00}{0.00}{\heifont}{0}{看门人和众人：}
\BookLine{38.50}{22.14}{}{0}{忒瑞西阿斯！}
\BookLine{55.50}{0.00}{\heifont}{0}{忒瑞西阿斯：}
\BookLine{64.00}{22.14}{}{0}{{\kaifont （拿手杖敲打地面）}去把国王叫来！灾难和瘟疫}
\BookLine{72.50}{22.14}{}{0}{要来了！去把那个什么什么国王叫来！}
\BookLine{89.50}{0.00}{\heifont}{0}{阿格莱亚：}
\BookLine{98.00}{22.14}{}{0}{这边请，伟大的忒瑞西阿斯。}
\BookLine{115.00}{7.38}{\kaifont}{0}{她匆匆走进宫殿，门开了，阿德墨托斯走了出来，更多}
\BookLine{123.50}{0.00}{\kaifont}{0}{的仆人围了过来。}
\BookLine{140.50}{0.00}{\heifont}{0}{阿德墨托斯：}
\BookLine{149.00}{22.14}{}{0}{出什么事了？阿格莱亚，这是谁？}
\BookLine{166.00}{0.00}{\heifont}{0}{阿格莱亚：}
\BookLine{174.50}{22.14}{}{0}{{\kaifont （悄悄地）}是从德尔斐来的，忒瑞西阿斯。}
\newpage
% 原书页码：53
\null
\BookNumber{53}
\Vertical{129}{11}{}{7}{10}第\\一\\幕\EndVertical
\BookLine{30.00}{0.00}{\heifont}{0}{阿德墨托斯：}
\BookLine{38.50}{22.14}{}{0}{忒瑞西阿斯！}
\BookLine{55.50}{0.00}{\heifont}{0}{阿格莱亚：}
\BookLine{64.00}{22.14}{}{0}{{\kaifont （手指前额）}如山峰一样古老，阿德墨托斯}
\BookLine{72.50}{22.14}{}{0}{国王。}
\BookLine{89.50}{0.00}{\heifont}{0}{阿德墨托斯：}
\BookLine{98.00}{22.14}{}{0}{欢迎您，欢迎您，尊贵的忒瑞西阿斯。您是我}
\BookLine{106.50}{22.14}{}{0}{父亲的老朋友了。欢迎来到费莱。我是塞萨利}
\BookLine{115.00}{22.14}{}{0}{国王，阿德墨托斯。}
\BookLine{132.00}{0.00}{\heifont}{0}{忒瑞西阿斯：}
\BookLine{140.50}{22.14}{}{0}{{\kaifont （舞着手杖）}后边去，往后站。又挤又乱的……}
\BookLine{149.00}{22.14}{}{0}{你，带着耳朵呢吗？}
\BookLine{166.00}{0.00}{\heifont}{0}{阿德墨托斯：}
\BookLine{174.50}{22.14}{}{0}{是的，忒瑞西阿斯。}
\newpage
% 原书页码：54
\null
\BookNumber{54}
\Vertical{4}{169}{}{6.2}{8}阿\\尔\\刻\\提\\斯\\之\\歌\EndVertical
\BookLine{30.00}{0.00}{\heifont}{0}{忒瑞西阿斯：}
\BookLine{38.50}{22.14}{}{0}{那就把你的耳屎都清理干净，仔细聆听神的}
\BookLine{47.00}{22.14}{}{0}{教诲。}
\BookLine{64.00}{0.00}{\heifont}{0}{阿德墨托斯：}
\BookLine{72.50}{22.14}{}{0}{洗耳恭听，忒瑞西阿斯。}
\BookLine{89.50}{0.00}{\heifont}{0}{忒瑞西阿斯：}
\BookLine{98.00}{22.14}{}{0}{迈锡尼国王阿特柔斯，聆听神的教诲——}
\BookLine{115.00}{0.00}{\heifont}{0}{阿德墨托斯：}
\BookLine{123.50}{22.14}{}{0}{阿特柔斯？尊贵的忒瑞西阿斯，我是阿德墨托}
\BookLine{132.00}{22.14}{}{0}{斯，是——}
\BookLine{149.00}{0.00}{\heifont}{0}{忒瑞西阿斯：}
\BookLine{157.50}{22.14}{}{0}{阿德墨托斯？哦，那阿德墨托斯，你先住嘴，}
\BookLine{166.00}{22.14}{}{0}{听我把话说完。我从德尔斐的阿波罗神殿为你}
\BookLine{174.50}{22.14}{}{0}{带来了一个消息。荣耀，一份至高无上的荣耀}
\newpage
% 原书页码：55
\null
\BookNumber{55}
\Vertical{129}{11}{}{7}{10}第\\一\\幕\EndVertical
\BookLine{30.00}{22.14}{}{0}{将会降临塞萨利。孩子，这儿是塞萨利吧？{\kaifont （把}}
\BookLine{38.50}{22.14}{}{0}{{\kaifont 手放在男孩的头上，男孩点了点头）}一份至高}
\BookLine{47.00}{22.14}{}{0}{无上的荣耀，一场危难将会降临塞萨利。}
\BookLine{64.00}{0.00}{\heifont}{0}{阿德墨托斯：}
\BookLine{72.50}{22.14}{}{0}{您说危难？忒瑞西阿斯？}
\BookLine{89.50}{0.00}{\heifont}{0}{忒瑞西阿斯：}
\BookLine{98.00}{22.14}{}{0}{荣耀，还有危难。危难就是荣耀，傻瓜。不对，}
\BookLine{106.50}{22.14}{}{0}{荣耀就是危难。在塞萨利，所有事情不都是要}
\BookLine{115.00}{22.14}{}{0}{先告诉你吗？}
\BookLine{132.00}{0.00}{\heifont}{0}{阿德墨托斯：}
\BookLine{140.50}{22.14}{}{0}{稍等片刻，忒瑞西阿斯。给我的消息同样应该}
\BookLine{149.00}{22.14}{}{0}{让我未来的王后知悉。阿格莱亚，把公主叫来。}
\BookLine{166.00}{7.38}{\kaifont}{0}{阿格莱亚匆匆走进宫殿。}
\newpage
% 原书页码：56
\null
\BookNumber{56}
\Vertical{4}{169}{}{6.2}{8}阿\\尔\\刻\\提\\斯\\之\\歌\EndVertical
\BookLine{30.00}{0.00}{\heifont}{0}{阿德墨托斯：}
\BookLine{38.50}{22.14}{}{0}{我要在今天迎娶伊俄尔科斯国王佩利阿斯的女}
\BookLine{47.00}{22.14}{}{0}{儿阿尔刻提斯，忒瑞西阿斯，您将是最尊贵的}
\BookLine{55.50}{22.14}{}{0}{客人。请先休息片刻，忒瑞西阿斯……}
\BookLine{72.50}{0.00}{\heifont}{0}{忒瑞西阿斯：}
\BookLine{81.00}{22.14}{}{0}{婚礼我天天见，让王后快点。听我说，克里特}
\BookLine{89.50}{22.14}{}{0}{国王米诺斯……孩子，他叫什么名字来着？}
\BookLine{106.50}{7.38}{\kaifont}{0}{男孩耳语。}
\BookLine{123.50}{0.00}{\heifont}{0}{忒瑞西阿斯：}
\BookLine{132.00}{22.14}{}{0}{好吧，这有什么关系？王后来了吗？}
\BookLine{149.00}{7.38}{\kaifont}{0}{阿尔刻提斯从宫中走出来，急匆匆的，并充满好奇。}
\BookLine{166.00}{0.00}{\heifont}{0}{阿德墨托斯：}
\BookLine{174.50}{22.14}{}{0}{她来了。}
\newpage
% 原书页码：57
\null
\BookNumber{57}
\Vertical{129}{11}{}{7}{10}第\\一\\幕\EndVertical
\BookLine{30.00}{0.00}{\heifont}{0}{阿尔刻提斯：}
\BookLine{38.50}{22.14}{}{0}{尊贵的忒瑞西阿斯，伟大的忒瑞西阿斯！我父}
\BookLine{47.00}{22.14}{}{0}{亲的老朋友。我是阿尔刻提斯，伊俄尔科斯国}
\BookLine{55.50}{22.14}{}{0}{王佩利阿斯的女儿。}
\BookLine{72.50}{0.00}{\heifont}{0}{忒瑞西阿斯：}
\BookLine{81.00}{22.14}{}{0}{{\kaifont （挥舞着手杖驱逐着围在他身边的人）}后边！往}
\BookLine{89.50}{22.14}{}{0}{后去！野鹅，野鸭，呱呱乱叫。闭上你们的嘴，}
\BookLine{98.00}{22.14}{}{0}{安静一会儿。众神与人类之父宙斯命令了……}
\BookLine{106.50}{22.14}{}{0}{已经命令……孩子，他命令了什么来着？}
\BookLine{123.50}{7.38}{\kaifont}{0}{男孩耳语，忒瑞西阿斯用手杖打了他一下。}
\BookLine{140.50}{0.00}{\heifont}{0}{忒瑞西阿斯：}
\BookLine{149.00}{22.14}{}{0}{嗯，你们不必……他命令阿波罗——我的主}
\BookLine{157.50}{22.14}{}{0}{人——阿波罗从奥林匹斯山下来，在人间生活}
\BookLine{166.00}{22.14}{}{0}{一年，从这个至点到下个至点，作为一个人类}
\BookLine{174.50}{22.14}{}{0}{和大家一起生活。我已经把信息带到。孩子，}
\newpage
% 原书页码：58
\null
\BookNumber{58}
\Vertical{4}{169}{}{6.2}{8}阿\\尔\\刻\\提\\斯\\之\\歌\EndVertical
\BookLine{30.00}{22.14}{}{0}{我们上路吧。{\kaifont （他转身要走）}}
\BookLine{47.00}{0.00}{\heifont}{0}{阿尔刻提斯：}
\BookLine{55.50}{22.14}{}{0}{{\kaifont （当男孩还在对忒瑞西阿斯耳语时）}阿波罗要在}
\BookLine{64.00}{22.14}{}{0}{这地上生活？}
\BookLine{81.00}{0.00}{\heifont}{0}{忒瑞西阿斯：}
\BookLine{89.50}{22.14}{}{0}{{\kaifont （对男孩）}是的，是的。我又没聋。阿波罗，我}
\BookLine{98.00}{22.14}{}{0}{的主人，他选择了生活在这里。{\kaifont （用手杖敲敲脚}}
\BookLine{106.50}{22.14}{}{0}{{\kaifont 下的土地）}在这儿，作为塞萨利国王阿德墨托}
\BookLine{115.00}{22.14}{}{0}{斯的仆人。}
\BookLine{132.00}{0.00}{\heifont}{0}{阿德墨托斯：}
\BookLine{140.50}{22.14}{}{0}{在这儿？这儿？尊贵的忒瑞西阿斯？{\kaifont （飞快走到}}
\BookLine{149.00}{22.14}{}{0}{{\kaifont 他身边）}稍等一下，忒瑞西阿斯。我该怎么理}
\BookLine{157.50}{22.14}{}{0}{解这件事情？神圣的忒瑞西阿斯，您该不是说，}
\BookLine{166.00}{22.14}{}{0}{阿波罗要在这儿做我的仆人，每一天？每一天}
\BookLine{174.50}{22.14}{}{0}{都和我们生活在一起？}
\newpage
% 原书页码：59
\null
\BookNumber{59}
\Vertical{129}{11}{}{7}{10}第\\一\\幕\EndVertical
\BookLine{30.00}{7.38}{\kaifont}{0}{大家都屏气凝神地望着他，忒瑞西阿斯却垂下手臂，像}
\BookLine{38.50}{0.00}{\kaifont}{0}{是睡着了一样。突然他醒了过来，说起话来。}
\BookLine{55.50}{0.00}{\heifont}{0}{忒瑞西阿斯：}
\BookLine{64.00}{22.14}{}{0}{门外有四个牧人。他们在这一年之中会成为你}
\BookLine{72.50}{22.14}{}{0}{的仆人。给他们找些活儿干。他们其中之一就}
\BookLine{81.00}{22.14}{}{0}{是阿波罗。}
\BookLine{98.00}{0.00}{\heifont}{0}{阿德墨托斯：}
\BookLine{106.50}{22.14}{}{0}{{\kaifont （重复道）}其中一个是阿波罗？}
\BookLine{123.50}{0.00}{\heifont}{0}{忒瑞西阿斯：}
\BookLine{132.00}{22.14}{}{0}{四个牧人。其中一个是阿波罗。不用去想哪一}
\BookLine{140.50}{22.14}{}{0}{个是神，我不知道，你也不会知道。也别再问}
\BookLine{149.00}{22.14}{}{0}{我问题了，我没什么可回答的了。孩子，把牧}
\BookLine{157.50}{22.14}{}{0}{人们叫来。}
\BookLine{174.50}{7.38}{\kaifont}{0}{男孩走出去，静场。}
\newpage
% 原书页码：60
\null
\BookNumber{60}
\Vertical{4}{169}{}{6.2}{8}阿\\尔\\刻\\提\\斯\\之\\歌\EndVertical
\BookLine{30.00}{0.00}{\heifont}{0}{阿德墨托斯：}
\BookLine{38.50}{22.14}{}{0}{忒瑞西阿斯，我们是不是应该掩面下跪？}
\BookLine{55.50}{0.00}{\heifont}{0}{忒瑞西阿斯：}
\BookLine{64.00}{22.14}{}{0}{你没有听到我刚刚说的话吗？阿波罗在这里，}
\BookLine{72.50}{22.14}{}{0}{是一个普通人类。一个普通人，和普通的牧人}
\BookLine{81.00}{22.14}{}{0}{一样。像我这样就行了！}
\BookLine{98.00}{7.38}{\kaifont}{0}{男孩回来站到忒瑞西阿斯身边。四个牧人走了进来。他}
\BookLine{106.50}{0.00}{\kaifont}{0}{们满身尘土，肮脏且不修边幅，在人群面前显得十分害羞，}
\BookLine{115.00}{0.00}{\kaifont}{0}{他们卑躬屈膝地手扶前额，转向墙角站成一排，眼睛也不知}
\BookLine{123.50}{0.00}{\kaifont}{0}{该往哪儿看。有两人带着酒囊，每人都有一支手杖。忒瑞西}
\BookLine{132.00}{0.00}{\kaifont}{0}{阿斯粗声粗气地对他们说话。}
\BookLine{149.00}{0.00}{\heifont}{0}{忒瑞西阿斯：}
\BookLine{157.50}{22.14}{}{0}{{\kaifont （继续）}过来，别磨磨蹭蹭的。给你们的新主人}
\BookLine{166.00}{22.14}{}{0}{行礼。谁都看得出来你们喝了酒，工作开了个}
\BookLine{174.50}{22.14}{}{0}{好头啊。{\kaifont （挥舞着手杖）}要是我看到了肯定会揍}
\newpage
% 原书页码：61
\null
\BookNumber{61}
\Vertical{129}{11}{}{7}{10}第\\一\\幕\EndVertical
\BookLine{30.00}{22.14}{}{0}{你们。抬脚，往前走。孩子，四个都在这儿了}
\BookLine{38.50}{22.14}{}{0}{吗？嗯，国王成哑巴啦？}
\BookLine{55.50}{0.00}{\heifont}{0}{阿德墨托斯：}
\BookLine{64.00}{22.14}{}{0}{{\kaifont （整理了下情绪）}欢迎你们到塞萨利。欢迎你们}
\BookLine{72.50}{22.14}{}{0}{来参加宴会，今天是我的婚礼。明天我就会将}
\BookLine{81.00}{22.14}{}{0}{羊群和牛群分派给你们。长途跋涉辛苦了。欢}
\BookLine{89.50}{22.14}{}{0}{迎你们到塞萨利……忒瑞西阿斯，也欢迎你，}
\BookLine{98.00}{22.14}{}{0}{长途跋涉辛苦了。需要沐浴休息一下吗？}
\BookLine{115.00}{0.00}{\heifont}{0}{忒瑞西阿斯：}
\BookLine{123.50}{22.14}{}{0}{我还要继续长途跋涉。我的信息已经带到了。}
\BookLine{132.00}{22.14}{}{0}{孩子，带路走吧。}
\BookLine{149.00}{0.00}{\heifont}{0}{阿尔刻提斯：}
\BookLine{157.50}{22.14}{}{0}{{\kaifont （几步走到他面前，低声地）}神圣的忒瑞西阿}
\BookLine{166.00}{22.14}{}{0}{斯，没有什么消息带给阿尔刻提斯吗？}
\newpage
% 原书页码：62
\null
\BookNumber{62}
\Vertical{4}{169}{}{6.2}{8}阿\\尔\\刻\\提\\斯\\之\\歌\EndVertical
\BookLine{30.00}{0.00}{\heifont}{0}{忒瑞西阿斯：}
\BookLine{38.50}{22.14}{}{0}{这个女人是谁？}
\BookLine{55.50}{0.00}{\heifont}{0}{阿尔刻提斯：}
\BookLine{64.00}{22.14}{}{0}{我是阿尔刻提斯，伊俄尔科斯国王佩利阿斯的}
\BookLine{72.50}{22.14}{}{0}{女儿。我向德尔斐送去了许多祷告和奉献。}
\BookLine{89.50}{0.00}{\heifont}{0}{忒瑞西阿斯：}
\BookLine{98.00}{22.14}{}{0}{祷告和奉献，多得数不清。孩子，带路。}
\BookLine{115.00}{7.38}{\kaifont}{0}{男孩拉住他的袖子，扶着他的肩膀，想要对他耳语几句。}
\BookLine{132.00}{0.00}{\heifont}{0}{忒瑞西阿斯：}
\BookLine{140.50}{22.14}{}{0}{哦对，我是有一个消息，要带给一个女孩，或}
\BookLine{149.00}{22.14}{}{0}{是女人。带给伊俄卡斯忒，还是阿尔刻提斯，}
\BookLine{157.50}{22.14}{}{0}{或者是得伊阿尼拉，我也记不得是谁了，不过}
\BookLine{166.00}{22.14}{}{0}{这事儿我还记着。你这孩子，别老拽我！{\kaifont （拿着}}
\BookLine{174.50}{22.14}{}{0}{{\kaifont 手杖打男孩几下）}废物！放肆！}
\newpage
% 原书页码：63
\null
\BookNumber{63}
\Vertical{129}{11}{}{7}{10}第\\一\\幕\EndVertical
\BookLine{30.00}{7.38}{\kaifont}{0}{男孩摔倒在地，忒瑞西阿斯并没有住手。}
\BookLine{47.00}{0.00}{\heifont}{0}{男孩：}
\BookLine{55.50}{22.14}{}{0}{{\kaifont （喊叫着）}忒瑞西阿斯！救命！救命！阿德墨托}
\BookLine{64.00}{22.14}{}{0}{斯国王！}
\BookLine{81.00}{0.00}{\heifont}{0}{阿德墨托斯：}
\BookLine{89.50}{22.14}{}{0}{是啊，忒瑞西阿斯，这孩子也没有——}
\BookLine{106.50}{0.00}{\heifont}{0}{阿尔刻提斯：}
\BookLine{115.00}{22.14}{}{0}{只是个小小的错误，忒瑞西阿斯。求你放过他}
\BookLine{123.50}{22.14}{}{0}{吧，他知道错了。}
\BookLine{140.50}{0.00}{\heifont}{0}{忒瑞西阿斯：}
\BookLine{149.00}{22.14}{}{0}{{\kaifont （突然停了下来，盯着阿尔刻提斯）}我不管你叫}
\BookLine{157.50}{22.14}{}{0}{什么：伊俄卡斯忒，勒达，赫耳弥俄涅——}
\newpage
% 原书页码：64
\null
\BookNumber{64}
\Vertical{4}{169}{}{6.2}{8}阿\\尔\\刻\\提\\斯\\之\\歌\EndVertical
\BookLine{30.00}{0.00}{\heifont}{0}{阿尔刻提斯：}
\BookLine{38.50}{22.14}{}{0}{阿尔刻提斯。}
\BookLine{55.50}{0.00}{\heifont}{0}{忒瑞西阿斯：}
\BookLine{64.00}{22.14}{}{0}{我是有个消息，要带给某个女孩，但我忘了，}
\BookLine{72.50}{22.14}{}{0}{要么我已经说过了。嗯，我已经说过了，借着}
\BookLine{81.00}{22.14}{}{0}{雷声与闪电，借着神圣的祭坛——如果人类学}
\BookLine{89.50}{22.14}{}{0}{不会聆听，德尔斐要来有什么用？}
\BookLine{106.50}{7.38}{\kaifont}{0}{忒瑞西阿斯跟在男孩后面走了出去，阿德墨托斯跟上}
\BookLine{115.00}{0.00}{\kaifont}{0}{几步。}
\BookLine{132.00}{0.00}{\heifont}{0}{阿德墨托斯：}
\BookLine{140.50}{22.14}{}{0}{您说，您说会有一场危难是吗，忒瑞西阿斯？}
\BookLine{157.50}{0.00}{\heifont}{0}{忒瑞西阿斯：}
\BookLine{166.00}{22.14}{}{0}{{\kaifont （马上要走出门外）}当然有危难了，呆子。每次}
\BookLine{174.50}{22.14}{}{0}{他们{\kaifont （指着天空）}的来临就是一场危难。}
\newpage
% 原书页码：65
\null
\BookNumber{65}
\Vertical{129}{11}{}{7}{10}第\\一\\幕\EndVertical
\BookLine{30.00}{0.00}{\heifont}{0}{阿德墨托斯：}
\BookLine{38.50}{22.14}{}{0}{可我父亲说，阿波罗是爱着塞萨利的……}
\BookLine{55.50}{0.00}{\heifont}{0}{忒瑞西阿斯：}
\BookLine{64.00}{22.14}{}{0}{是的，爱爱爱。大家还是自个爱自个吧。你看}
\BookLine{72.50}{22.14}{}{0}{看我，五六百岁了，神很爱我，死都不让我死。}
\BookLine{81.00}{22.14}{}{0}{如果神不爱人类，我们肯定都很快乐；换个角}
\BookLine{89.50}{22.14}{}{0}{度同样成立：如果我们不爱神，我们肯定也很}
\BookLine{98.00}{22.14}{}{0}{快乐。}
\BookLine{115.00}{7.38}{\kaifont}{0}{忒瑞西阿斯带着男孩下场。一个尴尬的停顿。}
\BookLine{123.50}{7.38}{\kaifont}{0}{阿德墨托斯整理了一下情绪，带着实事求是却又威严的}
\BookLine{132.00}{0.00}{\kaifont}{0}{语气对四个牧人开始讲话。}
\BookLine{149.00}{0.00}{\heifont}{0}{阿德墨托斯：}
\BookLine{157.50}{22.14}{}{0}{再次欢迎你们来到塞萨利，来到费莱。{\kaifont （对看门}}
\BookLine{166.00}{22.14}{}{0}{{\kaifont 人）}请确保安顿好他们。{\kaifont （对牧人们）}很高兴你}
\BookLine{174.50}{22.14}{}{0}{们也成为了今天婚礼的客人。}
\newpage
% 原书页码：66
\null
\BookNumber{66}
\Vertical{4}{169}{}{6.2}{8}阿\\尔\\刻\\提\\斯\\之\\歌\EndVertical
\BookLine{30.00}{7.38}{\kaifont}{0}{看门人带领四个牧人沿着小路走向泉边，阿德墨托斯和}
\BookLine{38.50}{0.00}{\kaifont}{0}{舞台上的人们都好奇地看着。他们走过阿德墨托斯身边时，}
\BookLine{47.00}{0.00}{\kaifont}{0}{都像仆人一样怯生生的，到了泉边他们才放松下来，互相传}
\BookLine{55.50}{0.00}{\kaifont}{0}{递着酒囊，其中一人渐渐睡着了。阿德墨托斯并没有看向阿}
\BookLine{64.00}{0.00}{\kaifont}{0}{尔刻提斯，他对她说话，也是在对自己说。}
\BookLine{81.00}{0.00}{\heifont}{0}{阿德墨托斯：}
\BookLine{89.50}{22.14}{}{0}{{\kaifont （继续）}我也不知道该怎么理解这些事儿……我}
\BookLine{98.00}{22.14}{}{0}{自己也不过就是个牧人。阿尔刻提斯，费莱很}
\BookLine{106.50}{22.14}{}{0}{需要你。}
\BookLine{123.50}{7.38}{\kaifont}{0}{他伸出双手背到身后。她，沉思着愣在原地，毫无回应。}
\BookLine{140.50}{0.00}{\heifont}{0}{阿德墨托斯：}
\BookLine{149.00}{22.14}{}{0}{我得回去招呼客人了。{\kaifont （再次表达畏惧）}我真的}
\BookLine{157.50}{22.14}{}{0}{不知道，该怎么理解这些事情。{\kaifont （忽然笑了一}}
\BookLine{166.00}{22.14}{}{0}{{\kaifont 下）}阿尔刻提斯，在塞萨利有一个古老的习俗，}
\BookLine{174.50}{22.14}{}{0}{直到婚礼当晚，新郎都不能见到新娘的面容。}
\newpage
% 原书页码：67
\null
\BookNumber{67}
\Vertical{129}{11}{}{7}{10}第\\一\\幕\EndVertical
\BookLine{30.00}{22.14}{}{0}{这个习俗已经流传好几百年了，在伊俄尔科斯}
\BookLine{38.50}{22.14}{}{0}{也是这样吗？}
\BookLine{55.50}{0.00}{\heifont}{0}{阿尔刻提斯：}
\BookLine{64.00}{22.14}{}{0}{{\kaifont （低声）}是的，阿德墨托斯。}
\BookLine{81.00}{0.00}{\heifont}{0}{阿德墨托斯：}
\BookLine{89.50}{22.14}{}{0}{{\kaifont （活力满满地走过她身边，用手挡着自己的脸）}}
\BookLine{98.00}{22.14}{}{0}{很快——感谢神的恩赐——我就可以一直望着}
\BookLine{106.50}{22.14}{}{0}{你，直到我死去。}
\BookLine{123.50}{7.38}{\kaifont}{0}{走到王宫门口时，他忽然想起了什么。他仍然用手挡着}
\BookLine{132.00}{0.00}{\kaifont}{0}{自己的脸，缓缓地走到四个牧人在泉边所站立的地方。他下}
\BookLine{140.50}{0.00}{\kaifont}{0}{定决心般深深呼吸一口气，平静而直率地说道。}
\BookLine{157.50}{0.00}{\heifont}{0}{阿德墨托斯：}
\BookLine{166.00}{22.14}{}{0}{阿波罗，我父亲、我祖先、我土地的朋友，我}
\BookLine{174.50}{22.14}{}{0}{是一个虔诚的普通人，但我对真正的虔诚知晓}
\newpage
% 原书页码：68
\null
\BookNumber{68}
\Vertical{4}{169}{}{6.2}{8}阿\\尔\\刻\\提\\斯\\之\\歌\EndVertical
\BookLine{30.00}{22.14}{}{0}{不多。如果我不经意间有所冒犯，或是犯了错}
\BookLine{38.50}{22.14}{}{0}{误，希望这位“祂”，我的家族、我的人民的}
\BookLine{47.00}{22.14}{}{0}{朋友，能够原谅我。您在我人生最幸福的一}
\BookLine{55.50}{22.14}{}{0}{天到来，也请继续赐福于我和我的后裔吧……}
\BookLine{64.00}{22.14}{}{0}{{\kaifont （停顿）}我不太会说话，不过您能看透人心，请}
\BookLine{72.50}{22.14}{}{0}{鉴察我的心，或者，您可以在我脑中埋下愿望，}
\BookLine{81.00}{22.14}{}{0}{那些只有您能满足的愿望。}
\BookLine{98.00}{7.38}{\kaifont}{0}{他转过身迅速地走进宫殿。}
\BookLine{115.00}{0.00}{\heifont}{0}{阿尔刻提斯：}
\BookLine{123.50}{22.14}{}{0}{{\kaifont （一直在盯着牧人们看，带着期盼和渴望，又带}}
\BookLine{132.00}{22.14}{}{0}{{\kaifont 着怀疑和抗拒——尽管此刻他们已经从她身前}}
\BookLine{140.50}{22.14}{}{0}{{\kaifont 躲开了）}阿波罗在这儿吗？}
\BookLine{157.50}{0.00}{\heifont}{0}{阿格莱亚：}
\BookLine{166.00}{22.14}{}{0}{公主！}
\newpage
% 原书页码：69
\null
\BookNumber{69}
\Vertical{129}{11}{}{7}{10}第\\一\\幕\EndVertical
\BookLine{30.00}{0.00}{\heifont}{0}{阿尔刻提斯：}
\BookLine{38.50}{22.14}{}{0}{就是他们其中之一？那个老头就是德尔斐的}
\BookLine{47.00}{22.14}{}{0}{忒瑞西阿斯——那个颤颤巍巍、疯疯癫癫的}
\BookLine{55.50}{22.14}{}{0}{老头？}
\BookLine{72.50}{0.00}{\heifont}{0}{阿格莱亚：}
\BookLine{81.00}{22.14}{}{0}{{\kaifont （被惊讶到，坚定地）}不要怀疑这些事情，}
\BookLine{89.50}{22.14}{}{0}{公主。}
\BookLine{106.50}{0.00}{\heifont}{0}{阿尔刻提斯：}
\BookLine{115.00}{22.14}{}{0}{{\kaifont （朝着牧人们走了几步，突然做了决定）}让我和}
\BookLine{123.50}{22.14}{}{0}{他们单独待会儿。}
\BookLine{140.50}{7.38}{\kaifont}{0}{阿格莱亚对看门人做了个手势，两人下场。在接下来的}
\BookLine{149.00}{0.00}{\kaifont}{0}{谈话中，一个牧人睡着了，还在打鼾，其他几个则在她面前}
\BookLine{157.50}{0.00}{\kaifont}{0}{十分窘迫。他们互相递着酒囊，都不敢抬头看她。}
\newpage
% 原书页码：70
\null
\BookNumber{70}
\Vertical{4}{169}{}{6.2}{8}阿\\尔\\刻\\提\\斯\\之\\歌\EndVertical
\BookLine{30.00}{0.00}{\heifont}{0}{阿尔刻提斯：}
\BookLine{38.50}{22.14}{}{0}{您在这儿吗？我已经向您呼求成百上千次}
\BookLine{47.00}{22.14}{}{0}{了——向着天空、群星、太阳。我向德尔斐献}
\BookLine{55.50}{22.14}{}{0}{上了数不清的祷告。您现在是否，确确实实，}
\BookLine{64.00}{22.14}{}{0}{听到了我的声音？{\kaifont （静场，突然被鼾声打断）}有}
\BookLine{72.50}{22.14}{}{0}{人说，您不存在。有人说，神明都在很远的地}
\BookLine{81.00}{22.14}{}{0}{方，他们在奥林匹斯山上欢宴，饮酒，沉睡。}
\BookLine{89.50}{22.14}{}{0}{我已经把生命献给了您。您知道，我希望只为}
\BookLine{98.00}{22.14}{}{0}{您而活：去学习——由您来教导——生命的意}
\BookLine{106.50}{22.14}{}{0}{义。{\kaifont （没有回应）}人类是不是无法得到启示、教}
\BookLine{115.00}{22.14}{}{0}{训？我们是不是已经被抛弃了？{\kaifont （她等待着，场}}
\BookLine{123.50}{22.14}{}{0}{{\kaifont 上只有牧人们尴尬的沉默。她转身向王宫走去，}}
\BookLine{132.00}{22.14}{}{0}{{\kaifont 痛苦地自言自语）}那么我们只能自己寻找出路}
\BookLine{140.50}{22.14}{}{0}{了……生命不过是追寻这些那些过后的虚无，}
\BookLine{149.00}{22.14}{}{0}{是一种充满激情的空虚……}
\BookLine{166.00}{7.38}{\kaifont}{0}{第一个牧人——最脏的，四人中最不起眼的——站了起}
\BookLine{174.50}{0.00}{\kaifont}{0}{来。他谦卑而害羞地摸了摸自己的帽檐，说道。}
\newpage
% 原书页码：71
\null
\BookNumber{71}
\Vertical{129}{11}{}{7}{10}第\\一\\幕\EndVertical
\BookLine{30.00}{0.00}{\heifont}{0}{牧人：}
\BookLine{38.50}{22.14}{}{0}{公主，那个老头说我们之中有一个是神？他确}
\BookLine{47.00}{22.14}{}{0}{实这么说了？阿波罗？如果这样，夫人，我和}
\BookLine{55.50}{22.14}{}{0}{您一样惊讶。夫人，我们四人在过去三十天穿}
\BookLine{64.00}{22.14}{}{0}{越整个希腊，用同一个酒囊饮酒，用同一个碟}
\BookLine{72.50}{22.14}{}{0}{子吃饭，围着同一个火堆睡觉。如果我们之中}
\BookLine{81.00}{22.14}{}{0}{真有一位神明，我怎么会没发现呢？}
\BookLine{98.00}{7.38}{\kaifont}{0}{阿尔刻提斯带着希望和厌恶，朝他走了几步。}
\BookLine{115.00}{0.00}{\heifont}{0}{牧人：}
\BookLine{123.50}{22.14}{}{0}{照我看，夫人，我发誓我们只是普通的牧人罢}
\BookLine{132.00}{22.14}{}{0}{了，无知的牧人。但……但是有件事值得一}
\BookLine{140.50}{22.14}{}{0}{提，夫人，我们也不是平庸的牧人。怎么说}
\BookLine{149.00}{22.14}{}{0}{呢？那边那个家伙——在打鼾的那个家伙，没}
\BookLine{157.50}{22.14}{}{0}{有什么疾病是他治不好的，被蛇咬伤或是背部}
\BookLine{166.00}{22.14}{}{0}{骨折，不过我知道他不是神，公主。他旁边的}
\BookLine{174.50}{22.14}{}{0}{那个家伙，那个！{\kaifont （他走过去踢了那个牧人一}}
\newpage
% 原书页码：72
\null
\BookNumber{72}
\Vertical{4}{169}{}{6.2}{8}阿\\尔\\刻\\提\\斯\\之\\歌\EndVertical
\BookLine{30.00}{22.14}{}{0}{{\kaifont 脚）}你能不能别喝了？公主看着你呢。他从来}
\BookLine{38.50}{22.14}{}{0}{不会迷路，即使在最黑的夜里，他也能准确分}
\BookLine{47.00}{22.14}{}{0}{辨东西南北，很神奇。不过我知道他也不是太}
\BookLine{55.50}{22.14}{}{0}{阳之神。{\kaifont （小声说道）}另外，他有些很下流的习}
\BookLine{64.00}{22.14}{}{0}{惯，很下流。}
\BookLine{81.00}{0.00}{\heifont}{0}{阿尔刻提斯：}
\BookLine{89.50}{22.14}{}{0}{{\kaifont （毫不在意）}那个呢？}
\BookLine{106.50}{0.00}{\heifont}{0}{牧人：}
\BookLine{115.00}{22.14}{}{0}{那个？他是我们的歌手。}
\BookLine{132.00}{0.00}{\heifont}{0}{阿尔刻提斯：}
\BookLine{140.50}{22.14}{}{0}{啊！}
\BookLine{157.50}{0.00}{\heifont}{0}{牧人：}
\BookLine{166.00}{22.14}{}{0}{我跟您说，当他弹着竖琴唱起歌时——哦，公}
\BookLine{174.50}{22.14}{}{0}{主！每当那时我都对自己说“这是神啊”。即便}
\newpage
% 原书页码：73
\null
\BookNumber{73}
\Vertical{129}{11}{}{7}{10}第\\一\\幕\EndVertical
\BookLine{30.00}{22.14}{}{0}{没有什么值得快乐或忧伤的事情发生，他的音}
\BookLine{38.50}{22.14}{}{0}{乐也会让你感到满溢的欢喜和忧愁。他可以使}
\BookLine{47.00}{22.14}{}{0}{爱的记忆变得比爱本身更加甜蜜，但是公主，}
\BookLine{55.50}{22.14}{}{0}{他不是神。{\kaifont （她早就想反驳他，想要争辩一番）}}
\BookLine{64.00}{22.14}{}{0}{一个总是一副苦脸，随时要把自己喝死的人怎}
\BookLine{72.50}{22.14}{}{0}{么会是神呢？他在慢性自杀，您知道的，在我}
\BookLine{81.00}{22.14}{}{0}{们眼皮底下。神明们可不憎恨自己，公主。}
\BookLine{98.00}{0.00}{\heifont}{0}{阿尔刻提斯：}
\BookLine{106.50}{22.14}{}{0}{你呢？}
\BookLine{123.50}{0.00}{\heifont}{0}{牧人：}
\BookLine{132.00}{22.14}{}{0}{我？我是阿波罗？我不但不是阿波罗，我也不}
\BookLine{140.50}{22.14}{}{0}{相信阿波罗在这儿。}
\BookLine{157.50}{0.00}{\heifont}{0}{阿尔刻提斯：}
\BookLine{166.00}{22.14}{}{0}{忒瑞西阿斯，忒瑞西阿斯说了……}
\newpage
% 原书页码：74
\null
\BookNumber{74}
\Vertical{4}{169}{}{6.2}{8}阿\\尔\\刻\\提\\斯\\之\\歌\EndVertical
\BookLine{30.00}{0.00}{\heifont}{0}{牧人：}
\BookLine{38.50}{22.14}{}{0}{那个忒瑞西阿斯——那个路都走不稳还稀里糊}
\BookLine{47.00}{22.14}{}{0}{涂的老头？还能找到比他更差劲的送信人吗？}
\BookLine{55.50}{22.14}{}{0}{还能把神谕说得再不清不楚一点吗？}
\BookLine{72.50}{7.38}{\kaifont}{0}{阿尔刻提斯转身向王宫走去。她又转过身来看着牧人，}
\BookLine{81.00}{0.00}{\kaifont}{0}{像是在对自己说话。}
\BookLine{98.00}{0.00}{\heifont}{0}{阿尔刻提斯：}
\BookLine{106.50}{22.14}{}{0}{这样说来，我们真的很可悲。我们不仅得不到}
\BookLine{115.00}{22.14}{}{0}{帮助，还傻傻地相信我们可能会得到一些帮助。}
\BookLine{132.00}{0.00}{\heifont}{0}{牧人：}
\BookLine{140.50}{22.14}{}{0}{{\kaifont （向前走了几步到舞台上）}但如果神明们真的存}
\BookLine{149.00}{22.14}{}{0}{在，他们会如何对我们说话呢？他们会使用哪}
\BookLine{157.50}{22.14}{}{0}{种语言？和他们比起来，我们不过是病态的、}
\BookLine{166.00}{22.14}{}{0}{将死的、又聋又瞎还忙得团团转的小丑。为什}
\BookLine{174.50}{22.14}{}{0}{么有人会觉得神爱我们。您想得通吗？当爱人}
\newpage
% 原书页码：75
\null
\BookNumber{75}
\Vertical{129}{11}{}{7}{10}第\\一\\幕\EndVertical
\BookLine{30.00}{22.14}{}{0}{们之间隔着如此巨大的鸿沟，公主，这算哪门}
\BookLine{38.50}{22.14}{}{0}{子爱啊？{\kaifont （他慢慢走回自己的位置）}这种爱不会}
\BookLine{47.00}{22.14}{}{0}{快乐，简直不用怀疑。}
\BookLine{64.00}{0.00}{\heifont}{0}{阿尔刻提斯：}
\BookLine{72.50}{22.14}{}{0}{{\kaifont （恳切而尖锐地）}不，谁说不会快乐！}
\BookLine{89.50}{0.00}{\heifont}{0}{牧人：}
\BookLine{98.00}{22.14}{}{0}{{\kaifont （带着同样的情绪）}是。如果他真让我们看到他}
\BookLine{106.50}{22.14}{}{0}{的荣光，那会杀死我们的。{\kaifont （停顿）}对于今早发}
\BookLine{115.00}{22.14}{}{0}{生的事，我有些想法，或许有别的方法，去填}
\BookLine{123.50}{22.14}{}{0}{补这个鸿沟。或许神明们会找到一种方法，让}
\BookLine{132.00}{22.14}{}{0}{他们所爱之人变得更好，更接近他们。如果忒}
\BookLine{140.50}{22.14}{}{0}{瑞西阿斯说得没错，那阿波罗确实在塞萨利。}
\BookLine{149.00}{22.14}{}{0}{也或许，那个老糊涂捎错了信儿，或许信息本}
\BookLine{157.50}{22.14}{}{0}{来是，阿波罗在这里，分散在了每一个人身}
\BookLine{166.00}{22.14}{}{0}{上——我们四个牧人，还有其他人！比如说阿}
\BookLine{174.50}{22.14}{}{0}{德墨托斯，我才见过他几小时，我必须要说，}
\newpage
% 原书页码：76
\null
\BookNumber{76}
\Vertical{4}{169}{}{6.2}{8}阿\\尔\\刻\\提\\斯\\之\\歌\EndVertical
\BookLine{30.00}{22.14}{}{0}{公主，一开始，我对阿德墨托斯非常失望，他}
\BookLine{38.50}{22.14}{}{0}{看上去一点也没有什么不凡之处。您见过赫拉}
\BookLine{47.00}{22.14}{}{0}{克勒斯吧，公主？}
\BookLine{64.00}{0.00}{\heifont}{0}{阿尔刻提斯：}
\BookLine{72.50}{22.14}{}{0}{{\kaifont （轻轻点了点头）}见过。}
\BookLine{89.50}{0.00}{\heifont}{0}{牧人：}
\BookLine{98.00}{22.14}{}{0}{{\kaifont （突然回想起来）}啊，对。他追求过您。很有个}
\BookLine{106.50}{22.14}{}{0}{样儿！赫拉克勒斯，宙斯与阿尔克墨涅之子，您}
\BookLine{115.00}{22.14}{}{0}{一眼就能看出来。我见过不少比阿德墨托斯更强}
\BookLine{123.50}{22.14}{}{0}{的人，但渐渐地，我发现阿德墨托斯身上有一些}
\BookLine{132.00}{22.14}{}{0}{别的英雄没有的东西……这个世界在改变，慢}
\BookLine{140.50}{22.14}{}{0}{慢地改变。除非有不一样的男人、女人们出现，}
\BookLine{149.00}{22.14}{}{0}{公主，这个世界怎么会变得更好呢？}
\BookLine{166.00}{7.38}{\kaifont}{0}{阿德墨托斯上场，右肩上披着像阿波罗一样的浅蓝色的}
\BookLine{174.50}{0.00}{\kaifont}{0}{斗篷，他站在最上面一级台阶观望着。}
\newpage
% 原书页码：77
\null
\BookNumber{77}
\Vertical{129}{11}{}{7}{10}第\\一\\幕\EndVertical
\BookLine{30.00}{0.00}{\heifont}{0}{牧人：}
\BookLine{38.50}{22.14}{}{0}{这不就是，或许，就是那些伤感的爱人们，{\kaifont （指}}
\BookLine{47.00}{22.14}{}{0}{{\kaifont 指天空）}在想方设法，在我刚说的鸿沟上面搭}
\BookLine{55.50}{22.14}{}{0}{建的一座桥梁？{\kaifont （意识到阿德墨托斯的出现，谦}}
\BookLine{64.00}{22.14}{}{0}{{\kaifont 卑地低下头，用手触摸着前额，默念着）}我们}
\BookLine{72.50}{22.14}{}{0}{祝您幸福，长长久久。{\kaifont （回到自己的位置上）}}
\BookLine{89.50}{7.38}{\kaifont}{0}{太阳渐渐落山，宫殿背后的天空被染上了色彩。}
\BookLine{106.50}{0.00}{\heifont}{0}{阿德墨托斯：}
\BookLine{115.00}{22.14}{}{0}{{\kaifont （板着脸）}阿尔刻提斯公主，阿格莱亚已经告知}
\BookLine{123.50}{22.14}{}{0}{了我您离开的意愿。您，当然有离开的自由，}
\BookLine{132.00}{22.14}{}{0}{我们不会强留。塞萨利没有奴隶，公主，它的}
\BookLine{140.50}{22.14}{}{0}{王后也不会成为奴隶。我已经命令您的马夫和}
\BookLine{149.00}{22.14}{}{0}{女仆为您打点好旅途上需要的一切。{\kaifont （向下走一}}
\BookLine{157.50}{22.14}{}{0}{{\kaifont 步）}在您离开之前，我想说一件事，说这些不}
\BookLine{166.00}{22.14}{}{0}{是为了让您可怜我，也没有劝您回头的打算。}
\BookLine{174.50}{22.14}{}{0}{我说这些话是因为我们都是成年人了，对心中}
\newpage
% 原书页码：78
\null
\BookNumber{78}
\Vertical{4}{169}{}{6.2}{8}阿\\尔\\刻\\提\\斯\\之\\歌\EndVertical
\BookLine{30.00}{22.14}{}{0}{所想的事情不必掩藏。{\kaifont （轻微的停顿）}能够同时}
\BookLine{38.50}{22.14}{}{0}{驾驭狮子和野猪，这件事情直到现在都让我觉}
\BookLine{47.00}{22.14}{}{0}{得不可思议。但为什么能做到这一点，原因我}
\BookLine{55.50}{22.14}{}{0}{毫不怀疑：我爱您。我今后永远都见不到您了，}
\BookLine{64.00}{22.14}{}{0}{我永远不会再娶妻，我也永远不会是现在的我。}
\BookLine{72.50}{22.14}{}{0}{从今天开始，我不得不意识到，这个世界总是}
\BookLine{81.00}{22.14}{}{0}{有所残缺的。我是一个普通的人，但那充满在}
\BookLine{89.50}{22.14}{}{0}{我心中的爱，不是普通的爱。如果这样的爱都}
\BookLine{98.00}{22.14}{}{0}{得不到回应，那生命本身就成了一场谎言。人}
\BookLine{106.50}{22.14}{}{0}{们只需要随波逐流地生活着，所谓正义、荣誉、}
\BookLine{115.00}{22.14}{}{0}{爱，不过是短暂的、瞬间的一种妄想。愿您旅}
\BookLine{123.50}{22.14}{}{0}{途一切平安，阿尔刻提斯公主。}
\BookLine{140.50}{7.38}{\kaifont}{0}{阿尔刻提斯站在那里，慢慢抬起头，但仍不抬眼。她伸}
\BookLine{149.00}{0.00}{\kaifont}{0}{出了她的手。}
\BookLine{166.00}{0.00}{\heifont}{0}{阿尔刻提斯：}
\BookLine{174.50}{22.14}{}{0}{阿德墨托斯。{\kaifont （轻微停顿）}阿德墨托斯，再向我}
\newpage
% 原书页码：79
\null
\BookNumber{79}
\Vertical{129}{11}{}{7}{10}第\\一\\幕\EndVertical
\BookLine{30.00}{22.14}{}{0}{求一次婚吧。}
\BookLine{47.00}{7.38}{\kaifont}{0}{阿德墨托斯一个箭步向阿尔刻提斯走去，她再次伸出手，}
\BookLine{55.50}{0.00}{\kaifont}{0}{让他停下。}
\BookLine{72.50}{22.14}{}{0}{让我去爱你所爱，成为你的塞萨利的王后。让}
\BookLine{81.00}{22.14}{}{0}{我为你生儿育女，即使经历痛苦，让我在盛大}
\BookLine{89.50}{22.14}{}{0}{的宴会上挽着你走过，让我在你失意时安慰你，}
\BookLine{98.00}{22.14}{}{0}{让我在你远行归来时准备好沐浴的热水，让我}
\BookLine{106.50}{22.14}{}{0}{好好照料你的家，阿德墨托斯，像你所期待}
\BookLine{115.00}{22.14}{}{0}{的一样。让我为你、你的孩子、你的人民而生}
\BookLine{123.50}{22.14}{}{0}{活——让我每一刻都为你活着，每一刻都可以}
\BookLine{132.00}{22.14}{}{0}{为你死去。}
\BookLine{149.00}{0.00}{\heifont}{0}{阿德墨托斯：}
\BookLine{157.50}{22.14}{}{0}{{\kaifont （开心地，大声地）}不，不，阿尔刻提斯！应该}
\BookLine{166.00}{22.14}{}{0}{是我……}
\newpage
% 原书页码：80
\null
\BookNumber{80}
\Vertical{4}{169}{}{6.2}{8}阿\\尔\\刻\\提\\斯\\之\\歌\EndVertical
\BookLine{30.00}{7.38}{\kaifont}{0}{他拉起她伸出的手，他们并没有拥抱。}
\BookLine{47.00}{22.14}{}{0}{阿尔刻提斯，你是否愿意成为塞萨利国王——}
\BookLine{55.50}{22.14}{}{0}{阿德墨托斯的妻子？}
\BookLine{72.50}{0.00}{\heifont}{0}{阿尔刻提斯：}
\BookLine{81.00}{22.14}{}{0}{我愿交出我的所有，阿德墨托斯……}
\BookLine{98.00}{7.38}{\kaifont}{0}{他们走进宫殿。}
\newpage
% 原书页码：81
\null
\ActOrnament[二]
\newpage
% 原书页码：82
\null
\newpage
% 原书页码：83
\null
\BookNumber{83}
\Vertical{129}{11}{}{7}{10}第\\二\\幕\EndVertical
\BookLine{30.00}{7.38}{\kaifont}{0}{同一场景，十二年之后。仍然是黎明第一道光，仍然是}
\BookLine{38.50}{0.00}{\kaifont}{0}{同一个看门人绕过宫殿走出来。他并没有敲鼓，轻声而缓慢}
\BookLine{47.00}{0.00}{\kaifont}{0}{地说着话，情绪十分低落。}
\BookLine{64.00}{0.00}{\heifont}{0}{看门人：}
\BookLine{72.50}{22.14}{}{0}{黎明时分，照例巡查。热情诚挚的阿德墨托斯，}
\BookLine{81.00}{22.14}{}{0}{塞萨利之王，千金万马的拥有者。{\kaifont （低着头原地}}
\BookLine{89.50}{22.14}{}{0}{{\kaifont 站了一会儿，开始对观众说话）}事情要多糟糕}
\BookLine{98.00}{22.14}{}{0}{有多糟糕，非常糟糕。{\kaifont （门半掩着，他走到右}}
\BookLine{106.50}{22.14}{}{0}{{\kaifont 边门口，对外面的人们轻声说着）}没有，朋友}
\BookLine{115.00}{22.14}{}{0}{们，没什么新消息。国王还活着，他又撑过了}
\BookLine{123.50}{22.14}{}{0}{一晚，王后一直守着他，紧紧握着他的手。没}
\BookLine{132.00}{22.14}{}{0}{有新消息，没什么变化……把那条狗牵走，我}
\BookLine{140.50}{22.14}{}{0}{们这儿需要安静。对，国王喝了点掺着公牛血}
\BookLine{149.00}{22.14}{}{0}{的酒，{\kaifont （走回到场中央，对观众说道）}但是不得}
\BookLine{157.50}{22.14}{}{0}{不承认，国王已经处在死亡边缘。你们应该记}
\BookLine{166.00}{22.14}{}{0}{得那几个牧人，十二年前，忒瑞西阿斯带他们}
\BookLine{174.50}{22.14}{}{0}{到这儿来，说他们中间有一个人是阿波罗——}
\newpage
% 原书页码：84
\null
\BookNumber{84}
\Vertical{4}{169}{}{6.2}{8}阿\\尔\\刻\\提\\斯\\之\\歌\EndVertical
\BookLine{30.00}{22.14}{}{0}{阿波罗会作为仆人在这儿生活一年，你们还记}
\BookLine{38.50}{22.14}{}{0}{得吗？}
\BookLine{55.50}{22.14}{}{0}{嗯，一年结束之后他们走了吗？没有，他们四}
\BookLine{64.00}{22.14}{}{0}{个还待在这儿。阿波罗确实是这四人之一？没}
\BookLine{72.50}{22.14}{}{0}{人知道。别想这些事了，这种事想多了会变成}
\BookLine{81.00}{22.14}{}{0}{神经病的。不过朋友们，那一年确实很神奇。}
\BookLine{89.50}{22.14}{}{0}{我没法解释清楚：对于国家繁荣程度来说，那}
\BookLine{98.00}{22.14}{}{0}{一年和其他年份没太大区别。对于来访的，路}
\BookLine{106.50}{22.14}{}{0}{过的人来说，一切看起来都是老样子。但对我}
\BookLine{115.00}{22.14}{}{0}{们这些生活在其中的人来说，一切都不一样了。}
\BookLine{123.50}{22.14}{}{0}{人类生活的本质从没改变过，是我们看待生活}
\BookLine{132.00}{22.14}{}{0}{的眼光发生了改变。阿波罗确实就在这里。{\kaifont （停}}
\BookLine{140.50}{22.14}{}{0}{{\kaifont 顿）}嗯？阿德墨托斯国王怎么生病了？我不是}
\BookLine{149.00}{22.14}{}{0}{说了吗，那四个牧人留了下来。他们是好人，}
\BookLine{157.50}{22.14}{}{0}{好工人——虽然他们操着一口希腊南方的乡村}
\BookLine{166.00}{22.14}{}{0}{口音。其中还有两人——呃！嗜酒如命。有一}
\BookLine{174.50}{22.14}{}{0}{天晚上，太阳落山的时候他们就坐在那儿喝酒，}
\newpage
% 原书页码：85
\null
\BookNumber{85}
\Vertical{129}{11}{}{7}{10}第\\二\\幕\EndVertical
\BookLine{30.00}{22.14}{}{0}{那儿，就在那儿，一边喝一边开始吵架。他们}
\BookLine{38.50}{22.14}{}{0}{叫着嚷着，阿德墨托斯国王出来想要制止。结}
\BookLine{47.00}{22.14}{}{0}{果，天哪，其中一个牧人拿着刀，一不小心刺}
\BookLine{55.50}{22.14}{}{0}{到了国王，从这儿{\kaifont （指着喉咙）}一刀划到腰上。}
\BookLine{64.00}{22.14}{}{0}{伤口大得能把手伸进去。这已经是几星期前的}
\BookLine{72.50}{22.14}{}{0}{事儿了，但伤口还是没见好。伤口发炎然后结}
\BookLine{81.00}{22.14}{}{0}{痂，又流水又结痂的，国王看来是难逃此劫了。}
\BookLine{89.50}{22.14}{}{0}{没人打算怪罪那个牧人，因为他可能就是阿波}
\BookLine{98.00}{22.14}{}{0}{罗啊——你明白吗？这或许就是阿波罗的旨}
\BookLine{106.50}{22.14}{}{0}{意——想想看！现在那个牧人就在门外——他}
\BookLine{115.00}{22.14}{}{0}{一直在那儿——跪在地上，把脸埋在地上，希}
\BookLine{123.50}{22.14}{}{0}{望死的是他自己。}
\BookLine{140.50}{7.38}{\kaifont}{0}{门边有人在嘀咕，声音越来越大，听上去很激动。}
\BookLine{157.50}{0.00}{\heifont}{0}{看门人：}
\BookLine{166.00}{22.14}{}{0}{这又是怎么了？}
\newpage
% 原书页码：86
\null
\BookNumber{86}
\Vertical{4}{169}{}{6.2}{8}阿\\尔\\刻\\提\\斯\\之\\歌\EndVertical
\BookLine{30.00}{7.38}{\kaifont}{0}{阿格莱亚急匆匆地从宫殿里出来。}
\BookLine{47.00}{0.00}{\heifont}{0}{阿格莱亚：}
\BookLine{55.50}{22.14}{}{0}{出什么事了？看门人，这么吵闹你也不管管？}
\BookLine{72.50}{0.00}{\heifont}{0}{看门人：}
\BookLine{81.00}{22.14}{}{0}{{\kaifont （一瘸一拐地朝门口走去）}我也刚听见。}
\BookLine{98.00}{0.00}{\heifont}{0}{阿格莱亚：}
\BookLine{106.50}{22.14}{}{0}{{\kaifont （站在门边）}你们啰唆个不停在说什么？不知道}
\BookLine{115.00}{22.14}{}{0}{我们这儿需要安静吗？}
\BookLine{132.00}{0.00}{\heifont}{0}{看门人：}
\BookLine{140.50}{22.14}{}{0}{你们讲不讲理啊？你们有没有心肝啊？什么？}
\BookLine{149.00}{22.14}{}{0}{什么？}
\BookLine{166.00}{0.00}{\heifont}{0}{阿格莱亚：}
\BookLine{174.50}{22.14}{}{0}{什么？什么信使？告诉他去宫殿前门，一定要}
\newpage
% 原书页码：87
\null
\BookNumber{87}
\Vertical{129}{11}{}{7}{10}第\\二\\幕\EndVertical
\BookLine{30.00}{22.14}{}{0}{安静。什么信使？从哪儿来的？}
\BookLine{47.00}{7.38}{\kaifont}{0}{阿尔刻提斯上场，她正从宫殿里走出来。}
\BookLine{64.00}{0.00}{\heifont}{0}{阿尔刻提斯：}
\BookLine{72.50}{22.14}{}{0}{我们这儿要安静，要安静。阿格莱亚，这么吵}
\BookLine{81.00}{22.14}{}{0}{闹你也不管管吗？}
\BookLine{98.00}{0.00}{\heifont}{0}{阿格莱亚：}
\BookLine{106.50}{22.14}{}{0}{阿尔刻提斯王后，他们说来了一位信使。}
\BookLine{123.50}{0.00}{\heifont}{0}{阿尔刻提斯：}
\BookLine{132.00}{22.14}{}{0}{什么信使？}
\BookLine{149.00}{7.38}{\kaifont}{0}{有人从门缝递给他一小片长方形的金叶子，他转交给阿}
\BookLine{157.50}{0.00}{\kaifont}{0}{尔刻提斯。}
\newpage
% 原书页码：88
\null
\BookNumber{88}
\Vertical{4}{169}{}{6.2}{8}阿\\尔\\刻\\提\\斯\\之\\歌\EndVertical
\BookLine{30.00}{0.00}{\heifont}{0}{看门人：}
\BookLine{38.50}{22.14}{}{0}{他们说夜里来了一位信使，留下了这个。}
\BookLine{55.50}{7.38}{\kaifont}{0}{阿尔刻提斯拿在手里看了看，看门人也凑上去仔细看了}
\BookLine{64.00}{0.00}{\kaifont}{0}{看，严肃地说道。}
\BookLine{81.00}{0.00}{\heifont}{0}{看门人：}
\BookLine{89.50}{22.14}{}{0}{这是金的，夫人。您看，那些标志——是南方}
\BookLine{98.00}{22.14}{}{0}{的文字。那是太阳——祭坛，还有桂冠……阿}
\BookLine{106.50}{22.14}{}{0}{尔刻提斯王后，这是从德尔斐送来的，这是从}
\BookLine{115.00}{22.14}{}{0}{德尔斐送来的神谕！}
\BookLine{132.00}{7.38}{\kaifont}{0}{阿尔刻提斯立刻行起了仪式：她把叶子放到额前，然后}
\BookLine{140.50}{0.00}{\kaifont}{0}{放到胸前，然后放到唇边。}
\BookLine{157.50}{0.00}{\heifont}{0}{阿尔刻提斯：}
\BookLine{166.00}{22.14}{}{0}{信使在哪里？}
\newpage
% 原书页码：89
\null
\BookNumber{89}
\Vertical{129}{11}{}{7}{10}第\\二\\幕\EndVertical
\BookLine{30.00}{7.38}{\kaifont}{0}{看门人走到此刻微微开着的门边，小声与外面的人交谈}
\BookLine{38.50}{0.00}{\kaifont}{0}{起来。}
\BookLine{55.50}{0.00}{\heifont}{0}{看门人：}
\BookLine{64.00}{22.14}{}{0}{{\kaifont （转身对阿尔刻提斯）}他们说，他几个小时之前}
\BookLine{72.50}{22.14}{}{0}{走了。}
\BookLine{89.50}{0.00}{\heifont}{0}{阿尔刻提斯：}
\BookLine{98.00}{22.14}{}{0}{谁能看懂这些字？看门人，去把那四个牧人}
\BookLine{106.50}{22.14}{}{0}{叫来！}
\BookLine{123.50}{0.00}{\heifont}{0}{看门人：}
\BookLine{132.00}{22.14}{}{0}{他们就在这儿，夫人，四个人都在这儿了。}
\BookLine{149.00}{7.38}{\kaifont}{0}{门口又传来说话的声音。看门人开门让四个牧人进来。}
\BookLine{157.50}{0.00}{\kaifont}{0}{第一幕曾与阿尔刻提斯长谈过的那个牧人俯伏在地上，脸朝}
\BookLine{166.00}{0.00}{\kaifont}{0}{着地面，双手抱着自己的头。}
\newpage
% 原书页码：90
\null
\BookNumber{90}
\Vertical{4}{169}{}{6.2}{8}阿\\尔\\刻\\提\\斯\\之\\歌\EndVertical
\BookLine{30.00}{0.00}{\heifont}{0}{阿尔刻提斯：}
\BookLine{38.50}{22.14}{}{0}{你们有人认识南方文字吗？}
\BookLine{55.50}{0.00}{\heifont}{0}{牧人：}
\BookLine{64.00}{22.14}{}{0}{{\kaifont （抬起头）}我认识，一点，阿尔刻提斯王后。}
\BookLine{81.00}{0.00}{\heifont}{0}{阿尔刻提斯：}
\BookLine{89.50}{22.14}{}{0}{你先起来。站起来。}
\BookLine{106.50}{7.38}{\kaifont}{0}{他站了起来，她把叶子放在他手里，他看了一下，然后}
\BookLine{115.00}{0.00}{\kaifont}{0}{把叶子依次放在自己额前、胸前和唇边，敬畏地说道。}
\BookLine{132.00}{0.00}{\heifont}{0}{牧人：}
\BookLine{140.50}{22.14}{}{0}{这是从德尔斐的阿波罗神殿送来的信息。}
\BookLine{157.50}{7.38}{\kaifont}{0}{几个镇上的人从门口挤了进来，牧人艰难地辨认并读出}
\BookLine{166.00}{0.00}{\kaifont}{0}{文字。}
\newpage
% 原书页码：91
\null
\BookNumber{91}
\Vertical{129}{11}{}{7}{10}第\\二\\幕\EndVertical
\BookLine{30.00}{0.00}{\heifont}{0}{牧人：}
\BookLine{38.50}{22.14}{}{0}{平安……与长寿……归于热情诚挚的阿德墨托}
\BookLine{47.00}{22.14}{}{0}{斯，塞萨利之王，千金万马的拥有者。}
\BookLine{64.00}{0.00}{\heifont}{0}{阿尔刻提斯：}
\BookLine{72.50}{22.14}{}{0}{长寿？}
\BookLine{89.50}{0.00}{\heifont}{0}{牧人：}
\BookLine{98.00}{22.14}{}{0}{平安与长寿归于热情诚挚的阿德墨托斯。}
\BookLine{115.00}{0.00}{\heifont}{0}{阿尔刻提斯：}
\BookLine{123.50}{22.14}{}{0}{等等。{\kaifont （对看门人）}把门关起来。}
\BookLine{140.50}{7.38}{\kaifont}{0}{看门人把镇上的人赶下台去，他看了阿尔刻提斯一眼，}
\BookLine{149.00}{0.00}{\kaifont}{0}{把其他三个牧人也叫了下去，然后关上门。阿尔刻提斯走到}
\BookLine{157.50}{0.00}{\kaifont}{0}{牧人身旁，他正在努力辨认叶子上的字。}
\newpage
% 原书页码：92
\null
\BookNumber{92}
\Vertical{4}{169}{}{6.2}{8}阿\\尔\\刻\\提\\斯\\之\\歌\EndVertical
\BookLine{30.00}{0.00}{\heifont}{0}{牧人：}
\BookLine{38.50}{22.14}{}{0}{“阿德墨托斯国王，不，不会，死。”}
\BookLine{55.50}{0.00}{\heifont}{0}{阿格莱亚：}
\BookLine{64.00}{22.14}{}{0}{伟大的阿波罗！伟大的阿波罗！}
\BookLine{81.00}{0.00}{\heifont}{0}{牧人：}
\BookLine{89.50}{22.14}{}{0}{“不会死，如果，如果，因为，不，如果，另一}
\BookLine{98.00}{22.14}{}{0}{个，如果另一个人，第二个人，想要，{\kaifont （他指了}}
\BookLine{106.50}{22.14}{}{0}{{\kaifont 指自己的心）}希望，渴望，想要去死。”夫人，}
\BookLine{115.00}{22.14}{}{0}{我不认识最后一个字。}
\BookLine{132.00}{0.00}{\heifont}{0}{阿尔刻提斯：}
\BookLine{140.50}{22.14}{}{0}{我明白了。}
\BookLine{157.50}{0.00}{\heifont}{0}{牧人：}
\BookLine{166.00}{22.14}{}{0}{{\kaifont （费力地，终于想到）}“在他的位置，代替他。”}
\BookLine{174.50}{22.14}{}{0}{伟大的阿波罗！}
\newpage
% 原书页码：93
\null
\BookNumber{93}
\Vertical{129}{11}{}{7}{10}第\\二\\幕\EndVertical
\BookLine{30.00}{0.00}{\heifont}{0}{阿格莱亚：}
\BookLine{38.50}{22.14}{}{0}{这是什么意思？什么意思？}
\BookLine{55.50}{0.00}{\heifont}{0}{看门人：}
\BookLine{64.00}{22.14}{}{0}{我。我应该去死。{\kaifont （走开）}我的剑，我的剑呢？}
\BookLine{81.00}{0.00}{\heifont}{0}{阿格莱亚：}
\BookLine{89.50}{22.14}{}{1}{{\kaifont （明白过来）}不。{\kaifont （对看门人坚决地说道）}不，这}
\BookLine{98.00}{22.14}{}{0}{是给我的信息。我是看着他出生的，我也应该}
\BookLine{106.50}{22.14}{}{0}{为了他去死。我要去投河自尽。}
\BookLine{123.50}{0.00}{\heifont}{0}{看门人：}
\BookLine{132.00}{22.14}{}{0}{{\kaifont （回答阿格莱亚）}这不是女人该干的事情。}
\BookLine{149.00}{0.00}{\heifont}{0}{牧人：}
\BookLine{157.50}{22.14}{}{0}{{\kaifont （跪下，他用双拳捂住自己的眼睛）}我早就想要}
\BookLine{166.00}{22.14}{}{0}{去死了。}
\newpage
% 原书页码：94
\null
\BookNumber{94}
\Vertical{4}{169}{}{6.2}{8}阿\\尔\\刻\\提\\斯\\之\\歌\EndVertical
\BookLine{30.00}{0.00}{\heifont}{0}{阿格莱亚：}
\BookLine{38.50}{22.14}{}{0}{阿尔刻提斯王后，请吩咐吧，我一定办到。要}
\BookLine{47.00}{22.14}{}{0}{怎么为阿德墨托斯国王去死？}
\BookLine{64.00}{0.00}{\heifont}{0}{看门人：}
\BookLine{72.50}{22.14}{}{0}{{\kaifont （自告奋勇，站在阿格莱亚和阿尔刻提斯中间）}}
\BookLine{81.00}{22.14}{}{0}{阿尔刻提斯王后，下命令给我吧。很明显这个}
\BookLine{89.50}{22.14}{}{0}{信息就是给我的。}
\BookLine{106.50}{0.00}{\heifont}{0}{牧人：}
\BookLine{115.00}{22.14}{}{0}{{\kaifont （再次全身俯伏在地）}是我伤害了他，这是我应}
\BookLine{123.50}{22.14}{}{0}{得的！}
\BookLine{140.50}{0.00}{\heifont}{0}{阿尔刻提斯：}
\BookLine{149.00}{22.14}{}{0}{{\kaifont （一直一动不动，沉默着没有说话）}你们怎么能}
\BookLine{157.50}{22.14}{}{0}{替他去死呢？{\kaifont （停顿）}你们以为把自己扔进河里}
\BookLine{166.00}{22.14}{}{0}{就行了吗？一刀刺进心里就行了吗？不，神明}
\BookLine{174.50}{22.14}{}{0}{们不会让一切来得这么容易。}
\newpage
% 原书页码：95
\null
\BookNumber{95}
\Vertical{129}{11}{}{7}{10}第\\二\\幕\EndVertical
\BookLine{30.00}{0.00}{\heifont}{0}{阿格莱亚：}
\BookLine{38.50}{22.14}{}{0}{{\kaifont （啜泣着）}阿尔刻提斯王后，要怎么为阿德墨托}
\BookLine{47.00}{22.14}{}{0}{斯国王去死？}
\BookLine{64.00}{0.00}{\heifont}{0}{牧人：}
\BookLine{72.50}{22.14}{}{0}{我知道，我知道。{\kaifont （转身向门口走去）}平安与长}
\BookLine{81.00}{22.14}{}{0}{寿归于阿德墨托斯国王！}
\BookLine{98.00}{0.00}{\heifont}{0}{阿尔刻提斯：}
\BookLine{106.50}{22.14}{}{0}{{\kaifont （对牧人）}站住！}
\BookLine{123.50}{7.38}{\kaifont}{0}{大家都看着她，她平静地说道。}
\BookLine{140.50}{0.00}{\heifont}{0}{阿尔刻提斯：}
\BookLine{149.00}{22.14}{}{0}{阿格莱亚，去把我结婚时穿的那条裙子拿出来。}
\BookLine{166.00}{0.00}{\heifont}{0}{阿格莱亚：}
\BookLine{174.50}{22.14}{}{0}{{\kaifont （突然明白过来，盯着她，惊恐地小声说）}不，}
\newpage
% 原书页码：96
\null
\BookNumber{96}
\Vertical{4}{169}{}{6.2}{8}阿\\尔\\刻\\提\\斯\\之\\歌\EndVertical
\BookLine{30.00}{22.14}{}{0}{不，您不能……}
\BookLine{47.00}{0.00}{\heifont}{0}{阿尔刻提斯：}
\BookLine{55.50}{22.14}{}{0}{告诉孩子们我等下过来看他们。告诉国王外面}
\BookLine{64.00}{22.14}{}{0}{太阳很暖和，他应该出来晒晒太阳。}
\BookLine{81.00}{7.38}{\kaifont}{0}{大家彼此争论起来。}
\BookLine{98.00}{0.00}{\heifont}{0}{阿格莱亚：}
\BookLine{106.50}{22.14}{}{0}{不，不，阿尔刻提斯王后。我们已经老了，我}
\BookLine{115.00}{22.14}{}{0}{们活得差不多了。}
\BookLine{132.00}{0.00}{\heifont}{0}{看门人：}
\BookLine{140.50}{22.14}{}{0}{夫人，您看看我——我是个老人。}
\BookLine{157.50}{0.00}{\heifont}{0}{阿格莱亚：}
\BookLine{166.00}{22.14}{}{0}{您不能这么做，他也不会同意的。您是王后，}
\BookLine{174.50}{22.14}{}{0}{您是孩子们的母亲。}
\newpage
% 原书页码：97
\null
\BookNumber{97}
\Vertical{129}{11}{}{7}{10}第\\二\\幕\EndVertical
\BookLine{30.00}{0.00}{\heifont}{0}{阿尔刻提斯：}
\BookLine{38.50}{22.14}{}{0}{{\kaifont （同时用言语和手势制止他们的发言）}首先！首}
\BookLine{47.00}{22.14}{}{0}{先我命令你们，这件事情不许对任何人说起。}
\BookLine{55.50}{22.14}{}{0}{阿格莱亚，你听到了吗？}
\BookLine{72.50}{0.00}{\heifont}{0}{阿格莱亚：}
\BookLine{81.00}{22.14}{}{0}{是的，阿尔刻提斯王后。}
\BookLine{98.00}{0.00}{\heifont}{0}{阿尔刻提斯：}
\BookLine{106.50}{22.14}{}{0}{看门人！}
\BookLine{123.50}{0.00}{\heifont}{0}{看门人：}
\BookLine{132.00}{22.14}{}{0}{是的，阿尔刻提斯王后。}
\BookLine{149.00}{0.00}{\heifont}{0}{阿尔刻提斯：}
\BookLine{157.50}{22.14}{}{0}{直到明天，不能对任何人说。无论这里发生了}
\BookLine{166.00}{22.14}{}{0}{什么，你们都不许惊讶，不许悲伤。现在让我}
\BookLine{174.50}{22.14}{}{0}{和他单独待一会儿。{\kaifont （指着牧人）}}
\newpage
% 原书页码：98
\null
\BookNumber{98}
\Vertical{4}{169}{}{6.2}{8}阿\\尔\\刻\\提\\斯\\之\\歌\EndVertical
\BookLine{30.00}{7.38}{\kaifont}{0}{阿格莱亚走进宫殿，看门人下场。牧人站在那里，一只}
\BookLine{38.50}{0.00}{\kaifont}{0}{手放在额前沉思着，突然喊了起来。}
\BookLine{55.50}{0.00}{\heifont}{0}{牧人：}
\BookLine{64.00}{22.14}{}{0}{阿德墨托斯国王——起来吧！起来吧！}
\BookLine{81.00}{0.00}{\heifont}{0}{阿尔刻提斯：}
\BookLine{89.50}{22.14}{}{0}{你救不了他的。}
\BookLine{106.50}{0.00}{\heifont}{0}{牧人：}
\BookLine{115.00}{22.14}{}{0}{我早就想要去死了。}
\BookLine{132.00}{0.00}{\heifont}{0}{阿尔刻提斯：}
\BookLine{140.50}{22.14}{}{0}{是，或许你能做到，但你没法把这件事完全做}
\BookLine{149.00}{22.14}{}{0}{好。你希望去死，是的，但不是因为你爱阿德}
\BookLine{157.50}{22.14}{}{0}{墨托斯，只是为了让心中的罪恶感得到解脱。}
\BookLine{166.00}{22.14}{}{0}{这件事应该是为了爱而做，不是为了赎罪，是}
\BookLine{174.50}{22.14}{}{0}{为了爱。}
\newpage
% 原书页码：99
\null
\BookNumber{99}
\Vertical{129}{11}{}{7}{10}第\\二\\幕\EndVertical
\BookLine{30.00}{0.00}{\heifont}{0}{牧人：}
\BookLine{38.50}{22.14}{}{0}{{\kaifont （有些恼怒）}我，也爱着阿德墨托斯。}
\BookLine{55.50}{0.00}{\heifont}{0}{阿尔刻提斯：}
\BookLine{64.00}{22.14}{}{0}{有谁不爱他呢？但你的死太轻微了。你渴望去}
\BookLine{72.50}{22.14}{}{0}{死，而我，恐惧、害怕、憎恶死亡。我必须为}
\BookLine{81.00}{22.14}{}{0}{了阿德墨托斯而死，{\kaifont （望着天空）}在这日光下。}
\BookLine{89.50}{22.14}{}{0}{只有这样，他才能够重生。你能为他做到这}
\BookLine{98.00}{22.14}{}{0}{些吗？}
\BookLine{115.00}{7.38}{\kaifont}{0}{牧人沉默了，她继续自言自语般地说着。}
\BookLine{132.00}{22.14}{}{0}{我现在知道了要做什么，该怎么做，但我不明}
\BookLine{140.50}{22.14}{}{0}{白为什么……为什么是我。你……你能让我明}
\BookLine{149.00}{22.14}{}{0}{白，为什么我必须死去。}
\BookLine{166.00}{0.00}{\heifont}{0}{牧人：}
\BookLine{174.50}{22.14}{}{0}{我？}
\newpage
% 原书页码：100
\null
\BookNumber{100}
\Vertical{4}{169}{}{6.2}{8}阿\\尔\\刻\\提\\斯\\之\\歌\EndVertical
\BookLine{30.00}{0.00}{\heifont}{0}{阿尔刻提斯：}
\BookLine{38.50}{22.14}{}{0}{是的，你，你从德尔斐来，阿波罗就在你……}
\BookLine{55.50}{0.00}{\heifont}{0}{牧人：}
\BookLine{64.00}{22.14}{}{0}{{\kaifont （真诚地回绝）}夫人，公主，我告诉过您很多}
\BookLine{72.50}{22.14}{}{0}{很多次了。就算阿波罗在这儿，也绝对和我没}
\BookLine{81.00}{22.14}{}{0}{关系。}
\BookLine{98.00}{0.00}{\heifont}{0}{阿尔刻提斯：}
\BookLine{106.50}{22.14}{}{0}{{\kaifont （温和而轻柔地）}是你让我明白了，我的婚姻是}
\BookLine{115.00}{22.14}{}{0}{阿波罗的旨意。}
\BookLine{132.00}{0.00}{\heifont}{0}{牧人：}
\BookLine{140.50}{22.14}{}{0}{{\kaifont （有些恼怒）}不，不。您当时问了问题，我像其}
\BookLine{149.00}{22.14}{}{0}{他人一样回答您。您知道，我不过是个普普通}
\BookLine{157.50}{22.14}{}{0}{通的人。}
\BookLine{174.50}{7.38}{\kaifont}{0}{阿尔刻提斯转过身去，面对宫殿大门，沮丧至极，片刻}
\newpage
% 原书页码：101
\null
\BookNumber{101}
\Vertical{129}{11}{}{7}{10}第\\二\\幕\EndVertical
\BookLine{30.00}{0.00}{\kaifont}{0}{又转过身来对牧人说话。}
\BookLine{47.00}{0.00}{\heifont}{0}{阿尔刻提斯：}
\BookLine{55.50}{22.14}{}{0}{那就再次作为一个普普通通的人回答我的问题，}
\BookLine{64.00}{22.14}{}{0}{告诉我为什么阿波罗要让我死。}
\BookLine{81.00}{0.00}{\heifont}{0}{牧人：}
\BookLine{89.50}{22.14}{}{0}{{\kaifont （仍然带着一些恼怒）}那么作为一个普普通通的}
\BookLine{98.00}{22.14}{}{0}{人，我现在回答您——像其他人一样回答您：}
\BookLine{106.50}{22.14}{}{0}{德尔斐的旨意说，我们其中一人必须去死，我}
\BookLine{115.00}{22.14}{}{0}{也做好准备去死了，为什么一定要明白为什么}
\BookLine{123.50}{22.14}{}{0}{呢？}
\BookLine{140.50}{0.00}{\heifont}{0}{阿尔刻提斯：}
\BookLine{149.00}{22.14}{}{0}{如果不明白原因，那我们的生命和动物有什么}
\BookLine{157.50}{22.14}{}{0}{区别？}
\newpage
% 原书页码：102
\null
\BookNumber{102}
\Vertical{4}{169}{}{6.2}{8}阿\\尔\\刻\\提\\斯\\之\\歌\EndVertical
\BookLine{30.00}{0.00}{\heifont}{0}{牧人：}
\BookLine{38.50}{22.14}{}{0}{不！公主，明白，意味着您要看到整个事情。}
\BookLine{47.00}{22.14}{}{0}{我们人类有谁看到过整个故事呢，有谁看到过}
\BookLine{55.50}{22.14}{}{0}{故事的结局呢？如果您让我为阿德墨托斯去}
\BookLine{64.00}{22.14}{}{0}{死，我不知道等我死后会发生什么，但是我心}
\BookLine{72.50}{22.14}{}{0}{甘情愿。因为我一直明白这世上有两种死亡：}
\BookLine{81.00}{22.14}{}{0}{一种是终局；另一种是延续，因着随之而来的}
\BookLine{89.50}{22.14}{}{0}{一切而有了分量。我知道，如果这死亡的命运}
\BookLine{98.00}{22.14}{}{0}{是从德尔斐赐下给我的，那这死就不会是白白}
\BookLine{106.50}{22.14}{}{0}{的。如果众神真的存在，我想这就是他们的启}
\BookLine{115.00}{22.14}{}{0}{示：众神所做的一切还未曾揭示——而我们看}
\BookLine{123.50}{22.14}{}{0}{到的只是那广大之中极小的一部分。让我接受}
\BookLine{132.00}{22.14}{}{0}{这死亡吧，公主。因为这能够拯救我脱离另外}
\BookLine{140.50}{22.14}{}{0}{一种死亡，那正是我所恐惧的，也是所有人恐}
\BookLine{149.00}{22.14}{}{0}{惧的——死，不过是生命的停止，坟墓中的}
\BookLine{157.50}{22.14}{}{0}{尘埃。}
\BookLine{174.50}{7.38}{\kaifont}{0}{阿尔刻提斯在他的言辞中找到了所有答案。她的情绪有}
\newpage
% 原书页码：103
\null
\BookNumber{103}
\Vertical{129}{11}{}{7}{10}第\\二\\幕\EndVertical
\BookLine{30.00}{0.00}{\kaifont}{0}{了变化，轻声而迅速地说道。}
\BookLine{47.00}{0.00}{\heifont}{0}{阿尔刻提斯：}
\BookLine{55.50}{22.14}{}{0}{不，牧人，活下去——为了阿德墨托斯，为了}
\BookLine{64.00}{22.14}{}{0}{我，为了我的孩子们活下去。你是伊皮伊尼兹}
\BookLine{72.50}{22.14}{}{0}{的朋友呢，虽然因为你伤害了他父亲，他现在}
\BookLine{81.00}{22.14}{}{0}{恨你，你让他伤心坏了。他一直把你当作好朋}
\BookLine{89.50}{22.14}{}{0}{友，是你教会了他游泳、钓鱼。在离开之前，}
\BookLine{98.00}{22.14}{}{0}{我会告诉伊皮伊尼兹，你给了我多么大的恩惠。}
\BookLine{115.00}{7.38}{\kaifont}{0}{牧人默默走了出去，阿尔刻提斯拿金叶子又行了一次}
\BookLine{123.50}{0.00}{\kaifont}{0}{仪式：额前，胸前，唇边。阿格莱亚从宫殿里走了出来。}
\BookLine{140.50}{0.00}{\heifont}{0}{阿格莱亚：}
\BookLine{149.00}{22.14}{}{0}{{\kaifont （低着头）}裙子已经准备好了，阿尔刻提斯}
\BookLine{157.50}{22.14}{}{0}{王后。}
\newpage
% 原书页码：104
\null
\BookNumber{104}
\Vertical{4}{169}{}{6.2}{8}阿\\尔\\刻\\提\\斯\\之\\歌\EndVertical
\BookLine{30.00}{0.00}{\heifont}{0}{阿尔刻提斯：}
\BookLine{38.50}{22.14}{}{0}{阿格莱亚，我现在要做这件事，所以我不能，}
\BookLine{47.00}{22.14}{}{0}{绝不能去见孩子们。他们的小脑瓜闻起来太美}
\BookLine{55.50}{22.14}{}{0}{好了。你明白吗？}
\BookLine{72.50}{0.00}{\heifont}{0}{阿格莱亚：}
\BookLine{81.00}{22.14}{}{0}{明白，阿尔刻提斯王后。}
\BookLine{98.00}{0.00}{\heifont}{0}{阿尔刻提斯：}
\BookLine{106.50}{22.14}{}{0}{你告诉伊皮伊尼兹，就说我说的，让他去见那}
\BookLine{115.00}{22.14}{}{0}{个刺伤了他父亲的牧人，让他原谅他，并感谢}
\BookLine{123.50}{22.14}{}{0}{他给予我的恩惠。}
\BookLine{140.50}{0.00}{\heifont}{0}{阿格莱亚：}
\BookLine{149.00}{22.14}{}{0}{我会告诉他。}
\BookLine{166.00}{0.00}{\heifont}{0}{阿尔刻提斯：}
\BookLine{174.50}{22.14}{}{0}{你帮我剪下一束头发来，不要告诉任何人，把}
\newpage
% 原书页码：105
\null
\BookNumber{105}
\Vertical{129}{11}{}{7}{10}第\\二\\幕\EndVertical
\BookLine{30.00}{22.14}{}{0}{它放在祭坛上，{\kaifont （指着门）}就在路的那边，阿波}
\BookLine{38.50}{22.14}{}{0}{罗的丛林里。阿格莱亚，等我走了之后，告诉}
\BookLine{47.00}{22.14}{}{0}{阿德墨托斯国王，我希望他再娶一位妻子。}
\BookLine{64.00}{0.00}{\heifont}{0}{阿格莱亚：}
\BookLine{72.50}{22.14}{}{0}{王后啊！}
\BookLine{89.50}{0.00}{\heifont}{0}{阿尔刻提斯：}
\BookLine{98.00}{22.14}{}{0}{一个男人需要这样的安慰。但是，阿格莱亚！}
\BookLine{106.50}{22.14}{}{0}{人们都说继母会对前妻留下的孩子非常不好。}
\BookLine{115.00}{22.14}{}{0}{你要陪着他们，看好他们。阿格莱亚——要时}
\BookLine{123.50}{22.14}{}{0}{不时地，对他提起我。{\kaifont （快要说不下去了）}对他}
\BookLine{132.00}{22.14}{}{0}{提起我！}
\BookLine{149.00}{7.38}{\kaifont}{0}{阿尔刻提斯跑进宫殿里。阿格莱亚正要跟上去，听到门}
\BookLine{157.50}{0.00}{\kaifont}{0}{口有人在兴高采烈地讨论着什么。一些镇上的人走进了院子，}
\BookLine{166.00}{0.00}{\kaifont}{0}{罗德普，宫里的一个年轻侍女，一边掩饰着欢笑，一边从门}
\BookLine{174.50}{0.00}{\kaifont}{0}{口跑向宫殿大门。}
\newpage
% 原书页码：106
\null
\BookNumber{106}
\Vertical{4}{169}{}{6.2}{8}阿\\尔\\刻\\提\\斯\\之\\歌\EndVertical
\BookLine{30.00}{0.00}{\heifont}{0}{阿格莱亚：}
\BookLine{38.50}{22.14}{}{0}{有什么事吗？外面吵吵闹闹是怎么回事，罗}
\BookLine{47.00}{22.14}{}{0}{德普！}
\BookLine{64.00}{0.00}{\heifont}{0}{罗德普：}
\BookLine{72.50}{22.14}{}{0}{他马上就要到了！}
\BookLine{89.50}{0.00}{\heifont}{0}{阿格莱亚：}
\BookLine{98.00}{22.14}{}{0}{谁？}
\BookLine{115.00}{0.00}{\heifont}{0}{罗德普和镇上的人们：}
\BookLine{123.50}{22.14}{}{0}{赫拉克勒斯来了！他已经从山谷下来了！宙斯}
\BookLine{132.00}{22.14}{}{0}{之子，赫拉克勒斯！}
\BookLine{149.00}{0.00}{\heifont}{0}{阿格莱亚：}
\BookLine{157.50}{22.14}{}{0}{你们都闭嘴。快从院子里出去。你们难道忘了}
\BookLine{166.00}{22.14}{}{0}{这个家里有人生病吗？罗德普，快给我进去，}
\BookLine{174.50}{22.14}{}{0}{闭上你的嘴！}
\newpage
% 原书页码：107
\null
\BookNumber{107}
\Vertical{129}{11}{}{7}{10}第\\二\\幕\EndVertical
\BookLine{30.00}{7.38}{\kaifont}{0}{镇上的人们下场。看门人打开宫殿的大门，罗德普刚好}
\BookLine{38.50}{0.00}{\kaifont}{0}{溜了进去。看门人身后跟着两个仆人，被他们抬着的卧榻挡}
\BookLine{47.00}{0.00}{\kaifont}{0}{着。他对阿格莱亚说话，而她正焦虑地低着头。}
\BookLine{64.00}{0.00}{\heifont}{0}{看门人：}
\BookLine{72.50}{22.14}{}{0}{赫拉克勒斯！大日子！}
\BookLine{89.50}{0.00}{\heifont}{0}{阿格莱亚：}
\BookLine{98.00}{22.14}{}{0}{{\kaifont （讽刺着）}是啊，是啊，每年都来。他为什么每}
\BookLine{106.50}{22.14}{}{0}{年都要来，为什么？}
\BookLine{123.50}{0.00}{\heifont}{0}{看门人：}
\BookLine{132.00}{22.14}{}{0}{这有什么，多尊贵的朋友！}
\BookLine{149.00}{0.00}{\heifont}{0}{阿格莱亚：}
\BookLine{157.50}{22.14}{}{0}{是啊，尊贵的朋友有时候也该为自己的缺席补}
\BookLine{166.00}{22.14}{}{0}{上点礼物。}
\newpage
% 原书页码：108
\null
\BookNumber{108}
\Vertical{4}{169}{}{6.2}{8}阿\\尔\\刻\\提\\斯\\之\\歌\EndVertical
\BookLine{30.00}{0.00}{\heifont}{0}{看门人：}
\BookLine{38.50}{22.14}{}{0}{阿格莱亚，你这是怎么回事！}
\BookLine{55.50}{0.00}{\heifont}{0}{阿格莱亚：}
\BookLine{64.00}{22.14}{}{0}{{\kaifont （暴躁地）}我知道怎么回事。}
\BookLine{81.00}{7.38}{\kaifont}{0}{两个仆人从宫殿里出来，抬着一张卧榻和一个脚凳。}
\BookLine{98.00}{0.00}{\heifont}{0}{看门人：}
\BookLine{106.50}{22.14}{}{0}{{\kaifont （对仆人）}过来，放这儿，你们知道的。}
\BookLine{123.50}{7.38}{\kaifont}{0}{看门人转过身去扶阿德墨托斯，阿格莱亚指挥了起来。}
\BookLine{140.50}{0.00}{\heifont}{0}{阿格莱亚：}
\BookLine{149.00}{22.14}{}{0}{这儿，放这儿！{\kaifont （照着太阳找到位置）}转过来。}
\BookLine{157.50}{22.14}{}{0}{{\kaifont （从他们手里拿过垫子）}这样，对了。王后喜欢}
\BookLine{166.00}{22.14}{}{0}{坐在这儿。}
\newpage
% 原书页码：109
\null
\BookNumber{109}
\Vertical{129}{11}{}{7}{10}第\\二\\幕\EndVertical
\BookLine{30.00}{7.38}{\kaifont}{0}{阿德墨托斯由人搀扶着走了出来，双手分别搭在看门人}
\BookLine{38.50}{0.00}{\kaifont}{0}{和一个守卫的肩上。}
\BookLine{55.50}{0.00}{\heifont}{0}{阿德墨托斯：}
\BookLine{64.00}{22.14}{}{0}{已经快要正午了。{\kaifont （对看门人说）}你已经守了一}
\BookLine{72.50}{22.14}{}{0}{夜，去休息吧。}
\BookLine{89.50}{0.00}{\heifont}{0}{看门人：}
\BookLine{98.00}{22.14}{}{0}{啊，陛下，像我们这些老年人，睡得都很少。}
\BookLine{106.50}{22.14}{}{0}{轻点，好一点。}
\BookLine{123.50}{0.00}{\heifont}{0}{阿德墨托斯：}
\BookLine{132.00}{22.14}{}{0}{我没看到王后呢，阿格莱亚。}
\BookLine{149.00}{0.00}{\heifont}{0}{阿格莱亚：}
\BookLine{157.50}{22.14}{}{0}{{\kaifont （忙着整理垫子）}陛下，她马上就过来。您放}
\BookLine{166.00}{22.14}{}{0}{心吧。}
\newpage
% 原书页码：110
\null
\BookNumber{110}
\Vertical{4}{169}{}{6.2}{8}阿\\尔\\刻\\提\\斯\\之\\歌\EndVertical
\BookLine{30.00}{7.38}{\kaifont}{0}{阿格莱亚匆忙地走回宫中。}
\BookLine{47.00}{0.00}{\heifont}{0}{看门人：}
\BookLine{55.50}{22.14}{}{0}{陛下，等您舒服地坐好了，我有些消息要告诉}
\BookLine{64.00}{22.14}{}{0}{您。估计您都不敢相信。{\kaifont （斜起眼看看太阳）}是}
\BookLine{72.50}{22.14}{}{0}{的，陛下，今天会非常热。我们已经进入夏至}
\BookLine{81.00}{22.14}{}{0}{时节了。据海员们说，太阳现在很低，在这些}
\BookLine{89.50}{22.14}{}{0}{日子，那个祂离我们非常近。他们就是这么}
\BookLine{98.00}{22.14}{}{0}{说的。}
\BookLine{115.00}{0.00}{\heifont}{0}{阿德墨托斯：}
\BookLine{123.50}{22.14}{}{0}{所以是什么消息？}
\BookLine{140.50}{0.00}{\heifont}{0}{看门人：}
\BookLine{149.00}{22.14}{}{0}{一位尊贵的朋友，陛下，要来见您。一位，好}
\BookLine{157.50}{22.14}{}{0}{朋友。}
\newpage
% 原书页码：111
\null
\BookNumber{111}
\Vertical{129}{11}{}{7}{10}第\\二\\幕\EndVertical
\BookLine{30.00}{0.00}{\heifont}{0}{阿德墨托斯：}
\BookLine{38.50}{22.14}{}{0}{赫拉克勒斯！}
\BookLine{55.50}{0.00}{\heifont}{0}{看门人：}
\BookLine{64.00}{22.14}{}{0}{是的，宙斯之子赫拉克勒斯！}
\BookLine{81.00}{0.00}{\heifont}{0}{阿德墨托斯：}
\BookLine{89.50}{22.14}{}{0}{今天！今天真是个好日子！}
\BookLine{106.50}{0.00}{\heifont}{0}{看门人：}
\BookLine{115.00}{22.14}{}{0}{整个镇上都传开了，他走得也真是够慢的了。}
\BookLine{123.50}{22.14}{}{0}{每到一个村庄，或者农场，人们都献上好酒，}
\BookLine{132.00}{22.14}{}{0}{给他戴上花环。他呢，喝得醉醺醺的，哦，陛}
\BookLine{140.50}{22.14}{}{0}{下，今天可不能让他拥抱您啊。您知道他那力}
\BookLine{149.00}{22.14}{}{0}{气，还不得把您害死。}
\BookLine{166.00}{0.00}{\heifont}{0}{阿德墨托斯：}
\BookLine{174.50}{22.14}{}{0}{王后知道吗？}
\newpage
% 原书页码：112
\null
\BookNumber{112}
\Vertical{4}{169}{}{6.2}{8}阿\\尔\\刻\\提\\斯\\之\\歌\EndVertical
\BookLine{30.00}{0.00}{\heifont}{0}{看门人：}
\BookLine{38.50}{22.14}{}{0}{女仆们应该已经告诉她了。}
\BookLine{55.50}{0.00}{\heifont}{0}{阿德墨托斯：}
\BookLine{64.00}{22.14}{}{0}{他要是换个时间来就好了。}
\BookLine{81.00}{0.00}{\heifont}{0}{看门人：}
\BookLine{89.50}{22.14}{}{0}{赫拉克勒斯想什么时候来就什么时候来。陛下，}
\BookLine{98.00}{22.14}{}{0}{小王子们都激动坏了，伊皮伊尼兹肯定又会让}
\BookLine{106.50}{22.14}{}{0}{他讲一遍杀死复仇女神的狮子和打扫奥革阿斯}
\BookLine{115.00}{22.14}{}{0}{牛棚的故事。}
\BookLine{132.00}{0.00}{\heifont}{0}{阿德墨托斯：}
\BookLine{140.50}{22.14}{}{0}{去把阿格莱亚叫来。}
\BookLine{157.50}{7.38}{\kaifont}{0}{阿格莱亚从宫殿里走出来。}
\newpage
% 原书页码：113
\null
\BookNumber{113}
\Vertical{129}{11}{}{7}{10}第\\二\\幕\EndVertical
\BookLine{30.00}{0.00}{\heifont}{0}{看门人：}
\BookLine{38.50}{22.14}{}{0}{她来了。}
\BookLine{55.50}{0.00}{\heifont}{0}{阿格莱亚：}
\BookLine{64.00}{22.14}{}{0}{陛下，您有什么吩咐？}
\BookLine{81.00}{7.38}{\kaifont}{0}{阿格莱亚背着国王打手势让看门人先离开。看门人从左}
\BookLine{89.50}{0.00}{\kaifont}{0}{侧下场。}
\BookLine{106.50}{0.00}{\heifont}{0}{阿德墨托斯：}
\BookLine{115.00}{22.14}{}{0}{阿格莱亚，来。你听说了吗？}
\BookLine{132.00}{7.38}{\kaifont}{0}{看门人刚要下场，很紧张地看了阿格莱亚一眼，听到她}
\BookLine{140.50}{0.00}{\kaifont}{0}{的回答后才算放心。}
\BookLine{157.50}{0.00}{\heifont}{0}{阿格莱亚：}
\BookLine{166.00}{22.14}{}{0}{听说了，宙斯之子，要来看望我们。}
\newpage
% 原书页码：114
\null
\BookNumber{114}
\Vertical{4}{169}{}{6.2}{8}阿\\尔\\刻\\提\\斯\\之\\歌\EndVertical
\BookLine{30.00}{0.00}{\heifont}{0}{阿德墨托斯：}
\BookLine{38.50}{22.14}{}{0}{{\kaifont （淡然地）}阿格莱亚，这是我最后的日子了。我}
\BookLine{47.00}{22.14}{}{0}{知道，我感觉到了。}
\BookLine{64.00}{0.00}{\heifont}{0}{阿格莱亚：}
\BookLine{72.50}{22.14}{}{0}{陛下，这种事情没人会知道的。您还是少说几}
\BookLine{81.00}{22.14}{}{0}{句吧。您需要保存体力。}
\BookLine{98.00}{0.00}{\heifont}{0}{阿德墨托斯：}
\BookLine{106.50}{22.14}{}{0}{太阳落山之前我就会死去。如果可行的话，我}
\BookLine{115.00}{22.14}{}{0}{希望你能办到，拖上个一两天，不要让赫拉克}
\BookLine{123.50}{22.14}{}{0}{勒斯得知我死亡的消息。你就说我们有个仆人}
\BookLine{132.00}{22.14}{}{0}{死了，就像什么都没发生过一样照常生活。王}
\BookLine{140.50}{22.14}{}{0}{后会告诉你该怎么做。你知道赫拉克勒斯一直}
\BookLine{149.00}{22.14}{}{0}{以来的生活，总是担着沉重的工作，一个接一}
\BookLine{157.50}{22.14}{}{0}{个。他高高兴兴地来，我不想让他看到一片哀}
\BookLine{166.00}{22.14}{}{0}{伤。尤其是赫拉克勒斯，总是出生入死，他很}
\BookLine{174.50}{22.14}{}{0}{讨厌葬礼啊哀哭啊这些事情。}
\newpage
% 原书页码：115
\null
\BookNumber{115}
\Vertical{129}{11}{}{7}{10}第\\二\\幕\EndVertical
\BookLine{30.00}{0.00}{\heifont}{0}{阿格莱亚：}
\BookLine{38.50}{22.14}{}{0}{{\kaifont （又整理起垫子）}陛下，我不傻，您说的这些我}
\BookLine{47.00}{22.14}{}{0}{都明白。您好好躺下，闭上眼睛休息，阳光会}
\BookLine{55.50}{22.14}{}{0}{带给您力量的。}
\BookLine{72.50}{0.00}{\heifont}{0}{阿德墨托斯：}
\BookLine{81.00}{22.14}{}{0}{再给我一天的力气就好，好来欢迎赫拉克勒斯。}
\BookLine{89.50}{22.14}{}{0}{你去看看王后在哪里呢。}
\BookLine{106.50}{0.00}{\heifont}{0}{阿格莱亚：}
\BookLine{115.00}{22.14}{}{0}{她来了。}
\BookLine{132.00}{7.38}{\kaifont}{0}{阿格莱亚从左侧下场。阿尔刻提斯从宫殿里走出来。她}
\BookLine{140.50}{0.00}{\kaifont}{0}{站在最高一级的台阶上，穿着第一幕的那条裙子。阿德墨托}
\BookLine{149.00}{0.00}{\kaifont}{0}{斯回头看到她，没有说话。}
\BookLine{166.00}{0.00}{\heifont}{0}{阿德墨托斯：}
\BookLine{174.50}{22.14}{}{0}{这是咱们结婚时你穿的衣服！}
\newpage
% 原书页码：116
\null
\BookNumber{116}
\Vertical{4}{169}{}{6.2}{8}阿\\尔\\刻\\提\\斯\\之\\歌\EndVertical
\BookLine{30.00}{7.38}{\kaifont}{0}{阿尔刻提斯微笑着，手指放在唇边示意他不要说话，并}
\BookLine{38.50}{0.00}{\kaifont}{0}{走到他身边，轻声说着。}
\BookLine{55.50}{0.00}{\heifont}{0}{阿尔刻提斯：}
\BookLine{64.00}{22.14}{}{0}{妻子们都留着它等着葬礼的时候穿，我穿上它}
\BookLine{72.50}{22.14}{}{0}{是为了好好活着。}
\BookLine{89.50}{0.00}{\heifont}{0}{阿德墨托斯：}
\BookLine{98.00}{22.14}{}{0}{亲爱的，坐这儿来。我刚还跟阿格莱亚说……}
\BookLine{115.00}{7.38}{\kaifont}{0}{阿尔刻提斯坐在阿德墨托斯身边的凳子上，把手放在他}
\BookLine{123.50}{0.00}{\kaifont}{0}{的手边。}
\BookLine{140.50}{0.00}{\heifont}{0}{阿尔刻提斯：}
\BookLine{149.00}{22.14}{}{0}{是的，我知道，等赫拉克勒斯来了……}
\BookLine{166.00}{0.00}{\heifont}{0}{阿德墨托斯：}
\BookLine{174.50}{22.14}{}{0}{他马上就到了。}
\newpage
% 原书页码：117
\null
\BookNumber{117}
\Vertical{129}{11}{}{7}{10}第\\二\\幕\EndVertical
\BookLine{30.00}{0.00}{\heifont}{0}{阿尔刻提斯：}
\BookLine{38.50}{22.14}{}{0}{一定是在路上耽搁了，只有危险的时候他才跑}
\BookLine{47.00}{22.14}{}{0}{得飞快，说不定他要今晚才能到，或者要等明}
\BookLine{55.50}{22.14}{}{0}{天了。不过我们不着急。{\kaifont （起身走到舞台中央，}}
\BookLine{64.00}{22.14}{}{0}{{\kaifont 站在那里倾听着）}我听到山谷里传来的欢呼声}
\BookLine{72.50}{22.14}{}{0}{了。{\kaifont （走回到凳子旁边，重复着）}我们不着急。}
\BookLine{81.00}{22.14}{}{0}{{\kaifont （站在舞台中央，望着太阳仿佛要寻求帮助）}这}
\BookLine{89.50}{22.14}{}{0}{样的阳光能够治愈人，夏至日的阳光，你感觉}
\BookLine{98.00}{22.14}{}{0}{到了吗？}
\BookLine{115.00}{0.00}{\heifont}{0}{阿德墨托斯：}
\BookLine{123.50}{22.14}{}{0}{这阳光能治愈人，但只能治愈别人。我已经不}
\BookLine{132.00}{22.14}{}{0}{需要这些了，不再需要什么希望。我的一生很}
\BookLine{140.50}{22.14}{}{0}{短暂，但短短一个钟头也足以成为完美的时光。}
\BookLine{149.00}{22.14}{}{0}{我得到过了，最完美的时光。一个拥有过幸福}
\BookLine{157.50}{22.14}{}{0}{的人，就不会成为时间的囚徒。来，在最后这}
\BookLine{166.00}{22.14}{}{0}{一天，让我们把还没说完的话都好好说完。}
\newpage
% 原书页码：118
\null
\BookNumber{118}
\Vertical{4}{169}{}{6.2}{8}阿\\尔\\刻\\提\\斯\\之\\歌\EndVertical
\BookLine{30.00}{0.00}{\heifont}{0}{阿尔刻提斯：}
\BookLine{38.50}{22.14}{}{0}{{\kaifont （走回他身边，温柔地微笑着）}最后这一天？你}
\BookLine{47.00}{22.14}{}{0}{还是好好休息，留着点力气吧。听我说就好了。}
\BookLine{55.50}{22.14}{}{0}{我来把我们两个的话都说了。我不知道是谁给}
\BookLine{64.00}{22.14}{}{0}{我起过一个别名：沉默的阿尔刻提斯。我想应}
\BookLine{72.50}{22.14}{}{0}{该是你，不过是赫拉克勒斯告诉了全希腊的人。}
\BookLine{89.50}{0.00}{\heifont}{0}{阿德墨托斯：}
\BookLine{98.00}{22.14}{}{0}{你从前一直是沉默不语的吗？}
\BookLine{115.00}{0.00}{\heifont}{0}{阿尔刻提斯：}
\BookLine{123.50}{22.14}{}{0}{我？不是的，我还是个小女孩的时候，哈，我}
\BookLine{132.00}{22.14}{}{0}{可是个很喜欢争论和反驳别人的姑娘，这你是}
\BookLine{140.50}{22.14}{}{0}{知道的。有的时候我也会受够了自己的沉默，}
\BookLine{149.00}{22.14}{}{0}{无聊而沉默的阿尔刻提斯！有些时候，我也想}
\BookLine{157.50}{22.14}{}{0}{像其他的女人那样，心里想到什么就随口说什}
\BookLine{166.00}{22.14}{}{0}{么。其他的妻子们通常会怎么说？我想她们大}
\BookLine{174.50}{22.14}{}{0}{概都会这样说：阿德墨托斯，我有没有告诉过}
\newpage
% 原书页码：119
\null
\BookNumber{119}
\Vertical{129}{11}{}{7}{10}第\\二\\幕\EndVertical
\BookLine{30.00}{22.14}{}{0}{你——让我再好好看看你的眼睛——我有没有}
\BookLine{38.50}{22.14}{}{0}{告诉过你，我爱你胜过我的生命？}
\BookLine{55.50}{7.38}{\kaifont}{0}{她正向他靠过去，一阵疼痛突然击中了她。她站了起来，}
\BookLine{64.00}{0.00}{\kaifont}{0}{有些晕眩，用右手扶着自己的头，左手则抚着阿德墨托斯的}
\BookLine{72.50}{0.00}{\kaifont}{0}{伤口。阿德墨托斯并没有注意到，他闭上眼睛，一边笑一边}
\BookLine{81.00}{0.00}{\kaifont}{0}{小声念叨着。}
\BookLine{98.00}{0.00}{\heifont}{0}{阿德墨托斯：}
\BookLine{106.50}{22.14}{}{0}{这话听着可不像，不，听着不像我的沉默的阿}
\BookLine{115.00}{22.14}{}{0}{尔刻提斯啊。现在轮到替我说话了，该说什}
\BookLine{123.50}{22.14}{}{0}{么呢？}
\BookLine{140.50}{0.00}{\heifont}{0}{阿尔刻提斯：}
\BookLine{149.00}{22.14}{}{0}{{\kaifont （疼痛减轻了一些，她静静地说）}我会说，我，}
\BookLine{157.50}{22.14}{}{0}{阿德墨托斯，无法选择自己的出生和死亡。生}
\BookLine{166.00}{22.14}{}{0}{命是神明所赐的礼物，一份神奇的礼物，而赐}
\BookLine{174.50}{22.14}{}{0}{下生命的神明，也会赐予阿尔刻提斯，或是我}
\newpage
% 原书页码：120
\null
\BookNumber{120}
\Vertical{4}{169}{}{6.2}{8}阿\\尔\\刻\\提\\斯\\之\\歌\EndVertical
\BookLine{30.00}{22.14}{}{0}{死亡。}
\BookLine{47.00}{0.00}{\heifont}{0}{阿德墨托斯：}
\BookLine{55.50}{22.14}{}{0}{{\kaifont （有点惊奇）}你的死亡，阿尔刻提斯？}
\BookLine{72.50}{0.00}{\heifont}{0}{阿尔刻提斯：}
\BookLine{81.00}{22.14}{}{0}{难道不是同一种力量吗，赐予生命的也……你}
\BookLine{89.50}{22.14}{}{0}{哪里不舒服吗，阿德墨托斯？}
\BookLine{106.50}{0.00}{\heifont}{0}{阿德墨托斯：}
\BookLine{115.00}{22.14}{}{0}{{\kaifont （很惊讶地把手放到一边）}没有，我不知道为什}
\BookLine{123.50}{22.14}{}{0}{么，不过我的疼痛好像……}
\BookLine{140.50}{0.00}{\heifont}{0}{阿尔刻提斯：}
\BookLine{149.00}{22.14}{}{0}{举起手试试！}
\BookLine{166.00}{7.38}{\kaifont}{0}{他缓缓地举起那只不能动弹的左手，痛苦转移到了阿尔}
\BookLine{174.50}{0.00}{\kaifont}{0}{刻提斯的身上，她抓着衣服，痛苦地弯下了腰。}
\newpage
% 原书页码：121
\null
\BookNumber{121}
\Vertical{129}{11}{}{7}{10}第\\二\\幕\EndVertical
\BookLine{30.00}{0.00}{\heifont}{0}{阿德墨托斯：}
\BookLine{38.50}{22.14}{}{0}{这光芒！{\kaifont （他盯着自己的前方，有了希望）}不，}
\BookLine{47.00}{22.14}{}{0}{我不该想这些事情了。我已经不需要这些了，}
\BookLine{55.50}{22.14}{}{0}{我们已经品尝了生命的富足，我们已经不是时}
\BookLine{64.00}{22.14}{}{0}{间的囚徒了。}
\BookLine{81.00}{7.38}{\kaifont}{0}{剧烈的疼痛过去了，阿尔刻提斯小心翼翼地动弹一下，}
\BookLine{89.50}{0.00}{\kaifont}{0}{又坐了下来，脸贴脸地对阿德墨托斯说话。}
\BookLine{106.50}{0.00}{\heifont}{0}{阿尔刻提斯：}
\BookLine{115.00}{22.14}{}{0}{是啊，我们，有过那么多快乐。}
\BookLine{132.00}{0.00}{\heifont}{0}{阿德墨托斯：}
\BookLine{140.50}{22.14}{}{0}{{\kaifont （突然抓住了她的手，激动地）}你之前讨厌我！}
\BookLine{157.50}{0.00}{\heifont}{0}{阿尔刻提斯：}
\BookLine{166.00}{22.14}{}{0}{{\kaifont （把手抽出来）}没有，从来没有。}
\newpage
% 原书页码：122
\null
\BookNumber{122}
\Vertical{4}{169}{}{6.2}{8}阿\\尔\\刻\\提\\斯\\之\\歌\EndVertical
\BookLine{30.00}{0.00}{\heifont}{0}{阿德墨托斯：}
\BookLine{38.50}{22.14}{}{0}{有个年轻人，一而再地去参加那场狮子与野}
\BookLine{47.00}{22.14}{}{0}{猪的比赛……他又来了——那个塞萨利的大}
\BookLine{55.50}{22.14}{}{0}{傻瓜。}
\BookLine{72.50}{7.38}{\kaifont}{0}{两人都笑了起来。}
\BookLine{89.50}{0.00}{\heifont}{0}{阿尔刻提斯：}
\BookLine{98.00}{22.14}{}{0}{啊，人生多奇妙啊！}
\BookLine{115.00}{0.00}{\heifont}{0}{阿德墨托斯：}
\BookLine{123.50}{22.14}{}{0}{你不讨厌我吗，当我驾着那两只该死的野兽绕}
\BookLine{132.00}{22.14}{}{0}{过城墙的时候？}
\BookLine{149.00}{0.00}{\heifont}{0}{阿尔刻提斯：}
\BookLine{157.50}{22.14}{}{0}{不，我当时更受煎熬啊。因为我已经爱上了那}
\BookLine{166.00}{22.14}{}{0}{个塞萨利来的，一脸严肃的年轻人。你是唯}
\BookLine{174.50}{22.14}{}{0}{一一个两次来参加比试的追求者。连赫拉克勒}
\newpage
% 原书页码：123
\null
\BookNumber{123}
\Vertical{129}{11}{}{7}{10}第\\二\\幕\EndVertical
\BookLine{30.00}{22.14}{}{0}{斯都放弃了，当时，我不知道我的人生将会}
\BookLine{38.50}{22.14}{}{0}{怎样。}
\BookLine{55.50}{0.00}{\heifont}{0}{阿德墨托斯：}
\BookLine{64.00}{22.14}{}{0}{{\kaifont （骄傲地，激动地）}我，我知道。}
\BookLine{81.00}{0.00}{\heifont}{0}{阿尔刻提斯：}
\BookLine{89.50}{22.14}{}{0}{亲爱的阿德墨托斯……你知道，你娶了一个自}
\BookLine{98.00}{22.14}{}{0}{私又倔强的女孩。}
\BookLine{115.00}{7.38}{\kaifont}{0}{阿尔刻提斯做了个手势示意他不要再说下去。}
\BookLine{132.00}{0.00}{\heifont}{0}{阿德墨托斯：}
\BookLine{140.50}{22.14}{}{0}{{\kaifont （更加热烈）}我们的爱！我们之间的爱！我们在}
\BookLine{149.00}{22.14}{}{0}{夜晚说的那些悄悄话！伊皮伊尼兹出生的时候，}
\BookLine{157.50}{22.14}{}{0}{我差点就失去你了。阿尔刻提斯！这所有的一}
\BookLine{166.00}{22.14}{}{0}{切，人生啊……我真想一直这样活下去，和你}
\BookLine{174.50}{22.14}{}{0}{一起，在你身边。}
\newpage
% 原书页码：124
\null
\BookNumber{124}
\Vertical{4}{169}{}{6.2}{8}阿\\尔\\刻\\提\\斯\\之\\歌\EndVertical
\BookLine{30.00}{0.00}{\heifont}{0}{阿尔刻提斯：}
\BookLine{38.50}{22.14}{}{0}{嘘！这些话我们人类可不能大声说出来。}
\BookLine{55.50}{0.00}{\heifont}{0}{阿德墨托斯：}
\BookLine{64.00}{22.14}{}{0}{{\kaifont （试着让双脚落地）}我不知道怎么了……我的膝}
\BookLine{72.50}{22.14}{}{0}{盖竟然没有在发抖了。{\kaifont （充满着喜悦的希望）}阿}
\BookLine{81.00}{22.14}{}{0}{尔刻提斯！也许，也许我真的会活下去！}
\BookLine{98.00}{0.00}{\heifont}{0}{阿尔刻提斯：}
\BookLine{106.50}{22.14}{}{0}{{\kaifont （带着同样的喜悦，望着天空）}无论是死，是}
\BookLine{115.00}{22.14}{}{0}{生，我们都被眷顾着，指引着，懂得着。哦，}
\BookLine{123.50}{22.14}{}{0}{阿德墨托斯，静静地躺在那里吧，静静地。}
\BookLine{140.50}{0.00}{\heifont}{0}{阿德墨托斯：}
\BookLine{149.00}{22.14}{}{0}{我简直不敢……相信，不敢去奢望……}
\BookLine{166.00}{0.00}{\heifont}{0}{阿尔刻提斯：}
\BookLine{174.50}{22.14}{}{0}{{\kaifont （疼痛再次袭击了她，她向宫殿走去）}阿德墨托}
\newpage
% 原书页码：125
\null
\BookNumber{125}
\Vertical{129}{11}{}{7}{10}第\\二\\幕\EndVertical
\BookLine{30.00}{22.14}{}{0}{斯，如果德尔斐有人送来消息给我，说要让我}
\BookLine{38.50}{22.14}{}{0}{为了我的孩子，或是塞萨利，或是——我的丈}
\BookLine{47.00}{22.14}{}{0}{夫而死，我不会犹豫。}
\BookLine{64.00}{0.00}{\heifont}{0}{阿德墨托斯：}
\BookLine{72.50}{22.14}{}{0}{不，不，没有人希望别人为自己而死。每一个}
\BookLine{81.00}{22.14}{}{0}{人都要承受他自己的死亡。}
\BookLine{98.00}{0.00}{\heifont}{0}{阿尔刻提斯：}
\BookLine{106.50}{22.14}{}{0}{{\kaifont （掩藏着自己的疼痛）}你说的是什么话呀，阿德}
\BookLine{115.00}{22.14}{}{0}{墨托斯！想想那些战士们——成千上万的战士}
\BookLine{123.50}{22.14}{}{0}{们——都是为他人而死的啊。我们这些女人，}
\BookLine{132.00}{22.14}{}{0}{这些可怜又怯懦的战士，许许多多也都为了自}
\BookLine{140.50}{22.14}{}{0}{己的丈夫和孩子们而死去。}
\BookLine{157.50}{7.38}{\kaifont}{0}{她踉踉跄跄地走向宫殿大门。}
\newpage
% 原书页码：126
\null
\BookNumber{126}
\Vertical{4}{169}{}{6.2}{8}阿\\尔\\刻\\提\\斯\\之\\歌\EndVertical
\BookLine{30.00}{0.00}{\heifont}{0}{阿德墨托斯：}
\BookLine{38.50}{22.14}{}{0}{如果是神明们决意让妻子为了丈夫去死，我会}
\BookLine{47.00}{22.14}{}{0}{很看不起他们的。{\kaifont （起身，惊讶地呼喊着）}你}
\BookLine{55.50}{22.14}{}{0}{看！那光芒！}
\BookLine{72.50}{0.00}{\heifont}{0}{阿尔刻提斯：}
\BookLine{81.00}{22.14}{}{0}{阿格莱亚！阿格莱亚！}
\BookLine{98.00}{0.00}{\heifont}{0}{阿德墨托斯：}
\BookLine{106.50}{22.14}{}{0}{{\kaifont （起身跑向她）}阿尔刻提斯，你看起来不太对}
\BookLine{115.00}{22.14}{}{0}{劲！阿格莱亚！阿格莱亚！}
\BookLine{132.00}{7.38}{\kaifont}{0}{阿格莱亚从宫门出来跑向她。}
\BookLine{149.00}{0.00}{\heifont}{0}{阿尔刻提斯：}
\BookLine{157.50}{22.14}{}{0}{把我扶到床上去。}
\newpage
% 原书页码：127
\null
\BookNumber{127}
\Vertical{129}{11}{}{7}{10}第\\二\\幕\EndVertical
\BookLine{30.00}{0.00}{\heifont}{0}{阿德墨托斯：}
\BookLine{38.50}{22.14}{}{0}{你生病了，你是不是病了，阿尔刻提斯？}
\BookLine{55.50}{0.00}{\heifont}{0}{阿尔刻提斯：}
\BookLine{64.00}{22.14}{}{0}{{\kaifont （转过脸）}我把我的生命给你。要快乐，要}
\BookLine{72.50}{22.14}{}{0}{快乐。}
\BookLine{89.50}{7.38}{\kaifont}{0}{她晕倒在他们的臂弯，大家把她抬进了宫殿。右边门口}
\BookLine{98.00}{0.00}{\kaifont}{0}{传来人群的欢呼声，断断续续的歌声，接着就是赫拉克勒斯}
\BookLine{106.50}{0.00}{\kaifont}{0}{洪亮的嗓音。}
\BookLine{123.50}{0.00}{\heifont}{0}{赫拉克勒斯：}
\BookLine{132.00}{22.14}{}{0}{我的好哥们儿阿德墨——托斯在哪里？}
\BookLine{149.00}{0.00}{\heifont}{0}{居民：}
\BookLine{157.50}{22.14}{}{0}{赫拉克勒斯来啦！赫拉克勒斯来啦！}
\newpage
% 原书页码：128
\null
\BookNumber{128}
\Vertical{4}{169}{}{6.2}{8}阿\\尔\\刻\\提\\斯\\之\\歌\EndVertical
\BookLine{30.00}{0.00}{\heifont}{0}{赫拉克勒斯：}
\BookLine{38.50}{22.14}{}{0}{阿尔刻——提斯！圣洁的阿尔刻——提斯在哪}
\BookLine{47.00}{22.14}{}{0}{里？沉默的阿尔刻提斯！热情诚——挚的阿德}
\BookLine{55.50}{22.14}{}{0}{墨托斯！}
\BookLine{72.50}{7.38}{\kaifont}{0}{看门人走了进来。}
\BookLine{89.50}{0.00}{\heifont}{0}{居民：}
\BookLine{98.00}{22.14}{}{0}{赫拉克勒斯来啦！赫拉克勒斯！}
\BookLine{115.00}{0.00}{\heifont}{0}{看门人：}
\BookLine{123.50}{22.14}{}{0}{神明保佑！现在可怎么办？}
\BookLine{140.50}{7.38}{\kaifont}{0}{镇上的人不断从小门涌进来，一边还叫喊着：“宙斯之子}
\BookLine{149.00}{0.00}{\kaifont}{0}{赫拉克勒斯万岁！”赫拉克勒斯上场——醉醺醺的，兴高采}
\BookLine{157.50}{0.00}{\kaifont}{0}{烈的，头上戴着花环，手里拿着酒壶。两个村民大笑着，抬}
\BookLine{166.00}{0.00}{\kaifont}{0}{着赫拉克勒斯沉甸甸的大棍子。}
\newpage
% 原书页码：129
\null
\BookNumber{129}
\Vertical{129}{11}{}{7}{10}第\\二\\幕\EndVertical
\BookLine{30.00}{0.00}{\heifont}{0}{赫拉克勒斯：}
\BookLine{38.50}{22.14}{}{0}{阿尔刻提斯，伊俄尔科斯最美的女儿！阿德墨}
\BookLine{47.00}{22.14}{}{0}{托斯，我最最亲爱的朋友！他们都在哪儿！}
\BookLine{64.00}{0.00}{\heifont}{0}{居民：}
\BookLine{72.50}{22.14}{}{0}{赫拉克勒斯，猛兽的毁灭者！赫拉克勒斯，人}
\BookLine{81.00}{22.14}{}{0}{类的好朋友！赫拉克勒斯万岁！}
\BookLine{98.00}{0.00}{\heifont}{0}{看门人：}
\BookLine{106.50}{22.14}{}{0}{伟大的赫拉克勒斯，对于您的到来，费莱致以}
\BookLine{115.00}{22.14}{}{0}{一千次的欢迎！}
\BookLine{132.00}{0.00}{\heifont}{0}{赫拉克勒斯：}
\BookLine{140.50}{22.14}{}{0}{我的好朋友伊皮伊尼兹在哪里，勇敢的猎手，}
\BookLine{149.00}{22.14}{}{0}{能干的渔夫？伊皮伊尼兹，我要跟你比试摔跤。}
\BookLine{157.50}{22.14}{}{0}{我以神的雷电起誓，这次可不会再让你摔我一}
\BookLine{166.00}{22.14}{}{0}{大跤了！}
\newpage
% 原书页码：130
\null
\BookNumber{130}
\Vertical{4}{169}{}{6.2}{8}阿\\尔\\刻\\提\\斯\\之\\歌\EndVertical
\BookLine{30.00}{0.00}{\heifont}{0}{看门人：}
\BookLine{38.50}{22.14}{}{0}{他们马上就过来，赫拉克勒斯。他们全都高兴}
\BookLine{47.00}{22.14}{}{0}{坏了。}
\BookLine{64.00}{7.38}{\kaifont}{0}{宫殿里传来一阵哀号和哭喊。}
\BookLine{81.00}{0.00}{\heifont}{0}{赫拉克勒斯：}
\BookLine{89.50}{22.14}{}{0}{我听到有哭声。老人家，他们在哭什么？}
\BookLine{106.50}{0.00}{\heifont}{0}{看门人：}
\BookLine{115.00}{22.14}{}{0}{您说哭声，赫拉克勒斯！那是妇人和姑娘们在}
\BookLine{123.50}{22.14}{}{0}{为了您的到来欢呼尖叫呢。}
\BookLine{140.50}{7.38}{\kaifont}{0}{阿格莱亚从宫殿里匆匆走出来。}
\BookLine{157.50}{0.00}{\heifont}{0}{阿格莱亚：}
\BookLine{166.00}{22.14}{}{0}{神圣的赫拉克勒斯，向您致以一千次的欢迎。}
\BookLine{174.50}{22.14}{}{0}{国王和王后马上就来，他们太……太开心了，}
\newpage
% 原书页码：131
\null
\BookNumber{131}
\Vertical{129}{11}{}{7}{10}第\\二\\幕\EndVertical
\BookLine{30.00}{22.14}{}{0}{您知道的。哦，宙斯之子，见到您是多么荣}
\BookLine{38.50}{22.14}{}{0}{幸啊。}
\BookLine{55.50}{0.00}{\heifont}{0}{赫拉克勒斯：}
\BookLine{64.00}{22.14}{}{0}{{\kaifont （大声地）}我不是什么宙斯之子！}
\BookLine{81.00}{0.00}{\heifont}{0}{阿格莱亚：}
\BookLine{89.50}{22.14}{}{0}{{\kaifont （捂上耳朵）}赫拉克勒斯，您在乱说什么呀？}
\BookLine{106.50}{0.00}{\heifont}{0}{赫拉克勒斯：}
\BookLine{115.00}{22.14}{}{0}{我是安菲特律翁和阿尔克墨涅的儿子。我是一}
\BookLine{123.50}{22.14}{}{0}{个普通人，阿格莱亚。我做起事情来和别人一}
\BookLine{132.00}{22.14}{}{0}{样费力。}
\BookLine{149.00}{0.00}{\heifont}{0}{阿格莱亚：}
\BookLine{157.50}{22.14}{}{0}{哦，愿众神不要降灾给这个家！您居然说出这}
\BookLine{166.00}{22.14}{}{0}{么坏的话来，赫拉克勒斯，您一定是喝醉了。}
\newpage
% 原书页码：132
\null
\BookNumber{132}
\Vertical{4}{169}{}{6.2}{8}阿\\尔\\刻\\提\\斯\\之\\歌\EndVertical
\BookLine{30.00}{0.00}{\heifont}{0}{赫拉克勒斯：}
\BookLine{38.50}{22.14}{}{0}{我告诉你，我只是一个人，一个普通的人。}
\BookLine{55.50}{0.00}{\heifont}{0}{阿格莱亚：}
\BookLine{64.00}{22.14}{}{0}{神也好，人也好，赫拉克勒斯，让我瞧瞧，您}
\BookLine{72.50}{22.14}{}{0}{怎么脏成这样了！这还是我们全希腊最帅的男}
\BookLine{81.00}{22.14}{}{0}{人吗？凭着不朽的众神起誓，我都快认不出您}
\BookLine{89.50}{22.14}{}{0}{来了！现在听好了，您还记得阿格莱亚总是为}
\BookLine{98.00}{22.14}{}{0}{您准备的热水澡吗？他们会把您这“普通人”的}
\BookLine{106.50}{22.14}{}{0}{样子给洗掉的。我为您准备了精油，还记得那}
\BookLine{115.00}{22.14}{}{0}{些精油吗？}
\BookLine{132.00}{0.00}{\heifont}{0}{赫拉克勒斯：}
\BookLine{140.50}{22.14}{}{0}{{\kaifont （自信地）}首先，阿格莱亚，首先！}
\BookLine{157.50}{7.38}{\kaifont}{0}{他做了一个喝酒的姿势。}
\newpage
% 原书页码：133
\null
\BookNumber{133}
\Vertical{129}{11}{}{7}{10}第\\二\\幕\EndVertical
\BookLine{30.00}{0.00}{\heifont}{0}{阿格莱亚：}
\BookLine{38.50}{22.14}{}{0}{您还没忘了我们的好酒，是不是？马上就}
\BookLine{47.00}{22.14}{}{0}{准备！}
\BookLine{64.00}{0.00}{\heifont}{0}{赫拉克勒斯：}
\BookLine{72.50}{22.14}{}{0}{{\kaifont （突然吼了起来）}我在来的路上听说阿德墨托斯}
\BookLine{81.00}{22.14}{}{0}{被人刺伤了。谁伤害的他，阿格莱亚？}
\BookLine{98.00}{0.00}{\heifont}{0}{阿格莱亚：}
\BookLine{106.50}{22.14}{}{0}{哦，这事儿已经过去了，赫拉克勒斯。他已经}
\BookLine{115.00}{22.14}{}{0}{好起来了，跟你我一样。}
\BookLine{132.00}{7.38}{\kaifont}{0}{罗德普和另一个姑娘从宫里走出来，扛着酒罐，拿着几}
\BookLine{140.50}{0.00}{\kaifont}{0}{个杯子。}
\BookLine{149.00}{22.14}{}{0}{过来，坐这儿，好好放松一下。}
\BookLine{166.00}{7.38}{\kaifont}{0}{赫拉克勒斯想抓住其中一个女孩，女孩躲开了。他转而}
\newpage
% 原书页码：134
\null
\BookNumber{134}
\Vertical{4}{169}{}{6.2}{8}阿\\尔\\刻\\提\\斯\\之\\歌\EndVertical
\BookLine{30.00}{0.00}{\kaifont}{0}{去追另一个，一把抓住了她。}
\BookLine{47.00}{0.00}{\heifont}{0}{赫拉克勒斯：}
\BookLine{55.50}{22.14}{}{0}{你叫什么名字，我的小心肝？}
\BookLine{72.50}{0.00}{\heifont}{0}{阿格莱亚：}
\BookLine{81.00}{22.14}{}{0}{{\kaifont （生气地）}赫拉克勒斯！}
\BookLine{98.00}{7.38}{\kaifont}{0}{女孩跑掉，赫拉克勒斯追上去，跌了一跤，摔倒在地上。}
\BookLine{115.00}{0.00}{\heifont}{0}{赫拉克勒斯：}
\BookLine{123.50}{22.14}{}{0}{哦！我受伤了！真他妈的混账！我的膝盖！我}
\BookLine{132.00}{22.14}{}{0}{的膝盖！}
\BookLine{149.00}{0.00}{\heifont}{0}{阿格莱亚：}
\BookLine{157.50}{22.14}{}{0}{赫拉克勒斯！我可一点都不同情您。您忘了这}
\BookLine{166.00}{22.14}{}{0}{是哪儿了吗？不朽的天神啊，阿尔刻提斯王后}
\BookLine{174.50}{22.14}{}{0}{要是看到您这副模样，该怎么想您啊！}
\newpage
% 原书页码：135
\null
\BookNumber{135}
\Vertical{129}{11}{}{7}{10}第\\二\\幕\EndVertical
\BookLine{30.00}{0.00}{\heifont}{0}{赫拉克勒斯：}
\BookLine{38.50}{22.14}{}{0}{{\kaifont （慢慢起身，坐下）}我都走了二十天了。他们}
\BookLine{47.00}{22.14}{}{0}{人呢？我的朋友们，阿德墨托斯和阿尔刻提}
\BookLine{55.50}{22.14}{}{0}{斯呢？}
\BookLine{72.50}{0.00}{\heifont}{0}{阿格莱亚：}
\BookLine{81.00}{22.14}{}{0}{{\kaifont （严肃地）}好了，赫拉克勒斯，您是这个家的老}
\BookLine{89.50}{22.14}{}{0}{朋友了，是不是？}
\BookLine{106.50}{0.00}{\heifont}{0}{赫拉克勒斯：}
\BookLine{115.00}{22.14}{}{0}{当然是了！}
\BookLine{132.00}{0.00}{\heifont}{0}{阿格莱亚：}
\BookLine{140.50}{22.14}{}{0}{我们有什么就可以说什么，对不对？感谢神明，}
\BookLine{149.00}{22.14}{}{0}{我们不需要像对其他客人一样对您掩藏什么。}
\BookLine{166.00}{0.00}{\heifont}{0}{赫拉克勒斯：}
\BookLine{174.50}{22.14}{}{0}{我就是他们俩的兄长！哥哥！}
\newpage
% 原书页码：136
\null
\BookNumber{136}
\Vertical{4}{169}{}{6.2}{8}阿\\尔\\刻\\提\\斯\\之\\歌\EndVertical
\BookLine{30.00}{0.00}{\heifont}{0}{阿格莱亚：}
\BookLine{38.50}{22.14}{}{0}{我要跟您说的是，这个家里有一位夫人，一个}
\BookLine{47.00}{22.14}{}{0}{孤儿……这个家有位夫人……}
\BookLine{64.00}{0.00}{\heifont}{0}{赫拉克勒斯：}
\BookLine{72.50}{22.14}{}{0}{怎么了？死了？}
\BookLine{89.50}{0.00}{\heifont}{0}{阿格莱亚：}
\BookLine{98.00}{22.14}{}{0}{{\kaifont （手指放在唇间）}您也知道国王和王后对我们这}
\BookLine{106.50}{22.14}{}{0}{些仆人非常好……}
\BookLine{123.50}{0.00}{\heifont}{0}{赫拉克勒斯：}
\BookLine{132.00}{22.14}{}{0}{死了？一个孤儿？}
\BookLine{149.00}{0.00}{\heifont}{0}{阿格莱亚：}
\BookLine{157.50}{22.14}{}{0}{是的，她没有双亲。好了，他们应该很快就到}
\BookLine{166.00}{22.14}{}{0}{了，等您沐浴过后。但是现在，此刻，您知道，}
\BookLine{174.50}{22.14}{}{0}{他们像朋友一样，充满怜悯地陪着那个可怜的}
\newpage
% 原书页码：137
\null
\BookNumber{137}
\Vertical{129}{11}{}{7}{10}第\\二\\幕\EndVertical
\BookLine{30.00}{22.14}{}{0}{姑娘。这些事情您都很明白的，对吗？}
\BookLine{47.00}{0.00}{\heifont}{0}{赫拉克勒斯：}
\BookLine{55.50}{22.14}{}{0}{不，不，阿格莱亚，很多事情我都不明白。但}
\BookLine{64.00}{22.14}{}{0}{谁又全都明白呢？{\kaifont （停顿）}阿尔刻提斯明白，对}
\BookLine{72.50}{22.14}{}{0}{不对，阿格莱亚？}
\BookLine{89.50}{0.00}{\heifont}{0}{阿格莱亚：}
\BookLine{98.00}{22.14}{}{0}{是的，是的，赫拉克勒斯。}
\BookLine{115.00}{0.00}{\heifont}{0}{赫拉克勒斯：}
\BookLine{123.50}{22.14}{}{0}{那我也对你实话实说吧，阿格莱亚。我来这}
\BookLine{132.00}{22.14}{}{0}{儿——走了二十天来这儿——是要问阿尔刻提}
\BookLine{140.50}{22.14}{}{0}{斯一个问题。}
\BookLine{157.50}{0.00}{\heifont}{0}{阿格莱亚：}
\BookLine{166.00}{22.14}{}{0}{一个问题，赫拉克勒斯？}
\newpage
% 原书页码：138
\null
\BookNumber{138}
\Vertical{4}{169}{}{6.2}{8}阿\\尔\\刻\\提\\斯\\之\\歌\EndVertical
\BookLine{30.00}{7.38}{\kaifont}{0}{赫拉克勒斯两次指了指天空暗示自己的问题。}
\BookLine{47.00}{0.00}{\heifont}{0}{阿格莱亚：}
\BookLine{55.50}{22.14}{}{0}{关于，众神？}
\BookLine{72.50}{7.38}{\kaifont}{0}{他点点头，阿格莱亚回过神来，干脆地说。}
\BookLine{89.50}{0.00}{\heifont}{0}{阿格莱亚：}
\BookLine{98.00}{22.14}{}{0}{好吧，不过您还是要先去沐浴。洗完澡之后，}
\BookLine{106.50}{22.14}{}{0}{您看起来会像个十七岁的小伙子。我还要给您}
\BookLine{115.00}{22.14}{}{0}{一顶漂亮的花环，还有香水！}
\BookLine{132.00}{7.38}{\kaifont}{0}{赫拉克勒斯起身，跟在她身后登上宫殿的台阶。}
\BookLine{149.00}{0.00}{\heifont}{0}{赫拉克勒斯：}
\BookLine{157.50}{22.14}{}{0}{我洗澡的时候，你再给我讲讲阿波罗到塞萨利}
\BookLine{166.00}{22.14}{}{0}{的事情吧？}
\newpage
% 原书页码：139
\null
\BookNumber{139}
\Vertical{129}{11}{}{7}{10}第\\二\\幕\EndVertical
\BookLine{30.00}{0.00}{\heifont}{0}{阿格莱亚：}
\BookLine{38.50}{22.14}{}{0}{好的，没问题，我再给您讲一遍。}
\BookLine{55.50}{0.00}{\heifont}{0}{赫拉克勒斯：}
\BookLine{64.00}{22.14}{}{0}{{\kaifont （突然打断了她）}一直没人搞清楚哪个是阿}
\BookLine{72.50}{22.14}{}{0}{波罗？}
\BookLine{89.50}{7.38}{\kaifont}{0}{阿格莱亚摇摇头。}
\BookLine{98.00}{22.14}{}{0}{一整年时间，没人搞清楚哪个是阿波罗，包括}
\BookLine{106.50}{22.14}{}{0}{阿尔刻提斯？}
\BookLine{123.50}{7.38}{\kaifont}{0}{她又摇了摇头，赫拉克勒斯拍拍自己的脑门。}
\BookLine{132.00}{22.14}{}{0}{阿格莱亚，谁能搞明白他们？我们估计永远搞}
\BookLine{140.50}{22.14}{}{0}{不明白。我每次一想起他们啊，我就开始发抖，}
\BookLine{149.00}{22.14}{}{0}{头晕。}
\newpage
% 原书页码：140
\null
\BookNumber{140}
\Vertical{4}{169}{}{6.2}{8}阿\\尔\\刻\\提\\斯\\之\\歌\EndVertical
\BookLine{30.00}{0.00}{\heifont}{0}{阿格莱亚：}
\BookLine{38.50}{22.14}{}{0}{赫拉克勒斯！您太累了，跟我来……}
\BookLine{55.50}{7.38}{\kaifont}{0}{阿德墨托斯走了进来，赫拉克勒斯和阿格莱亚准备上台}
\BookLine{64.00}{0.00}{\kaifont}{0}{阶时，他正好站在最上面一级台阶上。}
\BookLine{81.00}{0.00}{\heifont}{0}{阿德墨托斯：}
\BookLine{89.50}{22.14}{}{0}{{\kaifont （大声地）}欢迎你，赫拉克勒斯，人类的朋友。}
\BookLine{98.00}{22.14}{}{0}{赫拉克勒斯，人类的大恩人！}
\BookLine{115.00}{0.00}{\heifont}{0}{赫拉克勒斯：}
\BookLine{123.50}{22.14}{}{0}{阿德墨托斯！我的老朋友！}
\BookLine{140.50}{7.38}{\kaifont}{0}{阿德墨托斯把双手放在赫拉克勒斯的肩上，他们互相盯}
\BookLine{149.00}{0.00}{\kaifont}{0}{着对方的眼睛看了好一阵子。}
\BookLine{166.00}{0.00}{\heifont}{0}{阿德墨托斯：}
\BookLine{174.50}{22.14}{}{0}{你从哪里来的，赫拉克勒斯？}
\newpage
% 原书页码：141
\null
\BookNumber{141}
\Vertical{129}{11}{}{7}{10}第\\二\\幕\EndVertical
\BookLine{30.00}{0.00}{\heifont}{0}{赫拉克勒斯：}
\BookLine{38.50}{22.14}{}{0}{从劳苦重担中来，阿德墨托斯……从劳苦重}
\BookLine{47.00}{22.14}{}{0}{担中。}
\BookLine{64.00}{0.00}{\heifont}{0}{阿格莱亚：}
\BookLine{72.50}{22.14}{}{0}{阿德墨托斯国王，赫拉克勒斯正准备去沐浴。}
\BookLine{81.00}{22.14}{}{0}{您可以等他待会儿出来了再聊。{\kaifont （看看天空）}天}
\BookLine{89.50}{22.14}{}{0}{色马上就晚了。}
\BookLine{106.50}{0.00}{\heifont}{0}{阿德墨托斯：}
\BookLine{115.00}{22.14}{}{0}{不行，我们先要一起喝上一碗！过来，坐下，}
\BookLine{123.50}{22.14}{}{0}{赫拉克勒斯。我病了一段时间，这个又呆又傻}
\BookLine{132.00}{22.14}{}{0}{的沙发，我之前总是坐在这儿晒太阳。现在跟}
\BookLine{140.50}{22.14}{}{0}{我说说，你最近在干什么大事。}
\BookLine{157.50}{7.38}{\kaifont}{0}{阿格莱亚和女孩们都回到宫里。}
\newpage
% 原书页码：142
\null
\BookNumber{142}
\Vertical{4}{169}{}{6.2}{8}阿\\尔\\刻\\提\\斯\\之\\歌\EndVertical
\BookLine{30.00}{0.00}{\heifont}{0}{赫拉克勒斯：}
\BookLine{38.50}{22.14}{}{0}{{\kaifont （急切地，睁大了眼睛）}阿德墨托斯，阿德墨托}
\BookLine{47.00}{22.14}{}{0}{斯——我杀了九头蛇。}
\BookLine{64.00}{0.00}{\heifont}{0}{阿德墨托斯：}
\BookLine{72.50}{22.14}{}{0}{{\kaifont （站起身，十分惊讶）}非凡的赫拉克勒斯！你}
\BookLine{81.00}{22.14}{}{0}{杀了九头蛇！凭着不朽的神明起誓，赫拉克勒}
\BookLine{89.50}{22.14}{}{0}{斯，你真是人类的好朋友。你杀了九头蛇！九}
\BookLine{98.00}{22.14}{}{0}{头蛇啊！}
\BookLine{115.00}{0.00}{\heifont}{0}{赫拉克勒斯：}
\BookLine{123.50}{22.14}{}{0}{{\kaifont （指了指自己，让阿德墨托斯靠得更近）}阿德墨}
\BookLine{132.00}{22.14}{}{0}{托斯！{\kaifont （又指了指自己）}很不容易！}
\BookLine{149.00}{0.00}{\heifont}{0}{阿德墨托斯：}
\BookLine{157.50}{22.14}{}{0}{我简直不敢相信！}
\newpage
% 原书页码：143
\null
\BookNumber{143}
\Vertical{129}{11}{}{7}{10}第\\二\\幕\EndVertical
\BookLine{30.00}{0.00}{\heifont}{0}{赫拉克勒斯：}
\BookLine{38.50}{22.14}{}{0}{{\kaifont （盯着他的脸，简直有些苦涩）}很不容易。{\kaifont （忽}}
\BookLine{47.00}{22.14}{}{0}{{\kaifont 然变得十分恳切）}阿德墨托斯，我是宙斯之}
\BookLine{55.50}{22.14}{}{0}{子吗？}
\BookLine{72.50}{0.00}{\heifont}{0}{阿德墨托斯：}
\BookLine{81.00}{22.14}{}{0}{赫拉克勒斯！所有人都知道，你是宙斯和阿尔}
\BookLine{89.50}{22.14}{}{0}{克墨涅的儿子。}
\BookLine{106.50}{0.00}{\heifont}{0}{赫拉克勒斯：}
\BookLine{115.00}{22.14}{}{0}{有一个人能告诉我答案。我来这儿就是要问阿}
\BookLine{123.50}{22.14}{}{0}{尔刻提斯这件事。阿尔刻提斯在哪儿？真正的}
\BookLine{132.00}{22.14}{}{0}{女王在哪儿？}
\BookLine{149.00}{0.00}{\heifont}{0}{阿德墨托斯：}
\BookLine{157.50}{22.14}{}{0}{阿格莱亚不是跟你说了吗！}
\newpage
% 原书页码：144
\null
\BookNumber{144}
\Vertical{4}{169}{}{6.2}{8}阿\\尔\\刻\\提\\斯\\之\\歌\EndVertical
\BookLine{30.00}{0.00}{\heifont}{0}{赫拉克勒斯：}
\BookLine{38.50}{22.14}{}{0}{哦，对。那个死去的孤儿是谁？是这个家最初}
\BookLine{47.00}{22.14}{}{0}{的那批人吗？}
\BookLine{64.00}{0.00}{\heifont}{0}{阿德墨托斯：}
\BookLine{72.50}{22.14}{}{0}{她称自己为仆人中的仆人。}
\BookLine{89.50}{0.00}{\heifont}{0}{赫拉克勒斯：}
\BookLine{98.00}{22.14}{}{0}{你们都很喜欢她是吗？}
\BookLine{115.00}{0.00}{\heifont}{0}{阿德墨托斯：}
\BookLine{123.50}{22.14}{}{0}{是的，我们都很爱她。}
\BookLine{140.50}{0.00}{\heifont}{0}{赫拉克勒斯：}
\BookLine{149.00}{22.14}{}{0}{你瞧，阿德墨托斯，人人都说我是宙斯之子，}
\BookLine{157.50}{22.14}{}{0}{所以我办起事情来应该比别人容易很多。如果}
\BookLine{166.00}{22.14}{}{0}{我拥有跟他一样的血液、心脏、肺腑{\kaifont （指着天}}
\BookLine{174.50}{22.14}{}{0}{{\kaifont 空）}，事情不就应该很容易吗？但是啊，阿德}
\newpage
% 原书页码：145
\null
\BookNumber{145}
\Vertical{129}{11}{}{7}{10}第\\二\\幕\EndVertical
\BookLine{30.00}{22.14}{}{0}{墨托斯，很不容易。九头蛇！{\kaifont （用手臂把阿德墨}}
\BookLine{38.50}{22.14}{}{0}{{\kaifont 托斯环起来，模仿着九头蛇做出恐怖的表情）}}
\BookLine{47.00}{22.14}{}{0}{我当时快要炸开了，鲜血跟泉水一样从我耳朵}
\BookLine{55.50}{22.14}{}{0}{里喷出来。如果我只是一个人——安菲特律翁}
\BookLine{64.00}{22.14}{}{0}{和阿尔克墨涅的儿子——那么，阿德墨托斯，}
\BookLine{72.50}{22.14}{}{0}{{\kaifont （用力戳了戳阿德墨托斯的脸）}我肯定是个非常}
\BookLine{81.00}{22.14}{}{0}{了不起的凡人！}
\BookLine{98.00}{0.00}{\heifont}{0}{阿德墨托斯：}
\BookLine{106.50}{22.14}{}{0}{神也好，人也好，赫拉克勒斯——是神也好，}
\BookLine{115.00}{22.14}{}{0}{人也好，所有人都会敬重你，感激你。}
\BookLine{132.00}{0.00}{\heifont}{0}{赫拉克勒斯：}
\BookLine{140.50}{22.14}{}{0}{但是，我想弄明白。有时候我觉得我就是宙斯}
\BookLine{149.00}{22.14}{}{0}{之子，但有时候我又……又是一只野兽。阿德}
\BookLine{157.50}{22.14}{}{0}{墨托斯，我觉得自己野蛮而残暴。每个月，信}
\BookLine{166.00}{22.14}{}{0}{使们从希腊各地送来消息，请我去干这个干那}
\BookLine{174.50}{22.14}{}{0}{个。我什么都不干了，在我解决这件事情之前，}
\newpage
% 原书页码：146
\null
\BookNumber{146}
\Vertical{4}{169}{}{6.2}{8}阿\\尔\\刻\\提\\斯\\之\\歌\EndVertical
\BookLine{30.00}{22.14}{}{0}{我什么事都不会再干了。就这几天，我已经预}
\BookLine{38.50}{22.14}{}{0}{见到了，有人会来找我，请我去……去下面的}
\BookLine{47.00}{22.14}{}{0}{世界，去冥界。阿德墨托斯，去把一个已死的}
\BookLine{55.50}{22.14}{}{0}{人带回来。{\kaifont （站起来，带着惊恐拒绝道）}不！不}
\BookLine{64.00}{22.14}{}{0}{行！这事我不会干！神也好，人也好，谁都别}
\BookLine{72.50}{22.14}{}{0}{指望我干这事！}
\BookLine{89.50}{0.00}{\heifont}{0}{阿德墨托斯：}
\BookLine{98.00}{22.14}{}{0}{{\kaifont （被震惊了）}不，赫拉克勒斯——这事没人干}
\BookLine{106.50}{22.14}{}{0}{过，也没人想过。}
\BookLine{123.50}{0.00}{\heifont}{0}{赫拉克勒斯：}
\BookLine{132.00}{22.14}{}{0}{我每次到这儿来都会看到它——冥界的入}
\BookLine{140.50}{22.14}{}{0}{口——我非常害怕，阿德墨托斯。}
\BookLine{157.50}{0.00}{\heifont}{0}{阿德墨托斯：}
\BookLine{166.00}{22.14}{}{0}{你不要再想这些事情了！}
\newpage
% 原书页码：147
\null
\BookNumber{147}
\Vertical{129}{11}{}{7}{10}第\\二\\幕\EndVertical
\BookLine{30.00}{0.00}{\heifont}{0}{赫拉克勒斯：}
\BookLine{38.50}{22.14}{}{0}{就算是宙斯之子也做不到这件事，对吗？所以，}
\BookLine{47.00}{22.14}{}{0}{你明白问题是什么吗？这难道不是个很要命的}
\BookLine{55.50}{22.14}{}{0}{问题吗？谁能回答呢？在这种你我都很无知的}
\BookLine{64.00}{22.14}{}{0}{问题上谁最有智慧呢？你知道是谁，就是那个}
\BookLine{72.50}{22.14}{}{0}{总是沉默，少言寡语的人。沉默，她总是沉默}
\BookLine{81.00}{22.14}{}{0}{着，对吗？}
\BookLine{98.00}{0.00}{\heifont}{0}{阿德墨托斯：}
\BookLine{106.50}{22.14}{}{0}{是的，她很沉默。}
\BookLine{123.50}{0.00}{\heifont}{0}{赫拉克勒斯：}
\BookLine{132.00}{22.14}{}{0}{但她会跟我说话的……会跟她的老朋友赫拉克}
\BookLine{140.50}{22.14}{}{0}{勒斯说话吧？她也会为我说话的，对吗？{\kaifont （灌一}}
\BookLine{149.00}{22.14}{}{0}{{\kaifont 大口酒）}你知道吗，阿德墨托斯，你了解你的}
\BookLine{157.50}{22.14}{}{0}{阿尔刻提斯是什么样的吗？}
\BookLine{174.50}{7.38}{\kaifont}{0}{阿德墨托斯痛苦地站了起来，又坐下。}
\newpage
% 原书页码：148
\null
\BookNumber{148}
\Vertical{4}{169}{}{6.2}{8}阿\\尔\\刻\\提\\斯\\之\\歌\EndVertical
\BookLine{30.00}{0.00}{\heifont}{0}{赫拉克勒斯：}
\BookLine{38.50}{22.14}{}{0}{不，你不知道。我去过很多地方，也追求过很}
\BookLine{47.00}{22.14}{}{0}{多人。希腊的这些王后啊，公主啊，我见过太}
\BookLine{55.50}{22.14}{}{0}{多了。我们认识的这些大人物的女儿，我们追}
\BookLine{64.00}{22.14}{}{0}{求过的姑娘都长大了。忒拜的俄狄浦斯有个女}
\BookLine{72.50}{22.14}{}{0}{儿，安提戈涅。还有佩涅洛佩，不久前嫁给了}
\BookLine{81.00}{22.14}{}{0}{伊萨卡的拉厄耳忒斯的儿子。斯巴达的勒达有}
\BookLine{89.50}{22.14}{}{0}{两个女儿，海伦和克吕泰涅斯特拉。我见过她}
\BookLine{98.00}{22.14}{}{0}{们每个人，也和她们都聊过天。可是在阿尔刻}
\BookLine{106.50}{22.14}{}{0}{提斯面前她们算什么呢？尘土，垃圾。}
\BookLine{123.50}{7.38}{\kaifont}{0}{阿格莱亚从宫殿里出来，手里拿着藤蔓编织的花环。看}
\BookLine{132.00}{0.00}{\kaifont}{0}{门人跟在后面，端着一罐精油。他们站在旁边聆听着。赫拉}
\BookLine{140.50}{0.00}{\kaifont}{0}{克勒斯起身，醉醺醺的要发狂。他转向阿德墨托斯，大喊着。}
\BookLine{157.50}{0.00}{\heifont}{0}{赫拉克勒斯：}
\BookLine{166.00}{22.14}{}{0}{你了解吗？你了解她的全部吗？}
\newpage
% 原书页码：149
\null
\BookNumber{149}
\Vertical{129}{11}{}{7}{10}第\\二\\幕\EndVertical
\BookLine{30.00}{0.00}{\heifont}{0}{阿德墨托斯：}
\BookLine{38.50}{22.14}{}{0}{{\kaifont （痛苦地抬起手臂）}我了解，赫拉克勒斯。}
\BookLine{55.50}{0.00}{\heifont}{0}{赫拉克勒斯：}
\BookLine{64.00}{22.14}{}{0}{你了解她的全部吗？她宽恕、原谅的能力？不，}
\BookLine{72.50}{22.14}{}{0}{你什么都不知道。阿尔刻提斯才是希腊的珍宝，}
\BookLine{81.00}{22.14}{}{0}{女人中的珍宝。}
\BookLine{98.00}{0.00}{\heifont}{0}{阿德墨托斯：}
\BookLine{106.50}{22.14}{}{0}{{\kaifont （往后退了三步）}赫拉克勒斯！阿尔刻提斯死}
\BookLine{115.00}{22.14}{}{0}{了。{\kaifont （停顿）}仆人中的仆人。原谅我，赫拉克}
\BookLine{123.50}{22.14}{}{0}{勒斯！}
\BookLine{140.50}{0.00}{\heifont}{0}{赫拉克勒斯：}
\BookLine{149.00}{22.14}{}{0}{{\kaifont （沉默了一阵。一股巨大的愤怒从心中升起，极}}
\BookLine{157.50}{22.14}{}{0}{{\kaifont 其痛苦地说）}阿德墨托斯……你算什么朋友。}
\BookLine{166.00}{22.14}{}{0}{凭着不朽的天神起誓，你是我最恨的敌人！}
\newpage
% 原书页码：150
\null
\BookNumber{150}
\Vertical{4}{169}{}{6.2}{8}阿\\尔\\刻\\提\\斯\\之\\歌\EndVertical
\BookLine{30.00}{7.38}{\kaifont}{0}{他一把扼住阿德墨托斯的咽喉，拖着丝毫没有抵抗的阿}
\BookLine{38.50}{0.00}{\kaifont}{0}{德墨托斯向后退了几步。}
\BookLine{55.50}{22.14}{}{0}{你要不是阿德墨托斯的话，我现在就杀了你！}
\BookLine{72.50}{0.00}{\heifont}{0}{阿格莱亚和看门人：}
\BookLine{81.00}{22.14}{}{0}{赫拉克勒斯，赫拉克勒斯！是阿尔刻提斯王后}
\BookLine{89.50}{22.14}{}{0}{要求这么做的！她让我们晚一些再告诉您！}
\BookLine{106.50}{0.00}{\heifont}{0}{赫拉克勒斯：}
\BookLine{115.00}{22.14}{}{0}{你们就没把我当朋友！我还以为我们是兄弟！}
\BookLine{132.00}{0.00}{\heifont}{0}{阿格莱亚：}
\BookLine{140.50}{22.14}{}{0}{{\kaifont （一直用拳头敲打赫拉克勒斯的后背）}赫拉克勒}
\BookLine{149.00}{22.14}{}{0}{斯！赫拉克勒斯！}
\BookLine{166.00}{0.00}{\heifont}{0}{赫拉克勒斯：}
\BookLine{174.50}{22.14}{}{0}{阿尔刻提斯死了，你们觉得用不着告诉我！你}
\newpage
% 原书页码：151
\null
\BookNumber{151}
\Vertical{129}{11}{}{7}{10}第\\二\\幕\EndVertical
\BookLine{30.00}{22.14}{}{0}{们还让我在这儿吹牛，喝酒，瞎胡闹！}
\BookLine{47.00}{0.00}{\heifont}{0}{阿格莱亚：}
\BookLine{55.50}{22.14}{}{0}{是她吩咐的，赫拉克勒斯。阿尔刻提斯下令我}
\BookLine{64.00}{22.14}{}{0}{们这么做。}
\BookLine{81.00}{0.00}{\heifont}{0}{赫拉克勒斯：}
\BookLine{89.50}{22.14}{}{0}{{\kaifont （沉思着，推着阿德墨托斯向后走了好几步）}你}
\BookLine{98.00}{22.14}{}{0}{们觉得我没有脑子，没有心肝，没有灵魂吗？}
\BookLine{106.50}{22.14}{}{0}{我连为朋友分担悲伤的资格都没有，就算有全}
\BookLine{115.00}{22.14}{}{0}{世界的感谢和赞美有什么用？}
\BookLine{132.00}{7.38}{\kaifont}{0}{停顿。}
\BookLine{149.00}{0.00}{\heifont}{0}{阿格莱亚：}
\BookLine{157.50}{22.14}{}{0}{只是为了看起来热情一些，赫拉克勒斯。她希}
\BookLine{166.00}{22.14}{}{0}{望如此！}
\newpage
% 原书页码：152
\null
\BookNumber{152}
\Vertical{4}{169}{}{6.2}{8}阿\\尔\\刻\\提\\斯\\之\\歌\EndVertical
\BookLine{30.00}{0.00}{\heifont}{0}{赫拉克勒斯：}
\BookLine{38.50}{22.14}{}{0}{客人才需要热情，真正的兄弟不需要。}
\BookLine{55.50}{0.00}{\heifont}{0}{阿德墨托斯：}
\BookLine{64.00}{22.14}{}{0}{{\kaifont （安静地）}请原谅我们，赫拉克勒斯。}
\BookLine{81.00}{0.00}{\heifont}{0}{赫拉克勒斯：}
\BookLine{89.50}{22.14}{}{0}{{\kaifont （放开阿德墨托斯）}阿尔刻提斯死了。阿尔刻提}
\BookLine{98.00}{22.14}{}{0}{斯已经在另一个世界了。{\kaifont （突然想到了什么，大}}
\BookLine{106.50}{22.14}{}{0}{{\kaifont 声地）}她在冥界！我的棍子呢？我的棍子！我}
\BookLine{115.00}{22.14}{}{0}{要去把她救出来！}
\BookLine{132.00}{0.00}{\heifont}{0}{阿格莱亚：}
\BookLine{140.50}{22.14}{}{0}{赫拉克勒斯！}
\BookLine{157.50}{0.00}{\heifont}{0}{阿德墨托斯：}
\BookLine{166.00}{22.14}{}{0}{不，赫拉克勒斯。你得活着！她是代替我死的，}
\BookLine{174.50}{22.14}{}{0}{她是为我而死的。不能再有人为此送命了。}
\newpage
% 原书页码：153
\null
\BookNumber{153}
\Vertical{129}{11}{}{7}{10}第\\二\\幕\EndVertical
\BookLine{30.00}{0.00}{\heifont}{0}{赫拉克勒斯：}
\BookLine{38.50}{22.14}{}{0}{我的棍子！我的棍子！阿德墨托斯，现在我要}
\BookLine{47.00}{22.14}{}{0}{告诉你一件事——我从没告诉过任何人。现}
\BookLine{55.50}{22.14}{}{0}{在，大家都应该知道。{\kaifont （向前走到小路的高处）}}
\BookLine{64.00}{22.14}{}{0}{我曾经差点强暴了阿尔刻提斯，极其粗野。是}
\BookLine{72.50}{22.14}{}{0}{的，我，赫拉克勒斯，宙斯之子，做了那样的}
\BookLine{81.00}{22.14}{}{0}{事情！但是有一位神明——某位神明——及时}
\BookLine{89.50}{22.14}{}{0}{介入进来，救了她也救了我。阿尔刻提斯原谅}
\BookLine{98.00}{22.14}{}{0}{了我。这怎么可能？有谁能理解呢？她没有对}
\BookLine{106.50}{22.14}{}{0}{任何人提起过这件事。我再次来到费莱时，在}
\BookLine{115.00}{22.14}{}{0}{她的表情中、眼神里看不到任何我曾经犯过罪}
\BookLine{123.50}{22.14}{}{0}{的痕迹。只有神明才能够理解，只有那样的微}
\BookLine{132.00}{22.14}{}{0}{笑可以解释。宽恕不是我们这些凡人能够做到}
\BookLine{140.50}{22.14}{}{0}{的。只有坚强的、圣洁的人才有能力宽恕。我}
\BookLine{149.00}{22.14}{}{0}{甚至希望她不要忘记那个邪恶的瞬间，不要。}
\BookLine{157.50}{22.14}{}{0}{因为能被她想起是我的幸福，因为她想起我时，}
\BookLine{166.00}{22.14}{}{0}{就已经宽恕了我。}
\newpage
% 原书页码：154
\null
\BookNumber{154}
\Vertical{4}{169}{}{6.2}{8}阿\\尔\\刻\\提\\斯\\之\\歌\EndVertical
\BookLine{30.00}{0.00}{\heifont}{0}{阿德墨托斯：}
\BookLine{38.50}{22.14}{}{0}{赫拉克勒斯，我和你一起去！}
\BookLine{55.50}{0.00}{\heifont}{0}{赫拉克勒斯：}
\BookLine{64.00}{22.14}{}{0}{留下来好好治理国家，阿德墨托斯。你的工作}
\BookLine{72.50}{22.14}{}{0}{我做不了，我的你也做不了。我能指望哪个神}
\BookLine{81.00}{22.14}{}{0}{明呢？肯定不会是我父亲——有和没有一个}
\BookLine{89.50}{22.14}{}{0}{样。你们在这儿都是崇拜哪一位神明？}
\BookLine{106.50}{0.00}{\heifont}{0}{阿格莱亚：}
\BookLine{115.00}{22.14}{}{0}{你忘了吗，赫拉克勒斯。这是阿波罗的领地啊。}
\BookLine{132.00}{0.00}{\heifont}{0}{赫拉克勒斯：}
\BookLine{140.50}{22.14}{}{0}{是，我想起来了。我和他没怎么打过交道，不}
\BookLine{149.00}{22.14}{}{0}{过……{\kaifont （望着天空）}阿波罗，我是赫拉克勒斯，}
\BookLine{157.50}{22.14}{}{0}{传说中宙斯和阿尔克墨涅的儿子。整个希腊都}
\BookLine{166.00}{22.14}{}{0}{知道您一直爱着这两位——阿德墨托斯和阿尔}
\BookLine{174.50}{22.14}{}{0}{刻提斯。您知道我打算做什么，您也知道我一}
\newpage
% 原书页码：155
\null
\BookNumber{155}
\Vertical{129}{11}{}{7}{10}第\\二\\幕\EndVertical
\BookLine{30.00}{22.14}{}{0}{个人做不到。请赐予我从未有过的力量吧，请}
\BookLine{38.50}{22.14}{}{0}{您来完成，或者说，您和我一起完成这件事。}
\BookLine{47.00}{22.14}{}{0}{如果我们能做到的话，人们就会拥有全新的认}
\BookLine{55.50}{22.14}{}{0}{识，他们会看到神明和人类合作起来可以做到}
\BookLine{64.00}{22.14}{}{0}{什么。}
\BookLine{81.00}{0.00}{\heifont}{0}{阿德墨托斯：}
\BookLine{89.50}{22.14}{}{0}{请等一下，赫拉克勒斯。{\kaifont （沿着小道走下来）}佩}
\BookLine{98.00}{22.14}{}{0}{索！佩索！这是赫拉克勒斯，这是阿尔刻提斯}
\BookLine{106.50}{22.14}{}{0}{和我，还有伊皮伊尼兹最亲爱的朋友。请让他}
\BookLine{115.00}{22.14}{}{0}{从这里通过，佩索——从这里去而复返。}
\BookLine{132.00}{0.00}{\heifont}{0}{赫拉克勒斯：}
\BookLine{140.50}{22.14}{}{0}{回宫殿里去吧，阿德墨托斯，你们都回去吧。}
\BookLine{149.00}{22.14}{}{0}{我必须独自做工，如同必须独自生活。}
\BookLine{166.00}{7.38}{\kaifont}{0}{场上一片黑暗，角色们纷纷退场，除了阿德墨托斯还站}
\BookLine{174.50}{0.00}{\kaifont}{0}{在宫殿门口，用斗篷挡着脸。低沉的鼓声响起，像是远方传}
\newpage
% 原书页码：156
\null
\BookNumber{156}
\Vertical{4}{169}{}{6.2}{8}阿\\尔\\刻\\提\\斯\\之\\歌\EndVertical
\BookLine{30.00}{0.00}{\kaifont}{0}{来的雷声。声渐响，鼓点急切且跳跃。赫拉克勒斯从洞穴的}
\BookLine{38.50}{0.00}{\kaifont}{0}{入口处消失，片刻后他再次出现，领着阿尔刻提斯。她的全}
\BookLine{47.00}{0.00}{\kaifont}{0}{身覆着深色的纱巾，在身后拖了很长。快要登上舞台时，赫}
\BookLine{55.50}{0.00}{\kaifont}{0}{拉克勒斯松开了手。她晃晃悠悠地摸索着向前走，像是在梦}
\BookLine{64.00}{0.00}{\kaifont}{0}{游。赫拉克勒斯夹起自己的棍子，独自一人离开。阿德墨托}
\BookLine{72.50}{0.00}{\kaifont}{0}{斯像是要避开眼前的强光一样用斗篷遮住自己。他慢慢靠近，}
\BookLine{81.00}{0.00}{\kaifont}{0}{揭开阿尔刻提斯的面纱。她把头靠在他的胸前，他将她带回}
\BookLine{89.50}{0.00}{\kaifont}{0}{宫殿。}
\newpage
% 原书页码：157
\null
\ActOrnament[三]
\newpage
% 原书页码：158
\null
\newpage
% 原书页码：159
\null
\BookNumber{159}
\Vertical{129}{11}{}{7}{10}第\\三\\幕\EndVertical
\BookLine{30.00}{7.38}{\kaifont}{0}{十二年之后。仍是黎明第一道光。阿尔刻提斯从舞台左}
\BookLine{38.50}{0.00}{\kaifont}{0}{侧上场。她已经老了，身形衰败，穿着破旧的衣服。她提着}
\BookLine{47.00}{0.00}{\kaifont}{0}{一个水罐，走到泉边去打水。路上传来一阵人群的吵闹。}
\BookLine{64.00}{0.00}{\heifont}{0}{居民：}
\BookLine{72.50}{22.14}{}{0}{{\kaifont （站在宫殿的门外）}阿吉斯国王，帮帮我们吧！}
\BookLine{81.00}{22.14}{}{0}{救救我们吧！我们需要国王的帮助……我们要}
\BookLine{89.50}{22.14}{}{0}{和国王谈谈！}
\BookLine{106.50}{7.38}{\kaifont}{0}{一名新来的年轻看门人匆匆忙忙走了进来，看起来很惊}
\BookLine{115.00}{0.00}{\kaifont}{0}{慌。他也提着羊皮灯笼和一只鼓，脸上满是灰尘。他把门开}
\BookLine{123.50}{0.00}{\kaifont}{0}{了一条小缝，透过门缝和外面的人说话。}
\BookLine{140.50}{0.00}{\heifont}{0}{看门人：}
\BookLine{149.00}{22.14}{}{0}{你们知道命令，敢踏进这门一步，就是死罪。}
\BookLine{166.00}{0.00}{\heifont}{0}{居民：}
\BookLine{174.50}{22.14}{}{0}{{\kaifont （从门里挤进来，把看门人甩在身后）}让他杀}
\newpage
% 原书页码：160
\null
\BookNumber{160}
\Vertical{4}{169}{}{6.2}{8}阿\\尔\\刻\\提\\斯\\之\\歌\EndVertical
\BookLine{30.00}{22.14}{}{0}{了我们吧！我们已经活不成了！国王得想想办}
\BookLine{38.50}{22.14}{}{0}{法啊！}
\BookLine{55.50}{0.00}{\heifont}{0}{看门人：}
\BookLine{64.00}{22.14}{}{0}{{\kaifont （害怕地望向宫殿）}国王已经在尽他全力了。国}
\BookLine{72.50}{22.14}{}{0}{王说了，瘟疫和灾难是众神降下的，只有众神}
\BookLine{81.00}{22.14}{}{0}{才能让它们停止。他还说了，你们要把死人火}
\BookLine{89.50}{22.14}{}{0}{化掉，只要他们一死就要尽快火化。用尘土蒙}
\BookLine{98.00}{22.14}{}{0}{起你们的脸，用尘土和灰烬铺满你们的全身，}
\BookLine{106.50}{22.14}{}{0}{国王说了，这样瘟疫就不会找上你们，会绕开}
\BookLine{115.00}{22.14}{}{0}{你们。}
\BookLine{132.00}{0.00}{\heifont}{0}{居民：}
\BookLine{140.50}{22.14}{}{0}{这些我们之前都听过了。我们要见阿尔刻提斯}
\BookLine{149.00}{22.14}{}{0}{王后。阿尔刻提斯王后，她能帮助我们！阿尔}
\BookLine{157.50}{22.14}{}{0}{刻提斯王后！}
\newpage
% 原书页码：161
\null
\BookNumber{161}
\Vertical{129}{11}{}{7}{10}第\\三\\幕\EndVertical
\BookLine{30.00}{0.00}{\heifont}{0}{看门人：}
\BookLine{38.50}{22.14}{}{0}{{\kaifont （愤怒地）}你们说什么，什么阿尔刻提斯王后？}
\BookLine{47.00}{22.14}{}{0}{这儿没有阿尔刻提斯王后。她是这儿的奴隶，}
\BookLine{55.50}{22.14}{}{0}{是最卑贱的奴隶。}
\BookLine{72.50}{7.38}{\kaifont}{0}{阿尔刻提斯从小路走上来，肩上扛着水罐。她站在那里}
\BookLine{81.00}{0.00}{\kaifont}{0}{听大家说话。看门人看到了她。}
\BookLine{98.00}{22.14}{}{0}{她来了！{\kaifont （充满攻击性情绪）}看看她！就是她把}
\BookLine{106.50}{22.14}{}{0}{瘟疫招来的。她就是瘟疫。就是她给塞萨利招}
\BookLine{115.00}{22.14}{}{0}{来了诅咒！}
\BookLine{132.00}{0.00}{\heifont}{0}{居民：}
\BookLine{140.50}{22.14}{}{0}{{\kaifont （先是一阵惊讶的沉默，随即爆发出反对的喊}}
\BookLine{149.00}{22.14}{}{0}{{\kaifont 声）}不，不是这样！不是阿尔刻提斯王后！他}
\BookLine{157.50}{22.14}{}{0}{刚说什么？她怎么会招来瘟疫？}
\newpage
% 原书页码：162
\null
\BookNumber{162}
\Vertical{4}{169}{}{6.2}{8}阿\\尔\\刻\\提\\斯\\之\\歌\EndVertical
\BookLine{30.00}{0.00}{\heifont}{0}{阿尔刻提斯：}
\BookLine{38.50}{22.14}{}{0}{{\kaifont （慢慢地转过头看着他）}我，招来瘟疫？是我招}
\BookLine{47.00}{22.14}{}{0}{来的瘟疫？}
\BookLine{64.00}{0.00}{\heifont}{0}{看门人：}
\BookLine{72.50}{22.14}{}{0}{她死过一次，不是吗？赫拉克勒斯把她从死亡}
\BookLine{81.00}{22.14}{}{0}{的世界带回来，不是吗？她把死亡一起带回来}
\BookLine{89.50}{22.14}{}{0}{了。和她有关系的人都难逃一死，她丈夫被杀}
\BookLine{98.00}{22.14}{}{0}{了，两个孩子被杀了。有个儿子跑得远远的，}
\BookLine{106.50}{22.14}{}{0}{十几年了没有任何消息。国王很快也会把她处}
\BookLine{115.00}{22.14}{}{0}{死或者流放的。}
\BookLine{132.00}{0.00}{\heifont}{0}{阿尔刻提斯：}
\BookLine{140.50}{22.14}{}{0}{不，塞萨利的灾难和我无关。你们可以自己评}
\BookLine{149.00}{22.14}{}{0}{判，如果真是因为我给塞萨利带来了这些不幸，}
\BookLine{157.50}{22.14}{}{0}{那就让他们杀了我，平息这些降在你们和后代}
\BookLine{166.00}{22.14}{}{0}{身上的灾祸。}
\newpage
% 原书页码：163
\null
\BookNumber{163}
\Vertical{129}{11}{}{7}{10}第\\三\\幕\EndVertical
\BookLine{30.00}{7.38}{\kaifont}{0}{阿尔刻提斯向舞台左侧走去，进入仆人们的院子。一束}
\BookLine{38.50}{0.00}{\kaifont}{0}{光渐渐照亮了站在屋顶上的阿波罗。}
\BookLine{55.50}{0.00}{\heifont}{0}{居民：}
\BookLine{64.00}{22.14}{}{0}{不！她才不会给塞萨利带来灾难！她是智慧的}
\BookLine{72.50}{22.14}{}{0}{阿尔刻提斯！}
\BookLine{89.50}{0.00}{\heifont}{0}{看门人：}
\BookLine{98.00}{22.14}{}{0}{{\kaifont （在反对的声音浪潮中高喊着）}回你们的家去}
\BookLine{106.50}{22.14}{}{0}{吧！你们都要死了！{\kaifont （对着站在前面的人说话）}}
\BookLine{115.00}{22.14}{}{0}{你们知道什么？一群无知的农民！一边儿去，}
\BookLine{123.50}{22.14}{}{0}{你们！我已经警告过你们了，别怪阿吉斯国王}
\BookLine{132.00}{22.14}{}{0}{不客气！}
\BookLine{149.00}{7.38}{\kaifont}{0}{镇上的人一边纷纷议论着，一边下场。看门人从左侧下}
\BookLine{157.50}{0.00}{\kaifont}{0}{场，一边走，一边摇着鼓。}
\newpage
% 原书页码：164
\null
\BookNumber{164}
\Vertical{4}{169}{}{6.2}{8}阿\\尔\\刻\\提\\斯\\之\\歌\EndVertical
\BookLine{30.00}{0.00}{\heifont}{0}{看门人：}
\BookLine{38.50}{22.14}{}{0}{黎明时分，照例巡查。塞萨利之王阿吉斯，千}
\BookLine{47.00}{22.14}{}{0}{金万马的拥有者。宫中万事顺遂。}
\BookLine{64.00}{7.38}{\kaifont}{0}{死神从洞穴中钻出来，慢慢爬到小路上，在门边四处}
\BookLine{72.50}{0.00}{\kaifont}{0}{嗅着。}
\BookLine{89.50}{0.00}{\heifont}{0}{阿波罗：}
\BookLine{98.00}{22.14}{}{0}{死神！}
\BookLine{115.00}{0.00}{\heifont}{0}{死神：}
\BookLine{123.50}{22.14}{}{0}{{\kaifont （没有注意到阿波罗，假装不在意）}啊，你又}
\BookLine{132.00}{22.14}{}{0}{来了！}
\BookLine{149.00}{0.00}{\heifont}{0}{阿波罗：}
\BookLine{157.50}{22.14}{}{0}{天越来越亮了，你在发抖呢。}
\newpage
% 原书页码：165
\null
\BookNumber{165}
\Vertical{129}{11}{}{7}{10}第\\三\\幕\EndVertical
\BookLine{30.00}{0.00}{\heifont}{0}{死神：}
\BookLine{38.50}{22.14}{}{0}{是，是，不过我有几个问题想问你。阿波罗大}
\BookLine{47.00}{22.14}{}{0}{人，我不明白，你干这些事有什么目的，死了}
\BookLine{55.50}{22.14}{}{0}{这么多人！我的地盘上现在都挤成一团，排起}
\BookLine{64.00}{22.14}{}{0}{好长的队啦！而且，我之前从没见过这么多孩}
\BookLine{72.50}{22.14}{}{0}{子！必须得承认，阿波罗大人，我不知道你这}
\BookLine{81.00}{22.14}{}{0}{是要干什么，你可是治愈的神啊，是生命与治}
\BookLine{89.50}{22.14}{}{0}{愈的神。结果你来了，带来的却是瘟疫和灾难。}
\BookLine{98.00}{22.14}{}{0}{你宠爱阿德墨托斯，爱他的家人，他的百姓，}
\BookLine{106.50}{22.14}{}{0}{然后，你把他们都杀了。}
\BookLine{123.50}{0.00}{\heifont}{0}{阿波罗：}
\BookLine{132.00}{22.14}{}{0}{{\kaifont （微笑）}我爱着阿尔刻提斯，但我杀了她——}
\BookLine{140.50}{22.14}{}{0}{一次。}
\BookLine{157.50}{0.00}{\heifont}{0}{死神：}
\BookLine{166.00}{22.14}{}{0}{多没风度啊，看看你干的事！}
\newpage
% 原书页码：166
\null
\BookNumber{166}
\Vertical{4}{169}{}{6.2}{8}阿\\尔\\刻\\提\\斯\\之\\歌\EndVertical
\BookLine{30.00}{0.00}{\heifont}{0}{阿波罗：}
\BookLine{38.50}{22.14}{}{0}{你的那堵墙修好了吗？就是赫拉克勒斯撞坏的}
\BookLine{47.00}{22.14}{}{0}{那堵墙？}
\BookLine{64.00}{0.00}{\heifont}{0}{死神：}
\BookLine{72.50}{22.14}{}{0}{撞坏？赫拉克勒斯？是你弄坏的。是你打破了}
\BookLine{81.00}{22.14}{}{0}{最古老的规矩和这世界的规律：生是生，死}
\BookLine{89.50}{22.14}{}{0}{是死。}
\BookLine{106.50}{0.00}{\heifont}{0}{阿波罗：}
\BookLine{115.00}{22.14}{}{0}{嗯——有一道微弱的光芒，照进了从未照耀到}
\BookLine{123.50}{22.14}{}{0}{的地方。}
\BookLine{140.50}{0.00}{\heifont}{0}{死神：}
\BookLine{149.00}{22.14}{}{0}{你坏了规矩，现在得来承担后果了。现在你的}
\BookLine{157.50}{22.14}{}{0}{这些幸福和引以为傲的智慧都不顶用了，因为}
\BookLine{166.00}{22.14}{}{0}{你还是没法让他们认识到你的存在。然后呢，}
\BookLine{174.50}{22.14}{}{0}{你就像一个被抛弃的情人一样，跑到爱人家里，}
\newpage
% 原书页码：167
\null
\BookNumber{167}
\Vertical{129}{11}{}{7}{10}第\\三\\幕\EndVertical
\BookLine{30.00}{22.14}{}{0}{把所有人都杀光。}
\BookLine{47.00}{0.00}{\heifont}{0}{阿波罗：}
\BookLine{55.50}{22.14}{}{0}{我已经让他们认识到我了。我已经开启了我的}
\BookLine{64.00}{22.14}{}{0}{故事。}
\BookLine{81.00}{0.00}{\heifont}{0}{死神：}
\BookLine{89.50}{22.14}{}{0}{你的所谓教训。}
\BookLine{106.50}{0.00}{\heifont}{0}{阿波罗：}
\BookLine{115.00}{22.14}{}{0}{对，我的教训——那些愿意为他人失去生命}
\BookLine{123.50}{22.14}{}{0}{的，我将使他们重新获得生命。}
\BookLine{140.50}{0.00}{\heifont}{0}{死神：}
\BookLine{149.00}{22.14}{}{0}{可你只带回来了一个，却把千千万万个送给}
\BookLine{157.50}{22.14}{}{0}{了我。}
\newpage
% 原书页码：168
\null
\BookNumber{168}
\Vertical{4}{169}{}{6.2}{8}阿\\尔\\刻\\提\\斯\\之\\歌\EndVertical
\BookLine{30.00}{0.00}{\heifont}{0}{阿波罗：}
\BookLine{38.50}{22.14}{}{0}{是的，必须让毁灭随之而来，只有这样，他们}
\BookLine{47.00}{22.14}{}{0}{才会记住这个故事。在那些能被永远记住的故}
\BookLine{55.50}{22.14}{}{0}{事里，死亡总是重头戏。}
\BookLine{72.50}{0.00}{\heifont}{0}{死神：}
\BookLine{81.00}{22.14}{}{0}{{\kaifont （发抖）}那道光芒！光芒！我对这些瘟疫什么的}
\BookLine{89.50}{22.14}{}{0}{不感兴趣。我只想等一个人类，她不会再从我}
\BookLine{98.00}{22.14}{}{0}{手里溜走了。}
\BookLine{115.00}{0.00}{\heifont}{0}{阿波罗：}
\BookLine{123.50}{22.14}{}{0}{阿尔刻提斯？你永远别想得到她。}
\BookLine{140.50}{0.00}{\heifont}{0}{死神：}
\BookLine{149.00}{22.14}{}{0}{她是个人！}
\BookLine{166.00}{0.00}{\heifont}{0}{阿波罗：}
\BookLine{174.50}{22.14}{}{0}{是的。}
\newpage
% 原书页码：169
\null
\BookNumber{169}
\Vertical{129}{11}{}{7}{10}第\\三\\幕\EndVertical
\BookLine{30.00}{0.00}{\heifont}{0}{死神：}
\BookLine{38.50}{22.14}{}{0}{她是，一个，人！}
\BookLine{55.50}{0.00}{\heifont}{0}{阿波罗：}
\BookLine{64.00}{22.14}{}{0}{人，只是个人。}
\BookLine{81.00}{0.00}{\heifont}{0}{死神：}
\BookLine{89.50}{22.14}{}{0}{你想干什么？你休想再把她从我这儿偷走一次。}
\BookLine{106.50}{0.00}{\heifont}{0}{阿波罗：}
\BookLine{115.00}{22.14}{}{0}{死神，太阳升起来了。你在发抖呢。}
\BookLine{132.00}{0.00}{\heifont}{0}{死神：}
\BookLine{140.50}{22.14}{}{0}{是，你告诉我答案，我马上就走。}
\BookLine{157.50}{0.00}{\heifont}{0}{阿波罗：}
\BookLine{166.00}{22.14}{}{0}{你要学会慢慢适应变化。}
\newpage
% 原书页码：170
\null
\BookNumber{170}
\Vertical{4}{169}{}{6.2}{8}阿\\尔\\刻\\提\\斯\\之\\歌\EndVertical
\BookLine{30.00}{0.00}{\heifont}{0}{死神：}
\BookLine{38.50}{22.14}{}{0}{我？}
\BookLine{55.50}{0.00}{\heifont}{0}{阿波罗：}
\BookLine{64.00}{22.14}{}{0}{有一道光已经照进你的王国了。}
\BookLine{81.00}{0.00}{\heifont}{0}{死神：}
\BookLine{89.50}{22.14}{}{0}{{\kaifont （迅速地回到他的洞中，尖叫着）}不会再有这种}
\BookLine{98.00}{22.14}{}{0}{事儿，不会。不会再有光进来，不会，不会有}
\BookLine{106.50}{22.14}{}{0}{什么变化！}
\BookLine{123.50}{7.38}{\kaifont}{0}{死神消失，阿波罗身上的光也褪去。二十一岁的伊皮伊}
\BookLine{132.00}{0.00}{\kaifont}{0}{尼兹和同样二十一岁的切瑞安德从宫殿里走出来。他们用斗}
\BookLine{140.50}{0.00}{\kaifont}{0}{篷遮住自己的鼻子，可以看到身上佩带着短剑。切瑞安德走}
\BookLine{149.00}{0.00}{\kaifont}{0}{在前面，带着好奇和畏惧四处看着，伊皮伊尼兹跟在他身后，}
\BookLine{157.50}{0.00}{\kaifont}{0}{闷闷不乐。}
\newpage
% 原书页码：171
\null
\BookNumber{171}
\Vertical{129}{11}{}{7}{10}第\\三\\幕\EndVertical
\BookLine{30.00}{0.00}{\heifont}{0}{切瑞安德：}
\BookLine{38.50}{22.14}{}{0}{这就是阿德墨托斯和阿尔刻提斯的宫殿啊！}
\BookLine{47.00}{22.14}{}{0}{阿波罗也在这儿待过对吧？他以前站在哪}
\BookLine{55.50}{22.14}{}{0}{里？——这儿，还是这儿？}
\BookLine{72.50}{7.38}{\kaifont}{0}{伊皮伊尼兹点头，一直低着头很少抬眼。}
\BookLine{89.50}{0.00}{\heifont}{0}{切瑞安德：}
\BookLine{98.00}{22.14}{}{0}{赫拉克勒斯把你母亲从死神那儿带了回来，那}
\BookLine{106.50}{22.14}{}{0}{是在哪儿？}
\BookLine{123.50}{7.38}{\kaifont}{0}{伊皮伊尼兹指了指洞穴。切瑞安德快步跑过去。}
\BookLine{140.50}{0.00}{\heifont}{0}{切瑞安德：}
\BookLine{149.00}{22.14}{}{0}{我竟然亲眼看到了这些！你看，伊皮伊尼兹！}
\BookLine{157.50}{22.14}{}{0}{你的老朋友佩索还在看着门呢。快跟他打个招}
\BookLine{166.00}{22.14}{}{0}{呼，我们喝一点泉水。这儿的泉水应该还没被}
\BookLine{174.50}{22.14}{}{0}{污染吧。}
\newpage
% 原书页码：172
\null
\BookNumber{172}
\Vertical{4}{169}{}{6.2}{8}阿\\尔\\刻\\提\\斯\\之\\歌\EndVertical
\BookLine{30.00}{0.00}{\heifont}{0}{伊皮伊尼兹：}
\BookLine{38.50}{22.14}{}{0}{佩索，你还记得我吗？他都不动弹。}
\BookLine{55.50}{0.00}{\heifont}{0}{切瑞安德：}
\BookLine{64.00}{22.14}{}{0}{你做一个祷告仪式，然后我们就可以喝了。}
\BookLine{81.00}{0.00}{\heifont}{0}{伊皮伊尼兹：}
\BookLine{89.50}{22.14}{}{0}{“生命之源——大地，空气……”不，切瑞安}
\BookLine{98.00}{22.14}{}{0}{德，我做不到。以前，父亲和母亲每天早上都}
\BookLine{106.50}{22.14}{}{0}{会带我们到这儿来，这个祷告我曾经做过上百}
\BookLine{115.00}{22.14}{}{0}{次。可是怎么会这样呢？神明就这样抛下我们}
\BookLine{123.50}{22.14}{}{0}{不管了？我的父亲被杀了，弟弟妹妹也被杀了，}
\BookLine{132.00}{22.14}{}{0}{而我在一个深夜被流放，只能长久与陌生人为}
\BookLine{140.50}{22.14}{}{0}{伍，我母亲成了奴隶，或许早已经死了。我的}
\BookLine{149.00}{22.14}{}{0}{国家在遭受瘟疫，阳光下到处都是死尸。我还}
\BookLine{157.50}{22.14}{}{0}{要向谁祷告？}
\newpage
% 原书页码：173
\null
\BookNumber{173}
\Vertical{129}{11}{}{7}{10}第\\三\\幕\EndVertical
\BookLine{30.00}{0.00}{\heifont}{0}{切瑞安德：}
\BookLine{38.50}{22.14}{}{0}{{\kaifont （沉默着做了决定）}我们这就离开费莱。现在就}
\BookLine{47.00}{22.14}{}{0}{从这扇门出去。关于我们计划的事情，我想你}
\BookLine{55.50}{22.14}{}{0}{还没准备好，我也不会帮你了。}
\BookLine{72.50}{0.00}{\heifont}{0}{伊皮伊尼兹：}
\BookLine{81.00}{22.14}{}{0}{只要不做这些祷告和仪式，我会做到答应自己}
\BookLine{89.50}{22.14}{}{0}{的事。为了正义和复仇，一剑刺下去！往他的}
\BookLine{98.00}{22.14}{}{0}{喉咙刺上一剑。是的，还要给他的女儿一剑，}
\BookLine{106.50}{22.14}{}{0}{劳德米娅也跑不了。}
\BookLine{123.50}{0.00}{\heifont}{0}{切瑞安德：}
\BookLine{132.00}{22.14}{}{0}{伊皮伊尼兹，我不会帮你的。我们可不是这么}
\BookLine{140.50}{22.14}{}{0}{计划的。要杀人还是要送死你随便，但我大老}
\BookLine{149.00}{22.14}{}{0}{远跑到这儿可不是为了加入你的野蛮屠杀的。}
\BookLine{166.00}{0.00}{\heifont}{0}{伊皮伊尼兹：}
\BookLine{174.50}{22.14}{}{0}{你想让我祷告？可我对着谁祷告？对着阿波}
\newpage
% 原书页码：174
\null
\BookNumber{174}
\Vertical{4}{169}{}{6.2}{8}阿\\尔\\刻\\提\\斯\\之\\歌\EndVertical
\BookLine{30.00}{22.14}{}{0}{罗？对着这个带着仇恨摧毁了整个家的阿}
\BookLine{38.50}{22.14}{}{0}{波罗？}
\BookLine{55.50}{7.38}{\kaifont}{0}{切瑞安德俯下身，把双手放在伊皮伊尼兹的肩上晃着他}
\BookLine{64.00}{0.00}{\kaifont}{0}{的肩膀，严肃而有些愤怒。}
\BookLine{81.00}{0.00}{\heifont}{0}{切瑞安德：}
\BookLine{89.50}{22.14}{}{0}{你以为你是第一个遭遇苦难的人吗？残暴，不}
\BookLine{98.00}{22.14}{}{0}{公，谋杀，羞辱，你以为这些事情只有在这儿}
\BookLine{106.50}{22.14}{}{0}{才会发生吗？你难道忘了吗，我们来这儿是为}
\BookLine{115.00}{22.14}{}{0}{了建立公义？是你自己说的，人类仅凭自身绝}
\BookLine{123.50}{22.14}{}{0}{不可能达到公平与正义，这是众神想要我们做}
\BookLine{132.00}{22.14}{}{0}{到的事情。}
\BookLine{149.00}{7.38}{\kaifont}{0}{他们坚定地互相望了望。}
\BookLine{166.00}{0.00}{\heifont}{0}{伊皮伊尼兹：}
\BookLine{174.50}{22.14}{}{0}{{\kaifont （安静地低着头）}有你这样的朋友真是太难}
\newpage
% 原书页码：175
\null
\BookNumber{175}
\Vertical{129}{11}{}{7}{10}第\\三\\幕\EndVertical
\BookLine{30.00}{22.14}{}{0}{得了。}
\BookLine{47.00}{7.38}{\kaifont}{0}{阿尔刻提斯从左侧上场，肩上扛着水罐。}
\BookLine{64.00}{0.00}{\heifont}{0}{切瑞安德：}
\BookLine{72.50}{22.14}{}{0}{你做好祷告的准备了吗？}
\BookLine{89.50}{7.38}{\kaifont}{0}{伊皮伊尼兹庄重地重复着仪式。不过一直没有出声。两}
\BookLine{98.00}{0.00}{\kaifont}{0}{人喝水，切瑞安德沿着小路走上去，十分兴奋。}
\BookLine{115.00}{0.00}{\heifont}{0}{切瑞安德：}
\BookLine{123.50}{22.14}{}{0}{我准备去敲门了。{\kaifont （看见阿尔刻提斯，匆忙走}}
\BookLine{132.00}{22.14}{}{0}{{\kaifont 回泉边，对伊皮伊尼兹说道）}这儿有个老妇人，}
\BookLine{140.50}{22.14}{}{0}{我们可以问问看有什么消息。}
\BookLine{157.50}{7.38}{\kaifont}{0}{阿尔刻提斯沿着小路走了几步，看到这两人后开始往回}
\BookLine{166.00}{0.00}{\kaifont}{0}{走。切瑞安德迎着她走上去。}
\newpage
% 原书页码：176
\null
\BookNumber{176}
\Vertical{4}{169}{}{6.2}{8}阿\\尔\\刻\\提\\斯\\之\\歌\EndVertical
\BookLine{30.00}{0.00}{\heifont}{0}{切瑞安德：}
\BookLine{38.50}{22.14}{}{0}{老婆婆，这是塞萨利国王阿吉斯的王宫吗？}
\BookLine{55.50}{7.38}{\kaifont}{0}{阿尔刻提斯的目光望向大门，切瑞安德回到伊皮伊尼兹}
\BookLine{64.00}{0.00}{\kaifont}{0}{身旁。}
\BookLine{81.00}{22.14}{}{0}{看上去是个奴隶。也许，是你的母亲也说不定。}
\BookLine{98.00}{0.00}{\heifont}{0}{伊皮伊尼兹：}
\BookLine{106.50}{22.14}{}{0}{{\kaifont （快步沿着小路走上去，凑上去看了看阿尔刻提}}
\BookLine{115.00}{22.14}{}{0}{{\kaifont 斯，直截了当地说道）}不是。}
\BookLine{132.00}{0.00}{\heifont}{0}{切瑞安德：}
\BookLine{140.50}{22.14}{}{0}{你确定吗？}
\BookLine{157.50}{0.00}{\heifont}{0}{伊皮伊尼兹：}
\BookLine{166.00}{22.14}{}{0}{{\kaifont （爽快地）}确定，肯定。}
\newpage
% 原书页码：177
\null
\BookNumber{177}
\Vertical{129}{11}{}{7}{10}第\\三\\幕\EndVertical
\BookLine{30.00}{0.00}{\heifont}{0}{切瑞安德：}
\BookLine{38.50}{22.14}{}{0}{老婆婆，跟我说说，这宫里有人染上瘟疫而死}
\BookLine{47.00}{22.14}{}{0}{吗？{\kaifont （她摇摇头）}阿吉斯国王和他的家人住在这}
\BookLine{55.50}{22.14}{}{0}{里，他有个小女儿劳德米娅，对吗？{\kaifont （她点点}}
\BookLine{64.00}{22.14}{}{0}{{\kaifont 头）}再跟我说说，他的这些守卫，是从色雷斯}
\BookLine{72.50}{22.14}{}{0}{带来的，还是本国人？{\kaifont （没有回答）}}
\BookLine{89.50}{0.00}{\heifont}{0}{伊皮伊尼兹：}
\BookLine{98.00}{22.14}{}{0}{说不定她是个聋子，或者是个哑巴。{\kaifont （大声地）}}
\BookLine{106.50}{22.14}{}{0}{他经常回自己的国家，回色雷斯吗？}
\BookLine{123.50}{0.00}{\heifont}{0}{阿尔刻提斯：}
\BookLine{132.00}{22.14}{}{0}{你们在这儿很危险，你们必须离开，马上离开}
\BookLine{140.50}{22.14}{}{0}{这儿。}
\BookLine{157.50}{0.00}{\heifont}{0}{伊皮伊尼兹：}
\BookLine{166.00}{22.14}{}{0}{哦，她能说话。老婆婆，跟我说说，阿德墨托}
\BookLine{174.50}{22.14}{}{0}{斯国王和阿尔刻提斯王后在这里的时候，你}
\newpage
% 原书页码：178
\null
\BookNumber{178}
\Vertical{4}{169}{}{6.2}{8}阿\\尔\\刻\\提\\斯\\之\\歌\EndVertical
\BookLine{30.00}{22.14}{}{0}{在吗？}
\BookLine{47.00}{0.00}{\heifont}{0}{阿尔刻提斯：}
\BookLine{55.50}{22.14}{}{0}{你们是谁？你们从哪儿来的？}
\BookLine{72.50}{0.00}{\heifont}{0}{伊皮伊尼兹：}
\BookLine{81.00}{22.14}{}{0}{你认识他们吗？你和他们说过话吗？{\kaifont （她点点}}
\BookLine{89.50}{22.14}{}{0}{{\kaifont 头）}阿德墨托斯国王死了？{\kaifont （她点点头）}阿尔刻}
\BookLine{98.00}{22.14}{}{0}{提斯王后还活着吗？{\kaifont （没有回应）}她还在宫里？}
\BookLine{115.00}{7.38}{\kaifont}{0}{阿尔刻提斯的目光再次望向宫门，伊皮伊尼兹很不耐烦。}
\BookLine{132.00}{0.00}{\heifont}{0}{伊皮伊尼兹：}
\BookLine{140.50}{22.14}{}{0}{凭着不朽的神明起誓，你既然会说话，快}
\BookLine{149.00}{22.14}{}{0}{说呀！}
\BookLine{166.00}{0.00}{\heifont}{0}{阿尔刻提斯：}
\BookLine{174.50}{22.14}{}{0}{你们必须马上离开。不过你告诉我，告诉我，}
\newpage
% 原书页码：179
\null
\BookNumber{179}
\Vertical{129}{11}{}{7}{10}第\\三\\幕\EndVertical
\BookLine{30.00}{22.14}{}{0}{你是谁？}
\BookLine{47.00}{0.00}{\heifont}{0}{切瑞安德：}
\BookLine{55.50}{22.14}{}{0}{{\kaifont （盯着阿尔刻提斯看了看，急切地拍着伊皮伊尼}}
\BookLine{64.00}{22.14}{}{0}{{\kaifont 兹的胳膊）}你再看看，仔细看看。你确定吗？}
\BookLine{81.00}{0.00}{\heifont}{0}{伊皮伊尼兹：}
\BookLine{89.50}{22.14}{}{0}{确定！这些山区的女人都长一个样。她们都是}
\BookLine{98.00}{22.14}{}{0}{这么少言寡语，神秘兮兮的。}
\BookLine{115.00}{0.00}{\heifont}{0}{阿尔刻提斯：}
\BookLine{123.50}{22.14}{}{0}{你们是怎么从那个门进来的？阿吉斯国王要是}
\BookLine{132.00}{22.14}{}{0}{知道了，一定会杀了你们。}
\BookLine{149.00}{0.00}{\heifont}{0}{切瑞安德：}
\BookLine{157.50}{22.14}{}{0}{不会的，我们来给阿吉斯国王送个信，他听到}
\BookLine{166.00}{22.14}{}{0}{应该会很高兴。}
\newpage
% 原书页码：180
\null
\BookNumber{180}
\Vertical{4}{169}{}{6.2}{8}阿\\尔\\刻\\提\\斯\\之\\歌\EndVertical
\BookLine{30.00}{0.00}{\heifont}{0}{阿尔刻提斯：}
\BookLine{38.50}{22.14}{}{0}{{\kaifont （不停地摇头）}年轻人，没什么好消息能救你们}
\BookLine{47.00}{22.14}{}{0}{的命。快走吧，快点上路朝北走。}
\BookLine{64.00}{0.00}{\heifont}{0}{伊皮伊尼兹：}
\BookLine{72.50}{22.14}{}{0}{我们带来的这个消息会让国王好好款待我们的。}
\BookLine{81.00}{22.14}{}{0}{我们到这儿来是要告诉他，他的敌人，他最大}
\BookLine{89.50}{22.14}{}{0}{的敌人死了。}
\BookLine{106.50}{0.00}{\heifont}{0}{切瑞安德：}
\BookLine{115.00}{22.14}{}{0}{不用现在说这个。}
\BookLine{132.00}{0.00}{\heifont}{0}{伊皮伊尼兹：}
\BookLine{140.50}{22.14}{}{0}{让他提心吊胆的那个人死了。}
\BookLine{157.50}{7.38}{\kaifont}{0}{阿尔刻提斯摇着头，从他身边走过，沿着小路往下。}
\newpage
% 原书页码：181
\null
\BookNumber{181}
\Vertical{129}{11}{}{7}{10}第\\三\\幕\EndVertical
\BookLine{30.00}{0.00}{\heifont}{0}{伊皮伊尼兹：}
\BookLine{38.50}{22.14}{}{0}{我们到这儿来是要告诉他，伊皮伊尼兹，阿德}
\BookLine{47.00}{22.14}{}{0}{墨托斯国王和阿尔刻提斯王后的儿子，已经死}
\BookLine{55.50}{22.14}{}{0}{了。他听了一定会非常非常高兴的。}
\BookLine{72.50}{7.38}{\kaifont}{0}{阿尔刻提斯在泉边突然停下来，她抬起头，把手放在心}
\BookLine{81.00}{0.00}{\kaifont}{0}{口，水罐滑落在一旁。}
\BookLine{98.00}{0.00}{\heifont}{0}{切瑞安德：}
\BookLine{106.50}{22.14}{}{0}{{\kaifont （走到她身边拿过她的水罐）}老婆婆，我来帮你}
\BookLine{115.00}{22.14}{}{0}{拿吧。}
\BookLine{132.00}{0.00}{\heifont}{0}{阿尔刻提斯：}
\BookLine{140.50}{22.14}{}{0}{{\kaifont （突然说道）}你就是伊皮伊尼兹！}
\BookLine{157.50}{0.00}{\heifont}{0}{切瑞安德：}
\BookLine{166.00}{22.14}{}{0}{{\kaifont （笑了一下）}不是我，老婆婆。伊皮伊尼兹已经}
\BookLine{174.50}{22.14}{}{0}{死了。}
\newpage
% 原书页码：182
\null
\BookNumber{182}
\Vertical{4}{169}{}{6.2}{8}阿\\尔\\刻\\提\\斯\\之\\歌\EndVertical
\BookLine{30.00}{0.00}{\heifont}{0}{阿尔刻提斯：}
\BookLine{38.50}{22.14}{}{0}{你有证据吗？}
\BookLine{55.50}{0.00}{\heifont}{0}{切瑞安德：}
\BookLine{64.00}{22.14}{}{0}{证据，有。}
\BookLine{81.00}{7.38}{\kaifont}{0}{他往罐子里装满水，沿着路往回走。}
\BookLine{98.00}{0.00}{\heifont}{0}{阿尔刻提斯：}
\BookLine{106.50}{22.14}{}{0}{我们一直以为有一天，他会悄悄地，或是假扮}
\BookLine{115.00}{22.14}{}{0}{成另一个人回来。我的眼睛已经不行了。{\kaifont （渴望}}
\BookLine{123.50}{22.14}{}{0}{{\kaifont 地）}你是伊皮伊尼兹！}
\BookLine{140.50}{0.00}{\heifont}{0}{切瑞安德：}
\BookLine{149.00}{22.14}{}{0}{不是，老婆婆，我不是。}
\BookLine{166.00}{0.00}{\heifont}{0}{阿尔刻提斯：}
\BookLine{174.50}{22.14}{}{0}{你认识他，你和他说过话？}
\newpage
% 原书页码：183
\null
\BookNumber{183}
\Vertical{129}{11}{}{7}{10}第\\三\\幕\EndVertical
\BookLine{30.00}{0.00}{\heifont}{0}{切瑞安德：}
\BookLine{38.50}{22.14}{}{0}{是的，我们常常一起聊天。我帮你把这个送到}
\BookLine{47.00}{22.14}{}{0}{哪里？}
\BookLine{64.00}{0.00}{\heifont}{0}{阿尔刻提斯：}
\BookLine{72.50}{22.14}{}{0}{{\kaifont （站在小路的最上方，盯着伊皮伊尼兹的脸看）}}
\BookLine{81.00}{22.14}{}{0}{让我看看。{\kaifont （接着又盯着切瑞安德的脸）}不要对}
\BookLine{89.50}{22.14}{}{0}{一个老人撒谎。他什么消息都没留下？}
\BookLine{106.50}{0.00}{\heifont}{0}{伊皮伊尼兹：}
\BookLine{115.00}{22.14}{}{0}{什么都没有。}
\BookLine{132.00}{0.00}{\heifont}{0}{阿尔刻提斯：}
\BookLine{140.50}{22.14}{}{0}{什么都没有？}
\BookLine{157.50}{0.00}{\heifont}{0}{伊皮伊尼兹：}
\BookLine{166.00}{22.14}{}{0}{只有这个，你看到这条腰带了吗？这是伊皮伊}
\BookLine{174.50}{22.14}{}{0}{尼兹出生前，阿尔刻提斯王后为阿德墨托斯国}
\newpage
% 原书页码：184
\null
\BookNumber{184}
\Vertical{4}{169}{}{6.2}{8}阿\\尔\\刻\\提\\斯\\之\\歌\EndVertical
\BookLine{30.00}{22.14}{}{0}{王织的。}
\BookLine{47.00}{0.00}{\heifont}{0}{阿尔刻提斯：}
\BookLine{55.50}{22.14}{}{0}{{\kaifont （看了一眼，突然哭了起来）}把这条腰带给你的}
\BookLine{64.00}{22.14}{}{0}{年轻人，他是什么模样？}
\BookLine{81.00}{0.00}{\heifont}{0}{伊皮伊尼兹：}
\BookLine{89.50}{22.14}{}{0}{{\kaifont （极力掩饰着自己的不耐烦，像是对一个聋子说}}
\BookLine{98.00}{22.14}{}{0}{{\kaifont 话一样）}老人家，阿尔刻提斯王后在这儿吗？}
\BookLine{106.50}{22.14}{}{0}{您就回答一下，在，还是不在。真让人受不了。}
\BookLine{115.00}{22.14}{}{0}{您能去叫一下阿尔刻提斯王后吗，能不能？}
\BookLine{132.00}{0.00}{\heifont}{0}{阿尔刻提斯：}
\BookLine{140.50}{22.14}{}{0}{我就是阿尔刻提斯。}
\BookLine{157.50}{0.00}{\heifont}{0}{切瑞安德：}
\BookLine{166.00}{22.14}{}{0}{{\kaifont （先是惊讶地沉默了一下，随即单膝跪地，望着}}
\BookLine{174.50}{22.14}{}{0}{{\kaifont 她的脸）}您……您是阿尔刻提斯？}
\newpage
% 原书页码：185
\null
\BookNumber{185}
\Vertical{129}{11}{}{7}{10}第\\三\\幕\EndVertical
\BookLine{30.00}{0.00}{\heifont}{0}{伊皮伊尼兹：}
\BookLine{38.50}{22.14}{}{0}{{\kaifont （愣在原地，一拳打在额头上）}我简直没有脸面}
\BookLine{47.00}{22.14}{}{0}{说出我的名字。{\kaifont （他也单膝跪地）}我是伊皮伊}
\BookLine{55.50}{22.14}{}{0}{尼兹。}
\BookLine{72.50}{0.00}{\heifont}{0}{阿尔刻提斯：}
\BookLine{81.00}{22.14}{}{0}{好，好啊。{\kaifont （目光从宫殿转了回来）}但现在太危}
\BookLine{89.50}{22.14}{}{0}{险了……太危险了……}
\BookLine{106.50}{0.00}{\heifont}{0}{伊皮伊尼兹：}
\BookLine{115.00}{22.14}{}{0}{请……原谅我。}
\BookLine{132.00}{0.00}{\heifont}{0}{阿尔刻提斯：}
\BookLine{140.50}{22.14}{}{0}{{\kaifont （拍了拍他的头）}你真是个又急躁又自我的年轻}
\BookLine{149.00}{22.14}{}{0}{人，就像我以前，是个又急躁又自我的姑娘。}
\BookLine{157.50}{22.14}{}{0}{你应该多继承你父亲一些。}
\newpage
% 原书页码：186
\null
\BookNumber{186}
\Vertical{4}{169}{}{6.2}{8}阿\\尔\\刻\\提\\斯\\之\\歌\EndVertical
\BookLine{30.00}{0.00}{\heifont}{0}{伊皮伊尼兹：}
\BookLine{38.50}{22.14}{}{0}{您的不幸……}
\BookLine{55.50}{0.00}{\heifont}{0}{阿尔刻提斯：}
\BookLine{64.00}{22.14}{}{0}{{\kaifont （打断他）}不，不，伊皮伊尼兹，不要说什么}
\BookLine{72.50}{22.14}{}{0}{不幸。}
\BookLine{89.50}{0.00}{\heifont}{0}{伊皮伊尼兹：}
\BookLine{98.00}{22.14}{}{0}{这些悲剧……}
\BookLine{115.00}{0.00}{\heifont}{0}{阿尔刻提斯：}
\BookLine{123.50}{22.14}{}{0}{不！当你遇到不幸的时候，要学会理解它。这}
\BookLine{132.00}{22.14}{}{0}{世上只有一种悲剧，那就是无知，对我们生命}
\BookLine{140.50}{22.14}{}{0}{的无知。那才是悲剧，和绝望。}
\BookLine{157.50}{0.00}{\heifont}{0}{切瑞安德：}
\BookLine{166.00}{22.14}{}{0}{{\kaifont （起身，询问道）}无知？}
\newpage
% 原书页码：187
\null
\BookNumber{187}
\Vertical{129}{11}{}{7}{10}第\\三\\幕\EndVertical
\BookLine{30.00}{0.00}{\heifont}{0}{阿尔刻提斯：}
\BookLine{38.50}{22.14}{}{0}{我曾经非常非常幸福，是的……难道我现在就}
\BookLine{47.00}{22.14}{}{0}{忘了吗？忘了赐予我幸福的那一位吗？正是曾}
\BookLine{55.50}{22.14}{}{0}{经给予我幸福的那一位，给予了我现在这一切。}
\BookLine{64.00}{22.14}{}{0}{我相信这些不幸也有其益处，这只是故事的一}
\BookLine{72.50}{22.14}{}{0}{部分，而我并没有看到所有……你们现在必须}
\BookLine{81.00}{22.14}{}{0}{马上离开。}
\BookLine{98.00}{0.00}{\heifont}{0}{伊皮伊尼兹：}
\BookLine{106.50}{22.14}{}{0}{{\kaifont （起身）}我们来这儿是要杀了阿吉斯国王，重新}
\BookLine{115.00}{22.14}{}{0}{夺回王位。}
\BookLine{132.00}{0.00}{\heifont}{0}{阿尔刻提斯：}
\BookLine{140.50}{22.14}{}{0}{{\kaifont （一边摇头一边小声说着）}不，不。}
\BookLine{157.50}{0.00}{\heifont}{0}{伊皮伊尼兹：}
\BookLine{166.00}{22.14}{}{0}{我们已经制订了万无一失的计划。就在}
\BookLine{174.50}{22.14}{}{0}{今晚……}
\newpage
% 原书页码：188
\null
\BookNumber{188}
\Vertical{4}{169}{}{6.2}{8}阿\\尔\\刻\\提\\斯\\之\\歌\EndVertical
\BookLine{30.00}{0.00}{\heifont}{0}{阿尔刻提斯：}
\BookLine{38.50}{22.14}{}{0}{不。不要这样，伊皮伊尼兹。这场灾难，这场}
\BookLine{47.00}{22.14}{}{0}{瘟疫已经够了。阿吉斯国王现在只关心一件事，}
\BookLine{55.50}{22.14}{}{0}{只担心一件事，他不在乎自己，只在乎他的孩}
\BookLine{64.00}{22.14}{}{0}{子，劳德米娅。神明会用他自己的方式解决这}
\BookLine{72.50}{22.14}{}{0}{件事的。}
\BookLine{89.50}{0.00}{\heifont}{0}{切瑞安德：}
\BookLine{98.00}{22.14}{}{0}{我们就照她说的办。伊皮伊尼兹，咱们现在}
\BookLine{106.50}{22.14}{}{0}{就走。}
\BookLine{123.50}{0.00}{\heifont}{0}{阿尔刻提斯：}
\BookLine{132.00}{22.14}{}{0}{十天之后，再回来。到北方去，用斗篷遮住你}
\BookLine{140.50}{22.14}{}{0}{的脸，到北方去。}
\BookLine{157.50}{0.00}{\heifont}{0}{切瑞安德：}
\BookLine{166.00}{22.14}{}{0}{把斗篷拿上，伊皮伊尼兹。}
\newpage
% 原书页码：189
\null
\BookNumber{189}
\Vertical{129}{11}{}{7}{10}第\\三\\幕\EndVertical
\BookLine{30.00}{7.38}{\kaifont}{0}{伊皮伊尼兹走到泉水边。}
\BookLine{47.00}{0.00}{\heifont}{0}{阿尔刻提斯：}
\BookLine{55.50}{22.14}{}{0}{年轻人，你的家乡在哪里？}
\BookLine{72.50}{0.00}{\heifont}{0}{切瑞安德：}
\BookLine{81.00}{22.14}{}{0}{埃维亚岛，就在德尔菲斯山下。}
\BookLine{98.00}{0.00}{\heifont}{0}{阿尔刻提斯：}
\BookLine{106.50}{22.14}{}{0}{你的母亲还在世吗？}
\BookLine{123.50}{0.00}{\heifont}{0}{切瑞安德：}
\BookLine{132.00}{22.14}{}{0}{是的，阿尔刻提斯王后。}
\BookLine{149.00}{0.00}{\heifont}{0}{阿尔刻提斯：}
\BookLine{157.50}{22.14}{}{0}{她知道我的名字？}
\newpage
% 原书页码：190
\null
\BookNumber{190}
\Vertical{4}{169}{}{6.2}{8}阿\\尔\\刻\\提\\斯\\之\\歌\EndVertical
\BookLine{30.00}{0.00}{\heifont}{0}{切瑞安德：}
\BookLine{38.50}{22.14}{}{0}{全希腊每一个孩子都知道您的名字！}
\BookLine{55.50}{0.00}{\heifont}{0}{阿尔刻提斯：}
\BookLine{64.00}{22.14}{}{0}{告诉她……阿尔刻提斯谢谢她。}
\BookLine{81.00}{7.38}{\kaifont}{0}{伊皮伊尼兹站在她身旁，她轻轻地摸了摸他。}
\BookLine{98.00}{0.00}{\heifont}{0}{阿尔刻提斯：}
\BookLine{106.50}{22.14}{}{0}{记住，我并没有活得不幸。曾经我十分悲伤和}
\BookLine{115.00}{22.14}{}{0}{绝望，但我得到了拯救。{\kaifont （指路）}从那片树林穿}
\BookLine{123.50}{22.14}{}{0}{过去，沿着河走。}
\BookLine{140.50}{7.38}{\kaifont}{0}{舞台上传来镇上居民们的声音，越来越多的人从门口闯}
\BookLine{149.00}{0.00}{\kaifont}{0}{进来，他们站在舞台边上叫喊着。看门人和守卫们从宫殿里}
\BookLine{157.50}{0.00}{\kaifont}{0}{出来，充满戒备。}
\BookLine{166.00}{7.38}{\kaifont}{0}{阿尔刻提斯示意两个年轻人藏在泉水旁边。守卫们挥舞}
\BookLine{174.50}{0.00}{\kaifont}{0}{着长矛逼居民们后退。}
\newpage
% 原书页码：191
\null
\BookNumber{191}
\Vertical{129}{11}{}{7}{10}第\\三\\幕\EndVertical
\BookLine{30.00}{0.00}{\heifont}{0}{居民：}
\BookLine{38.50}{22.14}{}{0}{水！我们要宫殿里的泉水！我们的水已经被污}
\BookLine{47.00}{22.14}{}{0}{染了。阿吉斯国王，帮帮我们！}
\BookLine{64.00}{0.00}{\heifont}{0}{守卫：}
\BookLine{72.50}{22.14}{}{0}{往后退！都往后退，出去，所有人都出去！}
\BookLine{89.50}{7.38}{\kaifont}{0}{阿吉斯国王从宫殿里走了出来，他穿着粗俗而华丽的衣}
\BookLine{98.00}{0.00}{\kaifont}{0}{服，年龄大概四十岁，浓密的黑色头发贴满了汗涔的前额。}
\BookLine{106.50}{0.00}{\kaifont}{0}{他的脸上也蒙着灰尘。}
\BookLine{123.50}{0.00}{\heifont}{0}{阿吉斯：}
\BookLine{132.00}{22.14}{}{0}{这些人是怎么进到院子里来的？}
\BookLine{149.00}{0.00}{\heifont}{0}{居民：}
\BookLine{157.50}{22.14}{}{0}{阿吉斯国王，帮帮我们！}
\newpage
% 原书页码：192
\null
\BookNumber{192}
\Vertical{4}{169}{}{6.2}{8}阿\\尔\\刻\\提\\斯\\之\\歌\EndVertical
\BookLine{30.00}{0.00}{\heifont}{0}{阿吉斯：}
\BookLine{38.50}{22.14}{}{0}{永远都是“帮帮我们”！我已经尽了全力。你}
\BookLine{47.00}{22.14}{}{0}{们谁让这些人进来的？}
\BookLine{64.00}{0.00}{\heifont}{0}{守卫：}
\BookLine{72.50}{22.14}{}{0}{{\kaifont （困惑地）}陛下，他们把大门撞坏了，我们拦不}
\BookLine{81.00}{22.14}{}{0}{住啊。}
\BookLine{98.00}{0.00}{\heifont}{0}{阿吉斯：}
\BookLine{106.50}{22.14}{}{0}{一群废物！你们忘了自己的任务吗？你们根本}
\BookLine{115.00}{22.14}{}{0}{就是害怕靠近他们，这才是问题。{\kaifont （对居民们说}}
\BookLine{123.50}{22.14}{}{0}{{\kaifont 话）}都往后站！{\kaifont （对其中一名守卫）}他们说水怎}
\BookLine{132.00}{22.14}{}{0}{么了？}
\BookLine{149.00}{0.00}{\heifont}{0}{守卫一：}
\BookLine{157.50}{22.14}{}{0}{陛下，他们说，镇上的水已经被污染了，想从}
\BookLine{166.00}{22.14}{}{0}{宫殿的泉眼里打水。}
\newpage
% 原书页码：193
\null
\BookNumber{193}
\Vertical{129}{11}{}{7}{10}第\\三\\幕\EndVertical
\BookLine{30.00}{0.00}{\heifont}{0}{阿吉斯：}
\BookLine{38.50}{22.14}{}{0}{知道了，你们可以从这儿打水。{\kaifont （捏起长袍的一}}
\BookLine{47.00}{22.14}{}{0}{{\kaifont 角遮住鼻子，指着一个爬到了自己脚边的人）}}
\BookLine{55.50}{22.14}{}{0}{你，往后退！费莱的百姓们，我有事情要告诉}
\BookLine{64.00}{22.14}{}{0}{大家。我以为自己已经倾尽全力，然而我今天}
\BookLine{72.50}{22.14}{}{0}{意识到，有一件事，我还没做。那个阿尔刻提}
\BookLine{81.00}{22.14}{}{0}{斯在哪里？}
\BookLine{98.00}{0.00}{\heifont}{0}{看门人：}
\BookLine{106.50}{22.14}{}{0}{陛下，她就在这儿。}
\BookLine{123.50}{0.00}{\heifont}{0}{阿吉斯：}
\BookLine{132.00}{22.14}{}{0}{{\kaifont （转过身，轻蔑地看着她，慢慢说道）}所以……}
\BookLine{140.50}{22.14}{}{0}{你就是……带来这些邪恶力量的人！费莱的百}
\BookLine{149.00}{22.14}{}{0}{姓们，今天是个大日子，因为我们终于找到这}
\BookLine{157.50}{22.14}{}{0}{场灾难的终极原因。如果可以证明是这个女人}
\BookLine{166.00}{22.14}{}{0}{把灾祸带来的，我们就用石头把她打死，或者}
\BookLine{174.50}{22.14}{}{0}{把她从这个国家永远流放。{\kaifont （对阿尔刻提斯说）}}
\newpage
% 原书页码：194
\null
\BookNumber{194}
\Vertical{4}{169}{}{6.2}{8}阿\\尔\\刻\\提\\斯\\之\\歌\EndVertical
\BookLine{30.00}{22.14}{}{0}{女人，你曾经死了，被埋葬，又活了过来，是}
\BookLine{38.50}{22.14}{}{0}{不是？}
\BookLine{55.50}{0.00}{\heifont}{0}{阿尔刻提斯：}
\BookLine{64.00}{22.14}{}{0}{是。}
\BookLine{81.00}{0.00}{\heifont}{0}{阿吉斯：}
\BookLine{89.50}{22.14}{}{0}{你和你的丈夫都相信赫拉克勒斯是在阿波罗的}
\BookLine{98.00}{22.14}{}{0}{帮助下完成的，所以这是阿波罗的旨意，是}
\BookLine{106.50}{22.14}{}{0}{不是？}
\BookLine{123.50}{0.00}{\heifont}{0}{阿尔刻提斯：}
\BookLine{132.00}{22.14}{}{0}{我们相信，这是真的。}
\BookLine{149.00}{0.00}{\heifont}{0}{阿吉斯：}
\BookLine{157.50}{22.14}{}{0}{你和你的丈夫都相信阿波罗出于对你和塞萨利}
\BookLine{166.00}{22.14}{}{0}{的爱，曾经在这儿生活过一年，是不是？}
\newpage
% 原书页码：195
\null
\BookNumber{195}
\Vertical{129}{11}{}{7}{10}第\\三\\幕\EndVertical
\BookLine{30.00}{0.00}{\heifont}{0}{阿尔刻提斯：}
\BookLine{38.50}{22.14}{}{0}{我们相信，这是真的。}
\BookLine{55.50}{0.00}{\heifont}{0}{阿吉斯：}
\BookLine{64.00}{22.14}{}{0}{那么，阿波罗的爱呢？费莱的百姓们，你们知}
\BookLine{72.50}{22.14}{}{0}{道，这是阿波罗的领地。如果阿波罗的确曾经}
\BookLine{81.00}{22.14}{}{0}{宠爱过这个女人和她的家人，如今岂不是可以}
\BookLine{89.50}{22.14}{}{0}{清楚看到，这爱已经变为仇恨？}
\BookLine{106.50}{7.38}{\kaifont}{0}{人群沉默了。}
\BookLine{123.50}{22.14}{}{0}{是不是？}
\BookLine{140.50}{7.38}{\kaifont}{0}{人群中有一些反对的议论。}
\BookLine{157.50}{22.14}{}{0}{你们说什么？你们搞不清楚？你们这群傻子！}
\BookLine{166.00}{22.14}{}{0}{你们难道忘了这场正在肆虐的瘟疫了吗？看}
\BookLine{174.50}{22.14}{}{0}{门人！}
\newpage
% 原书页码：196
\null
\BookNumber{196}
\Vertical{4}{169}{}{6.2}{8}阿\\尔\\刻\\提\\斯\\之\\歌\EndVertical
\BookLine{30.00}{0.00}{\heifont}{0}{看门人：}
\BookLine{38.50}{22.14}{}{0}{陛下？}
\BookLine{55.50}{0.00}{\heifont}{0}{阿吉斯：}
\BookLine{64.00}{22.14}{}{0}{在宫殿彻底闭锁之前，你看到过这场灾难的情}
\BookLine{72.50}{22.14}{}{0}{形吗？它是怎样袭击我们的？}
\BookLine{89.50}{0.00}{\heifont}{0}{看门人：}
\BookLine{98.00}{22.14}{}{0}{是的陛下，我看到了。}
\BookLine{115.00}{0.00}{\heifont}{0}{阿吉斯：}
\BookLine{123.50}{22.14}{}{0}{说说看！}
\BookLine{140.50}{0.00}{\heifont}{0}{看门人：}
\BookLine{149.00}{22.14}{}{0}{陛下，不只是一种——有三种疾病。第一种是}
\BookLine{157.50}{22.14}{}{0}{突然来临的，就好像……}
\BookLine{174.50}{7.38}{\kaifont}{0}{一名守卫从宫殿里出来，走到阿吉斯国王面前。}
\newpage
% 原书页码：197
\null
\BookNumber{197}
\Vertical{129}{11}{}{7}{10}第\\三\\幕\EndVertical
\BookLine{30.00}{0.00}{\heifont}{0}{阿吉斯：}
\BookLine{38.50}{22.14}{}{0}{{\kaifont （恼怒地）}好吧，怎么了？又怎么了？}
\BookLine{55.50}{0.00}{\heifont}{0}{守卫二：}
\BookLine{64.00}{22.14}{}{0}{陛下，您的女儿一直在拍打着门。她说她想出}
\BookLine{72.50}{22.14}{}{0}{来。她说她想和您待在一起。她一直在不停地}
\BookLine{81.00}{22.14}{}{0}{拍打着门。}
\BookLine{98.00}{0.00}{\heifont}{0}{阿吉斯：}
\BookLine{106.50}{22.14}{}{0}{{\kaifont （在女儿的事情上显出他的温柔，急切地）}告诉}
\BookLine{115.00}{22.14}{}{0}{她我马上就来陪她，现在还有些事要处理。告}
\BookLine{123.50}{22.14}{}{0}{诉她忍耐一下，我马上就来了。}
\BookLine{140.50}{0.00}{\heifont}{0}{守卫二：}
\BookLine{149.00}{22.14}{}{0}{是，陛下。}
\BookLine{166.00}{7.38}{\kaifont}{0}{他转身向宫门走去。}
\newpage
% 原书页码：198
\null
\BookNumber{198}
\Vertical{4}{169}{}{6.2}{8}阿\\尔\\刻\\提\\斯\\之\\歌\EndVertical
\BookLine{30.00}{0.00}{\heifont}{0}{阿吉斯：}
\BookLine{38.50}{22.14}{}{0}{等等！{\kaifont （心碎地）}告诉她忍耐一下。我等下会带}
\BookLine{47.00}{22.14}{}{0}{她去花园里走走。}
\BookLine{64.00}{0.00}{\heifont}{0}{守卫二：}
\BookLine{72.50}{22.14}{}{0}{是，陛下。}
\BookLine{89.50}{7.38}{\kaifont}{0}{护卫下场。}
\BookLine{106.50}{0.00}{\heifont}{0}{阿吉斯：}
\BookLine{115.00}{22.14}{}{0}{{\kaifont （对看门人）}三种？你说有三种疾病？}
\BookLine{132.00}{0.00}{\heifont}{0}{看门人：}
\BookLine{140.50}{22.14}{}{0}{第一种是突然来临的，就像闪电一样，陛下。}
\BookLine{157.50}{0.00}{\heifont}{0}{阿吉斯：}
\BookLine{166.00}{22.14}{}{0}{{\kaifont （颤抖着转向阿尔刻提斯）}就是你干的，你这个}
\BookLine{174.50}{22.14}{}{0}{老妇人！}
\newpage
% 原书页码：199
\null
\BookNumber{199}
\Vertical{129}{11}{}{7}{10}第\\三\\幕\EndVertical
\BookLine{30.00}{0.00}{\heifont}{0}{阿尔刻提斯：}
\BookLine{38.50}{22.14}{}{0}{不！不！}
\BookLine{55.50}{0.00}{\heifont}{0}{看门人：}
\BookLine{64.00}{22.14}{}{0}{{\kaifont （手放在自己的胃部）}感觉就像火烧一样。就是}
\BookLine{72.50}{22.14}{}{0}{这种病袭击了所有年轻人和孩子们。}
\BookLine{89.50}{0.00}{\heifont}{0}{阿吉斯：}
\BookLine{98.00}{22.14}{}{0}{{\kaifont （出离愤怒）}孩子！你这个蠢货！除了在这儿散}
\BookLine{106.50}{22.14}{}{0}{播各种关于疾病的预兆，你就没有别的事可做}
\BookLine{115.00}{22.14}{}{0}{了吗？我真该把你的舌头割下来！离开吧，你}
\BookLine{123.50}{22.14}{}{0}{们这些不朽的神明。不要看这些预兆，不要}
\BookLine{132.00}{22.14}{}{0}{听什么预兆。{\kaifont （对看门人）}我不想再听了。{\kaifont （绝}}
\BookLine{140.50}{22.14}{}{0}{{\kaifont 望而沮丧地朝着宫殿走去）}噢，走开吧，你们}
\BookLine{149.00}{22.14}{}{0}{这些人！谁才能拯救我们脱离这些可怕的夜}
\BookLine{157.50}{22.14}{}{0}{晚……这些困境……这些邪恶的力量？}
\BookLine{174.50}{7.38}{\kaifont}{0}{他无助地摇头，突然反应过来，指着阿尔刻提斯。}
\newpage
% 原书页码：200
\null
\BookNumber{200}
\Vertical{4}{169}{}{6.2}{8}阿\\尔\\刻\\提\\斯\\之\\歌\EndVertical
\BookLine{30.00}{22.14}{}{0}{你说！是因为神圣的阿波罗憎恨你，因为你他}
\BookLine{38.50}{22.14}{}{0}{才把这诅咒带到塞萨利来的！}
\BookLine{55.50}{7.38}{\kaifont}{0}{阿尔刻提斯静静地看着他，没有说话。}
\BookLine{72.50}{22.14}{}{0}{难道不是这样吗？你说啊！}
\BookLine{89.50}{0.00}{\heifont}{0}{阿尔刻提斯：}
\BookLine{98.00}{22.14}{}{0}{{\kaifont （不紧不慢）}阿吉斯国王，阿波罗来到这里的时}
\BookLine{106.50}{22.14}{}{0}{候，他的祭司忒瑞西阿斯曾说过，伟大的荣耀}
\BookLine{115.00}{22.14}{}{0}{总是伴随着巨大的危难。}
\BookLine{132.00}{0.00}{\heifont}{0}{阿吉斯：}
\BookLine{140.50}{22.14}{}{0}{危难？}
\BookLine{157.50}{0.00}{\heifont}{0}{阿尔刻提斯：}
\BookLine{166.00}{22.14}{}{0}{你此刻正在经历。这片土地仍然是温暖的——}
\BookLine{174.50}{22.14}{}{0}{炙热的——因为神的足迹。{\kaifont （像是在对自己说}}
\newpage
% 原书页码：201
\null
\BookNumber{201}
\Vertical{129}{11}{}{7}{10}第\\三\\幕\EndVertical
\BookLine{30.00}{22.14}{}{0}{{\kaifont 话，像是梦呓）}很长一段时间，我一直没能明}
\BookLine{38.50}{22.14}{}{0}{白。然而孤独和奴役让我渐渐清醒起来。阿吉}
\BookLine{47.00}{22.14}{}{0}{斯国王，在这儿，行事要小心。}
\BookLine{64.00}{0.00}{\heifont}{0}{阿吉斯：}
\BookLine{72.50}{22.14}{}{0}{{\kaifont （仰起脸，对她的警告毫不在意）}回答我的}
\BookLine{81.00}{22.14}{}{0}{问题。}
\BookLine{98.00}{0.00}{\heifont}{0}{阿尔刻提斯：}
\BookLine{106.50}{22.14}{}{0}{阿吉斯国王，神明和你我不同，不过有时我们}
\BookLine{115.00}{22.14}{}{0}{会和他们有几分相像。神明们不会爱过我们，}
\BookLine{123.50}{22.14}{}{0}{又恨恶起他们先前所爱的。就像你不会今天爱}
\BookLine{132.00}{22.14}{}{0}{着劳德米娅，明天却把她丢到大街上一样。}
\BookLine{149.00}{0.00}{\heifont}{0}{阿吉斯：}
\BookLine{157.50}{22.14}{}{0}{{\kaifont （愤怒地）}不许你提起她的名字，你，你这个带}
\BookLine{166.00}{22.14}{}{0}{来死亡与毁灭的人！{\kaifont （他向宫殿走去）}}
\newpage
% 原书页码：202
\null
\BookNumber{202}
\Vertical{4}{169}{}{6.2}{8}阿\\尔\\刻\\提\\斯\\之\\歌\EndVertical
\BookLine{30.00}{0.00}{\heifont}{0}{阿尔刻提斯：}
\BookLine{38.50}{22.14}{}{0}{我们向神明祈求健康，财富，自己的幸福。但}
\BookLine{47.00}{22.14}{}{0}{神明想要给予我们一些其他的东西，一些更好}
\BookLine{55.50}{22.14}{}{0}{的东西——对生命的理解。我们总是太快地拒}
\BookLine{64.00}{22.14}{}{0}{绝他们的赐予……不，不！阿波罗从没有放弃}
\BookLine{72.50}{22.14}{}{0}{过我……费莱的百姓，阿德墨托斯国王可曾待}
\BookLine{81.00}{22.14}{}{0}{你们有何不公？}
\BookLine{98.00}{0.00}{\heifont}{0}{居民：}
\BookLine{106.50}{22.14}{}{0}{没有，阿尔刻提斯王后！}
\BookLine{123.50}{0.00}{\heifont}{0}{阿尔刻提斯：}
\BookLine{132.00}{22.14}{}{0}{我呢？}
\BookLine{149.00}{0.00}{\heifont}{0}{居民：}
\BookLine{157.50}{22.14}{}{0}{没有，阿尔刻提斯王后！}
\newpage
% 原书页码：203
\null
\BookNumber{203}
\Vertical{129}{11}{}{7}{10}第\\三\\幕\EndVertical
\BookLine{30.00}{0.00}{\heifont}{0}{阿尔刻提斯：}
\BookLine{38.50}{22.14}{}{0}{你们觉得神明是为了惩罚我们的罪恶才降下这}
\BookLine{47.00}{22.14}{}{0}{些灾难吗？}
\BookLine{64.00}{0.00}{\heifont}{0}{居民：}
\BookLine{72.50}{22.14}{}{0}{不！不！}
\BookLine{89.50}{0.00}{\heifont}{0}{阿尔刻提斯：}
\BookLine{98.00}{22.14}{}{0}{就算赫拉克勒斯从死神那里带回来了什么，你}
\BookLine{106.50}{22.14}{}{0}{们觉得没有神明全权的准许，他能做到吗？}
\BookLine{123.50}{0.00}{\heifont}{0}{居民：}
\BookLine{132.00}{22.14}{}{0}{不能，不能！}
\BookLine{149.00}{0.00}{\heifont}{0}{阿吉斯：}
\BookLine{157.50}{22.14}{}{0}{那么，为什么会有这场瘟疫呢？}
\newpage
% 原书页码：204
\null
\BookNumber{204}
\Vertical{4}{169}{}{6.2}{8}阿\\尔\\刻\\提\\斯\\之\\歌\EndVertical
\BookLine{30.00}{0.00}{\heifont}{0}{阿尔刻提斯：}
\BookLine{38.50}{22.14}{}{0}{神明降下灾难……是让我们注意……让我们停}
\BookLine{47.00}{22.14}{}{0}{下来，睁开自己的眼睛……}
\BookLine{64.00}{0.00}{\heifont}{0}{阿吉斯：}
\BookLine{72.50}{22.14}{}{0}{让我们注意什么，阿尔刻提斯？}
\BookLine{89.50}{0.00}{\heifont}{0}{阿尔刻提斯：}
\BookLine{98.00}{22.14}{}{0}{{\kaifont （抬起头，像是在倾听）}我不知道，注意某种}
\BookLine{106.50}{22.14}{}{0}{启示。}
\BookLine{123.50}{7.38}{\kaifont}{0}{场上陷入片刻静待的状态，突然间站在路高处的守卫看}
\BookLine{132.00}{0.00}{\kaifont}{0}{到了两个年轻人，大喊道。}
\BookLine{149.00}{0.00}{\heifont}{0}{守卫一：}
\BookLine{157.50}{22.14}{}{0}{陛下，这儿有两个陌生人！}
\newpage
% 原书页码：205
\null
\BookNumber{205}
\Vertical{129}{11}{}{7}{10}第\\三\\幕\EndVertical
\BookLine{30.00}{0.00}{\heifont}{0}{阿吉斯：}
\BookLine{38.50}{22.14}{}{0}{{\kaifont （走上前去，保持一个安全距离）}你们是怎么进}
\BookLine{47.00}{22.14}{}{0}{来的？守卫！把他们围起来！}
\BookLine{64.00}{7.38}{\kaifont}{0}{守卫们围了过去。}
\BookLine{81.00}{22.14}{}{0}{放下武器！}
\BookLine{98.00}{0.00}{\heifont}{0}{伊皮伊尼兹：}
\BookLine{106.50}{22.14}{}{0}{阿吉斯国王，没门！}
\BookLine{123.50}{0.00}{\heifont}{0}{阿吉斯：}
\BookLine{132.00}{22.14}{}{0}{你们被困在那个洞里，还能怎么样？放下你们}
\BookLine{140.50}{22.14}{}{0}{的剑！}
\BookLine{157.50}{7.38}{\kaifont}{0}{阿尔刻提斯不住地摇头，嘴里默念着：“阿吉斯国王，阿}
\BookLine{166.00}{0.00}{\kaifont}{0}{吉斯国王……”}
\newpage
% 原书页码：206
\null
\BookNumber{206}
\Vertical{4}{169}{}{6.2}{8}阿\\尔\\刻\\提\\斯\\之\\歌\EndVertical
\BookLine{30.00}{0.00}{\heifont}{0}{阿吉斯：}
\BookLine{38.50}{22.14}{}{0}{{\kaifont （对阿尔刻提斯说）}是你把他们放进来的。你们}
\BookLine{47.00}{22.14}{}{0}{都没有往脸上蒙上灰尘。{\kaifont （对守卫）}杀了他们！}
\BookLine{55.50}{22.14}{}{0}{杀了他们！}
\BookLine{72.50}{7.38}{\kaifont}{0}{两个守卫怯生生地往下走。}
\BookLine{89.50}{22.14}{}{0}{懦夫！叛徒！照我说的去做！}
\BookLine{106.50}{7.38}{\kaifont}{0}{舞台上一名守卫突然疫病发作，他扔掉手中的长矛，喊}
\BookLine{115.00}{0.00}{\kaifont}{0}{了起来。}
\BookLine{132.00}{0.00}{\heifont}{0}{守卫三：}
\BookLine{140.50}{22.14}{}{0}{阿吉斯国王！水！是瘟疫！我觉得自己烧起来}
\BookLine{149.00}{22.14}{}{0}{了！救救我！帮帮我！我在燃烧！}
\BookLine{166.00}{7.38}{\kaifont}{0}{他断断续续地呕吐着。}
\newpage
% 原书页码：207
\null
\BookNumber{207}
\Vertical{129}{11}{}{7}{10}第\\三\\幕\EndVertical
\BookLine{30.00}{0.00}{\heifont}{0}{阿吉斯：}
\BookLine{38.50}{22.14}{}{0}{{\kaifont （向后退，人群都向后退着）}把他拖出去！用你}
\BookLine{47.00}{22.14}{}{0}{们的矛赶他走！}
\BookLine{64.00}{0.00}{\heifont}{0}{守卫三：}
\BookLine{72.50}{22.14}{}{0}{{\kaifont （身体剧烈地左右摇晃，一边呕吐一边喊着）}}
\BookLine{81.00}{22.14}{}{0}{水！水！我要死了！}
\BookLine{98.00}{0.00}{\heifont}{0}{阿吉斯：}
\BookLine{106.50}{22.14}{}{0}{{\kaifont （用手挡住自己的脸）}把他拖到路上去。}
\BookLine{123.50}{7.38}{\kaifont}{0}{一阵骚动。生病的守卫努力向门口爬去，其余的守卫和}
\BookLine{132.00}{0.00}{\kaifont}{0}{镇上的居民们都转过脸去，或是拿着长矛或是用脚踢，想把}
\BookLine{140.50}{0.00}{\kaifont}{0}{他快点赶走。生病的守卫将大家的注意力从伊皮伊尼兹和切}
\BookLine{149.00}{0.00}{\kaifont}{0}{瑞安德身上转移开了，场上一度只能听见大家喊叫的声音。}
\BookLine{166.00}{0.00}{\heifont}{0}{阿吉斯：}
\BookLine{174.50}{22.14}{}{0}{{\kaifont （尖叫着）}带上我的武器！我的马匹！回色雷斯}
\newpage
% 原书页码：208
\null
\BookNumber{208}
\Vertical{4}{169}{}{6.2}{8}阿\\尔\\刻\\提\\斯\\之\\歌\EndVertical
\BookLine{30.00}{22.14}{}{0}{去！回色雷斯去！{\kaifont （对阿尔刻提斯）}热情好客的}
\BookLine{38.50}{22.14}{}{0}{塞萨利我还给你，阿尔刻提斯王后！你自己统}
\BookLine{47.00}{22.14}{}{0}{治你的这些死人吧！}
\BookLine{64.00}{0.00}{\heifont}{0}{伊皮伊尼兹：}
\BookLine{72.50}{22.14}{}{0}{快，切瑞安德！}
\BookLine{89.50}{7.38}{\kaifont}{0}{他们冲上舞台。}
\BookLine{106.50}{0.00}{\heifont}{0}{切瑞安德：}
\BookLine{115.00}{22.14}{}{0}{给他一剑，伊皮伊尼兹！}
\BookLine{132.00}{0.00}{\heifont}{0}{伊皮伊尼兹：}
\BookLine{140.50}{22.14}{}{0}{阿吉斯，我是阿德墨托斯之子，伊皮伊尼兹！}
\BookLine{157.50}{0.00}{\heifont}{0}{阿吉斯：}
\BookLine{166.00}{22.14}{}{0}{你是谁？这是怎么了？守卫！}
\newpage
% 原书页码：209
\null
\BookNumber{209}
\Vertical{129}{11}{}{7}{10}第\\三\\幕\EndVertical
\BookLine{30.00}{7.38}{\kaifont}{0}{阿尔刻提斯摇着头，站在阿吉斯国王身前，伸出手阻止}
\BookLine{38.50}{0.00}{\kaifont}{0}{伊皮伊尼兹。}
\BookLine{55.50}{0.00}{\heifont}{0}{阿尔刻提斯：}
\BookLine{64.00}{22.14}{}{0}{不，伊皮伊尼兹！不要！}
\BookLine{81.00}{0.00}{\heifont}{0}{伊皮伊尼兹：}
\BookLine{89.50}{22.14}{}{0}{{\kaifont （对阿尔刻提斯的态度十分愤怒）}母亲！我们的}
\BookLine{98.00}{22.14}{}{0}{机会来了啊。是他杀了我的父亲！}
\BookLine{115.00}{0.00}{\heifont}{0}{阿尔刻提斯：}
\BookLine{123.50}{22.14}{}{0}{不要这样做！}
\BookLine{140.50}{0.00}{\heifont}{0}{阿吉斯：}
\BookLine{149.00}{22.14}{}{0}{守卫！{\kaifont （狂躁地左右踱步）}一群懦夫！守卫！}
\BookLine{166.00}{0.00}{\heifont}{0}{切瑞安德：}
\BookLine{174.50}{22.14}{}{0}{伊皮伊尼兹，照你母亲说的做。}
\newpage
% 原书页码：210
\null
\BookNumber{210}
\Vertical{4}{169}{}{6.2}{8}阿\\尔\\刻\\提\\斯\\之\\歌\EndVertical
\BookLine{30.00}{0.00}{\heifont}{0}{阿尔刻提斯：}
\BookLine{38.50}{22.14}{}{0}{{\kaifont （背对着阿吉斯和他说话）}对的，阿吉斯国王，}
\BookLine{47.00}{22.14}{}{0}{回你自己的国家去吧。}
\BookLine{64.00}{0.00}{\heifont}{0}{阿吉斯：}
\BookLine{72.50}{22.14}{}{0}{这是你儿子吗？阿尔刻提斯，回答我，是你儿}
\BookLine{81.00}{22.14}{}{0}{子吗？}
\BookLine{98.00}{7.38}{\kaifont}{0}{眼见复仇计划受挫，伊皮伊尼兹愤怒又沮丧，他跪在地}
\BookLine{106.50}{0.00}{\kaifont}{0}{上，一直在用刀柄敲打地面。}
\BookLine{123.50}{0.00}{\heifont}{0}{伊皮伊尼兹：}
\BookLine{132.00}{22.14}{}{0}{我要报仇！报仇！}
\BookLine{149.00}{0.00}{\heifont}{0}{阿尔刻提斯：}
\BookLine{157.50}{22.14}{}{0}{伊皮伊尼兹，记得你父亲说过的话：举刀杀人，}
\BookLine{166.00}{22.14}{}{0}{伤的是自己。}
\newpage
% 原书页码：211
\null
\BookNumber{211}
\Vertical{129}{11}{}{7}{10}第\\三\\幕\EndVertical
\BookLine{30.00}{0.00}{\heifont}{0}{伊皮伊尼兹：}
\BookLine{38.50}{22.14}{}{0}{{\kaifont （前额贴着地面哭泣）}他杀了我的父亲……我的}
\BookLine{47.00}{22.14}{}{0}{弟弟……妹妹……}
\BookLine{64.00}{7.38}{\kaifont}{0}{守卫二从宫殿里跑了出来。}
\BookLine{81.00}{0.00}{\heifont}{0}{守卫二：}
\BookLine{89.50}{22.14}{}{0}{阿吉斯国王！您的女儿，劳德米娅公主！她一}
\BookLine{98.00}{22.14}{}{0}{直在敲门在喊您，非常痛苦。阿吉斯国王，她}
\BookLine{106.50}{22.14}{}{0}{非常痛苦！}
\BookLine{123.50}{0.00}{\heifont}{0}{阿吉斯：}
\BookLine{132.00}{22.14}{}{0}{{\kaifont （举起双臂）}神明们，停手吧！劳德米娅！劳德}
\BookLine{140.50}{22.14}{}{0}{米娅！停手吧，你们这些不朽的众神！{\kaifont （阿吉斯}}
\BookLine{149.00}{22.14}{}{0}{{\kaifont 跑进宫殿）}}
\BookLine{166.00}{7.38}{\kaifont}{0}{阿尔刻提斯站在伊皮伊尼兹身边，把手放在他的肩膀上。}
\newpage
% 原书页码：212
\null
\BookNumber{212}
\Vertical{4}{169}{}{6.2}{8}阿\\尔\\刻\\提\\斯\\之\\歌\EndVertical
\BookLine{30.00}{0.00}{\heifont}{0}{阿尔刻提斯：}
\BookLine{38.50}{22.14}{}{0}{一个人一旦从复仇中得到快乐，他就不会再有}
\BookLine{47.00}{22.14}{}{0}{其他快乐了，这是你父亲说过的话。{\kaifont （转向居}}
\BookLine{55.50}{22.14}{}{0}{{\kaifont 民们，冷静地）}朋友们，回家去，拿上你们的}
\BookLine{64.00}{22.14}{}{0}{篮子和水罐，然后到南边大门外的硫矿那里}
\BookLine{72.50}{22.14}{}{0}{去。伊皮伊尼兹，你还记得小时候常常去玩的}
\BookLine{81.00}{22.14}{}{0}{那片矿区吗？那儿有很多工人们制铁时用的硫}
\BookLine{89.50}{22.14}{}{0}{黄……取一些拿到路上去烧，然后撒到死者们}
\BookLine{98.00}{22.14}{}{0}{的身上。{\kaifont （对守卫们）}你们去帮着一起干吧，这}
\BookLine{106.50}{22.14}{}{0}{儿已经没有什么事了。伊皮伊尼兹，站起来。}
\BookLine{115.00}{22.14}{}{0}{你来带领他们。}
\BookLine{132.00}{0.00}{\heifont}{0}{伊皮伊尼兹：}
\BookLine{140.50}{22.14}{}{0}{{\kaifont （起身）}是的，母亲。{\kaifont （带着平静的威严，对守}}
\BookLine{149.00}{22.14}{}{0}{{\kaifont 卫和镇上的居民）}大家跟我来。}
\BookLine{166.00}{7.38}{\kaifont}{0}{他们准备离开，切瑞安德返回来，轻柔地，带着敬畏地}
\BookLine{174.50}{0.00}{\kaifont}{0}{问道。}
\newpage
% 原书页码：213
\null
\BookNumber{213}
\Vertical{129}{11}{}{7}{10}第\\三\\幕\EndVertical
\BookLine{30.00}{0.00}{\heifont}{0}{切瑞安德：}
\BookLine{38.50}{22.14}{}{0}{阿尔刻提斯王后……您说的那个启示，阿波罗}
\BookLine{47.00}{22.14}{}{0}{神给我们的启示。阿吉斯国王的这个决定，就}
\BookLine{55.50}{22.14}{}{0}{是那个启示吗？}
\BookLine{72.50}{0.00}{\heifont}{0}{阿尔刻提斯：}
\BookLine{81.00}{22.14}{}{0}{{\kaifont （抬起头像听着）}不……启示还没有显明呢。}
\BookLine{98.00}{0.00}{\heifont}{0}{切瑞安德：}
\BookLine{106.50}{22.14}{}{0}{{\kaifont （带着年轻人的一脸热忱）}您就是那个启示啊！}
\BookLine{115.00}{22.14}{}{0}{您就是神明的信息和启示，阿尔刻提斯王后！}
\BookLine{132.00}{0.00}{\heifont}{0}{阿尔刻提斯：}
\BookLine{140.50}{22.14}{}{0}{{\kaifont （近乎麻木，摇着头，轻声地）}不，不……}
\BookLine{157.50}{7.38}{\kaifont}{0}{切瑞安德跟在伊皮伊尼兹身后走了出去。阿吉斯国王从}
\BookLine{166.00}{0.00}{\kaifont}{0}{宫殿里出来，痛苦地哭号着。}
\newpage
% 原书页码：214
\null
\BookNumber{214}
\Vertical{4}{169}{}{6.2}{8}阿\\尔\\刻\\提\\斯\\之\\歌\EndVertical
\BookLine{30.00}{0.00}{\heifont}{0}{阿吉斯：}
\BookLine{38.50}{22.14}{}{0}{她死了！劳德米娅死了！才十二岁就死了……}
\BookLine{55.50}{7.38}{\kaifont}{0}{他用拳头捶着柱子和墙，向下走几步又返回。}
\BookLine{72.50}{22.14}{}{0}{才十二岁。她搂着我的脖子，痛得要命：“父}
\BookLine{81.00}{22.14}{}{0}{亲，帮帮我，父亲，帮帮我啊！”她的头发，}
\BookLine{89.50}{22.14}{}{0}{她那么像她的母亲。她的眼睛，她的眼睛。{\kaifont （看}}
\BookLine{98.00}{22.14}{}{0}{{\kaifont 见阿尔刻提斯）}你！是你造成的这一切！都是}
\BookLine{106.50}{22.14}{}{0}{你干的！}
\BookLine{123.50}{0.00}{\heifont}{0}{阿尔刻提斯：}
\BookLine{132.00}{22.14}{}{0}{{\kaifont （在阿吉斯说话的间隙一直在小声说着，像是在}}
\BookLine{140.50}{22.14}{}{0}{{\kaifont 恳求）}阿吉斯……阿吉斯……阿吉斯……}
\BookLine{157.50}{0.00}{\heifont}{0}{阿吉斯：}
\BookLine{166.00}{22.14}{}{0}{{\kaifont （抓住她的手放在自己的额前和胸前）}让我也染}
\BookLine{174.50}{22.14}{}{0}{上瘟疫吧，让它把我们都杀了吧。“父亲，帮}
\newpage
% 原书页码：215
\null
\BookNumber{215}
\Vertical{129}{11}{}{7}{10}第\\三\\幕\EndVertical
\BookLine{30.00}{22.14}{}{0}{帮我！”她是我的一切啊。现在她死了，死了，}
\BookLine{38.50}{22.14}{}{0}{死了！}
\BookLine{55.50}{0.00}{\heifont}{0}{阿尔刻提斯：}
\BookLine{64.00}{22.14}{}{0}{阿吉斯……阿吉斯……}
\BookLine{81.00}{0.00}{\heifont}{0}{阿吉斯：}
\BookLine{89.50}{22.14}{}{0}{你，赫拉克勒斯曾把你从死神那儿带回来，你}
\BookLine{98.00}{22.14}{}{0}{有办法。{\kaifont （突然想到了什么）}赫拉克勒斯把你带}
\BookLine{106.50}{22.14}{}{0}{回来了。在哪里？是这里吗？{\kaifont （跌跌撞撞地沿}}
\BookLine{115.00}{22.14}{}{0}{{\kaifont 着路走下去）}在那儿是不是？{\kaifont （匆忙地走下去）}}
\BookLine{123.50}{22.14}{}{0}{告诉我，阿尔刻提斯，赫拉克勒斯是怎么做到}
\BookLine{132.00}{22.14}{}{0}{的？在下面都发生了什么？}
\BookLine{149.00}{7.38}{\kaifont}{0}{阿尔刻提斯沉默地摇摇头。}
\BookLine{166.00}{0.00}{\heifont}{0}{阿吉斯：}
\BookLine{174.50}{22.14}{}{0}{告诉我他是怎么做到的，我也可以。劳德米娅，}
\newpage
% 原书页码：216
\null
\BookNumber{216}
\Vertical{4}{169}{}{6.2}{8}阿\\尔\\刻\\提\\斯\\之\\歌\EndVertical
\BookLine{30.00}{22.14}{}{0}{我要来救你。阿尔刻提斯，你回答我！}
\BookLine{47.00}{0.00}{\heifont}{0}{阿尔刻提斯：}
\BookLine{55.50}{22.14}{}{0}{阿吉斯，我那时候什么也没看到，没听到。}
\BookLine{72.50}{0.00}{\heifont}{0}{阿吉斯：}
\BookLine{81.00}{22.14}{}{0}{你说谎！}
\BookLine{98.00}{0.00}{\heifont}{0}{阿尔刻提斯：}
\BookLine{106.50}{22.14}{}{0}{阿吉斯，你听我说，我有话对你说。}
\BookLine{123.50}{0.00}{\heifont}{0}{阿吉斯：}
\BookLine{132.00}{22.14}{}{0}{你说啊！说啊！}
\BookLine{149.00}{0.00}{\heifont}{0}{阿尔刻提斯：}
\BookLine{157.50}{22.14}{}{0}{“父亲，帮帮我！”}
\newpage
% 原书页码：217
\null
\BookNumber{217}
\Vertical{129}{11}{}{7}{10}第\\三\\幕\EndVertical
\BookLine{30.00}{0.00}{\heifont}{0}{阿吉斯：}
\BookLine{38.50}{22.14}{}{0}{不要嘲弄我了。}
\BookLine{55.50}{0.00}{\heifont}{0}{阿尔刻提斯：}
\BookLine{64.00}{22.14}{}{0}{我不是在嘲弄。阿吉斯国王，劳德米娅说这些}
\BookLine{72.50}{22.14}{}{0}{是什么意思？}
\BookLine{89.50}{0.00}{\heifont}{0}{阿吉斯：}
\BookLine{98.00}{22.14}{}{0}{她觉得很疼，疼得要命！}
\BookLine{115.00}{0.00}{\heifont}{0}{阿尔刻提斯：}
\BookLine{123.50}{22.14}{}{0}{是，但这不是全部。她还想要表达什么？}
\BookLine{140.50}{0.00}{\heifont}{0}{阿吉斯：}
\BookLine{149.00}{22.14}{}{0}{还有什么？}
\BookLine{166.00}{0.00}{\heifont}{0}{阿尔刻提斯：}
\BookLine{174.50}{22.14}{}{0}{阿吉斯国王，死亡本身的痛苦，只是痛苦的一}
\newpage
% 原书页码：218
\null
\BookNumber{218}
\Vertical{4}{169}{}{6.2}{8}阿\\尔\\刻\\提\\斯\\之\\歌\EndVertical
\BookLine{30.00}{22.14}{}{0}{部分，并非全部。死亡带给人最大的痛苦不是}
\BookLine{38.50}{22.14}{}{0}{分离，尽管这也会让人非常非常悲伤。我死过}
\BookLine{47.00}{22.14}{}{0}{一次。阿吉斯国王，死亡带给人最大的痛苦是}
\BookLine{55.50}{22.14}{}{0}{什么？}
\BookLine{72.50}{0.00}{\heifont}{0}{阿吉斯：}
\BookLine{81.00}{22.14}{}{0}{告诉我！}
\BookLine{98.00}{0.00}{\heifont}{0}{阿尔刻提斯：}
\BookLine{106.50}{22.14}{}{0}{是发现自己从未活过的绝望，是发现生命毫无}
\BookLine{115.00}{22.14}{}{0}{意义的绝望。生命没有任何意义，快乐也好，}
\BookLine{123.50}{22.14}{}{0}{不快乐也好，都成了空白。“父亲，帮帮我！”}
\BookLine{140.50}{0.00}{\heifont}{0}{阿吉斯：}
\BookLine{149.00}{22.14}{}{0}{她曾经多么爱我，阿尔刻提斯。}
\BookLine{166.00}{0.00}{\heifont}{0}{阿尔刻提斯：}
\BookLine{174.50}{22.14}{}{0}{是的。}
\newpage
% 原书页码：219
\null
\BookNumber{219}
\Vertical{129}{11}{}{7}{10}第\\三\\幕\EndVertical
\BookLine{30.00}{0.00}{\heifont}{0}{阿吉斯：}
\BookLine{38.50}{22.14}{}{0}{她曾经多么爱我。}
\BookLine{55.50}{0.00}{\heifont}{0}{阿尔刻提斯：}
\BookLine{64.00}{22.14}{}{0}{是的，但爱不是全部。}
\BookLine{81.00}{0.00}{\heifont}{0}{阿吉斯：}
\BookLine{89.50}{22.14}{}{0}{是。对她，对我来说都是的。我不会听你的。}
\BookLine{106.50}{0.00}{\heifont}{0}{阿尔刻提斯：}
\BookLine{115.00}{22.14}{}{0}{爱不是全部的意义。它是一个启示，让我们明}
\BookLine{123.50}{22.14}{}{0}{白一切是有意义的。它只是一个启示，让我们}
\BookLine{132.00}{22.14}{}{0}{明白一切是有意义的。劳德米娅感到很绝望，}
\BookLine{140.50}{22.14}{}{0}{希望你能帮助她。这就是死亡——死就是绝}
\BookLine{149.00}{22.14}{}{0}{望。除非你赋予它意义，否则她的生命不过是}
\BookLine{157.50}{22.14}{}{0}{空白与空虚。}
\newpage
% 原书页码：220
\null
\BookNumber{220}
\Vertical{4}{169}{}{6.2}{8}阿\\尔\\刻\\提\\斯\\之\\歌\EndVertical
\BookLine{30.00}{0.00}{\heifont}{0}{阿吉斯：}
\BookLine{38.50}{22.14}{}{0}{我能赋予它什么意义？}
\BookLine{55.50}{0.00}{\heifont}{0}{阿尔刻提斯：}
\BookLine{64.00}{22.14}{}{0}{{\kaifont （平静地）}你真是一个粗俗、残忍、无知的人。}
\BookLine{72.50}{22.14}{}{0}{{\kaifont （短暂的静场）}是你杀了“我的”劳德米娅。三}
\BookLine{81.00}{22.14}{}{0}{次。无情地。你自己都不知道多少次杀死了她。}
\BookLine{98.00}{0.00}{\heifont}{0}{阿吉斯：}
\BookLine{106.50}{22.14}{}{0}{不！}
\BookLine{123.50}{0.00}{\heifont}{0}{阿尔刻提斯：}
\BookLine{132.00}{22.14}{}{0}{你什么都不明白。回你自己的国家去吧。在那}
\BookLine{140.50}{22.14}{}{0}{儿，只有在那儿，你还可以为劳德米娅做些}
\BookLine{149.00}{22.14}{}{0}{事情。}
\BookLine{166.00}{7.38}{\kaifont}{0}{阿吉斯沿着小路走上来，走过阿尔刻提斯身边，径直走}
\BookLine{174.50}{0.00}{\kaifont}{0}{向宫殿大门。}
\newpage
% 原书页码：221
\null
\BookNumber{221}
\Vertical{129}{11}{}{7}{10}第\\三\\幕\EndVertical
\BookLine{30.00}{0.00}{\heifont}{0}{阿尔刻提斯：}
\BookLine{38.50}{22.14}{}{0}{阿吉斯国王，所有的死者……{\kaifont （她指着地下世}}
\BookLine{47.00}{22.14}{}{0}{{\kaifont 界的入口）}那里成千上万的死者，在等着我们}
\BookLine{55.50}{22.14}{}{0}{向他们表明，他们不是白白地、无益地死去。}
\BookLine{72.50}{0.00}{\heifont}{0}{阿吉斯：}
\BookLine{81.00}{22.14}{}{0}{我能赋予劳德米娅的生命什么意义呢？}
\BookLine{98.00}{0.00}{\heifont}{0}{阿尔刻提斯：}
\BookLine{106.50}{22.14}{}{0}{从今天开始，你会慢慢明白的。}
\BookLine{123.50}{0.00}{\heifont}{0}{阿吉斯：}
\BookLine{132.00}{22.14}{}{0}{{\kaifont （他的头抵着宫殿的门）}不。}
\BookLine{149.00}{0.00}{\heifont}{0}{阿尔刻提斯：}
\BookLine{157.50}{22.14}{}{0}{我也是慢慢才明白的。即使是我。你会明白的，}
\BookLine{166.00}{22.14}{}{0}{阿吉斯国王……从劳德米娅遭受的痛苦之中，}
\BookLine{174.50}{22.14}{}{0}{你会明白的。}
\newpage
% 原书页码：222
\null
\BookNumber{222}
\Vertical{4}{169}{}{6.2}{8}阿\\尔\\刻\\提\\斯\\之\\歌\EndVertical
\BookLine{30.00}{7.38}{\kaifont}{0}{被彻底击垮的阿吉斯走进宫殿。他慢慢退场后，灯光渐}
\BookLine{38.50}{0.00}{\kaifont}{0}{渐照在阿波罗身上，他从宫殿里走出，慢慢走下来。他穿着}
\BookLine{47.00}{0.00}{\kaifont}{0}{斗篷，兜帽被取了下来落在肩上，露出头顶的花环。他先是}
\BookLine{55.50}{0.00}{\kaifont}{0}{站在门边同阿尔刻提斯说话，随后慢慢跟在她身后走着。场}
\BookLine{64.00}{0.00}{\kaifont}{0}{上只剩阿尔刻提斯一人，她闭着眼睛朝左边走了几步。她的}
\BookLine{72.50}{0.00}{\kaifont}{0}{头一直低着，像是背负着重担，她的身体缩成一团，像是已}
\BookLine{81.00}{0.00}{\kaifont}{0}{经很老了。她转向右边，慢慢向通往大路的门边走去，半睁}
\BookLine{89.50}{0.00}{\kaifont}{0}{着她的眼睛。}
\BookLine{106.50}{0.00}{\heifont}{0}{阿波罗：}
\BookLine{115.00}{22.14}{}{0}{再走几步，阿尔刻提斯。穿过这道门……穿过}
\BookLine{123.50}{22.14}{}{0}{那条路……走到我的林子里去。}
\BookLine{140.50}{0.00}{\heifont}{0}{阿尔刻提斯：}
\BookLine{149.00}{22.14}{}{0}{太远了……太高了……}
\BookLine{166.00}{0.00}{\heifont}{0}{阿波罗：}
\BookLine{174.50}{22.14}{}{0}{再走一步。这不是在上山，你不需要抬脚。}
\newpage
% 原书页码：223
\null
\BookNumber{223}
\Vertical{129}{11}{}{7}{10}第\\三\\幕\EndVertical
\BookLine{30.00}{0.00}{\heifont}{0}{阿尔刻提斯：}
\BookLine{38.50}{22.14}{}{0}{太远了。我在这儿找个坟墓就好了。}
\BookLine{55.50}{0.00}{\heifont}{0}{阿波罗：}
\BookLine{64.00}{22.14}{}{0}{阿尔刻提斯，你不会去到坟墓里的。}
\BookLine{81.00}{0.00}{\heifont}{0}{阿尔刻提斯：}
\BookLine{89.50}{22.14}{}{0}{哦，不，我只想找到我的坟墓。}
\BookLine{106.50}{0.00}{\heifont}{0}{阿波罗：}
\BookLine{115.00}{22.14}{}{0}{坟墓意味着终结，可你的生命不会这样终结。}
\BookLine{123.50}{22.14}{}{0}{在许许多多人里面，你是第一个无须面对终结}
\BookLine{132.00}{22.14}{}{0}{的人。再走一步，阿尔刻提斯。}
\BookLine{149.00}{0.00}{\heifont}{0}{阿尔刻提斯：}
\BookLine{157.50}{22.14}{}{0}{我会有子孙吗？后代的后代？}
\newpage
% 原书页码：224
\null
\BookNumber{224}
\Vertical{4}{169}{}{6.2}{8}阿\\尔\\刻\\提\\斯\\之\\歌\EndVertical
\BookLine{30.00}{0.00}{\heifont}{0}{阿波罗：}
\BookLine{38.50}{22.14}{}{0}{多得数不清。}
\BookLine{55.50}{0.00}{\heifont}{0}{阿尔刻提斯：}
\BookLine{64.00}{22.14}{}{0}{好……我忘了，他叫什么名字？}
\BookLine{81.00}{0.00}{\heifont}{0}{阿波罗：}
\BookLine{89.50}{22.14}{}{0}{阿德墨托斯。}
\BookLine{106.50}{0.00}{\heifont}{0}{阿尔刻提斯：}
\BookLine{115.00}{22.14}{}{0}{对，还有我想侍奉的，明亮如光的那一位？}
\BookLine{132.00}{0.00}{\heifont}{0}{阿波罗：}
\BookLine{140.50}{22.14}{}{0}{阿波罗。}
\BookLine{157.50}{0.00}{\heifont}{0}{阿尔刻提斯：}
\BookLine{166.00}{22.14}{}{0}{对……{\kaifont （走到门边）}许许多多的日子……都被}
\BookLine{174.50}{22.14}{}{0}{眷顾着……{\kaifont （抬起头，闭上眼睛）}这些幸福，}
\newpage
% 原书页码：225
\null
\BookNumber{225}
\Vertical{129}{11}{}{7}{10}第\\三\\幕\EndVertical
\BookLine{30.00}{22.14}{}{0}{我该感谢谁呢？}
\BookLine{47.00}{0.00}{\heifont}{0}{阿波罗：}
\BookLine{55.50}{22.14}{}{0}{朋友不会问彼此这样的问题。}
\BookLine{72.50}{7.38}{\kaifont}{0}{她走了出去。}
\BookLine{81.00}{7.38}{\kaifont}{0}{阿波罗提高了自己的嗓音，像是要保证她在墙外仍然能}
\BookLine{89.50}{0.00}{\kaifont}{0}{听到自己的声音。}
\BookLine{106.50}{22.14}{}{0}{那些真正爱过的人，不会问彼此这样的问题。}
\BookLine{115.00}{22.14}{}{0}{阿尔刻提斯。}
\BookLine{132.00}{7.38}{\kaifont}{0}{——剧终——}
\newpage
% 原书页码：226
\null
\newpage
% 原书页码：227
\null
\WorkOrnament{醉\\酒\\的\\女\\神}
\newpage
% 原书页码：228
\null
\BookNumber{228}
\Vertical{15}{28}{}{14}{21}人\\物\EndVertical
\BookLine{112.00}{7.38}{\kaifont}{0}{克洛托——命运三女神之一}
\BookLine{120.50}{7.38}{\kaifont}{0}{拉刻西斯——命运三女神之一}
\BookLine{129.00}{7.38}{\kaifont}{0}{阿特洛波斯——命运三女神之一}
\BookLine{137.50}{7.38}{\kaifont}{0}{阿波罗——太阳神}
\newpage
% 原书页码：229
\null
\BookNumber{229}
\Vertical{15}{28}{}{14}{21}转\\场\EndVertical
\BookLine{92.00}{7.38}{\kaifont}{0}{在《阿尔刻提斯之歌》的三幕戏结束后，幕布降下。阿波}
\BookLine{100.50}{0.00}{\kaifont}{0}{罗走到幕前。}
\BookLine{117.50}{0.00}{\heifont}{0}{阿波罗：}
\BookLine{126.00}{22.14}{}{0}{{\kaifont （对观众）}等一下！等一下！还有一件事。我们}
\BookLine{134.50}{22.14}{}{0}{现在是在希腊，在这里我们可不希望观众刚看}
\BookLine{143.00}{22.14}{}{0}{完一个讲述人类艰难悲惨命运的故事就立刻回}
\BookLine{151.50}{22.14}{}{0}{家。这里有个习俗：我们会要求诗人们写一出}
\BookLine{160.00}{22.14}{}{0}{短小的羊人剧来娱乐大家，让大家发笑。{\kaifont （悄悄}}
\BookLine{168.50}{22.14}{}{0}{{\kaifont 地）}我们认为悲剧视角不能单独存在，容易显}
\BookLine{177.00}{22.14}{}{0}{得太过。就像你们中间有人{\kaifont （指着观众席）}说}
\BookLine{185.50}{22.14}{}{0}{过：无论死亡或是阳光，都经不起人一再地凝}
\newpage
% 原书页码：230
\null
\BookNumber{230}
\Vertical{4}{169}{}{6.2}{8}阿\\尔\\刻\\提\\斯\\之\\歌\EndVertical
\BookLine{30.00}{22.14}{}{0}{视。以及，我们会要求这出羊人剧包含一些元}
\BookLine{38.50}{22.14}{}{0}{素——关于之前这出戏的——一些其他的视}
\BookLine{47.00}{22.14}{}{0}{角。所以——应该给大家演点什么呢？讲讲忒}
\BookLine{55.50}{22.14}{}{0}{瑞西阿斯——那个六百岁的老头，被过分宠爱}
\BookLine{64.00}{22.14}{}{0}{的老头？神明们时不时地让他变年轻，时不时}
\BookLine{72.50}{22.14}{}{0}{地把他变成一个女人？还有宙斯和赫拉因为这}
\BookLine{81.00}{22.14}{}{0}{个事情吵架？{\kaifont （笑了起来，不过努力控制着自}}
\BookLine{89.50}{22.14}{}{0}{{\kaifont 己）}不，不——不太合适，太粗俗了。希腊人}
\BookLine{98.00}{22.14}{}{0}{是能包含这种不太得体的故事的，不过{\kaifont （又无}}
\BookLine{106.50}{22.14}{}{0}{{\kaifont 法自抑地笑了起来）}它是……不，不——不能}
\BookLine{115.00}{22.14}{}{0}{在这儿演。要么我们讲讲阿尔刻提斯王后的姐}
\BookLine{123.50}{22.14}{}{0}{妹？她的父亲佩利阿斯就是个蠢老头，两个姐}
\BookLine{132.00}{22.14}{}{0}{姐也不太聪明。一个家总是这样——只能有一}
\BookLine{140.50}{22.14}{}{0}{个聪明人。我们也可以讲讲美狄亚是怎么给这}
\BookLine{149.00}{22.14}{}{0}{两个姐姐吹耳边风，假装告诉她们如何能使父}
\BookLine{157.50}{22.14}{}{0}{亲返老还童……{\kaifont （突然改变了声调）}不！是有}
\BookLine{166.00}{22.14}{}{0}{些人，只有听到残忍的事情才能笑出来。但这}
\BookLine{174.50}{22.14}{}{0}{个故事残忍至极，等你们回到家，想起自己在}
\newpage
% 原书页码：231
\null
\BookNumber{231}
\Vertical{129}{11}{}{7}{10}醉\\酒\\的\\女\\神\EndVertical
\BookLine{30.00}{22.14}{}{0}{为这样的故事感到惊奇，也会挺羞愧的。还有}
\BookLine{38.50}{22.14}{}{0}{一个故事：不算特别有趣。今晚看戏时你有没}
\BookLine{47.00}{22.14}{}{0}{有好奇过，阿德墨托斯国王的生命是如何得到}
\BookLine{55.50}{22.14}{}{0}{延续的？那几个厉害的女人，行事怪异的姐妹，}
\BookLine{64.00}{22.14}{}{0}{命运三女神——她们也会被骗吗？我们来讲讲}
\BookLine{72.50}{22.14}{}{0}{这件事是如何发生的。}
\BookLine{89.50}{7.38}{\kaifont}{0}{他脱掉外面的长袍，里面是一身后厨小伙计的打扮，他}
\BookLine{98.00}{0.00}{\kaifont}{0}{冲着台侧喊着：“我的帽子！帽子！”}
\BookLine{115.00}{22.14}{}{0}{在这出小戏里，我还是阿波罗，不过假扮成了}
\BookLine{123.50}{22.14}{}{0}{一个后厨的小伙计。}
\BookLine{140.50}{7.38}{\kaifont}{0}{从台侧递出来一顶锥形草帽和一条腰带，腰带上挂着洋}
\BookLine{149.00}{0.00}{\kaifont}{0}{葱之类的东西——表示这是个在厨房干活的伙计。他又冲着}
\BookLine{157.50}{0.00}{\kaifont}{0}{台侧喊：“我的木壶！”一边喊一边戴上帽子，系上腰带。}
\newpage
% 原书页码：232
\null
\BookNumber{232}
\Vertical{4}{169}{}{6.2}{8}阿\\尔\\刻\\提\\斯\\之\\歌\EndVertical
\BookLine{30.00}{0.00}{\heifont}{0}{阿波罗：}
\BookLine{38.50}{22.14}{}{0}{我讨厌乔装打扮，讨厌喝醉酒的人——}
\BookLine{55.50}{7.38}{\kaifont}{0}{从台侧递出来一根绳子，上面挂着三只壶。}
\BookLine{72.50}{0.00}{\heifont}{0}{阿波罗：}
\BookLine{81.00}{22.14}{}{0}{看到我脖子上挂的这几个壶了吗？我讨厌撒谎，}
\BookLine{89.50}{22.14}{}{0}{讨厌耍花招——但是我要耍花招去干的这件事}
\BookLine{98.00}{22.14}{}{0}{情，就算是全能的宙斯来，也只能靠骗：我要}
\BookLine{106.50}{22.14}{}{0}{去延长一个人的生命。}
\BookLine{123.50}{7.38}{\kaifont}{0}{他朝着幕布后方或是台侧喊道：“都准备好了吗？”又转}
\BookLine{132.00}{0.00}{\kaifont}{0}{过身来对着观众。}
\BookLine{149.00}{22.14}{}{0}{好了，羊人剧的表演准备开始，为这出严肃的}
\BookLine{157.50}{22.14}{}{0}{三幕剧《阿尔刻提斯之歌》收尾。}
\BookLine{174.50}{7.38}{\kaifont}{0}{幕布升起，他依然站在舞台前侧的柱子旁边。}
\newpage
% 原书页码：233
\null
\BookNumber{233}
\Vertical{129}{11}{}{7}{10}醉\\酒\\的\\女\\神\EndVertical
\BookLine{30.00}{7.38}{\kaifont}{0}{三位命运女神正坐在一张长凳上，凳子被宽大的褶摆所}
\BookLine{38.50}{0.00}{\kaifont}{0}{覆盖。她们戴着老妇人的面具，风格怪诞却不失威严。克洛}
\BookLine{47.00}{0.00}{\kaifont}{0}{托拿着纺锤，拉刻西斯腿上放着生命之线，阿特洛波斯拿着}
\BookLine{55.50}{0.00}{\kaifont}{0}{剪刀。道具设计要让每一样工具都显得很大，接近两腿之宽。}
\BookLine{64.00}{0.00}{\kaifont}{0}{她们工作时身体前后晃动着，从右至左传递着生命之线。观}
\BookLine{72.50}{0.00}{\kaifont}{0}{众看着她们沉默了好一阵子，只有克洛托偶尔轻声哼两句。}
\BookLine{89.50}{0.00}{\heifont}{0}{克洛托：}
\BookLine{98.00}{22.14}{}{0}{什么东西，出生的时候四条腿，后来变成两条}
\BookLine{106.50}{22.14}{}{0}{腿？别说！别说！}
\BookLine{123.50}{0.00}{\heifont}{0}{拉刻西斯：}
\BookLine{132.00}{22.14}{}{0}{{\kaifont （无聊）}你知道答案！}
\BookLine{149.00}{0.00}{\heifont}{0}{克洛托：}
\BookLine{157.50}{22.14}{}{0}{让我假装一下我不知道。}
\newpage
% 原书页码：234
\null
\BookNumber{234}
\Vertical{4}{169}{}{6.2}{8}阿\\尔\\刻\\提\\斯\\之\\歌\EndVertical
\BookLine{30.00}{0.00}{\heifont}{0}{阿特洛波斯：}
\BookLine{38.50}{22.14}{}{0}{没有新谜语了。咱们全都猜过了。}
\BookLine{55.50}{0.00}{\heifont}{0}{拉刻西斯：}
\BookLine{64.00}{22.14}{}{0}{没有谜语的人生真无聊！克洛托，你想一个}
\BookLine{72.50}{22.14}{}{0}{谜语。}
\BookLine{89.50}{0.00}{\heifont}{0}{克洛托：}
\BookLine{98.00}{22.14}{}{0}{安静，等等，给我一分钟想一想……那是什}
\BookLine{106.50}{22.14}{}{0}{么……那是什么？}
\BookLine{123.50}{7.38}{\kaifont}{0}{乔装打扮的阿波罗上场。}
\BookLine{140.50}{0.00}{\heifont}{0}{阿波罗：}
\BookLine{149.00}{22.14}{}{0}{{\kaifont （对观众）}这就是伟大的三姐妹——命运女神。}
\BookLine{157.50}{22.14}{}{0}{克洛托编织生命之线，拉刻西斯丈量每一根线}
\BookLine{166.00}{22.14}{}{0}{的长度，阿特洛波斯把它们剪断。这份决定我}
\BookLine{174.50}{22.14}{}{0}{们命运的工作千篇一律，她们干得非常无聊，}
\newpage
% 原书页码：235
\null
\BookNumber{235}
\Vertical{129}{11}{}{7}{10}醉\\酒\\的\\女\\神\EndVertical
\BookLine{30.00}{22.14}{}{0}{就像其他无聊的人一样，她们也通过游戏取}
\BookLine{38.50}{22.14}{}{0}{乐——解谜题，猜谜语。她们也不能玩牌，因}
\BookLine{47.00}{22.14}{}{0}{为双手都被命运之线的工作给占着。}
\BookLine{64.00}{0.00}{\heifont}{0}{阿特洛波斯：}
\BookLine{72.50}{22.14}{}{0}{姐姐！你的胳膊！你干活儿的时候别撞我呀。}
\BookLine{89.50}{0.00}{\heifont}{0}{拉刻西斯：}
\BookLine{98.00}{22.14}{}{0}{没办法——线太太太长了！我没遇到过这么长}
\BookLine{106.50}{22.14}{}{0}{的线。}
\BookLine{123.50}{0.00}{\heifont}{0}{克洛托：}
\BookLine{132.00}{22.14}{}{0}{又长，颜色又暗，又脏兮兮的！这么多年就是}
\BookLine{140.50}{22.14}{}{0}{个奴隶！}
\BookLine{157.50}{0.00}{\heifont}{0}{拉刻西斯：}
\BookLine{166.00}{22.14}{}{0}{是啊！{\kaifont （对阿特洛波斯）}剪了吧，妹妹。{\kaifont （阿特}}
\BookLine{174.50}{22.14}{}{0}{{\kaifont 洛波斯剪线——咔！）}这个，剪了。这条是蓝}
\newpage
% 原书页码：236
\null
\BookNumber{236}
\Vertical{4}{169}{}{6.2}{8}阿\\尔\\刻\\提\\斯\\之\\歌\EndVertical
\BookLine{30.00}{22.14}{}{0}{线——蓝线代表勇敢；很勇敢，很短暂。}
\BookLine{47.00}{0.00}{\heifont}{0}{阿特洛波斯：}
\BookLine{55.50}{22.14}{}{0}{看得出来！{\kaifont （咔）}}
\BookLine{72.50}{0.00}{\heifont}{0}{拉刻西斯：}
\BookLine{81.00}{22.14}{}{0}{你差点剪到那条紫色的线，阿特洛波斯。}
\BookLine{98.00}{0.00}{\heifont}{0}{阿特洛波斯：}
\BookLine{106.50}{22.14}{}{0}{这条？紫色是代表国王？}
\BookLine{123.50}{0.00}{\heifont}{0}{拉刻西斯：}
\BookLine{132.00}{22.14}{}{0}{是的，你干活儿的时候小心点，亲爱的。这是}
\BookLine{140.50}{22.14}{}{0}{塞萨利国王阿德墨托斯的生命线。}
\BookLine{157.50}{0.00}{\heifont}{0}{阿波罗：}
\BookLine{166.00}{22.14}{}{0}{{\kaifont （在一旁）}哎呀！}
\newpage
% 原书页码：237
\null
\BookNumber{237}
\Vertical{129}{11}{}{7}{10}醉\\酒\\的\\女\\神\EndVertical
\BookLine{30.00}{0.00}{\heifont}{0}{拉刻西斯：}
\BookLine{38.50}{22.14}{}{0}{我标记得很清楚。他日落的时候就会死去。}
\BookLine{55.50}{0.00}{\heifont}{0}{阿波罗：}
\BookLine{64.00}{22.14}{}{0}{{\kaifont （对观众）}不！不！}
\BookLine{81.00}{0.00}{\heifont}{0}{拉刻西斯：}
\BookLine{89.50}{22.14}{}{0}{那塞萨利可要满城哀号了。大地震动，人们撕}
\BookLine{98.00}{22.14}{}{0}{碎衣裳……不是现在，亲爱的。还有一个小}
\BookLine{106.50}{22.14}{}{0}{时呢。}
\BookLine{123.50}{0.00}{\heifont}{0}{阿波罗：}
\BookLine{132.00}{22.14}{}{0}{{\kaifont （在一旁）}做点什么！快做点什么，伪装着阿}
\BookLine{140.50}{22.14}{}{0}{波罗！}
\BookLine{157.50}{7.38}{\kaifont}{0}{他做出疯狂奔跑的动作，但人仍在原地。他突然停下来，}
\BookLine{166.00}{0.00}{\kaifont}{0}{对观众说话。}
\newpage
% 原书页码：238
\null
\BookNumber{238}
\Vertical{4}{169}{}{6.2}{8}阿\\尔\\刻\\提\\斯\\之\\歌\EndVertical
\BookLine{30.00}{22.14}{}{0}{在场的有谁没听过那个故事——为什么我那么}
\BookLine{38.50}{22.14}{}{0}{喜欢阿德墨托斯国王和阿尔刻提斯王后？}
\BookLine{55.50}{7.38}{\kaifont}{0}{他坐到地上，伸出自己的食指，继续说着。}
\BookLine{72.50}{22.14}{}{0}{是十年以前吧？我对时间没什么概念，我是太}
\BookLine{81.00}{22.14}{}{0}{阳神，走到哪里哪里就是白天。差不多是十年}
\BookLine{89.50}{22.14}{}{0}{前，我的父亲，咱们大家的父亲，对我非常生}
\BookLine{98.00}{22.14}{}{0}{气。我干了什么吗？}
\BookLine{115.00}{7.38}{\kaifont}{0}{他来回摆弄着手指。}
\BookLine{132.00}{22.14}{}{0}{先别问这个，就忘了吧……总之，他对我降下}
\BookLine{140.50}{22.14}{}{0}{了一个惩罚。他命令我到地上来，在人类当中}
\BookLine{149.00}{22.14}{}{0}{生活一年。——我，作为一个普通人，一个仆}
\BookLine{157.50}{22.14}{}{0}{人。我选择了到塞萨利去，成为阿德墨托斯国}
\BookLine{166.00}{22.14}{}{0}{王的牧人。我半隐藏了自己的身份，大家知道}
\BookLine{174.50}{22.14}{}{0}{点什么，但也不太清楚。我就像一个普通人一}
\newpage
% 原书页码：239
\null
\BookNumber{239}
\Vertical{129}{11}{}{7}{10}醉\\酒\\的\\女\\神\EndVertical
\BookLine{30.00}{22.14}{}{0}{样生活，在他们身边，就像此刻在你们身边，}
\BookLine{38.50}{22.14}{}{0}{在公义的人身边，也在不义的人身边。每一天，}
\BookLine{47.00}{22.14}{}{0}{国王都会对我和其他牧人下指令，每一天，王}
\BookLine{55.50}{22.14}{}{0}{后也会关心我们的生活是否顺利，家人是否安}
\BookLine{64.00}{22.14}{}{0}{好。我渐渐喜欢上了阿德墨托斯国王和阿尔刻}
\BookLine{72.50}{22.14}{}{0}{提斯王后，因为他们，我渐渐喜欢上了人类。}
\BookLine{81.00}{22.14}{}{0}{现在阿德墨托斯要死了。{\kaifont （起身）}不行！我有我}
\BookLine{89.50}{22.14}{}{0}{的计划，我要阻止这件事情，要做点什么。做}
\BookLine{98.00}{22.14}{}{0}{点什么，伪装者阿波罗。}
\BookLine{115.00}{7.38}{\kaifont}{0}{他又开始在原地疯狂跑步，大声地抱怨着。}
\BookLine{132.00}{22.14}{}{0}{哦，我的后背！哎呀，哎呀。他们打我，更惨}
\BookLine{140.50}{22.14}{}{0}{的是我要迟到了，我又要挨打了。}
\BookLine{157.50}{0.00}{\heifont}{0}{拉刻西斯：}
\BookLine{166.00}{22.14}{}{0}{这个哭哭啼啼的人是谁？}
\newpage
% 原书页码：240
\null
\BookNumber{240}
\Vertical{4}{169}{}{6.2}{8}阿\\尔\\刻\\提\\斯\\之\\歌\EndVertical
\BookLine{30.00}{0.00}{\heifont}{0}{阿波罗：}
\BookLine{38.50}{22.14}{}{0}{别拦着我，我都没有机会说说话。我已经迟到}
\BookLine{47.00}{22.14}{}{0}{了，而且，我的差事还是个秘密，什么都不}
\BookLine{55.50}{22.14}{}{0}{能说。}
\BookLine{72.50}{0.00}{\heifont}{0}{阿特洛波斯：}
\BookLine{81.00}{22.14}{}{0}{拉刻西斯，用你的线把他围住。这个傻子能有}
\BookLine{89.50}{22.14}{}{0}{什么秘密？我们才是知道所有秘密的人。}
\BookLine{106.50}{7.38}{\kaifont}{0}{放在三姐妹腿上的线，观众是看不到的。拉刻西斯站了}
\BookLine{115.00}{0.00}{\kaifont}{0}{起来，她的手在头顶挥舞转了三个大圈，像是要丢出一个套}
\BookLine{123.50}{0.00}{\kaifont}{0}{索，随即套索圈住了舞台上的东西。阿波罗摆出被捉住的姿}
\BookLine{132.00}{0.00}{\kaifont}{0}{势。拉刻西斯每次用力一拉，阿波罗就被拽得离她更近。在}
\BookLine{140.50}{0.00}{\kaifont}{0}{接下来的一段台词中，拉刻西斯把线的另一头举在空中，一}
\BookLine{149.00}{0.00}{\kaifont}{0}{会儿把阿波罗高高抬起来，几乎要勒死他，一会儿又把他摔}
\BookLine{157.50}{0.00}{\kaifont}{0}{在地上。}
\newpage
% 原书页码：241
\null
\BookNumber{241}
\Vertical{129}{11}{}{7}{10}醉\\酒\\的\\女\\神\EndVertical
\BookLine{30.00}{0.00}{\heifont}{0}{阿波罗：}
\BookLine{38.50}{22.14}{}{0}{夫人们，美丽的夫人们，放过我吧。要是我迟}
\BookLine{47.00}{22.14}{}{0}{到的话，整个奥林匹斯山都要闹起来了。阿佛}
\BookLine{55.50}{22.14}{}{0}{洛狄忒会担心得发疯的——哦，我说得太多}
\BookLine{64.00}{22.14}{}{0}{了。我的任务是立刻前往，什么都别说——尤}
\BookLine{72.50}{22.14}{}{0}{其是不要和女人说。这件事男人不会感兴趣。}
\BookLine{81.00}{22.14}{}{0}{夫人们，放我走吧。}
\BookLine{98.00}{0.00}{\heifont}{0}{阿特洛波斯：}
\BookLine{106.50}{22.14}{}{0}{姐姐，把你的线拽紧。}
\BookLine{123.50}{0.00}{\heifont}{0}{阿波罗：}
\BookLine{132.00}{22.14}{}{0}{我要窒息了，你要把我勒死了。}
\BookLine{149.00}{0.00}{\heifont}{0}{拉刻西斯：}
\BookLine{157.50}{22.14}{}{0}{{\kaifont （强迫他）}别发牢骚了，现在就把你的秘密说}
\BookLine{166.00}{22.14}{}{0}{出来。}
\newpage
% 原书页码：242
\null
\BookNumber{242}
\Vertical{4}{169}{}{6.2}{8}阿\\尔\\刻\\提\\斯\\之\\歌\EndVertical
\BookLine{30.00}{0.00}{\heifont}{0}{阿波罗：}
\BookLine{38.50}{22.14}{}{0}{我不能说。我不敢。}
\BookLine{55.50}{0.00}{\heifont}{0}{阿特洛波斯：}
\BookLine{64.00}{22.14}{}{0}{再紧一点，姐姐。孩子，要么说秘密，要么就}
\BookLine{72.50}{22.14}{}{0}{窒息。{\kaifont （做把他勒住的手势）}}
\BookLine{89.50}{0.00}{\heifont}{0}{阿波罗：}
\BookLine{98.00}{22.14}{}{0}{噢！噢！等一下！你们要是能放我走，我说}
\BookLine{106.50}{22.14}{}{0}{一半。}
\BookLine{123.50}{0.00}{\heifont}{0}{阿特洛波斯：}
\BookLine{132.00}{22.14}{}{0}{全说出来，要么我们就把你挂在这根绳子上。}
\BookLine{149.00}{0.00}{\heifont}{0}{阿波罗：}
\BookLine{157.50}{22.14}{}{0}{我说，我说。不过——{\kaifont （担忧地看看自己）}你们}
\BookLine{166.00}{22.14}{}{0}{发誓，指着冥河起誓，你们不会告诉任何人，}
\BookLine{174.50}{22.14}{}{0}{指着遗忘河起誓，你们会把它忘掉。}
\newpage
% 原书页码：243
\null
\BookNumber{243}
\Vertical{129}{11}{}{7}{10}醉\\酒\\的\\女\\神\EndVertical
\BookLine{30.00}{0.00}{\heifont}{0}{拉刻西斯：}
\BookLine{38.50}{22.14}{}{0}{我们只有一种誓言——指着黄泉起誓。不过我}
\BookLine{47.00}{22.14}{}{0}{们不会发誓的——更何况是为你一个哭哭啼啼的}
\BookLine{55.50}{22.14}{}{0}{小奴隶。你现在就把知道的事情告诉我们，否}
\BookLine{64.00}{22.14}{}{0}{则你立刻就能见到这几条河。}
\BookLine{81.00}{0.00}{\heifont}{0}{阿波罗：}
\BookLine{89.50}{22.14}{}{0}{我害怕，我真的是不敢说。我……嗯……我}
\BookLine{98.00}{22.14}{}{0}{带着……来的……那些瓶子……哦，夫人们，}
\BookLine{106.50}{22.14}{}{0}{放我走吧，放我走吧。}
\BookLine{123.50}{0.00}{\heifont}{0}{克洛托和阿特洛波斯：}
\BookLine{132.00}{22.14}{}{0}{使劲拉，姐姐。}
\BookLine{149.00}{0.00}{\heifont}{0}{阿波罗：}
\BookLine{157.50}{22.14}{}{0}{不要！不要！我说！我带的酒是要送给……送}
\BookLine{166.00}{22.14}{}{0}{给阿佛洛狄忒。每十天她就要喝这些酒……来}
\BookLine{174.50}{22.14}{}{0}{恢复自己的美貌。}
\newpage
% 原书页码：244
\null
\BookNumber{244}
\Vertical{4}{169}{}{6.2}{8}阿\\尔\\刻\\提\\斯\\之\\歌\EndVertical
\BookLine{30.00}{0.00}{\heifont}{0}{阿特洛波斯：}
\BookLine{38.50}{22.14}{}{0}{骗子！傻子！她喝神仙的酒，吃神仙的食物，}
\BookLine{47.00}{22.14}{}{0}{和大家都一样。}
\BookLine{64.00}{0.00}{\heifont}{0}{阿波罗：}
\BookLine{72.50}{22.14}{}{0}{{\kaifont （悄悄地）}但她却是最美的？……这是赫菲斯}
\BookLine{81.00}{22.14}{}{0}{托斯的礼物，来自狄奥尼索斯的葡萄园，这些}
\BookLine{89.50}{22.14}{}{0}{酿酒的葡萄是在阿波罗的注视下成熟的——那}
\BookLine{98.00}{22.14}{}{0}{位从不说谎的阿波罗。}
\BookLine{115.00}{0.00}{\heifont}{0}{三姐妹：}
\BookLine{123.50}{22.14}{}{0}{{\kaifont （心照不宣地，带着雀跃的期待）}姐妹们！}
\BookLine{140.50}{0.00}{\heifont}{0}{阿特洛波斯：}
\BookLine{149.00}{22.14}{}{0}{{\kaifont （甜甜地）}好孩子，把那些瓶子拿过来。}
\BookLine{166.00}{0.00}{\heifont}{0}{阿波罗：}
\BookLine{174.50}{22.14}{}{0}{{\kaifont （害怕地）}这可不行！夫人！我已经把秘密都告}
\newpage
% 原书页码：245
\null
\BookNumber{245}
\Vertical{129}{11}{}{7}{10}醉\\酒\\的\\女\\神\EndVertical
\BookLine{30.00}{22.14}{}{0}{诉你们了！这可不行！}
\BookLine{47.00}{0.00}{\heifont}{0}{阿特洛波斯：}
\BookLine{55.50}{22.14}{}{0}{拉刻西斯，你肯定，能从那堆线里面找到这个}
\BookLine{64.00}{22.14}{}{0}{不值一提的奴隶的生命线吧——一根黄色的长}
\BookLine{72.50}{22.14}{}{0}{长的线？}
\BookLine{89.50}{0.00}{\heifont}{0}{阿波罗：}
\BookLine{98.00}{22.14}{}{0}{{\kaifont （跪在地上）}饶了我吧！}
\BookLine{115.00}{0.00}{\heifont}{0}{阿特洛波斯：}
\BookLine{123.50}{22.14}{}{0}{{\kaifont （对拉刻西斯）}看，就在这儿——土黄土黄的，}
\BookLine{132.00}{22.14}{}{0}{因为不诚实，所以缠成一坨；因为固执，所以}
\BookLine{140.50}{22.14}{}{0}{线很硬；因为愚蠢，所以粗细不均。亲爱的，}
\BookLine{149.00}{22.14}{}{0}{把线递给我。}
\BookLine{166.00}{0.00}{\heifont}{0}{阿波罗：}
\BookLine{174.50}{22.14}{}{0}{{\kaifont （前额贴着地面）}哦，真希望我从没有出}
\newpage
% 原书页码：246
\null
\BookNumber{246}
\Vertical{4}{169}{}{6.2}{8}阿\\尔\\刻\\提\\斯\\之\\歌\EndVertical
\BookLine{30.00}{22.14}{}{0}{生过！}
\BookLine{47.00}{0.00}{\heifont}{0}{拉刻西斯：}
\BookLine{55.50}{22.14}{}{1}{{\kaifont （递给阿特洛波斯）}给。{\kaifont （打了个手势）}我本来是}
\BookLine{64.00}{22.14}{}{0}{想给剪成五段的。}
\BookLine{81.00}{0.00}{\heifont}{0}{阿波罗：}
\BookLine{89.50}{22.14}{}{0}{{\kaifont （起身，把酒瓶递给她们，啜泣着）}给，你们拿}
\BookLine{98.00}{22.14}{}{0}{走！都拿走！反正我怎么都是死。阿佛洛狄忒}
\BookLine{106.50}{22.14}{}{0}{也会杀了我的，我的命就这么长了。}
\BookLine{123.50}{7.38}{\kaifont}{0}{三姐妹拿起瓶子，态度强硬。}
\BookLine{140.50}{0.00}{\heifont}{0}{阿特洛波斯：}
\BookLine{149.00}{22.14}{}{0}{一个字都别说了。把嘴闭上，我们不想听你说}
\BookLine{157.50}{22.14}{}{0}{话了。}
\BookLine{174.50}{7.38}{\kaifont}{0}{阿波罗身上的绳索被松开，他头朝地摔下来，肩膀撞}
\BookLine{183.00}{0.00}{\kaifont}{0}{到地上。三姐妹对着酒壶喝了起来，她们一边喝一边开心地}
\BookLine{191.50}{0.00}{\kaifont}{0}{喊着。}
\newpage
% 原书页码：247
\null
\BookNumber{247}
\Vertical{129}{11}{}{7}{10}醉\\酒\\的\\女\\神\EndVertical
\BookLine{30.00}{0.00}{\heifont}{0}{三姐妹：}
\BookLine{38.50}{22.14}{}{0}{姐妹们！}
\BookLine{55.50}{0.00}{\heifont}{0}{拉刻西斯：}
\BookLine{64.00}{22.14}{}{0}{妹妹，我看起来怎么样？}
\BookLine{81.00}{0.00}{\heifont}{0}{阿特洛波斯：}
\BookLine{89.50}{22.14}{}{0}{噢，我简直想吃了你，我呢？}
\BookLine{106.50}{0.00}{\heifont}{0}{克洛托：}
\BookLine{115.00}{22.14}{}{0}{姐姐，我看着怎么样？}
\BookLine{132.00}{0.00}{\heifont}{0}{拉刻西斯：}
\BookLine{140.50}{22.14}{}{0}{美！太美了！我呢？}
\newpage
% 原书页码：248
\null
\BookNumber{248}
\Vertical{4}{169}{}{6.2}{8}阿\\尔\\刻\\提\\斯\\之\\歌\EndVertical
\BookLine{30.00}{0.00}{\heifont}{0}{阿特洛波斯：}
\BookLine{38.50}{22.14}{}{0}{可惜这山上一面镜子都没有，一小片水池也没}
\BookLine{47.00}{22.14}{}{0}{有，没法让我们知道谁是最美的。}
\BookLine{64.00}{0.00}{\heifont}{0}{拉刻西斯：}
\BookLine{72.50}{22.14}{}{0}{{\kaifont （做梦一般，把手放在自己的脸上）}我感觉……}
\BookLine{81.00}{22.14}{}{0}{克罗诺斯在追我，要在一个黑暗的角落把我}
\BookLine{89.50}{22.14}{}{0}{抓住。}
\BookLine{106.50}{0.00}{\heifont}{0}{阿特洛波斯：}
\BookLine{115.00}{22.14}{}{0}{波塞冬要发疯了——在大陆上横冲直撞，想要}
\BookLine{123.50}{22.14}{}{0}{把我吞掉。}
\BookLine{140.50}{0.00}{\heifont}{0}{克洛托：}
\BookLine{149.00}{22.14}{}{0}{我自己的父亲——谁又能怪他呢？——他已}
\BookLine{157.50}{22.14}{}{0}{经不是他自己。}
\newpage
% 原书页码：249
\null
\BookNumber{249}
\Vertical{129}{11}{}{7}{10}醉\\酒\\的\\女\\神\EndVertical
\BookLine{30.00}{0.00}{\heifont}{0}{阿特洛波斯：}
\BookLine{38.50}{22.14}{}{0}{{\kaifont （轻声地）}他也不是一点用处都没有。而且，他}
\BookLine{47.00}{22.14}{}{0}{也不丑。{\kaifont （对克洛托）}问问他，他看到了什么。}
\BookLine{64.00}{0.00}{\heifont}{0}{拉刻西斯：}
\BookLine{72.50}{22.14}{}{0}{问问他，我们之中谁是最美的。}
\BookLine{89.50}{0.00}{\heifont}{0}{克洛托：}
\BookLine{98.00}{22.14}{}{0}{孩子！孩子！你说话可以。我是说，你……你}
\BookLine{106.50}{22.14}{}{0}{可以喙{\kaifont （说）}。拉刻塞{\kaifont （拉刻西）}斯，你跟他说}
\BookLine{115.00}{22.14}{}{0}{{\kaifont （说）}，我涩头{\kaifont （舌头）}不好使了。}
\BookLine{132.00}{0.00}{\heifont}{0}{拉刻西斯：}
\BookLine{140.50}{22.14}{}{0}{孩子，你看看我们！你来说说，谁是最美的。}
\BookLine{157.50}{7.38}{\kaifont}{0}{阿波罗一直脸朝着地面，他站起身来，盯着三姐妹看了}
\BookLine{166.00}{0.00}{\kaifont}{0}{看。他表现得好像自己要瞎了一样，捂住眼睛又松开手，盯}
\BookLine{174.50}{0.00}{\kaifont}{0}{着她们每个人看了半天。}
\newpage
% 原书页码：250
\null
\BookNumber{250}
\Vertical{4}{169}{}{6.2}{8}阿\\尔\\刻\\提\\斯\\之\\歌\EndVertical
\BookLine{30.00}{0.00}{\heifont}{0}{阿波罗：}
\BookLine{38.50}{22.14}{}{0}{我干了什么呀？这些美丽的光辉！我干了什}
\BookLine{47.00}{22.14}{}{0}{么？——你——还有你——还有你！如果你们}
\BookLine{55.50}{22.14}{}{0}{想，就杀了我吧，但我真的说不出来你们哪一}
\BookLine{64.00}{22.14}{}{0}{位更美。{\kaifont （跪在地上）}哦，夫人们——如果更}
\BookLine{72.50}{22.14}{}{0}{多的美貌并没有让你们变得更加残忍，那就放}
\BookLine{81.00}{22.14}{}{0}{我走吧，让我躲起来。阿佛洛狄忒肯定会知道}
\BookLine{89.50}{22.14}{}{0}{这件事。让我逃回克里特去，回去做我从前的}
\BookLine{98.00}{22.14}{}{0}{工作。}
\BookLine{115.00}{0.00}{\heifont}{0}{阿特洛波斯：}
\BookLine{123.50}{22.14}{}{0}{好孩子，你以前是做什么的？}
\BookLine{140.50}{0.00}{\heifont}{0}{阿波罗：}
\BookLine{149.00}{22.14}{}{0}{我在集市上跟我父亲一起工作，我讲故事，说}
\BookLine{157.50}{22.14}{}{0}{谜语。}
\BookLine{174.50}{7.38}{\kaifont}{0}{三姐妹惊呆了，她们简直要尖叫起来。}
\newpage
% 原书页码：251
\null
\BookNumber{251}
\Vertical{129}{11}{}{7}{10}醉\\酒\\的\\女\\神\EndVertical
\BookLine{30.00}{0.00}{\heifont}{0}{三姐妹：}
\BookLine{38.50}{22.14}{}{0}{什么？你刚才说什么？}
\BookLine{55.50}{0.00}{\heifont}{0}{阿波罗：}
\BookLine{64.00}{22.14}{}{0}{讲故事，说谜语。美丽的夫人们也喜欢谜}
\BookLine{72.50}{22.14}{}{0}{语吗？}
\BookLine{89.50}{0.00}{\heifont}{0}{三姐妹：}
\BookLine{98.00}{22.14}{}{0}{{\kaifont （来回摇晃着身体，互相拍打着对方）}姐妹们，}
\BookLine{106.50}{22.14}{}{0}{我们喜欢谜语吗？}
\BookLine{123.50}{0.00}{\heifont}{0}{阿特洛波斯：}
\BookLine{132.00}{22.14}{}{0}{哦，他知道的也就是那些老谜语了。噢！跛了}
\BookLine{140.50}{22.14}{}{0}{的马……大脚趾……}
\BookLine{157.50}{0.00}{\heifont}{0}{拉刻西斯：}
\BookLine{166.00}{22.14}{}{0}{天上的云……赫拉的眼睫毛……}
\newpage
% 原书页码：252
\null
\BookNumber{252}
\Vertical{4}{169}{}{6.2}{8}阿\\尔\\刻\\提\\斯\\之\\歌\EndVertical
\BookLine{30.00}{0.00}{\heifont}{0}{克洛托：}
\BookLine{38.50}{22.14}{}{0}{{\kaifont （喋喋不休）}什么东西，出生的时候四条腿？}
\BookLine{55.50}{0.00}{\heifont}{0}{阿特洛波斯：}
\BookLine{64.00}{22.14}{}{0}{海豚……埃奇纳……}
\BookLine{81.00}{0.00}{\heifont}{0}{阿波罗：}
\BookLine{89.50}{22.14}{}{0}{这些大家都知道了！我有一些新的谜语——}
\BookLine{106.50}{0.00}{\heifont}{0}{三姐妹：}
\BookLine{115.00}{22.14}{}{0}{{\kaifont （又一次尖叫）}新谜语！}
\BookLine{132.00}{0.00}{\heifont}{0}{阿波罗：}
\BookLine{140.50}{22.14}{}{0}{{\kaifont （缓缓地）}是什么必须——}
\BookLine{157.50}{7.38}{\kaifont}{0}{他停了下来，三姐妹完全被他吸引住了。}
\newpage
% 原书页码：253
\null
\BookNumber{253}
\Vertical{129}{11}{}{7}{10}醉\\酒\\的\\女\\神\EndVertical
\BookLine{30.00}{0.00}{\heifont}{0}{拉刻西斯：}
\BookLine{38.50}{22.14}{}{0}{然后，孩子，然后。是什么必须——}
\BookLine{55.50}{0.00}{\heifont}{0}{阿波罗：}
\BookLine{64.00}{22.14}{}{0}{不过——要有赏罚我才玩。这样！如果我}
\BookLine{72.50}{22.14}{}{0}{输了……}
\BookLine{89.50}{0.00}{\heifont}{0}{克洛托：}
\BookLine{98.00}{22.14}{}{0}{如果你酥{\kaifont （输）}了，你必西{\kaifont （须）}告诉{\kaifont （诉）}我}
\BookLine{106.50}{22.14}{}{0}{们，谁四{\kaifont （是）}最美滴{\kaifont （的）}。}
\BookLine{123.50}{0.00}{\heifont}{0}{阿波罗：}
\BookLine{132.00}{22.14}{}{0}{不！不！我不敢！}
\BookLine{149.00}{0.00}{\heifont}{0}{拉刻西斯：}
\BookLine{157.50}{22.14}{}{0}{{\kaifont （尖声地）}你必须！}
\newpage
% 原书页码：254
\null
\BookNumber{254}
\Vertical{4}{169}{}{6.2}{8}阿\\尔\\刻\\提\\斯\\之\\歌\EndVertical
\BookLine{30.00}{0.00}{\heifont}{0}{阿波罗：}
\BookLine{38.50}{22.14}{}{0}{要是我赢了呢？}
\BookLine{55.50}{0.00}{\heifont}{0}{阿特洛波斯：}
\BookLine{64.00}{22.14}{}{0}{赢？白痴！蠢货！奴才！没人赢过我们！}
\BookLine{81.00}{0.00}{\heifont}{0}{阿波罗：}
\BookLine{89.50}{22.14}{}{0}{但是，要是我赢了呢？}
\BookLine{106.50}{0.00}{\heifont}{0}{拉刻西斯：}
\BookLine{115.00}{22.14}{}{0}{他知不知道我们是谁！}
\BookLine{132.00}{0.00}{\heifont}{0}{阿波罗：}
\BookLine{140.50}{22.14}{}{0}{但是，要是我赢了呢？}
\BookLine{157.50}{0.00}{\heifont}{0}{克洛托：}
\BookLine{166.00}{22.14}{}{0}{这个傻子还敢嗦{\kaifont （说）}赢！}
\newpage
% 原书页码：255
\null
\BookNumber{255}
\Vertical{129}{11}{}{7}{10}醉\\酒\\的\\女\\神\EndVertical
\BookLine{30.00}{0.00}{\heifont}{0}{阿波罗：}
\BookLine{38.50}{22.14}{}{0}{要是我赢了，你们必须满足我一个愿望。一个}
\BookLine{47.00}{22.14}{}{0}{愿望，随便什么愿望。}
\BookLine{64.00}{0.00}{\heifont}{0}{拉刻西斯：}
\BookLine{72.50}{22.14}{}{0}{行，行啊。哦，你可真无聊！快说谜语吧，是}
\BookLine{81.00}{22.14}{}{0}{什么必须——}
\BookLine{98.00}{0.00}{\heifont}{0}{阿波罗：}
\BookLine{106.50}{22.14}{}{0}{指着黄泉起誓！}
\BookLine{123.50}{0.00}{\heifont}{0}{克洛托和拉刻西斯：}
\BookLine{132.00}{22.14}{}{0}{我发誓！指着黄泉起誓！指着黄泉起誓！}
\BookLine{149.00}{0.00}{\heifont}{0}{阿波罗：}
\BookLine{157.50}{22.14}{}{0}{{\kaifont （对阿特洛波斯）}你也发誓。}
\newpage
% 原书页码：256
\null
\BookNumber{256}
\Vertical{4}{169}{}{6.2}{8}阿\\尔\\刻\\提\\斯\\之\\歌\EndVertical
\BookLine{30.00}{0.00}{\heifont}{0}{阿特洛波斯：}
\BookLine{38.50}{22.14}{}{0}{{\kaifont （思考了一下，大声地）}指着黄泉起誓！}
\BookLine{55.50}{0.00}{\heifont}{0}{阿波罗：}
\BookLine{64.00}{22.14}{}{0}{那么，准备好了吗？}
\BookLine{81.00}{0.00}{\heifont}{0}{拉刻西斯：}
\BookLine{89.50}{22.14}{}{0}{等一下！一下。{\kaifont （靠近阿特洛波斯，悄悄地）}太}
\BookLine{98.00}{22.14}{}{0}{阳快落山了。你别忘了阿德——你知道，生命}
\BookLine{106.50}{22.14}{}{0}{线，阿德——}
\BookLine{123.50}{0.00}{\heifont}{0}{阿特洛波斯：}
\BookLine{132.00}{22.14}{}{0}{什么？什么阿德？你在嘀嘀咕咕什么呢，}
\BookLine{140.50}{22.14}{}{0}{傻子？}
\BookLine{157.50}{0.00}{\heifont}{0}{拉刻西斯：}
\BookLine{166.00}{22.14}{}{0}{{\kaifont （大声）}你别忘了塞萨利国王，阿德墨托斯的生}
\BookLine{174.50}{22.14}{}{0}{命线。日落时分，阿特洛波斯，你是聋了吗？}
\newpage
% 原书页码：257
\null
\BookNumber{257}
\Vertical{129}{11}{}{7}{10}醉\\酒\\的\\女\\神\EndVertical
\BookLine{30.00}{0.00}{\heifont}{0}{阿特洛波斯：}
\BookLine{38.50}{22.14}{}{0}{噢，你别再啰里啰唆了，管好你自己的事吧。}
\BookLine{47.00}{22.14}{}{0}{我当然没聋！孩子，接着说你的谜语！}
\BookLine{64.00}{0.00}{\heifont}{0}{阿波罗：}
\BookLine{72.50}{22.14}{}{0}{好！你们的猜谜时间，就是背诵缪斯女神和她}
\BookLine{81.00}{22.14}{}{0}{们母亲姓名的时间。}
\BookLine{98.00}{0.00}{\heifont}{0}{拉刻西斯：}
\BookLine{106.50}{22.14}{}{0}{嗯！九位女神，一位母亲。行，开始！}
\BookLine{123.50}{0.00}{\heifont}{0}{阿波罗：}
\BookLine{132.00}{22.14}{}{0}{对每个生命来说都是必须——却只能拯救其中}
\BookLine{140.50}{22.14}{}{0}{之一，是什么？}
\BookLine{157.50}{7.38}{\kaifont}{0}{三姐妹闭上眼睛前后摇晃着身体，嘴里嘟囔着谜题。阿}
\BookLine{166.00}{0.00}{\kaifont}{0}{波罗突然开始向缪斯女神们唱起祈祷的歌谣来。}
\newpage
% 原书页码：258
\null
\BookNumber{258}
\Vertical{4}{169}{}{6.2}{8}阿\\尔\\刻\\提\\斯\\之\\歌\EndVertical
\BookLine{30.00}{0.00}{\heifont}{0}{阿波罗：}
\BookLine{38.50}{22.14}{}{0}{摩涅莫绪涅，九女神的母亲；}
\BookLine{47.00}{22.14}{}{0}{波吕许谟尼亚，众神的香气——}
\BookLine{64.00}{0.00}{\heifont}{0}{拉刻西斯：}
\BookLine{72.50}{22.14}{}{0}{{\kaifont （尖叫道）}别唱了！这不公平！你这样我们怎么}
\BookLine{81.00}{22.14}{}{0}{思考？}
\BookLine{98.00}{0.00}{\heifont}{0}{克洛托：}
\BookLine{106.50}{22.14}{}{0}{姐姐，你把耳朵捂上啊。}
\BookLine{123.50}{0.00}{\heifont}{0}{阿特洛波斯：}
\BookLine{132.00}{22.14}{}{0}{不公平！{\kaifont （嘀囔着）}对每个生命来说都是}
\BookLine{140.50}{22.14}{}{0}{必须——}
\BookLine{157.50}{7.38}{\kaifont}{0}{她们把手指插进耳朵里。}
\newpage
% 原书页码：259
\null
\BookNumber{259}
\Vertical{129}{11}{}{7}{10}醉\\酒\\的\\女\\神\EndVertical
\BookLine{30.00}{0.00}{\heifont}{0}{阿波罗：}
\BookLine{38.50}{22.14}{}{0}{埃拉托，爱的声音；}
\BookLine{47.00}{22.14}{}{0}{欧忒耳珀，请帮助我。}
\BookLine{55.50}{22.14}{}{0}{卡利俄珀，灵魂的窃贼；}
\BookLine{64.00}{22.14}{}{0}{乌拉尼亚，被众星包裹；}
\BookLine{72.50}{22.14}{}{0}{克利俄，回望的目光；}
\BookLine{81.00}{22.14}{}{0}{欧忒耳珀，请帮助我。}
\BookLine{89.50}{22.14}{}{0}{忒耳西科瑞，美丽脚踝的守护者；}
\BookLine{98.00}{22.14}{}{0}{塔利亚，带来漫长的欢乐；}
\BookLine{106.50}{22.14}{}{0}{墨尔波墨涅，令人害怕却也喜爱；}
\BookLine{115.00}{22.14}{}{0}{欧忒耳珀，请帮助我。}
\BookLine{123.50}{22.14}{}{0}{{\kaifont （突然大声地）}惩罚！惩罚！}
\BookLine{140.50}{7.38}{\kaifont}{0}{克洛托和阿特洛波斯将脸埋在拉刻西斯的肩上，抱怨}
\BookLine{149.00}{0.00}{\kaifont}{0}{叹息。}
\BookLine{166.00}{0.00}{\heifont}{0}{拉刻西斯：}
\BookLine{174.50}{22.14}{}{0}{{\kaifont （有气无力）}答案是什么？}
\newpage
% 原书页码：260
\null
\BookNumber{260}
\Vertical{4}{169}{}{6.2}{8}阿\\尔\\刻\\提\\斯\\之\\歌\EndVertical
\BookLine{30.00}{0.00}{\heifont}{0}{阿波罗：}
\BookLine{38.50}{22.14}{}{0}{{\kaifont （把帽子丢在一边，带着胜利的语气）}是我！太}
\BookLine{47.00}{22.14}{}{0}{阳神阿波罗！}
\BookLine{64.00}{0.00}{\heifont}{0}{三姐妹：}
\BookLine{72.50}{22.14}{}{0}{阿波罗！是你？}
\BookLine{89.50}{0.00}{\heifont}{0}{拉刻西斯：}
\BookLine{98.00}{22.14}{}{0}{{\kaifont （十分生气）}呸！你能救谁的命？}
\BookLine{115.00}{0.00}{\heifont}{0}{阿波罗：}
\BookLine{123.50}{22.14}{}{0}{我赢得的！我的愿望！一个人的性命！塞萨利}
\BookLine{132.00}{22.14}{}{0}{之王，阿德墨托斯的性命。}
\BookLine{149.00}{7.38}{\kaifont}{0}{三姐妹之中爆发了一阵令人害怕的喧闹声。}
\BookLine{166.00}{0.00}{\heifont}{0}{三姐妹：}
\BookLine{174.50}{22.14}{}{0}{骗子！不可能！这你别想了！}
\newpage
% 原书页码：261
\null
\BookNumber{261}
\Vertical{129}{11}{}{7}{10}醉\\酒\\的\\女\\神\EndVertical
\BookLine{30.00}{0.00}{\heifont}{0}{阿波罗：}
\BookLine{38.50}{22.14}{}{0}{指着黄泉起誓。}
\BookLine{55.50}{0.00}{\heifont}{0}{三姐妹：}
\BookLine{64.00}{22.14}{}{0}{这会违反所有的法则。宙斯会审判的，骗子。}
\BookLine{81.00}{0.00}{\heifont}{0}{阿波罗：}
\BookLine{89.50}{22.14}{}{0}{{\kaifont （警告）}指着黄泉起誓。}
\BookLine{106.50}{0.00}{\heifont}{0}{三姐妹：}
\BookLine{115.00}{22.14}{}{0}{宙斯！我们要去找宙斯说这件事，他会决定的。}
\BookLine{132.00}{0.00}{\heifont}{0}{阿波罗：}
\BookLine{140.50}{22.14}{}{0}{宙斯也是指着黄泉起誓，并且会遵守自己的}
\BookLine{149.00}{22.14}{}{0}{誓言。}
\BookLine{166.00}{7.38}{\kaifont}{0}{突然静场。}
\newpage
% 原书页码：262
\null
\BookNumber{262}
\Vertical{4}{169}{}{6.2}{8}阿\\尔\\刻\\提\\斯\\之\\歌\EndVertical
\BookLine{30.00}{0.00}{\heifont}{0}{阿特洛波斯：}
\BookLine{38.50}{22.14}{}{0}{{\kaifont （下了决心，但带着不祥的语气）}你可以实}
\BookLine{47.00}{22.14}{}{0}{现你的愿望——阿德墨托斯国王的性命。}
\BookLine{55.50}{22.14}{}{0}{不过——}
\BookLine{72.50}{0.00}{\heifont}{0}{阿波罗：}
\BookLine{81.00}{22.14}{}{0}{{\kaifont （带着胜利的语气）}我救了阿德墨托斯的命！}
\BookLine{98.00}{0.00}{\heifont}{0}{三姐妹：}
\BookLine{106.50}{22.14}{}{0}{不过——}
\BookLine{123.50}{0.00}{\heifont}{0}{阿波罗：}
\BookLine{132.00}{22.14}{}{0}{我救了阿德墨托斯的命！你还要说什么？}
\BookLine{149.00}{0.00}{\heifont}{0}{阿特洛波斯：}
\BookLine{157.50}{22.14}{}{0}{需要有人替他去死。}
\newpage
% 原书页码：263
\null
\BookNumber{263}
\Vertical{129}{11}{}{7}{10}醉\\酒\\的\\女\\神\EndVertical
\BookLine{30.00}{0.00}{\heifont}{0}{阿波罗：}
\BookLine{38.50}{22.14}{}{0}{{\kaifont （轻松地）}哦——随便选个奴隶就好。在你们那}
\BookLine{47.00}{22.14}{}{0}{些线里找一根灰色、油腻的线就好了，圣洁的}
\BookLine{55.50}{22.14}{}{0}{拉刻西斯。}
\BookLine{72.50}{0.00}{\heifont}{0}{拉刻西斯：}
\BookLine{81.00}{22.14}{}{0}{{\kaifont （十分愤怒）}什么？你让我们拿走别人的性命？}
\BookLine{98.00}{0.00}{\heifont}{0}{阿特洛波斯：}
\BookLine{106.50}{22.14}{}{0}{你让我们杀人？}
\BookLine{123.50}{0.00}{\heifont}{0}{克洛托：}
\BookLine{132.00}{22.14}{}{0}{阿波罗觉得我们都是刽子手？}
\BookLine{149.00}{0.00}{\heifont}{0}{阿波罗：}
\BookLine{157.50}{22.14}{}{0}{{\kaifont （感到有些畏惧）}那，伟大的命运女神们，这件}
\BookLine{166.00}{22.14}{}{0}{事怎么办呢？}
\newpage
% 原书页码：264
\null
\BookNumber{264}
\Vertical{4}{169}{}{6.2}{8}阿\\尔\\刻\\提\\斯\\之\\歌\EndVertical
\BookLine{30.00}{0.00}{\heifont}{0}{拉刻西斯：}
\BookLine{38.50}{22.14}{}{0}{我——一个杀手？{\kaifont （张开双臂，庄严地）}我的左}
\BookLine{47.00}{22.14}{}{0}{手是人生际遇，我的右手是自然规律。}
\BookLine{64.00}{0.00}{\heifont}{0}{阿波罗：}
\BookLine{72.50}{22.14}{}{0}{那，仁慈的命运女神们，这件事怎么办呢？}
\BookLine{89.50}{0.00}{\heifont}{0}{拉刻西斯：}
\BookLine{98.00}{22.14}{}{0}{必须有人为阿德墨托斯国王献出自己的生命，}
\BookLine{106.50}{22.14}{}{0}{出于自己的选择和意愿。对于这样的死亡，我}
\BookLine{115.00}{22.14}{}{0}{们无法掌控，无论是人生际遇还是自然规律，}
\BookLine{123.50}{22.14}{}{0}{都无法控制出于自由意志的奉献。必须有人选}
\BookLine{132.00}{22.14}{}{0}{择代替塞萨利国王阿德墨托斯去死。}
\BookLine{149.00}{0.00}{\heifont}{0}{阿波罗：}
\BookLine{157.50}{22.14}{}{0}{{\kaifont （用手捂住自己的脸）}不！不！我知道要发生什}
\BookLine{166.00}{22.14}{}{0}{么了！{\kaifont （大喊道）}阿尔刻提斯！阿尔刻提斯！}
\newpage
% 原书页码：265
\null
\BookNumber{265}
\Vertical{129}{11}{}{7}{10}醉\\酒\\的\\女\\神\EndVertical
\BookLine{30.00}{7.38}{\kaifont}{0}{他跌跌撞撞地跑下场。}
\BookLine{47.00}{7.38}{\kaifont}{0}{——剧终——}
\end{document}
